\documentclass[11pt,twoside]{book}

% ── Packages ──────────────────────────────────────────────────────────
\usepackage[utf8]{inputenc}
\usepackage[T1]{fontenc}
\usepackage{amsmath,amssymb,amsthm}
\usepackage{mathrsfs}
\usepackage{enumerate}
\usepackage{hyperref}
\usepackage{cleveref}
\usepackage{listings}
\usepackage{xcolor}
\usepackage{geometry}
\usepackage{fancyhdr}
\usepackage{titlesec}
\usepackage{epigraph}
% \usepackage{microtype}  % requires scalable fonts

\geometry{
  papersize={7in,10in},
  margin=1in,
  inner=1.2in,
  outer=0.8in
}

% ── Theorem environments ──────────────────────────────────────────────
\theoremstyle{plain}
\newtheorem{theorem}{Theorem}[chapter]
\newtheorem{lemma}[theorem]{Lemma}
\newtheorem{proposition}[theorem]{Proposition}
\newtheorem{corollary}[theorem]{Corollary}

\theoremstyle{definition}
\newtheorem{definition}[theorem]{Definition}
\newtheorem{example}[theorem]{Example}
\newtheorem{remark}[theorem]{Remark}

\theoremstyle{remark}
\newtheorem*{notation}{Notation}
\newtheorem*{convention}{Convention}

% ── Lean 4 code listings ─────────────────────────────────────────────
\lstdefinelanguage{Lean4}{
  morekeywords={theorem,lemma,def,structure,class,instance,namespace,
    section,end,open,import,variable,where,by,have,let,show,calc,
    match,with,if,then,else,do,return,fun,forall,exists,Prop,Type,
    sorry,exact,apply,intro,rw,simp,linarith,omega,induction,cases,
    constructor,refine,obtain,rcases,ext,congr,rfl,trivial,
    noncomputable,axiom,opaque,set_option,attribute},
  sensitive=true,
  morecomment=[l]{--},
  morecomment=[s]{/-}{-/},
  morestring=[b]",
  literate=
    {→}{{$\to$}}1
    {←}{{$\leftarrow$}}1
    {↔}{{$\leftrightarrow$}}1
    {∀}{{$\forall$}}1
    {∃}{{$\exists$}}1
    {λ}{{$\lambda$}}1
    {≤}{{$\leq$}}1
    {≥}{{$\geq$}}1
    {≠}{{$\neq$}}1
    {∧}{{$\wedge$}}1
    {∨}{{$\vee$}}1
    {¬}{{$\neg$}}1
    {∘}{{$\circ$}}1
    {⊆}{{$\subseteq$}}1
    {ℕ}{{$\mathbb{N}$}}1
    {ℤ}{{$\mathbb{Z}$}}1
    {ℝ}{{$\mathbb{R}$}}1
}

\lstset{
  language=Lean4,
  basicstyle=\small\ttfamily,
  keywordstyle=\color{blue}\bfseries,
  commentstyle=\color{gray}\itshape,
  stringstyle=\color{red!70!black},
  breaklines=true,
  frame=single,
  backgroundcolor=\color{gray!5},
  numbers=left,
  numberstyle=\tiny\color{gray},
  xleftmargin=2em,
  framexleftmargin=1.5em,
  captionpos=b,
  aboveskip=1em,
  belowskip=1em
}

% ── Headers ───────────────────────────────────────────────────────────
\pagestyle{fancy}
\fancyhf{}
\fancyhead[LE]{\small\itshape\nouppercase{\leftmark}}
\fancyhead[RO]{\small\itshape\nouppercase{\rightmark}}
\fancyfoot[C]{\thepage}
\renewcommand{\headrulewidth}{0.4pt}

% ── Custom commands ───────────────────────────────────────────────────
\newcommand{\IS}{\mathcal{I}}          % Ideatic space
\newcommand{\compose}{\circ}           % Composition operator
\newcommand{\void}{\mathbf{0}}         % Void element
\newcommand{\depth}{\operatorname{d}}  % Depth function
\newcommand{\res}{\sim}                % Resonance relation
\newcommand{\interp}[2]{#1 \compose #2}  % Interpretation
\newcommand{\Rep}{\mathcal{R}}         % Repertoire
\newcommand{\Sig}{\mathcal{S}}         % Signal space
\newcommand{\Pay}{\pi}                 % Payoff

% ── Title ─────────────────────────────────────────────────────────────
\title{
  \vspace{-2cm}
  {\LARGE\textsc{Signalling Games in}}\\[0.3cm]
  {\Huge\textbf{Ideatic Diffusion Theory}}\\[0.5cm]
  {\large\textit{A Formal Theory of Meaning, Interpretation, and Strategic Communication}}\\[1cm]
  {\normalsize With 364 Machine-Verified Theorems in Lean 4}
}
\author{
  \textsc{Halley Young}\\[0.3cm]
  \small\textit{University of Pennsylvania}
}
\date{\today}

% ══════════════════════════════════════════════════════════════════════
\begin{document}

\frontmatter
\maketitle

% ── Epigraph ──────────────────────────────────────────────────────────
\newpage
\thispagestyle{empty}
\vspace*{3cm}
\begin{center}
\textit{``The limits of my language mean the limits of my world.''}\\[0.3cm]
--- Ludwig Wittgenstein, \textit{Tractatus Logico-Philosophicus} (5.6)
\vspace{2cm}

\textit{``A sign is something which stands to somebody for something\\
in some respect or capacity.''}\\[0.3cm]
--- Charles Sanders Peirce, \textit{Collected Papers} (2.228)
\vspace{2cm}

\textit{``Culture is public because meaning is.''}\\[0.3cm]
--- Clifford Geertz, \textit{The Interpretation of Cultures}
\end{center}

% ── Preface ───────────────────────────────────────────────────────────
\chapter*{Preface}
\addcontentsline{toc}{chapter}{Preface}

This book develops a formal theory of signalling games within the framework
of Ideatic Diffusion Theory (IDT). Where classical game theory treats
signals as tokens in an abstract communication channel, we treat them
as \emph{ideas}---elements of an ideatic space equipped with composition,
depth, and resonance. The central insight is that \emph{interpretation
is composition}: when a receiver with background idea $r$ encounters
signal $s$, they do not simply ``decode'' it but rather \emph{compose}
it with their existing understanding, producing the interpretation
$r \compose s$.

This compositional view of interpretation has profound consequences for
game theory. Payoffs depend not on the signal itself but on what
receivers \emph{make of it}---how it composes with their existing repertoire
of ideas. A signal that is profound to one receiver may be trivial to
another, not because the signal differs but because their repertoires do.
This is the formal kernel of hermeneutic theory: meaning is not in the
text but in the encounter between text and reader.

The book is organized in four parts. Part~I (\Cref{ch:hermeneutic,ch:signalling,ch:costly,ch:activation,ch:misinterp})
establishes the interpretation framework: repertoires, signalling problems,
costly signalling, resonance activation, and misinterpretation. Part~II
(\Cref{ch:equilibrium,ch:persuasion,ch:lossy,ch:echo,ch:multireceiver})
develops equilibrium theory, persuasion bounds, information loss, echo
chambers, and multi-receiver signalling. Part~III
(\Cref{ch:marketplace,ch:matching,ch:coalition,ch:auction,ch:curation})
studies populations: competition, matching, coalitions, auctions, and
repertoire curation. Part~IV
(\Cref{ch:cooperative,ch:bargaining,ch:network,ch:social,ch:synthesis})
addresses cooperation, bargaining, network diffusion, social choice,
and offers a synthesis.

Every theorem in this book has been formally verified in Lean~4,
building on the \texttt{IDT\_Foundations} library. There are no
\texttt{sorry} axioms. The 364 machine-checked theorems constitute,
to our knowledge, the first comprehensive formal treatment of signalling
games in an algebraic theory of ideas.

\vspace{1cm}
\noindent\textit{Philadelphia, 2026}

% ── Table of Contents ─────────────────────────────────────────────────
\tableofcontents

% ══════════════════════════════════════════════════════════════════════
\mainmatter

% ── Part I: The Interpretation Framework ──────────────────────────────
\part{The Interpretation Framework}
\label{part:interpretation}

% ══════════════════════════════════════════════════════════════════════
% Chapter 1: Hermeneutic Foundations
% ══════════════════════════════════════════════════════════════════════

\chapter{Hermeneutic Foundations: Repertoires, Interpretation, and the Hermeneutic Bound}
\label{ch:hermeneutic}

\epigraph{``The limits of my language mean the limits of my world.''}{Ludwig Wittgenstein, \textit{Tractatus Logico-Philosophicus} (5.6)}

% ── Introduction ──────────────────────────────────────────────────────

Every human mind is a repertoire: a stock of ideas accumulated through
enculturation.  When a signal arrives---a gesture, an utterance, a ritual
act---the receiver does not passively absorb it.  They \emph{interpret}
it by composing it with their existing ideas.  This chapter builds the
foundational structures on which the entire sequel rests.

In the language of Bourdieu~\cite{bourdieu1984}, a repertoire is the
cognitive component of \emph{habitus}---the durable, transposable
dispositions through which agents perceive and act in the social world.
In the language of Saussure~\cite{saussure1916}, each element of the
repertoire is a signified, and interpretation is the operation that maps
a signifier to a new signified by composing it with an existing one.  In
the language of Peirce~\cite{peirce1931}, interpretation is the
generation of an \emph{interpretant} from the encounter of a sign with a
prior sign.

The payoff in a signalling game is \emph{never} about the signal
itself---it is always about the interpretations the signal produces in
other minds.  This insight, already present in Gadamer's
\emph{Truth and Method}~\cite{gadamer1960}, receives its first fully
formal treatment here.

\medskip
\noindent\textbf{Chapter overview.}  We introduce four core concepts:
\emph{repertoire} (\S\ref{sec:ch1:core}), \emph{interpretation} and the
\emph{hermeneutic bound} (\S\ref{sec:ch1:hermeneutic-bound}), \emph{silence}
(\S\ref{sec:ch1:silence}), and \emph{cultural consensus}
(\S\ref{sec:ch1:consensus}).
We close with the \emph{telephone game} (\S\ref{sec:ch1:telephone}) and
structural invariants (\S\ref{sec:ch1:invariants}).  All sixteen theorems
are machine-verified in Lean~4 with zero \texttt{sorry} axioms.

% ── §1. Core Definitions ─────────────────────────────────────────────
\section{Core Definitions}
\label{sec:ch1:core}

We work within an ideatic space $(\IS, \compose, \void, \depth, \res)$
as established in the IDT Foundations.

\begin{definition}[Repertoire]
\label{def:repertoire}
A \emph{repertoire} is a finite list of ideas:
\[
  \Rep = [r_1, r_2, \ldots, r_n] \in \IS^{*}.
\]
In Bourdieu's terms, this is the \emph{habitus} made formal---the
totality of schemas through which an agent perceives and interprets the
social world.
\end{definition}

\begin{definition}[Interpretation]
\label{def:interpretation}
When a receiver with repertoire $\Rep = [r_1, \ldots, r_n]$ encounters
signal~$s$, they produce the list of interpretations
\[
  \operatorname{interpret}(\Rep, s) 
  = [r_1 \compose s,\; r_2 \compose s,\; \ldots,\; r_n \compose s].
\]
This is Gadamer's \emph{Horizontverschmelzung} (fusion of horizons):
each background idea merges with the incoming signal.
\end{definition}

In Lean~4, the definition is:

\begin{lstlisting}[caption={Interpretation in Lean~4},label=lst:interpret]
def interpret (rep : Repertoire I) (signal : I) : List I :=
  rep.map (fun r => IdeaticSpace.compose r signal)
\end{lstlisting}

\begin{definition}[Maximum Interpretation Depth]
\label{def:maxinterpdepth}
The \emph{maximum interpretation depth} of a list of ideas is
\[
  \operatorname{maxInterpDepth}([a_1, \ldots, a_n])
  = \max_{i} \depth(a_i).
\]
\end{definition}

\begin{definition}[Depth Sum]
\label{def:depthsum}
The \emph{depth sum} of a list of ideas is
$\operatorname{depthSum}([a_1, \ldots, a_n]) = \sum_i \depth(a_i)$.
\end{definition}

% ── §2. The Hermeneutic Bound ────────────────────────────────────────
\section{The Hermeneutic Bound}
\label{sec:ch1:hermeneutic-bound}

``You can only understand what you are prepared to understand.''
The following two theorems make this aphorism precise.

\begin{theorem}[Total Hermeneutic Bound]
\label{thm:total-hermeneutic-bound}
For any repertoire $\Rep$ and signal~$s$,
\[
  \operatorname{depthSum}(\operatorname{interpret}(\Rep, s))
  \;\leq\;
  \operatorname{depthSum}(\Rep) + |\Rep| \cdot \depth(s).
\]
\end{theorem}

\begin{proof}[Proof sketch]
By induction on $|\Rep|$.  The base case ($\Rep = []$) is trivial.
For $\Rep = r :: \Rep'$, we use the depth-composition bound
$\depth(r \compose s) \leq \depth(r) + \depth(s)$ and the
inductive hypothesis on~$\Rep'$.
The full proof is verified in \texttt{IDT\_Signal\_Ch01.lean},
theorem \texttt{total\_hermeneutic\_bound}.
\end{proof}

\begin{remark}[Mass Media]
This theorem explains why mass media (one signal, many receivers) has
bounded total interpretive impact per receiver.  A richer repertoire
absorbs more, but the signal's contribution scales only linearly with
audience size.
\end{remark}

\begin{theorem}[MaxDepth Interpretation Bound]
\label{thm:max-interp-depth}
For any repertoire $\Rep$ and signal~$s$,
\[
  \operatorname{maxInterpDepth}(\operatorname{interpret}(\Rep, s))
  \;\leq\;
  \operatorname{maxInterpDepth}(\Rep) + \depth(s).
\]
\end{theorem}

\begin{proof}[Proof sketch]
By induction on $|\Rep|$.  At each step, we bound
$\depth(r \compose s) \leq \depth(r) + \depth(s)$ and propagate
through the maximum.  See theorem \texttt{max\_interp\_depth\_bound}.
\end{proof}

\begin{remark}
Even the most sophisticated listener adds only $\depth(s)$ to their
best existing understanding.  No single signal can vault a receiver
beyond their current horizon by more than the signal's intrinsic
complexity.
\end{remark}

% ── §3. Silence and the Unsaid ───────────────────────────────────────
\section{Silence and the Unsaid}
\label{sec:ch1:silence}

\begin{theorem}[Silence is Transparent]
\label{thm:silence-transparent}
The void signal leaves the mind unchanged:
\[
  \operatorname{interpret}(\Rep, \void) = \Rep.
\]
\end{theorem}

\begin{proof}[Proof sketch]
By induction on $\Rep$.  For each $r \in \Rep$, the void-right identity
$r \compose \void = r$ yields the result immediately.
See \texttt{silence\_is\_transparent}.
\end{proof}

\begin{remark}[Ritual Silence]
Quaker meetings, Buddhist \emph{noble silence}, the minute of silence
for the dead---all exploit this principle.  By saying nothing, the
receiver's existing framework remains undisturbed.  The void signal is
the formal foundation of taboo: what must not be said leaves the mind in
its current state.
\end{remark}

\begin{theorem}[The Blank Receiver]
\label{thm:blank-receiver}
A receiver whose only idea is $\void$ interprets any signal as the
signal itself:
\[
  \operatorname{interpret}([\void], s) = [s].
\]
\end{theorem}

\begin{proof}[Proof sketch]
Direct computation: $\void \compose s = s$ by the void-left identity.
See \texttt{blank\_receiver}.
\end{proof}

\begin{remark}[The Objective Observer]
The impossible ideal of the ``objective observer'' requires total
cultural amnesia: a repertoire consisting solely of~$\void$.  Only then
does interpretation reduce to transparent reception.
\end{remark}

\begin{theorem}[Negative Capability / Keats's Theorem]
\label{thm:negative-capability}
If $\void \in \Rep$, then $s \in \operatorname{interpret}(\Rep, s)$.
\end{theorem}

\begin{proof}[Proof sketch]
Since $\void \in \Rep$, the map sends $\void \mapsto \void \compose s = s$,
so $s$ appears in the output.
See \texttt{negative\_capability}.
\end{proof}

\begin{remark}
Keats described \emph{negative capability} as the capacity for ``being
in uncertainties, Mysteries, doubts, without any irritable reaching after
fact \& reason'' \cite{keats1817}.  In our framework, a mind retaining a
slot of pure receptivity ($\void$ in the repertoire) captures the signal
in its original form alongside culturally mediated versions.
\end{remark}

% ── §4. Cultural Consensus ───────────────────────────────────────────
\section{Cultural Consensus}
\label{sec:ch1:consensus}

The next three theorems establish that resonance in backgrounds implies
resonance in interpretations---the structural foundation of cultural
consensus.

\begin{theorem}[Durkheimian Consensus]
\label{thm:durkheimian-consensus}
If $r_1 \res r_2$ and both receive signal~$s$, then
$r_1 \compose s \;\res\; r_2 \compose s$.
\end{theorem}

\begin{proof}[Proof sketch]
Direct application of the resonance-compose-right lemma from IDT
Foundations.  See \texttt{durkheimian\_consensus}.
\end{proof}

\begin{remark}
Two members of the same culture (shared myths, shared categories)
\emph{always} produce resonant interpretations of the same event.
Cultural consensus is a structural consequence of shared
background---Durkheim's \emph{conscience collective}~\cite{durkheim1893}.
\end{remark}

\begin{theorem}[L\'evi-Strauss's Mythical Transformation]
\label{thm:mythical-transformation}
For any receiver~$r$, if $s_1 \res s_2$ then
$r \compose s_1 \;\res\; r \compose s_2$.
\end{theorem}

\begin{proof}[Proof sketch]
By resonance-compose-left.  See \texttt{mythical\_transformation}.
\end{proof}

\begin{remark}
Ritual repetition-with-variation works because the participant's
invariant background ensures that structurally similar performances
produce resonant experiences.  L\'evi-Strauss's mythical
``transformations'' are perceived as ``the same story told differently''
precisely because the listener's background resonates with both
variants~\cite{levistrauss1962}.
\end{remark}

\begin{theorem}[Intercultural Solidarity / Turner's Communitas]
\label{thm:turners-communitas}
If $r_1 \res r_2$ and $s_1 \res s_2$, then
\[
  r_1 \compose s_1 \;\res\; r_2 \compose s_2.
\]
\end{theorem}

\begin{proof}[Proof sketch]
By the full resonance-compose lemma.  See \texttt{turners\_communitas}.
\end{proof}

\begin{remark}
Members of allied cultures witnessing similar rituals produce resonant
interpretations.  This is the formal basis of what Turner called
\emph{communitas}~\cite{turner1969}: the spontaneous solidarity that
arises when structurally similar participants undergo similar experiences.
\end{remark}

% ── §5. The Telephone Game ───────────────────────────────────────────
\section{The Telephone Game: Cultural Transmission Chains}
\label{sec:ch1:telephone}

\begin{theorem}[Bartlett's Serial Reproduction]
\label{thm:bartlett}
Two-stage interpretation through different minds equals single-stage
interpretation through the composed mind:
\[
  r_B \compose (r_A \compose s)
  \;=\;
  (r_B \compose r_A) \compose s.
\]
\end{theorem}

\begin{proof}[Proof sketch]
This is the associativity axiom of the ideatic space, applied to the
triple $(r_B, r_A, s)$.
See \texttt{bartletts\_serial\_reproduction}.
\end{proof}

\begin{remark}
Bartlett's (1932) serial reproduction experiments~\cite{bartlett1932}
showed that stories change predictably through chains of retelling.  This
theorem shows \emph{why}: the chain of interpreters is structurally
equivalent to a single composite interpreter.  Oral tradition scholars
can study the ``net effect'' of a transmission chain without tracking
each individual retelling.
\end{remark}

\begin{theorem}[Telephone Depth Bound]
\label{thm:telephone-depth}
The depth after two-stage interpretation satisfies
\[
  \depth(r_B \compose (r_A \compose s))
  \;\leq\;
  \depth(r_B) + \depth(r_A) + \depth(s).
\]
\end{theorem}

\begin{proof}[Proof sketch]
Rewrite using associativity, then apply depth-compose-le twice.
See \texttt{telephone\_depth}.
\end{proof}

\begin{theorem}[Resonant Transmission Lines]
\label{thm:resonant-transmission}
If $r_{A_1} \res r_{A_2}$ and $r_{B_1} \res r_{B_2}$, then
\[
  r_{B_1} \compose (r_{A_1} \compose s)
  \;\res\;
  r_{B_2} \compose (r_{A_2} \compose s).
\]
\end{theorem}

\begin{proof}[Proof sketch]
Apply associativity to both sides, then use resonance-compose twice.
See \texttt{resonant\_transmission\_lines}.
\end{proof}

\begin{remark}
Independent lines of ritual transmission (e.g., different griots in West
Africa, different monastic orders transmitting the same texts) produce
resonant outcomes if corresponding specialists resonate.  Ritual
orthodoxy can be maintained across geographically separated communities
without central authority.
\end{remark}

% ── §6. Structural Invariants ────────────────────────────────────────
\section{Structural Invariants of Interpretation}
\label{sec:ch1:invariants}

\begin{theorem}[Void is Unpersuasive]
\label{thm:void-unpersuasive}
$\depth(r \compose \void) = \depth(r)$.  Silence contributes zero
marginal depth.
\end{theorem}

\begin{proof}[Proof sketch]
By the depth-compose-void identity.  See \texttt{void\_is\_unpersuasive}.
\end{proof}

\begin{remark}[Game-theoretic significance]
The null message (babbling) contributes zero information.  This is
\emph{why} babbling equilibria always exist: silence costs nothing and
changes nothing.
\end{remark}

\begin{theorem}[Interpretation Conserves Count]
\label{thm:conserves-count}
$|\operatorname{interpret}(\Rep, s)| = |\Rep|$.
\end{theorem}

\begin{proof}[Proof sketch]
By the length-preserving property of \texttt{List.map}.
See \texttt{interpretation\_conserves\_count}.
\end{proof}

\begin{theorem}[Durkheim's Division of Labor]
\label{thm:division-of-labor}
$\operatorname{interpret}(\Rep_1 \mathbin{+\!+} \Rep_2, s)
 = \operatorname{interpret}(\Rep_1, s) \mathbin{+\!+}
   \operatorname{interpret}(\Rep_2, s)$.
\end{theorem}

\begin{proof}[Proof sketch]
By \texttt{List.map\_append}.  See \texttt{durkheims\_division\_of\_labor}.
\end{proof}

\begin{remark}
A polymath's interpretation equals the union of specialist
interpretations.  Interdisciplinary understanding is \emph{not}
synergistic---it is the disjoint union of specialist understandings.
This formalizes Durkheim's cognitive division of labor~\cite{durkheim1893}.
\end{remark}

\begin{theorem}[Silence Preserves DepthSum]
\label{thm:silence-depthsum}
$\operatorname{depthSum}(\operatorname{interpret}(\Rep, \void))
 = \operatorname{depthSum}(\Rep)$.
\end{theorem}

\begin{proof}[Proof sketch]
Follows immediately from Theorem~\ref{thm:silence-transparent}.
See \texttt{silence\_preserves\_depthSum}.
\end{proof}

\begin{theorem}[Three-Stage Telephone]
\label{thm:three-stage}
\[
  r_C \compose (r_B \compose (r_A \compose s))
  \;=\;
  ((r_C \compose r_B) \compose r_A) \compose s.
\]
\end{theorem}

\begin{proof}[Proof sketch]
Three applications of associativity.
See \texttt{three\_stage\_telephone}.
\end{proof}

\begin{remark}
Arbitrarily long chains of cultural transmission collapse into a single
composite interpreter.  This is why we can meaningfully speak of ``the
Western tradition'' as a single interpretive lens, despite it being the
product of thousands of individual acts of interpretation.
\end{remark}

% ── Summary ──────────────────────────────────────────────────────────
\section*{Chapter Summary}

This chapter established the foundational apparatus of the signalling-game
framework within IDT\@.  We introduced \emph{repertoires} as formal models
of culturally constituted minds, and \emph{interpretation} as
compositional fusion of horizons.  The \emph{Total Hermeneutic Bound}
(Theorem~\ref{thm:total-hermeneutic-bound}) shows that no signal can
inject more complexity into a mind than the signal's depth times the
repertoire's size.  \emph{Silence is Transparent}
(Theorem~\ref{thm:silence-transparent}) provides the formal foundation
for babbling equilibria.  The \emph{Durkheimian Consensus}
(Theorem~\ref{thm:durkheimian-consensus}) and \emph{Turner's Communitas}
(Theorem~\ref{thm:turners-communitas}) establish that resonant
backgrounds guarantee resonant interpretations, providing the formal
basis of cultural consensus.  The \emph{Telephone Game} theorems
(Theorems~\ref{thm:bartlett}--\ref{thm:resonant-transmission}) show that
sequential interpretation chains collapse into single composite
interpreters via associativity.

All sixteen theorems are verified in
\texttt{FormalAnthropology/IDT\_Signal\_Ch01.lean}.

% ══════════════════════════════════════════════════════════════════════
% Chapter 2: The Signalling Problem
% ══════════════════════════════════════════════════════════════════════

\chapter{The Signalling Problem}
\label{ch:signalling}

\epigraph{``Words are also actions, and actions are a kind of words.''}{Ralph Waldo Emerson}

% ── Introduction ──────────────────────────────────────────────────────

In every human society, actors choose signals strategically: a politician
picks rhetoric, a suitor offers gifts, a priest designs liturgy.  The
sender's fundamental problem is that they cannot control how receivers
will \emph{interpret} the signal---each receiver composes the signal with
their own repertoire.

This chapter formalises the core game-theoretic structure that governs
strategic communication.  Following Spence~\cite{spence1973} and
Crawford \& Sobel~\cite{crawford1982}, we model the signalling problem
as a game in which a sender chooses a signal, a receiver interprets it
through their repertoire, and both receive payoffs from the resulting
interpretations.  The crucial departure from classical game theory is
that payoffs depend not on the signal itself but on what receivers
\emph{make of it}---how it composes with their existing repertoire of
ideas.

\medskip
\noindent\textbf{Chapter overview.}  We define the \emph{signalling game}
(\S\ref{sec:ch2:game}), establish the \emph{babbling equilibrium}
(\S\ref{sec:ch2:babbling}), analyse \emph{audience structure}
(\S\ref{sec:ch2:audience}), prove \emph{depth bounds} on outcomes
(\S\ref{sec:ch2:depth}), study \emph{resonance in game outcomes}
(\S\ref{sec:ch2:resonance}), and explore \emph{strategic properties}
(\S\ref{sec:ch2:strategic}).  All eighteen theorems are machine-verified.

% ── §1. The Signalling Game ──────────────────────────────────────────
\section{The Signalling Game}
\label{sec:ch2:game}

\begin{definition}[Signalling Game]
\label{def:signalling-game}
A \emph{signalling game} $G = (\Rep, \Pay_S, \Pay_R)$ consists of:
\begin{itemize}
  \item A receiver repertoire $\Rep \in \IS^*$;
  \item A sender payoff function $\Pay_S : \IS^* \to \mathbb{N}$;
  \item A receiver payoff function $\Pay_R : \IS^* \to \mathbb{N}$.
\end{itemize}
The payoff depends on the \emph{interpretation} (the list of composed
ideas), not the signal directly.
\end{definition}

\begin{lstlisting}[caption={SignallingGame in Lean~4},label=lst:signalling-game]
structure SignallingGame (I : Type*) [IdeaticSpace I] where
  repertoire : Repertoire I
  sender_payoff : List I -> Nat
  receiver_payoff : List I -> Nat
\end{lstlisting}

\begin{definition}[Outcome]
\label{def:outcome}
The \emph{outcome} of sending signal~$s$ in game~$G$ is
\[
  \operatorname{outcome}(G, s) = \operatorname{interpret}(G.\Rep, s).
\]
\end{definition}

\begin{definition}[Optimal Signal]
\label{def:optimal-signal}
A signal~$s$ is \emph{optimal} for the sender if no alternative yields
strictly higher payoff:
\[
  \forall s'.\; \Pay_S(\operatorname{outcome}(G, s'))
  \leq \Pay_S(\operatorname{outcome}(G, s)).
\]
\end{definition}

\begin{definition}[Constant Sender Payoff]
\label{def:constant-payoff}
A game has \emph{constant sender payoff} if
$\Pay_S(\mathbf{x}) = \Pay_S(\mathbf{y})$ for all interpretation
lists~$\mathbf{x}, \mathbf{y}$.
\end{definition}

% ── §2. Babbling Equilibrium ─────────────────────────────────────────
\section{Babbling Equilibrium}
\label{sec:ch2:babbling}

\begin{theorem}[Babbling Outcome]
\label{thm:babbling-outcome}
Sending void produces the original repertoire unchanged:
\[
  \operatorname{outcome}(G, \void) = G.\Rep.
\]
\end{theorem}

\begin{proof}[Proof sketch]
Unfold the definitions and apply Theorem~\ref{thm:silence-transparent}
from Chapter~1.  See \texttt{babbling\_outcome}.
\end{proof}

\begin{remark}
The ``babbling equilibrium'' of Crawford \& Sobel~\cite{crawford1982}
exists in every signalling game because silence is transparent---the
receiver's mind is left exactly as it was.
\end{remark}

\begin{theorem}[Constant Payoff Babbling Optimality]
\label{thm:constant-payoff}
If sender payoff is constant, every signal is optimal---including void.
\end{theorem}

\begin{proof}[Proof sketch]
If $\Pay_S$ is constant, $\Pay_S(\operatorname{outcome}(G,s'))
= \Pay_S(\operatorname{outcome}(G,s))$ for all $s,s'$, so the
optimality condition holds trivially.  See \texttt{constant\_payoff\_optimal}.
\end{proof}

\begin{remark}
When the sender genuinely does not care what the receiver thinks,
information transmission collapses entirely.  This is the cheapest of
cheap talk.
\end{remark}

\begin{theorem}[Outcome Length Invariance]
\label{thm:outcome-length}
Every signal produces the same number of interpretations:
\[
  |\operatorname{outcome}(G, s)| = |G.\Rep|.
\]
\end{theorem}

\begin{proof}[Proof sketch]
By Theorem~\ref{thm:conserves-count}.  See \texttt{outcome\_length}.
\end{proof}

\begin{remark}
You cannot make someone think \emph{more} thoughts by choosing a
cleverer signal.  The number of schemas activated is a fixed property of
the receiver, not the sender.
\end{remark}

\begin{theorem}[Singleton Repertoire Outcome]
\label{thm:singleton-outcome}
A receiver with a single idea~$r$ produces $[r \compose s]$.
\end{theorem}

\begin{proof}[Proof sketch]
Direct computation.  See \texttt{singleton\_outcome}.
\end{proof}

\begin{remark}
The fanatic (single-idea mind) has a perfectly predictable interpretation
of any signal.  Totalitarian receivers are, paradoxically, the easiest
to communicate with.
\end{remark}

% ── §3. Audience Structure ───────────────────────────────────────────
\section{Audience Structure}
\label{sec:ch2:audience}

\begin{theorem}[Merged Audience]
\label{thm:merged-audience}
Signalling to a combined audience equals the concatenation of separate
outcomes:
\[
  \operatorname{outcome}(\Rep_1 \mathbin{+\!+} \Rep_2, s)
  = \operatorname{outcome}(\Rep_1, s) \mathbin{+\!+}
    \operatorname{outcome}(\Rep_2, s).
\]
\end{theorem}

\begin{proof}[Proof sketch]
By \texttt{List.map\_append}.  See \texttt{merged\_audience\_outcome}.
\end{proof}

\begin{remark}
A political speech to a diverse audience (workers $\cup$ elites) produces
interpretations that are the disjoint union of what each subgroup would
have produced separately.  There is no ``synergistic'' interpretation
effect from mere co-presence.
\end{remark}

\begin{theorem}[Empty Audience]
\label{thm:empty-audience}
Signalling to no one produces nothing:
$\operatorname{outcome}([], s) = []$.
\end{theorem}

\begin{proof}[Proof sketch]
By definition.  See \texttt{empty\_audience}.
\end{proof}

\begin{remark}
If nobody is there to interpret, there are no interpretations.  A signal
without a receiver is not a communicative act---formalising Wittgenstein's
point that meaning requires use~\cite{wittgenstein1953}.
\end{remark}

\begin{theorem}[Blank Audience Transparency]
\label{thm:blank-audience}
A receiver whose repertoire is all void interprets any signal as
$n$~copies of the signal itself:
\[
  \operatorname{outcome}(\operatorname{replicate}(n, \void), s)
  = \operatorname{replicate}(n, s).
\]
\end{theorem}

\begin{proof}[Proof sketch]
By induction on~$n$, using $\void \compose s = s$ at each step.
See \texttt{blank\_audience}.
\end{proof}

% ── §4. Depth Bounds on Outcomes ─────────────────────────────────────
\section{Depth Bounds on Outcomes}
\label{sec:ch2:depth}

\begin{theorem}[Outcome Depth Bound]
\label{thm:outcome-depth-bound}
\[
  \operatorname{depthSum}(\operatorname{outcome}(G, s))
  \leq \operatorname{depthSum}(G.\Rep) + |G.\Rep| \cdot \depth(s).
\]
\end{theorem}

\begin{proof}[Proof sketch]
Direct application of the Total Hermeneutic Bound
(Theorem~\ref{thm:total-hermeneutic-bound}).
See \texttt{outcome\_depth\_bound}.
\end{proof}

\begin{remark}
The sender's capacity to inject complexity into the receiver's mind is
linearly bounded.  This is why propaganda targets breadth (many receivers)
rather than depth (one deep signal).
\end{remark}

\begin{theorem}[Void Signal Zero Cost]
\label{thm:void-zero-cost}
The void signal adds zero total depth to the outcome:
\[
  \operatorname{depthSum}(\operatorname{outcome}(G, \void))
  = \operatorname{depthSum}(G.\Rep).
\]
\end{theorem}

\begin{proof}[Proof sketch]
By Theorem~\ref{thm:silence-depthsum}.  See \texttt{void\_signal\_zero\_cost}.
\end{proof}

% ── §5. Resonance in Game Outcomes ───────────────────────────────────
\section{Resonance in Game Outcomes}
\label{sec:ch2:resonance}

\begin{theorem}[Resonant Repertoire Consensus]
\label{thm:resonant-repertoire}
If all elements in the repertoire resonate with a fixed idea~$a$, then
all interpretations of any signal also resonate with $a \compose s$.
\end{theorem}

\begin{proof}[Proof sketch]
For each $r \in \Rep$ with $r \res a$, by Theorem~\ref{thm:durkheimian-consensus}
we get $r \compose s \res a \compose s$.
See \texttt{resonant\_repertoire\_consensus}.
\end{proof}

\begin{remark}
A culturally homogeneous audience (everyone resonates with the same
archetype) produces interpretations that all resonate with that
archetype's interpretation.  Cultural consensus is self-reinforcing.
\end{remark}

\begin{theorem}[Resonant Signals Produce Resonant Outcomes]
\label{thm:resonant-signals-outcomes}
If $s_1 \res s_2$, then for each position~$i$, the interpretations
$r_i \compose s_1$ and $r_i \compose s_2$ resonate.
\end{theorem}

\begin{proof}[Proof sketch]
By induction on the repertoire, applying resonance-compose-left at each
step.  See \texttt{resonant\_signals\_resonant\_outcomes}.
\end{proof}

% ── §6. Strategic Properties ─────────────────────────────────────────
\section{Strategic Properties}
\label{sec:ch2:strategic}

\begin{theorem}[Self-Signalling]
\label{thm:self-signalling}
If $\mathit{sender\_idea} \in \Rep$, then
$\mathit{sender\_idea} \compose s \in \operatorname{outcome}(G, s)$.
\end{theorem}

\begin{proof}[Proof sketch]
The sender's idea is mapped to its composition with~$s$ by the
interpretation function.  See \texttt{self\_signalling}.
\end{proof}

\begin{remark}
An insider---someone whose perspective is represented in the audience---
can always produce resonant interpretations by sending signals resonant
with their own position.
\end{remark}

\begin{theorem}[Composition Signal]
\label{thm:composition-signal}
\[
  \operatorname{outcome}(\Rep, s_1 \compose s_2)
  = \operatorname{outcome}(\operatorname{interpret}(\Rep, s_1), s_2).
\]
\end{theorem}

\begin{proof}[Proof sketch]
By associativity of composition: for each~$r$,
$r \compose (s_1 \compose s_2) = (r \compose s_1) \compose s_2$.
See \texttt{composition\_signal}.
\end{proof}

\begin{remark}
A complex ritual (compound signal) has the same effect as first
``priming'' the audience with~$s_1$ and then delivering~$s_2$.  Ritual
structure exploits associativity.
\end{remark}

\begin{theorem}[Void Prepend]
\label{thm:void-prepend}
$\operatorname{outcome}(\Rep, \void \compose s)
 = \operatorname{outcome}(\Rep, s)$.
\end{theorem}

\begin{proof}[Proof sketch]
By void-right: $r \compose (\void \compose s) = r \compose s$ after
applying associativity and the void-right identity.
See \texttt{void\_prepend\_signal}.
\end{proof}

\begin{theorem}[Sequential Signal]
\label{thm:sequential-signal}
\[
  \operatorname{interpret}(\operatorname{interpret}(\Rep, s_1), s_2)
  = \operatorname{interpret}(\Rep, s_1 \compose s_2).
\]
\end{theorem}

\begin{proof}[Proof sketch]
By functoriality of map and associativity.
See \texttt{sequential\_signal}.
\end{proof}

\begin{remark}
Multi-stage rituals (initiation sequences, liturgical calendars) produce
the same interpretive effect as a single complex ritual.  Ceremonial
structure is algebraically decomposable.
\end{remark}

\begin{theorem}[Void Captures Signal]
\label{thm:void-captures}
If $\void \in G.\Rep$, then $s \in \operatorname{outcome}(G, s)$.
\end{theorem}

\begin{proof}[Proof sketch]
By Theorem~\ref{thm:negative-capability}.
See \texttt{void\_captures\_signal}.
\end{proof}

\begin{remark}
An ``open-minded'' receiver (one who retains $\void$ = pure receptivity)
always captures the sender's intended signal.
\end{remark}

\begin{theorem}[MaxDepth Outcome Bound]
\label{thm:max-outcome-depth}
\[
  \operatorname{maxInterpDepth}(\operatorname{outcome}(G, s))
  \leq \operatorname{maxInterpDepth}(G.\Rep) + \depth(s).
\]
\end{theorem}

\begin{proof}[Proof sketch]
By Theorem~\ref{thm:max-interp-depth}.
See \texttt{max\_outcome\_depth}.
\end{proof}

\begin{remark}
No matter how profound a signal, the deepest interpretation anyone can
produce is bounded by their existing depth plus the signal's depth.  The
receiver's capacity limits understanding, not the sender's intent---a
formalisation of ``you can't teach beyond the student's readiness.''
\end{remark}

\begin{theorem}[Babbling Payoff]
\label{thm:babbling-payoff}
\[
  \Pay_S(\operatorname{outcome}(G, \void))
  = \Pay_S(G.\Rep).
\]
\end{theorem}

\begin{proof}[Proof sketch]
By Theorem~\ref{thm:babbling-outcome}.
See \texttt{babbling\_payoff}.
\end{proof}

\begin{remark}
The babbling payoff equals the ``status quo'' payoff.  A signal is worth
sending only if it yields strictly more than silence.
\end{remark}

% ── Summary ──────────────────────────────────────────────────────────
\section*{Chapter Summary}

This chapter formalised the signalling problem within IDT\@.  The central
game-theoretic structure---a signalling game with repertoire-dependent
payoffs---was introduced, and the \emph{babbling equilibrium}
(Theorem~\ref{thm:babbling-outcome}) was established: silence always
leaves the receiver's mind unchanged.  We showed that outcomes have
invariant length (Theorem~\ref{thm:outcome-length}), bounded depth
(Theorem~\ref{thm:outcome-depth-bound}), and that resonant signals
produce resonant outcomes (Theorem~\ref{thm:resonant-signals-outcomes}).
The composition signal theorem (Theorem~\ref{thm:composition-signal})
reveals that compound signals decompose via associativity, providing the
algebraic foundation for analysing multi-stage rituals.

All eighteen theorems are verified in
\texttt{FormalAnthropology/IDT\_Signal\_Ch02.lean}.

% ══════════════════════════════════════════════════════════════════════
% Chapter 3: Costly Depth and Ritual Sacrifice
% ══════════════════════════════════════════════════════════════════════

\chapter{Costly Depth and Ritual Sacrifice}
\label{ch:costly}

\epigraph{``A signal that is cheap to produce is a signal that is easy
to fake.''}{Amotz Zahavi, \textit{The Handicap Principle}}

% ── Introduction ──────────────────────────────────────────────────────

Why do organisms---and human societies in particular---invest enormous
resources in signals that appear, on the surface, to be wasteful?  The
peacock drags a metabolically expensive tail through the underbrush; the
stotting gazelle leaps conspicuously in the presence of a predator
rather than fleeing at full speed; the young Maasai warrior hunts a lion
with a spear when safer alternatives exist.  In each case, the very
costliness of the signal is what makes it \emph{honest}: a weak peacock
cannot afford the tail, a slow gazelle cannot afford the leap, and a
cowardly warrior cannot afford the hunt.\footnote{Zahavi's original
  formulation (1975) was widely resisted until Grafen (1990) provided a
  formal game-theoretic model showing that handicap equilibria could be
  evolutionarily stable.  See Grafen, A., ``Biological signals as
  handicaps,'' \textit{Journal of Theoretical Biology}, 144(4), 1990,
  pp.~517--546.}

Sosis and Alcorta~\cite{sosis2003} extended this logic from biology to
religion, arguing that costly rituals---painful initiations, dietary
restrictions, time-consuming liturgical practices---function as
\emph{credibility-enhancing displays} (CREDs) that solve the free-rider
problem in cooperative groups.  Religious communities face a persistent
threat from individuals who would enjoy the benefits of group membership
(mutual aid, shared childcare, collective defence) without bearing the
costs of commitment.  Costly rituals screen out free-riders by imposing
costs that only genuinely committed individuals are willing to pay.

Henrich's~(2009) cultural evolution framework~\cite{henrich2009}
situates costly signalling within a broader theory of cultural group
selection: groups that evolve effective costly-signalling mechanisms
outcompete groups that do not, because they achieve higher levels of
internal cooperation.\footnote{See also Boyd and Richerson (1985) on the
  coevolution of genes and culture, and Norenzayan (2013) on ``Big
  Gods'' as a cultural evolutionary response to the scaling problem of
  cooperation.}  Irons~(2001) independently developed a costly
signalling account of religious commitment, emphasising that religious
costs are ``hard to fake'' precisely because they involve genuine
sacrifice---time, resources, pain, or opportunity cost.  Bulbulia~(2004)
further argued that religious costs function as \emph{credibility
enhancers}: the greater the cost borne by a signaller, the more credible
the signal of underlying commitment.\footnote{Bulbulia, J., ``The
  cognitive and evolutionary psychology of religion,''
  \textit{Biology and Philosophy}, 19(5), 2004, pp.~655--686.}

The central contribution of Ideatic Depth Theory to this literature is
the identification of \textbf{depth with cost}.  Where Zahavi spoke of
metabolic cost, Spence of educational cost, and Sosis of ritual cost,
IDT provides a unified formal measure: the \emph{depth} of a signal in
an ideatic space.  This identification is not merely notational
convenience.  It allows us to derive, as formal theorems, the key
properties that costly signalling theorists have posited as assumptions:
subadditivity of combined signals, linear scaling of iterated rituals,
the single-crossing property that separates types, and the bounded
impact of signals on audiences of finite sophistication.

Zahavi's (1975) handicap principle~\cite{zahavi1975} and Sosis's (2003)
costly signalling theory of ritual~\cite{sosis2003} propose that honest
communication requires \emph{costly} signals.  The depth of an idea---its
structural complexity---is a natural cost metric: deeper signals require
more cognitive resources to produce and maintain.

In classical signalling theory (Spence~\cite{spence1973}), the cost of
education serves as a separating device between high-ability and
low-ability workers.  In our framework, the cost of a signal is its
\emph{depth}: the structural complexity that must be generated and
maintained.  A simple slogan has low cost; a multi-layered ritual
performance has high cost.

The key anthropological insight is that rituals involving genuine
sacrifice---scarification, potlatch, fasting---are \emph{honest depth
signals}: they cannot be cheaply faked because the depth (complexity) is
intrinsic to the act itself.

\medskip
\noindent\textbf{Chapter overview.}  We identify depth with signal cost
(\S\ref{sec:ch3:cost}), study iterated signals and ritual repetition
(\S\ref{sec:ch3:iteration}), formalise the single-crossing property
(\S\ref{sec:ch3:crossing}), bound interpretation cost
(\S\ref{sec:ch3:interp-cost}), and analyse the interaction between
resonance and cost (\S\ref{sec:ch3:resonance-cost}).

\subsection*{Historical Context: From Veblen to Zahavi}

The intuition that costliness signals quality has deep roots in the
social sciences, well predating Zahavi's biological formalisation.
Thorstein Veblen's \textit{Theory of the Leisure Class}~(1899)
introduced the concept of ``conspicuous consumption''---the display of
wealth through deliberately wasteful expenditure---as a mechanism for
signalling social status.\footnote{Veblen, T., \textit{The Theory of
  the Leisure Class: An Economic Study of Institutions}, Macmillan,
  1899.  Veblen's analysis of ``conspicuous waste'' anticipates the
  costly signalling framework by nearly a century.}  The industrial
magnate who commissions a mansion far beyond his residential needs, the
aristocrat who maintains a retinue of idle servants---these are,
in IDT terms, high-depth signals whose cost is the point.

Marcel Mauss's \textit{The Gift}~(1925) extended this logic to
non-market societies.  The potlatch ceremonies of the Pacific Northwest
involved competitive gift-giving and, in extreme cases, competitive
destruction of wealth: canoes broken, blankets burned, coppers thrown
into the sea.  Mauss recognised that these acts of apparent waste were
in fact \emph{investments in social obligation}: the depth of the
sacrifice created debts that bound communities
together.\footnote{Mauss, M., \textit{Essai sur le don}, L'Ann\'{e}e
  Sociologique, 1925.  English translation: \textit{The Gift}, trans.\
  W.~D.~Halls, Routledge, 1990.}

What Veblen and Mauss lacked was a formal framework capable of deriving
the \emph{consequences} of costly signalling from its
\emph{structure}.  Zahavi~(1975) provided the biological version of
such a framework; Spence~(1973) provided the economic version.  IDT
now provides a framework general enough to encompass both, grounding
cost in the combinatorial structure of ideatic depth.  The theorems that
follow make precise the properties that Veblen, Mauss, Zahavi, and
Sosis described informally: why cost is subadditive, why empty signals
carry no cost, and why high-quality types bear cost more
efficiently than low-quality types.

% ── §1. Depth as Signal Cost ─────────────────────────────────────────
\section{Depth as Signal Cost}
\label{sec:ch3:cost}

\begin{definition}[Signal Cost]
\label{def:signal-cost}
The \emph{cost} of signal~$s$ is its depth:
$\operatorname{signalCost}(s) = \depth(s)$.
In Zahavi's handicap principle, this is the metabolic, social, or
cognitive cost of producing the signal.
\end{definition}

\begin{theorem}[Void is Costless]
\label{thm:void-costless}
$\operatorname{signalCost}(\void) = 0$.
\end{theorem}

\begin{proof}[Proof sketch]
By $\depth(\void) = 0$.  See \texttt{void\_costless}.
\end{proof}

\begin{remark}
Silence is always free.  Every signalling system must contend with the
zero-cost outside option of saying nothing.
\end{remark}

\begin{theorem}[Cost Subadditivity]
\label{thm:cost-subadditive}
\[
  \operatorname{signalCost}(s_1 \compose s_2)
  \leq \operatorname{signalCost}(s_1) + \operatorname{signalCost}(s_2).
\]
\end{theorem}

\begin{proof}[Proof sketch]
By the depth-compose bound.  See \texttt{cost\_subadditive}.
\end{proof}

\begin{remark}
A two-part ritual (fasting + scarification) costs at most as much as the
parts separately.  Compound signals can be \emph{cheaper} than their
parts (subadditivity), but never more expensive.  This formalises the
``economies of ritual combination'' observed in initiation ceremonies
that bundle multiple ordeals.
\end{remark}

\begin{theorem}[Void Right-Compose Cost]
\label{thm:void-right-compose-cost}
$\operatorname{signalCost}(s \compose \void) = \operatorname{signalCost}(s)$.
\end{theorem}

\begin{proof}[Proof sketch]
By $\depth(s \compose \void) = \depth(s)$.
See \texttt{void\_compose\_cost}.
\end{proof}

\begin{theorem}[Void Left-Compose Cost]
\label{thm:void-left-compose-cost}
$\operatorname{signalCost}(\void \compose s) = \operatorname{signalCost}(s)$.
\end{theorem}

\begin{proof}[Proof sketch]
By $\depth(\void \compose s) = \depth(s)$.
See \texttt{void\_left\_compose\_cost}.
\end{proof}

\begin{remark}
You cannot fake costly signals by ``adding silence.''  Prefixing or
appending a dramatic pause to a sacrifice does not change its cost.
\end{remark}

\begin{example}[The Potlatch of the Pacific Northwest]
Among the Kwakwaka'wakw (formerly called Kwakiutl) peoples of the
Pacific Northwest coast, the potlatch ceremony served as a paradigmatic
costly signal.  Chiefs competed for social rank by giving away---and in
extreme cases, deliberately destroying---enormous quantities of wealth:
blankets, canoes, eulachon oil, and prized copper shields known as
\textit{t'lakwa}.  The depth of the signal was measured not in symbolic
gestures but in real material sacrifice.  A chief who destroyed a copper
worth thousands of blankets was producing a signal of such depth that no
rival of lesser means could replicate it.  In IDT terms,
$\operatorname{signalCost}(\text{potlatch}) = \depth(\text{potlatch})$,
and this depth was precisely what made the signal unfakeable: one cannot
pretend to destroy wealth one does not possess.\footnote{Boas, F.,
  \textit{Kwakiutl Ethnography}, ed.\ H.~Codere, University of Chicago
  Press, 1966.  See also Codere, H., ``Fighting with property: a study
  of Kwakiutl potlatching and warfare,'' American Ethnological Society
  Monograph~18, 1950.}
\end{example}

\begin{example}[Scarification Among the Nuer]
Evans-Pritchard's ethnography of the Nuer of southern Sudan documented a
striking initiation rite: adolescent boys received six deep horizontal
cuts across the forehead, from ear to ear, using a small knife.  The
cuts were made without anaesthetic, and the resulting scars---called
\textit{gaar}---were permanent, visible markers of adult male
status.\footnote{Evans-Pritchard, E.~E., \textit{The Nuer: A
  Description of the Modes of Livelihood and Political Institutions of
  a Nilotic People}, Clarendon Press, 1940.}  The signal satisfies the
formal properties of costly depth: the pain and risk of infection
constitute genuine cost that cannot be faked (no one can acquire
\textit{gaar} scars without enduring the cutting), and the permanence of
the signal means that its depth is \emph{preserved over time}---it
cannot be retracted or diminished.  By Theorem~\ref{thm:void-costless},
void signals (mere words or gestures) have zero cost; the Nuer
initiation achieves honesty precisely by requiring non-zero depth.
\end{example}

\subsection*{Discussion: Why Depth Cannot Be Faked}

The examples above illustrate a crucial structural property of depth as
a cost metric: depth is \emph{intrinsic} to the signal, not extrinsic.
A potlatch chief cannot fake the destruction of a copper because the
copper must actually be broken; a Nuer initiate cannot fake the
scarification because the scars must actually be cut.  This
intrinsicality is what distinguishes depth-as-cost from purely
conventional costs (such as monetary fees, which can be borrowed,
embezzled, or counterfeited).  In the formal framework, this property
is captured by the fact that $\depth(s)$ is determined by the
compositional structure of~$s$ itself, not by any external accounting
mechanism.  A signal's cost is its complexity, and complexity cannot be
counterfeited.

\begin{theorem}[Trivial Composition is Costless]
\label{thm:trivial-compose-costless}
If $\operatorname{signalCost}(s_1) = 0$ and
$\operatorname{signalCost}(s_2) = 0$, then
$\operatorname{signalCost}(s_1 \compose s_2) = 0$.
\end{theorem}

\begin{proof}[Proof sketch]
By the depth-zero-closed property of ideatic spaces.
See \texttt{trivial\_compose\_costless}.
\end{proof}

\begin{remark}
Banality composed with banality remains banal.  You cannot bootstrap
costly signalling from costless components---depth must come from
somewhere.
\end{remark}

% ── §2. Iterated Signals ─────────────────────────────────────────────
\section{Iterated Signals and Ritual Repetition}
\label{sec:ch3:iteration}

\begin{theorem}[Iteration Cost Bound]
\label{thm:iteration-cost}
For $n$-fold repetition of signal~$s$:
\[
  \operatorname{signalCost}(\operatorname{composeN}(s, n))
  \leq n \cdot \operatorname{signalCost}(s).
\]
\end{theorem}

\begin{proof}[Proof sketch]
By induction on~$n$, using subadditivity at each step.
See \texttt{iteration\_cost\_bound}.
\end{proof}

\begin{remark}
Ritual repetition (daily prayer, weekly sabbath, annual pilgrimage) has
linearly bounded cumulative cost.  A devotee who prays 5 times daily
pays at most $5 \times$ the cost of a single prayer.  This linear
scaling is why religious systems can demand iterative commitment without
unbounded cost.
\end{remark}

\begin{theorem}[Zero-Cost Iteration]
\label{thm:zero-cost-iteration}
$\operatorname{signalCost}(\operatorname{composeN}(\void, n)) = 0$
for all~$n$.
\end{theorem}

\begin{proof}[Proof sketch]
$\operatorname{composeN}(\void, n) = \void$ by the void-composeN
identity, so $\depth = 0$.  See \texttt{zero\_cost\_iteration}.
\end{proof}

\begin{remark}
Empty ritual (void composed with void, repeated endlessly) never
generates genuine cost.  ``Going through the motions'' with no internal
depth is formally costless.
\end{remark}

\begin{remark}[Connection to Religious Studies: The Five Pillars of Islam]
The five daily prayers (\textit{salat}) of Islam exemplify iterated
costly signalling in a religious context.  Each individual prayer has
bounded cost: the time and cognitive effort required for the ablution
(\textit{wudu}), the physical prostrations (\textit{ruku} and
\textit{sujud}), and the recitation of Qur'anic passages.  By
Theorem~\ref{thm:iteration-cost}, the total annual cost of
\textit{salat} is linearly bounded by the number of iterations times
the cost of a single prayer.  What makes this system particularly
elegant from a signalling perspective is that the per-iteration cost is
calibrated to be bearable---Islam explicitly forbids making worship
unduly burdensome (\textit{la yukallifu Allahu nafsan illa wus'aha},
Qur'an 2:286)---while the cumulative commitment over a lifetime is
substantial enough to separate genuine believers from
opportunists.\footnote{Cf.\ Sosis, R., ``Why aren't we all
  Hutterites? Costly signaling theory and religious behavior,''
  \textit{Human Nature}, 14(2), 2003, pp.~91--127.}
\end{remark}

\begin{example}[Daily Prayer as Iterated Signal]
A Muslim who performs \textit{salat} five times daily for fifty years
executes approximately $5 \times 365 \times 50 = 91{,}250$ iterations
of the prayer signal.  By Theorem~\ref{thm:iteration-cost}, the total
signal cost over a lifetime is bounded above by
$91{,}250 \times \operatorname{signalCost}(\text{single prayer})$.
This linear bound has an important implication: it means that the
cumulative cost of religious observance grows \emph{predictably}, not
explosively.  A community can calibrate its ritual demands to be
sustainable across generations, which is exactly what we observe in
the major world religions.  By contrast, if signal cost were
superadditive (if each repetition made the next one more costly), no
religious system could sustain iterative practice over decades.
\end{example}

\begin{theorem}[Single Repetition Cost]
\label{thm:single-rep-cost}
$\operatorname{signalCost}(\operatorname{composeN}(s, 1))
 \leq \operatorname{signalCost}(s)$.
\end{theorem}

\begin{proof}[Proof sketch]
Specialise Theorem~\ref{thm:iteration-cost} with $n = 1$.
See \texttt{single\_rep\_cost}.
\end{proof}

% ── §3. The Single-Crossing Property ─────────────────────────────────
\section{The Single-Crossing Property}
\label{sec:ch3:crossing}

\begin{definition}[Costly Signal Model]
\label{def:costly-signal-model}
A \emph{costly signal model} consists of two cost functions:
\begin{itemize}
  \item $c_H : \IS \to \mathbb{N}$ (high-quality type), with
    $c_H(s) \leq \depth(s)$;
  \item $c_L : \IS \to \mathbb{N}$ (low-quality type), with
    $\depth(s) \leq c_L(s)$.
\end{itemize}
The high type can produce depth efficiently; the low type cannot.
\end{definition}

\begin{lstlisting}[caption={CostlySignalModel in Lean~4},label=lst:costly-signal]
structure CostlySignalModel (I : Type*) [IdeaticSpace I] where
  highTypeCost : I -> Nat
  lowTypeCost : I -> Nat
  high_efficient : forall s, highTypeCost s <= depth s
  low_costly : forall s, depth s <= lowTypeCost s
\end{lstlisting}

\begin{theorem}[High Type Cost Advantage]
\label{thm:high-type-advantage}
For any costly signal model and any signal~$s$:
$c_H(s) \leq c_L(s)$.
\end{theorem}

\begin{proof}[Proof sketch]
Chain the two constraints: $c_H(s) \leq \depth(s) \leq c_L(s)$.
See \texttt{high\_type\_advantage}.
\end{proof}

\begin{remark}
In Sosis's theory, genuine believers find costly rituals less onerous
than hypocrites.  The devout faster suffers less than the faker because
their internal ``type'' (depth of belief) makes the signal cheaper to
produce.  This is the formal foundation of ritual honesty.
\end{remark}

\begin{theorem}[Void is Free for All Types]
\label{thm:void-free-all}
$c_H(\void) \leq 0$ (and hence $c_H(\void) = 0$).
\end{theorem}

\begin{proof}[Proof sketch]
$c_H(\void) \leq \depth(\void) = 0$.
See \texttt{void\_free\_all\_types}.
\end{proof}

\begin{remark}
The babbling option is costless for all types.  Separation in costly
signalling must come from \emph{non-void} signals---only signals with
genuine depth can separate types.
\end{remark}

\subsection*{Discussion: The Ecology of Costly Signalling}

Why do different societies calibrate signal costs so differently?  The
Yanomam\"{o} demand chest-pounding duels; the Amish demand plain dress
and technological renunciation; the Benedictines demand a lifetime of
prayer, labour, and obedience.  From the perspective of IDT, these
are all costly signals, but their \emph{depth profiles} differ
dramatically.  The key insight from Sosis's comparative work is that
the optimal signal cost depends on the cooperative ecology of the
group: groups facing greater free-rider threats should evolve costlier
signals, and groups whose cooperation yields greater per-capita
benefits can afford to impose higher costs without losing members.

This prediction has been empirically tested.  Sosis and
Bressler~(2003) analysed the longevity of 277 communes founded in the
nineteenth-century United States, comparing religious communes (which
imposed costly ritual requirements) with secular communes (which
typically did not).\footnote{Sosis, R.\ and Bressler, E.~R.,
  ``Cooperation and commune longevity: a test of the costly signaling
  theory of religion,'' \textit{Cross-Cultural Research}, 37(2), 2003,
  pp.~211--239.}  They found that religious communes survived
significantly longer than secular ones, and that within the religious
category, communes imposing more costly requirements (dietary
restrictions, distinctive dress, behavioural taboos) survived longer
still.  Each additional costly requirement increased the predicted
lifespan of the commune.

\begin{example}[Kibbutzim and Religious Communes]
The contrast between religious and secular kibbutzim in
twentieth-century Israel provides a natural experiment in costly
signalling.  Religious kibbutzim required members to observe Shabbat,
keep kosher, and participate in daily prayer---all costly signals in
the IDT sense, with non-trivial depth.  Secular kibbutzim, by
contrast, relied primarily on ideological commitment, a signal whose
depth is harder to verify externally.  Consistent with costly
signalling theory, religious kibbutzim exhibited lower defection rates
and greater economic cooperation.\footnote{Sosis, R.\ and Ruffle,
  B.~J., ``Religious ritual and cooperation: testing for a
  relationship on Israeli religious and secular kibbutzim,''
  \textit{Current Anthropology}, 44(5), 2003, pp.~713--722.}
By Theorem~\ref{thm:high-type-advantage}, genuinely committed members
of religious kibbutzim bore these costs more efficiently than
uncommitted members, producing exactly the type-separation that
sustains long-term cooperation.
\end{example}

% ── §4. Interpretation Cost ──────────────────────────────────────────
\section{Interpretation Cost and Audience Effects}
\label{sec:ch3:interp-cost}

\begin{theorem}[Total Interpretation Cost Bound]
\label{thm:total-interp-cost}
\[
  \operatorname{depthSum}(\operatorname{interpret}(\Rep, s))
  \leq \operatorname{depthSum}(\Rep) + |\Rep| \cdot \operatorname{signalCost}(s).
\]
\end{theorem}

\begin{proof}[Proof sketch]
By the Total Hermeneutic Bound (Theorem~\ref{thm:total-hermeneutic-bound}),
noting that $\operatorname{signalCost}(s) = \depth(s)$.
See \texttt{total\_interpretation\_cost}.
\end{proof}

\begin{remark}
The ``impact'' of a costly signal on an audience is bounded by the
signal's cost times the audience size plus the audience's existing
complexity.  A potlatch (costly gift-giving) impresses proportionally to
its cost $\times$ audience.
\end{remark}

\begin{theorem}[Costless Signal Preserves Total Depth]
\label{thm:costless-preserves}
If $\operatorname{signalCost}(s) = 0$, then
$\operatorname{depthSum}(\operatorname{interpret}(\Rep, s))
 \leq \operatorname{depthSum}(\Rep)$.
\end{theorem}

\begin{proof}[Proof sketch]
Set $\depth(s) = 0$ in the hermeneutic bound to eliminate the linear
term.  See \texttt{costless\_preserves\_depth}.
\end{proof}

\begin{theorem}[Atomic Signal Bounded Cost]
\label{thm:atomic-cost}
If $s$ is atomic, then $\operatorname{signalCost}(s) \leq 1$.
\end{theorem}

\begin{proof}[Proof sketch]
By $\depth(s) \leq 1$ for atomic ideas.
See \texttt{atomic\_signal\_cost}.
\end{proof}

\begin{theorem}[Compound Atomic Cost]
\label{thm:compound-atomic-cost}
If $s_1, s_2$ are both atomic, then
$\operatorname{signalCost}(s_1 \compose s_2) \leq 2$.
\end{theorem}

\begin{proof}[Proof sketch]
By $\depth(s_1 \compose s_2) \leq \depth(s_1) + \depth(s_2) \leq 2$.
See \texttt{compound\_atomic\_cost}.
\end{proof}

\begin{remark}
The simplest possible compound ritual (two atomic acts) has strictly
bounded cost---complex enough to potentially separate types but simple
enough to be tractable.
\end{remark}

\begin{remark}[Connection to Evolutionary Biology]
The compound atomic cost theorem
(Theorem~\ref{thm:compound-atomic-cost}) has a direct biological
analogue in multi-component mating displays.  The peacock's courtship
involves both a visual component (the iridescent tail fan, with its
elaborate eyespot pattern) and an auditory component (a distinctive
rattling call produced by vibrating the tail feathers).  Each component
is approximately atomic in the IDT sense: a single, irreducible costly
act.  The total display cost is bounded by the sum of the component
costs, exactly as the theorem predicts.  Crucially, the two components
are not redundant---they are received through different sensory channels
and serve complementary verification functions.  The visual component
signals genetic quality (symmetry, pigmentation); the auditory component
signals current condition (vigour, muscular control).  The compound
signal is more informative than either component alone, while remaining
within the tractable cost bound of~$2$.\footnote{Loyau, A.\ et al.,
  ``Iridescent structurally based coloration of eyespots correlates
  with mating success in the peacock,'' \textit{Behavioral Ecology},
  18(6), 2007, pp.~1123--1131.}
\end{remark}

\begin{example}[Marriage Ceremonies as Compound Signals]
A traditional Hindu wedding (\textit{vivaha}) provides an instructive
example of compound costly signalling.  The ceremony involves multiple
individually costly acts: the exchange of garlands (\textit{jaimala}),
signalling mutual acceptance; the sacred fire ritual (\textit{agni
parikrama}), in which the couple circles the fire seven times while
reciting vows; the application of vermilion (\textit{sindoor}) to the
bride's hair parting; and the tying of the sacred thread
(\textit{mangalsutra}).  Each act is approximately atomic---a single
ritual gesture with bounded depth---and the total ceremony has bounded
compound cost by Theorem~\ref{thm:compound-atomic-cost} applied
iteratively.  The multi-component structure serves a dual purpose:
it increases the total signal cost (making the commitment harder to
fake) while distributing the cost across distinct symbolic registers
(visual, verbal, kinesthetic, material), thereby making the signal
legible to a diverse audience of witnesses with different interpretive
repertoires.\footnote{Pandey, R., \textit{Hindu Sa\d{m}sk\={a}ras:
  Socio-Religious Study of the Hindu Sacraments}, Motilal Banarsidass,
  1969.}
\end{example}

% ── §5. Resonance and Cost ───────────────────────────────────────────
\section{Resonance and Cost Interaction}
\label{sec:ch3:resonance-cost}

\begin{theorem}[Resonant Signals Have Comparable Outcomes]
\label{thm:resonant-comparable}
If $s_1 \res s_2$, then for each $r \in \Rep$,
$r \compose s_1 \res r \compose s_2$.
\end{theorem}

\begin{proof}[Proof sketch]
By resonance-compose-left applied to each repertoire element.
See \texttt{resonant\_signals\_comparable}.
\end{proof}

\begin{remark}
Two rituals that ``resonate'' produce similar effects on any audience,
even if one costs more than the other.  This is why cheap imitations of
costly rituals can sometimes ``work''---if they resonate with the
original, the audience's interpretations resonate.
\end{remark}

\begin{theorem}[Self-Resonant Interpretation]
\label{thm:self-resonant-interp}
For any repertoire~$\Rep$ and signal~$s$,
$\operatorname{interpret}(\Rep, s)$ is pairwise self-resonant.
\end{theorem}

\begin{proof}[Proof sketch]
Apply Theorem~\ref{thm:resonant-comparable} with $s_1 = s_2 = s$,
using reflexivity of resonance.
See \texttt{self\_resonant\_interpretation}.
\end{proof}

\begin{theorem}[Bounded Repertoire, Bounded Outcome]
\label{thm:bounded-repertoire}
If $\depth(r) \leq d$ for all $r \in \Rep$, then each interpretation
$x \in \operatorname{interpret}(\Rep, s)$ satisfies
$\depth(x) \leq d + \operatorname{signalCost}(s)$.
\end{theorem}

\begin{proof}[Proof sketch]
For each $r \in \Rep$,
$\depth(r \compose s) \leq \depth(r) + \depth(s) \leq d + \depth(s)$.
See \texttt{bounded\_repertoire\_bounded\_outcome}.
\end{proof}

\begin{remark}
Receivers with bounded sophistication cannot produce arbitrarily deep
interpretations.  Even the most costly signal only adds its own cost to
each receiver's maximum capacity.
\end{remark}

\begin{theorem}[Interpretation Filtration]
\label{thm:interp-filtration}
If $\Rep \subseteq \mathcal{F}_n$ (the depth-$n$ filtration), then
$\operatorname{interpret}(\Rep, s) \subseteq
 \mathcal{F}_{n + \operatorname{signalCost}(s)}$.
\end{theorem}

\begin{proof}[Proof sketch]
Each $r \in \Rep$ has $\depth(r) \leq n$, so
$\depth(r \compose s) \leq n + \depth(s)$.
See \texttt{interpretation\_filtration}.
\end{proof}

\begin{remark}
Costly signals push interpretations up the depth hierarchy by at most
their cost.  A depth-3 ritual received by depth-5 minds produces at most
depth-8 interpretations.  This is the formal ``budget'' of cultural
transformation through signalling.
\end{remark}

\begin{remark}[Open Questions in Costly Signalling]
Several questions remain open in the formal theory of costly depth
signals, and we flag them here as directions for future work.
First, can the \emph{single-crossing property}
(Theorem~\ref{thm:high-type-advantage}) be weakened?  Our current
formulation assumes a binary type space (high vs.\ low), but real
signalling environments feature a continuum of types.  Extending the
single-crossing result to a richer type space would require a graded
notion of cost efficiency, perhaps indexed by a continuous depth
parameter.  Second, what happens when signal cost is
\emph{stochastic}?  In practice, the cost of a ritual performance
varies across occasions: a fast during summer is more costly than a
fast during winter; a pilgrimage in wartime is more costly than in
peacetime.  A stochastic extension of IDT would model
$\operatorname{signalCost}(s)$ as a random variable and derive
probabilistic analogues of the theorems in this chapter.  Third, the
interaction between \emph{signal cost and audience size} deserves
further exploration: Theorem~\ref{thm:total-interp-cost} gives a
linear bound, but empirical evidence suggests that audience effects may
exhibit diminishing returns (a potlatch before 1000 witnesses is not
ten times more effective than one before 100).
\end{remark}

\begin{example}[The Olympic Games as Multi-Level Costly Signal]
The modern Olympic Games function as a costly signal at multiple
levels simultaneously, illustrating the compositional structure that
IDT formalises.  At the \emph{national} level, the host country bears
enormous costs: infrastructure construction, security, and
opportunity cost of diverted public funds.  These costs signal
national prosperity and organisational capacity.  At the
\emph{individual} level, each athlete bears years of training cost:
the depth of their athletic signal is measured in thousands of hours
of deliberate practice, dietary discipline, and foregone alternative
careers.  At the \emph{spectator} level, attendees bear costs of
travel, tickets, and time---signals of allegiance to their national
team.  By the subadditivity theorem
(Theorem~\ref{thm:cost-subadditive}), the total cost of the Olympic
signal is bounded by the sum of these component costs, and by the
interpretation filtration theorem
(Theorem~\ref{thm:interp-filtration}), the cultural impact of the
Games on any audience is bounded by the audience's existing
sophistication plus the signal's depth.\footnote{For a sociological
  analysis of the Olympics as costly signalling, see Roche, M.,
  \textit{Mega-Events and Modernity: Olympics and Expos in the Growth
  of Global Culture}, Routledge, 2000.}
\end{example}

% ── Summary ──────────────────────────────────────────────────────────
\section*{Chapter Summary}

This chapter connected IDT depth to signal cost theory.  The
identification $\operatorname{signalCost}(s) = \depth(s)$ allowed us to
prove that void is costless (Theorem~\ref{thm:void-costless}), cost is
subadditive (Theorem~\ref{thm:cost-subadditive}), and ritual repetition
has linearly bounded cumulative cost
(Theorem~\ref{thm:iteration-cost}).  The \emph{single-crossing property}
(Theorem~\ref{thm:high-type-advantage}) establishes that high-quality
types pay less for depth, providing the formal foundation of Zahavi's
handicap principle and Sosis's costly signalling theory of ritual.
Resonant signals produce comparable outcomes regardless of cost
difference (Theorem~\ref{thm:resonant-comparable}), and the
\emph{interpretation filtration} theorem
(Theorem~\ref{thm:interp-filtration}) shows that costly signals push
interpretations up the depth hierarchy by exactly their cost.

All eighteen theorems are verified in
\texttt{FormalAnthropology/IDT\_Signal\_Ch03.lean}.

% ══════════════════════════════════════════════════════════════════════
% Chapter 4: Resonance Activation and Cultural Priming
% ══════════════════════════════════════════════════════════════════════

\chapter{Resonance Activation and Cultural Priming}
\label{ch:activation}

\epigraph{``The question of relevance is the question of what is
activated.''}{Dan Sperber, \textit{Explaining Culture}}

% ── Introduction ──────────────────────────────────────────────────────

The concept of activation---which ideas a signal ``turns on'' in a
culturally constituted mind---lies at the intersection of several
major intellectual traditions.  Sperber and Wilson's \emph{relevance
theory}~\cite{sperber1995} articulates two foundational principles.
The \emph{cognitive principle of relevance} holds that human cognition
is geared toward the maximisation of relevance: cognitive systems
automatically allocate attention to inputs that promise the greatest
cognitive effect for the least processing effort.  The
\emph{communicative principle of relevance} adds that every ostensive
stimulus carries a presumption of its own optimal relevance---that the
speaker has made it worth the hearer's effort to
process.\footnote{Sperber and Wilson distinguish ``relevance'' from
mere ``informativeness.''  A stimulus is relevant not because it
carries a large quantity of information, but because it connects
efficiently with the hearer's existing cognitive environment---what we
formalise here as the repertoire.}  IDT's resonance relation gives
these principles a precise algebraic form: $r \res s$ holds exactly
when the signal~$s$ is relevant to the schema~$r$.

Cognitive linguistics provides a complementary perspective.  Lakoff's
\emph{Women, Fire, and Dangerous Things} (1987)~\cite{lakoff1987}
demonstrated that human categorisation is not a matter of classical
necessary-and-sufficient conditions but of radial categories organised
around prototypes and linked by family resemblance.  The schemas that
a signal activates are not an arbitrary set; they form a structured
neighbourhood in conceptual space, held together by what Lakoff calls
\emph{idealized cognitive models} (ICMs).  Fillmore's \emph{frame
semantics}~\cite{fillmore1982} makes a parallel claim at the level of
linguistic meaning: understanding a word requires activating an entire
frame---a structured background of knowledge against which the word
makes sense.\footnote{Fillmore's classic example: the word
``\emph{bachelor}'' activates a frame that includes the institution of
marriage, typical marital age, and social expectations.  Remove the
frame, and the word loses its meaning.  In IDT terms, the frame is
part of the repertoire, and the word resonates with it.}  The
activated subset of a repertoire (Definition~\ref{def:activated-subset})
is, in Fillmorean terms, the set of frames that a signal evokes.

From experimental psychology, Kahneman's \emph{Thinking, Fast and
Slow} (2011)~\cite{kahneman2011} documents the pervasive role of
\emph{priming}---the phenomenon whereby exposure to one stimulus
influences the response to a subsequent stimulus.  Subjects primed
with words associated with old age walk more slowly; subjects primed
with money-related images become more individualistic.  Kahneman
locates priming in ``System~1,'' the fast, automatic, associative mode
of thought.  IDT's activation framework provides an algebraic account
of exactly these effects: the primer~$s_1$ reshapes the repertoire
(via $\operatorname{interpret}(\Rep, s_1)$), and the subsequent
signal~$s_2$ then activates schemas within this reshaped
landscape.\footnote{The sequential interpretation associativity
theorem (Theorem~\ref{thm:sequential-activation}) shows that this
two-step process is equivalent to a single exposure to the compound
signal $s_1 \compose s_2$.  Priming, in algebraic terms, is
composition.}

What IDT contributes to these diverse literatures is a unified formal
language.  Sperber's relevance, Lakoff's ICMs, Fillmore's frames, and
Kahneman's priming are all, at bottom, claims about \emph{which
elements of a structured repertoire a signal activates}.  By
axiomatising resonance and deriving activation as a consequence, IDT
reveals the common algebraic skeleton beneath these superficially
different theories.

When a signal arrives at a mind, it does not activate the entire
repertoire uniformly.  It selectively \emph{activates} those ideas with
which it resonates---Sperber's (1996) relevance theory~\cite{sperber1996}
formalised.  A Catholic seeing a cross activates christological schemas.
A Hindu seeing the same cross might activate geometric symmetry schemas.
The signal is identical; the activation pattern differs because different
repertoire elements resonate with it.

This chapter builds the formal theory of activation: signals do not just
produce interpretations---they \emph{select which interpretations
matter} through resonance.  We connect this to Bourdieu's
\emph{habitus}~\cite{bourdieu1984}: the culturally acquired dispositions
that determine what is ``relevant'' are precisely the resonance structure
of the repertoire.

\medskip
\noindent\textbf{Chapter overview.}  We define activation
(\S\ref{sec:ch4:activation}), study its interaction with interpretation
(\S\ref{sec:ch4:interaction}), analyse sequential (double) activation
(\S\ref{sec:ch4:double}), establish depth bounds
(\S\ref{sec:ch4:depth-bounds}), and trace resonance chains
(\S\ref{sec:ch4:chains}).

\subsection*{Historical Context: From Associationism to Relevance Theory}

The idea that mental contents activate one another has a long
pedigree.  British associationism, from David Hume's \emph{Treatise of
Human Nature} (1739) through David Hartley's \emph{Observations on
Man} (1749), proposed that ideas are linked by resemblance, contiguity,
and cause-and-effect, and that the activation of one idea
automatically triggers associated ideas.  This philosophical tradition
was given cognitive-scientific form by Collins and
Loftus~\cite{collins1975}, whose \emph{spreading activation} model
posited that concepts are stored in a semantic network and that
activating one node sends activation spreading along weighted links to
neighbouring nodes.  The model explained priming effects, category
verification times, and typicality gradients.

Anthropologists have long employed activation concepts, if not always
under that name.  Malinowski's ``context of
situation''~\cite{malinowski1923} was an early recognition that the
meaning of an utterance depends on which aspects of the surrounding
context are cognitively active for the participants.  Bateson's theory
of ``frames''~\cite{bateson1972} proposed that metacommunicative
signals (``this is play'') activate interpretive frames that transform
the meaning of all subsequent signals.  Goffman~\cite{goffman1974}
extended Bateson's frame analysis into a comprehensive sociology of
everyday interaction, showing how \emph{frame shifts}---changes in
which schemas are active---are the fundamental unit of social
micro-dynamics.

IDT inherits and formalises this lineage.  The resonance
relation~$\res$ replaces the informal notion of ``associative link'';
the activated subset replaces the vague concept of ``what is salient'';
and the composition operation~$\compose$ replaces the hand-waving
appeal to ``contextual influence.''  The theorems that follow make
precise the claims that associationists, relevance theorists, and
frame analysts have long advanced informally.

% ── §1. Activation ───────────────────────────────────────────────────
\section{Activation}
\label{sec:ch4:activation}

\begin{definition}[Activation]
\label{def:activation}
An idea~$r$ is \emph{activated} by signal~$s$ if they resonate:
$\operatorname{activated}(r, s) \iff r \res s$.
Activation is the formal counterpart of Sperber's ``relevance'': a
signal is relevant to a schema if and only if it resonates with it.
\end{definition}

\begin{definition}[Activated Subset]
\label{def:activated-subset}
The \emph{activated subset} of repertoire~$\Rep$ under selection
predicate $\sigma$ is $\operatorname{activatedSubset}(\Rep, \sigma)
= \operatorname{filter}(\sigma, \Rep)$.
\end{definition}

\begin{definition}[Activated Interpretations]
\label{def:activated-interp}
The \emph{activated interpretations} compose only the selected elements
with the signal:
\[
  \operatorname{activatedInterpretations}(\Rep, s, \sigma)
  = \operatorname{map}(\lambda r.\; r \compose s,\;
    \operatorname{filter}(\sigma, \Rep)).
\]
This models selective attention: only relevant schemas participate in
interpretation.
\end{definition}

\begin{theorem}[Self-Activation]
\label{thm:self-activation}
Every idea activates itself: $\operatorname{activated}(r, r)$ for all~$r$.
\end{theorem}

\begin{proof}[Proof sketch]
By reflexivity of resonance.  See \texttt{self\_activation}.
\end{proof}

\begin{remark}
The formal basis of self-recognition.  Familiar symbols are immediately
meaningful because they resonate with what is already there.
\end{remark}

\begin{theorem}[Activation Symmetry]
\label{thm:activation-symm}
If $\operatorname{activated}(r, s)$, then $\operatorname{activated}(s, r)$.
\end{theorem}

\begin{proof}[Proof sketch]
By symmetry of resonance.  See \texttt{activation\_symm}.
\end{proof}

\begin{remark}
Relevance is reciprocal.  If a Christian schema is activated by a cross,
then the cross-schema would be activated by a Christian symbol.
Cultural exchange is inherently bidirectional.
\end{remark}

\begin{theorem}[Void Self-Activates]
\label{thm:void-self-activates}
$\operatorname{activated}(\void, \void)$.
\end{theorem}

\begin{proof}[Proof sketch]
By reflexivity: $\void \res \void$.
See \texttt{void\_self\_activates}.
\end{proof}

\begin{theorem}[Composition Preserves Activation (Right)]
\label{thm:activation-compose-right}
If $\operatorname{activated}(r, s)$, then
$\operatorname{activated}(r \compose c,\; s \compose c)$.
\end{theorem}

\begin{proof}[Proof sketch]
By resonance-compose-right.  See \texttt{activation\_compose\_right}.
\end{proof}

\begin{theorem}[Composition Preserves Activation (Left)]
\label{thm:activation-compose-left}
If $\operatorname{activated}(r, s)$, then
$\operatorname{activated}(c \compose r,\; c \compose s)$.
\end{theorem}

\begin{proof}[Proof sketch]
By resonance-compose-left.  See \texttt{activation\_compose\_left}.
\end{proof}

\begin{remark}
A ritual frame (opening prayer, spatial arrangement) that precedes two
resonant symbols does not destroy their mutual activation.  Context
preserves relevance.
\end{remark}

\begin{example}[The Cross as Multi-Activation Symbol]
\label{ex:cross-multi-activation}
The symbol of the cross provides a vivid illustration of how the same
signal activates radically different schemas depending on the
repertoire it encounters.  For a devout Catholic, the cross activates
a dense cluster of christological schemas: the crucifixion narrative,
the theology of redemption through suffering, the eucharistic
sacrifice, Marian devotion at the foot of the cross, and the
eschatological promise of resurrection.  These schemas are mutually
resonant, forming what Lakoff would call an idealized cognitive model
of salvific suffering.

For a mathematician, the identical visual form activates an entirely
disjoint set of schemas: Cartesian coordinate axes, the addition
operator in elementary arithmetic, the cross product of vectors, or
the symbol for the direct product of groups.  The resonance relation
$r \res s$ holds for a different subset of the repertoire, and the
activated interpretations accordingly diverge.

For a Red Cross worker, the cross activates schemas of humanitarian
aid, political neutrality, the Geneva Conventions, and emergency field
medicine.  The symbol's meaning is neither theological nor
mathematical but \emph{institutional}---it indexes a framework of
international law and organisational practice.

In each case, the signal is identical; what differs is the
repertoire~$\Rep$ and hence the activated subset.  The composition
preservation theorems (Theorems~\ref{thm:activation-compose-right}
and~\ref{thm:activation-compose-left}) guarantee that adding
contextual framing to the cross preserves these divergent activation
patterns rather than collapsing them.
\end{example}

\begin{example}[Priming in Advertising]
\label{ex:advertising-priming}
Consider a television advertisement for a luxury automobile.  Before
the car appears on screen, the viewer is shown a sequence of images:
a private jet landing on a sun-drenched runway, a penthouse overlooking
a glittering skyline, a well-dressed figure signing documents in a
boardroom of polished mahogany.  These images are not accidental; they
are primers.  Each image activates schemas of wealth, success,
aspiration, and social status.  By the time the automobile appears,
the viewer's repertoire has been reshaped: the schemas most recently
activated are those associated with luxury and achievement.  The car
is then interpreted not merely as a mode of transportation but as a
signifier of the lifestyle depicted in the preceding images.

The sequential interpretation associativity theorem
(Theorem~\ref{thm:sequential-activation}) explains why this works:
the priming sequence followed by the product reveal is algebraically
equivalent to a single compound signal that fuses aspiration-imagery
with automobile-imagery.  The advertiser's art lies in choosing primers
whose composition with the product signal maximises the desired
activation pattern.
\end{example}

\subsection*{Discussion: Activation as Cultural Attention}

The activation framework developed so far formalises what
anthropologists have long called \emph{cultural salience}---the
culturally variable tendency to notice, attend to, and find meaningful
certain features of the environment while ignoring others.  Mary
Douglas, in \emph{Purity and Danger} (1966)~\cite{douglas1966},
argued that every culture imposes a classification system that directs
attention to particular boundaries (pure/impure, sacred/profane,
inside/outside) and renders invisible everything that does not fit the
classificatory grid.  In IDT terms, Douglas's cultural classifications
are encoded in the repertoire~$\Rep$, and the boundaries she identifies
are precisely the regions of high resonance density: signals that fall
near a classificatory boundary activate many schemas simultaneously,
producing the sense of heightened significance---or danger---that
Douglas documented.

Conversely, signals that fall outside all cultural categories activate
nothing (the null activation theorem, Theorem~\ref{thm:null-activation})
and are experienced as meaningless noise.  Douglas's famous ``matter
out of place''---dirt as the by-product of classification---is, in
formal terms, a signal whose activated subset is empty relative to the
prevailing repertoire.  What counts as dirt, disorder, or anomaly is
not an objective property of the signal but a consequence of the
repertoire's resonance structure.

% ── §2. Activation and Interpretation ────────────────────────────────
\section{Activation and Interpretation Interaction}
\label{sec:ch4:interaction}

\begin{theorem}[Activated Interpretations Are a Sublist]
\label{thm:activated-sublist}
\[
  |\operatorname{activatedInterpretations}(\Rep, s, \sigma)|
  \leq |\operatorname{interpret}(\Rep, s)|.
\]
\end{theorem}

\begin{proof}[Proof sketch]
Filtering can only reduce list length; mapping preserves it.
See \texttt{activated\_interp\_sublist}.
\end{proof}

\begin{remark}
Attention is a filter, not a generator.  Selective attention can only
\emph{reduce} the set of interpretations, never add new ones.
\end{remark}

\begin{theorem}[Universal Activation = Full Interpretation]
\label{thm:universal-activation}
If $\sigma(r) = \mathtt{true}$ for all~$r$, then
$\operatorname{activatedInterpretations}(\Rep, s, \sigma)
 = \operatorname{interpret}(\Rep, s)$.
\end{theorem}

\begin{proof}[Proof sketch]
Filtering with the constant-true predicate is the identity.
See \texttt{universal\_activation}.
\end{proof}

\begin{remark}
A mind with zero attentional filters (the ``open mind'' ideal) interprets
everything.  Pure receptivity means every schema participates.
\end{remark}

\begin{theorem}[Null Activation = Empty Interpretation]
\label{thm:null-activation}
If $\sigma(r) = \mathtt{false}$ for all~$r$, then
$\operatorname{activatedInterpretations}(\Rep, s, \sigma) = []$.
\end{theorem}

\begin{proof}[Proof sketch]
Filtering with the constant-false predicate yields the empty list.
See \texttt{null\_activation}.
\end{proof}

\begin{remark}
Total cognitive shutdown---when nothing in your repertoire is deemed
relevant---produces zero interpretations.  This models culture shock or
the ``deer in headlights'' response to incomprehensible stimuli.
\end{remark}

\begin{example}[Culture Shock as Null Activation]
\label{ex:culture-shock}
An American businessman arrives in a remote village in the Papua New
Guinea highlands.  He is led to a clearing where men in elaborate
feathered headdresses are engaged in a ceremonial exchange of pearl
shells, pigs, and sweet potatoes, accompanied by stylised oratory,
rhythmic stamping, and the blowing of bamboo flutes.  Nothing in his
cognitive repertoire---business negotiation schemas, American social
etiquette, Western aesthetic categories---resonates with what he
observes.  The activation predicate $\sigma(r)$ returns
$\mathtt{false}$ for every element $r \in \Rep$.  The result, by
Theorem~\ref{thm:null-activation}, is an empty list of
interpretations: complete cognitive blankness.

This is the formal model of what Oberg (1960)~\cite{oberg1960} called
\emph{culture shock}---not merely unfamiliarity but the total failure
of one's interpretive apparatus.  The remedy, in IDT terms, is
repertoire expansion: by acquiring new schemas (through language
learning, participant observation, or explicit instruction), the
traveller gradually populates $\Rep$ with elements that resonate with
local signals, restoring the capacity for interpretation.
\end{example}

\begin{remark}[Connection to Artificial Intelligence: Attention Mechanisms]
\label{rem:attention-ai}
The activation framework bears a striking structural parallel with the
\emph{attention mechanism} in transformer neural
networks~\cite{vaswani2017}.  In a transformer, each input token is
projected into query, key, and value vectors; the attention score
between a query and a key determines how much the corresponding value
contributes to the output.  In IDT, a signal~$s$ plays the role of
the query, the repertoire elements~$r \in \Rep$ play the role of keys,
and the interpretations $r \compose s$ play the role of values.  The
resonance predicate $r \res s$ is a binary attention score:
it selects which values contribute and which are suppressed.

This analogy is more than metaphorical.  The activated interpretations
definition (Definition~\ref{def:activated-interp}) is a discrete,
binary version of the attention-weighted sum that transformers
compute.  The null activation theorem
(Theorem~\ref{thm:null-activation}) corresponds to the case where all
attention scores are zero---the model produces no output.  The
universal activation theorem
(Theorem~\ref{thm:universal-activation}) corresponds to uniform
attention---every key-value pair contributes equally.  Future work
might extend IDT's binary resonance to graded resonance, bringing the
formal framework even closer to the continuous attention scores used
in neural architectures.
\end{remark}

% ── §3. Double Activation ────────────────────────────────────────────
\section{Double Activation: Sequential Signals}
\label{sec:ch4:double}

\begin{theorem}[Sequential Interpretation Associativity]
\label{thm:sequential-activation}
\[
  \operatorname{interpret}(\operatorname{interpret}(\Rep, s_1), s_2)
  = \operatorname{interpret}(\Rep, s_1 \compose s_2).
\]
\end{theorem}

\begin{proof}[Proof sketch]
By functoriality of map and associativity of composition.
See \texttt{sequential\_activation}.
\end{proof}

\begin{remark}
Cultural priming works because sequential exposure to signals is
equivalent to a single compound signal.  Ritual sequences (opening chant
$\to$ sacrifice $\to$ closing prayer) produce the same interpretations
as a single complex ritual.
\end{remark}

\begin{theorem}[Resonant Sequential Identity]
\label{thm:resonant-sequential}
For any background~$r$ and signals $s_1, s_2$:
\[
  r \compose (s_1 \compose s_2) = (r \compose s_1) \compose s_2.
\]
\end{theorem}

\begin{proof}[Proof sketch]
By associativity.  See \texttt{resonant\_sequential\_identity}.
\end{proof}

\begin{theorem}[Void Priming is Identity]
\label{thm:void-priming}
\[
  \operatorname{interpret}(\operatorname{interpret}(\Rep, \void), s)
  = \operatorname{interpret}(\Rep, s).
\]
\end{theorem}

\begin{proof}[Proof sketch]
$\operatorname{interpret}(\Rep, \void) = \Rep$ by silence transparency,
so the second interpretation acts on the original repertoire.
See \texttt{void\_priming}.
\end{proof}

\begin{remark}
A ritual preceded by silence produces the same interpretations as the
ritual alone.  The ``pregnant pause'' before a sermon adds drama but not
information.
\end{remark}

\begin{theorem}[Post-Void is Identity]
\label{thm:post-void}
\[
  \operatorname{interpret}(\operatorname{interpret}(\Rep, s), \void)
  = \operatorname{interpret}(\Rep, s).
\]
\end{theorem}

\begin{proof}[Proof sketch]
By silence transparency applied to $\operatorname{interpret}(\Rep, s)$.
See \texttt{post\_void}.
\end{proof}

\subsection*{Discussion: Sequential Activation and Ritual Design}

The theorems of this section reveal that the careful ordering of
ritual elements is not arbitrary but reflects the algebraic structure
of composition.  Ritual designers---priests, shamans, liturgists,
and, in the modern world, advertisers and political
speechwriters---exploit sequential activation to shape the
interpretive landscape of their audiences.  The sequential
interpretation associativity theorem
(Theorem~\ref{thm:sequential-activation}) shows that the effect of a
ritual sequence can be analysed as a single compound signal, while
the void priming theorem (Theorem~\ref{thm:void-priming}) and the
post-void theorem (Theorem~\ref{thm:post-void}) establish that only
\emph{substantive} primers---those carrying non-trivial content---alter
the interpretive outcome.

\begin{example}[The Seder as Sequential Activation]
\label{ex:seder}
The Jewish Passover Seder follows a precise fifteen-step sequence:
\emph{Kadesh} (sanctification over wine), \emph{Urchatz} (hand
washing), \emph{Karpas} (dipping greens in salt water),
\emph{Yachatz} (breaking the middle matzah), \emph{Maggid}
(narrating the Exodus), and so on through \emph{Nirtzah} (acceptance).
Each step activates specific schemas---liberation, slavery, divine
intervention, communal identity, agricultural renewal---that prime
the participant for the next step.  The dipping of bitter herbs
(\emph{Maror}) in \emph{charoset} activates schemas of suffering and
sweetness simultaneously, and this compound activation primes the
participant for the meal (\emph{Shulchan Orech}), which is then
interpreted not merely as eating but as a re-enactment of the
celebratory feast of the freed Israelites.

The sequential associativity theorem tells us that this entire
fifteen-step cascade is algebraically equivalent to a single compound
signal $s_1 \compose s_2 \compose \cdots \compose s_{15}$.  But the
pedagogical genius of the Seder lies in \emph{decomposing} this
compound signal into steps, each of which is individually
interpretable and each of which prepares the cognitive ground for its
successor.
\end{example}

\begin{remark}[Connection to Pedagogy: Scaffolding]
\label{rem:scaffolding}
Vygotsky's concept of the \emph{zone of proximal development}
(ZPD)~\cite{vygotsky1978} and the pedagogical technique of
\emph{scaffolding} (Wood, Bruner, and Ross, 1976) describe the
practice of providing learners with structured support that enables
them to achieve tasks just beyond their current independent
capability.  The void priming theorem
(Theorem~\ref{thm:void-priming}) provides a formal criterion for
effective scaffolding: a preliminary lesson that activates nothing
new---that is, a void primer---adds no interpretive value.  Effective
scaffolding requires \emph{non-void primers} that genuinely activate
relevant schemas, reshaping the learner's repertoire so that the
target material falls within the activated subset.

The post-void theorem (Theorem~\ref{thm:post-void}) adds a
complementary lesson: trailing silence---review sessions that merely
repeat what has already been activated without introducing new
content---likewise adds nothing to the interpretive outcome.  The
implication for instructional design is that both the preamble and the
coda of a lesson must carry substantive content; purely ceremonial
framing is algebraically inert.
\end{remark}

% ── §4. Activation Depth Bounds ──────────────────────────────────────
\section{Activation Depth Bounds}
\label{sec:ch4:depth-bounds}

\begin{theorem}[Activated Interpretation Depth Bound]
\label{thm:activated-depth}
\[
  \operatorname{maxInterpDepth}(\operatorname{activatedInterpretations}(\Rep, s, \sigma))
  \leq \operatorname{maxInterpDepth}(\Rep) + \depth(s).
\]
\end{theorem}

\begin{proof}[Proof sketch]
By induction on $|\Rep|$.  Filtering can only select elements whose
depth is bounded by $\operatorname{maxInterpDepth}(\Rep)$; composition
adds at most $\depth(s)$.
See \texttt{activated\_depth\_bound}.
\end{proof}

\begin{remark}
Selective attention does not change the depth bound.  Relevance-filtering
is a horizontal cut, not a vertical one.
\end{remark}

\begin{theorem}[Activation-Filtered Count Bound]
\label{thm:activated-count}
\[
  |\operatorname{activatedInterpretations}(\Rep, s, \sigma)|
  \leq |\Rep|.
\]
\end{theorem}

\begin{proof}[Proof sketch]
By the length bound on filtered lists.
See \texttt{activated\_count\_bound}.
\end{proof}

\begin{remark}
You cannot hallucinate more schemas than you actually have, no matter how
stimulating the signal.
\end{remark}

\begin{example}[Selective Attention in Ethnography]
\label{ex:ethnographic-attention}
When Clifford Geertz observed the Balinese cockfight~\cite{geertz1973},
his ethnographic training activated schemas that ordinary Balinese
spectators---or casual Western tourists---lacked.  Where a tourist
might activate only schemas of animal cruelty or gambling, Geertz's
repertoire contained schemas of status rivalry, symbolic violence,
Weberian stratification, and Benthamite ``deep play.''  His trained
repertoire therefore produced a richer set of activated
interpretations.  But Theorem~\ref{thm:activated-count} guarantees
that even Geertz could not produce more interpretations than the
cardinality of his repertoire: $|\operatorname{activatedInterpretations}
(\Rep_{\text{Geertz}}, s_{\text{cockfight}}, \sigma)| \leq
|\Rep_{\text{Geertz}}|$.  The ethnographer's advantage is not
unbounded creativity but a \emph{larger and more finely differentiated
repertoire} through which more signals resonate.
\end{example}

\begin{remark}[Depth Bounds and Information Overload]
\label{rem:information-overload}
The activated interpretation depth bound
(Theorem~\ref{thm:activated-depth}) provides a formal explanation for
why information overload, though experientially overwhelming, is
\emph{cognitively bounded}.  In an age of constant digital
stimulation---news feeds, social media notifications, algorithmic
content streams---each individual signal can activate at most
$|\Rep|$ schemas, and the depth of any resulting interpretation is
bounded by $\operatorname{maxInterpDepth}(\Rep) + \depth(s)$.  The
subjective experience of overload arises not from unbounded depth of
processing but from the \emph{rapid succession} of signals, each
making finite but non-negligible demands on interpretive resources.

The so-called ``attention economy''---the competition among signals
for limited cognitive bandwidth---is thus formally bounded by
repertoire structure.  No amount of algorithmic optimisation can force
a mind to produce interpretations deeper or more numerous than its
repertoire permits.  The defense against information overload is not
faster processing but \emph{repertoire curation}: the deliberate
cultivation of schemas that resonate with signals one judges
important, and the pruning of schemas that resonate with noise.
\end{remark}

% ── §5. Resonance Chains ─────────────────────────────────────────────
\section{Resonance Chains and Activation Propagation}
\label{sec:ch4:chains}

\begin{theorem}[Resonant Backgrounds Produce Resonant Activations]
\label{thm:resonant-activations}
If $r_1 \res r_2$, then
$r_1 \compose s \;\res\; r_2 \compose s$.
\end{theorem}

\begin{proof}[Proof sketch]
By resonance-compose-right.
See \texttt{resonant\_backgrounds\_resonant\_activations}.
\end{proof}

\begin{remark}
Culturally similar people who are primed by the same stimulus produce
resonant responses.  This is Durkheim's collective effervescence: a
crowd sharing cultural schemas, exposed to the same ritual stimulus,
produces a harmonious collective response~\cite{durkheim1912}.
\end{remark}

\begin{theorem}[Activation Under Composition]
\label{thm:activation-composition}
If $\operatorname{activated}(r, s_1)$ and
$\operatorname{activated}(s_1, s_2)$, then
$\operatorname{activated}(r \compose s_1,\; s_1 \compose s_2)$.
\end{theorem}

\begin{proof}[Proof sketch]
By the full resonance-compose lemma.
See \texttt{activation\_composition\_chain}.
\end{proof}

\begin{remark}
Resonance ``propagates'' through compositional chains.  Multi-layered
rituals create cascading activations through the audience's schemas.
\end{remark}

\begin{theorem}[Void Is Universally Self-Activatable]
\label{thm:void-self-activatable}
$\operatorname{activated}(\void, \void)$.
\end{theorem}

\begin{proof}[Proof sketch]
By reflexivity of resonance.
See \texttt{void\_universally\_self\_activatable}.
\end{proof}

\begin{theorem}[Double Signal Depth Bound]
\label{thm:double-signal-depth}
\[
  \depth(r_2 \compose (r_1 \compose s))
  \leq \depth(r_2) + \depth(r_1) + \depth(s).
\]
\end{theorem}

\begin{proof}[Proof sketch]
Apply depth-compose-le twice.  See \texttt{double\_signal\_depth\_bound}.
\end{proof}

\begin{remark}
Cultural complexity does not compound super-linearly through chains of
interpretation.  A three-layer cultural transmission has bounded total
complexity.
\end{remark}

\begin{example}[Propaganda as Activation Chain]
\label{ex:propaganda}
The propaganda apparatus of totalitarian regimes provides a dark but
instructive illustration of activation chains.  Joseph Goebbels's
communication strategy in Nazi Germany involved the careful sequencing
of messages: initial signals activated schemas of national humiliation
(the Treaty of Versailles, economic hardship), which primed the
population for subsequent signals activating schemas of ethnic
scapegoating, which in turn primed for signals activating schemas of
militaristic renewal.  Each stage in the sequence depended on the
activation patterns established by its predecessors.

The activation composition theorem
(Theorem~\ref{thm:activation-composition}) shows that this cascade of
resonances propagates through compositional chains: if
$\operatorname{activated}(r, s_1)$ and
$\operatorname{activated}(s_1, s_2)$, then the compositions
$r \compose s_1$ and $s_1 \compose s_2$ are themselves mutually
activated.  But the double signal depth bound
(Theorem~\ref{thm:double-signal-depth}) provides a crucial
counterpoint: the total interpretive depth of such a chain is bounded
by the sum of the individual depths, not their product.  Propaganda
is dangerous, but its cognitive reach is formally limited by the
structure of the repertoires it targets.
\end{example}

\begin{remark}[Open Questions]
\label{rem:open-questions}
The activation framework developed in this chapter treats resonance as
a binary relation: a signal either resonates with a schema or it does
not.  Several natural extensions suggest themselves.

First, can activation be made \emph{quantitative} beyond the binary
resonates/does-not-resonate dichotomy?  A graded resonance
function $\rho(r, s) \in [0, 1]$ would allow modelling of
\emph{partial activation} and \emph{weak relevance}---phenomena that
are ubiquitous in everyday cognition.  A faint church bell in the
distance partially activates religious schemas without fully engaging
them; a foreign word that \emph{sounds like} a familiar word in one's
own language triggers tentative, low-confidence activation.  Graded
resonance connects naturally to fuzzy set theory~\cite{zadeh1965}
and to the graded category membership documented in cognitive
linguistics~\cite{rosch1975}.

Second, the current framework treats the repertoire as static during
a single activation event.  But activation itself may \emph{modify}
the repertoire: schemas that are frequently activated may become more
accessible (the ``frequency effect'' in psycholinguistics), while
schemas that are never activated may atrophy (the ``use it or lose it''
principle).  A dynamic activation theory would require enriching the
algebraic structure with a feedback mechanism from activation to
repertoire composition.

Third, the interaction between activation and social structure remains
largely unexplored.  When two minds with different repertoires are
exposed to the same signal, their divergent activation patterns may
generate communicative friction---or, under favourable conditions,
productive misunderstanding that drives cultural innovation.
Formalising this social dimension of activation is a task for future
chapters.
\end{remark}

% ── Summary ──────────────────────────────────────────────────────────
\section*{Chapter Summary}

This chapter formalised the theory of \emph{resonance activation}:
signals selectively activate repertoire elements through resonance,
producing filtered interpretations.  Self-activation
(Theorem~\ref{thm:self-activation}) and activation symmetry
(Theorem~\ref{thm:activation-symm}) establish the basic structure.
Universal activation recovers full interpretation
(Theorem~\ref{thm:universal-activation}), while null activation produces
cognitive shutdown (Theorem~\ref{thm:null-activation}).  Sequential
signals compose associatively
(Theorem~\ref{thm:sequential-activation}), and void priming is the
identity (Theorem~\ref{thm:void-priming}).  The activated depth bound
(Theorem~\ref{thm:activated-depth}) shows that selective attention does
not change the depth ceiling.  Resonance propagates through activation
chains (Theorem~\ref{thm:activation-composition}), providing the formal
basis for understanding how multi-layered rituals create cascading
effects through culturally constituted minds.

All eighteen theorems are verified in
\texttt{FormalAnthropology/IDT\_Signal\_Ch04.lean}.

% ══════════════════════════════════════════════════════════════════════
% Chapter 5: Misinterpretation and Cultural Distance
% ══════════════════════════════════════════════════════════════════════

\chapter{Misinterpretation and Cultural Distance}
\label{ch:misinterp}

\epigraph{``Understanding is always more than merely re-creating someone
else's meaning.''}{Hans-Georg Gadamer, \textit{Truth and Method}}

% ── Introduction ──────────────────────────────────────────────────────

The sender intends the receiver to compose signal~$s$ with
background~$r^*$ (the ``shared understanding'').  But the receiver
actually composes with~$r_{\mathrm{actual}}$---their real, possibly
different, cultural background.  The gap between
$r^* \compose s$ and $r_{\mathrm{actual}} \compose s$ is the formal
measure of \emph{misinterpretation}.

This chapter develops the formal theory of interpretive divergence.
Connecting to the Sapir--Whorf hypothesis~\cite{vygotsky1978}, which
holds that linguistic background shapes interpretation, and to Gadamer's
hermeneutic circle~\cite{gadamer1960}, which insists that understanding
requires pre-understanding, we show that cultural distance between
backgrounds governs interpretive distance between outcomes.

\medskip
\noindent\textbf{Chapter overview.}  We define the interpretation gap
(\S\ref{sec:ch5:gap}), analyse void backgrounds
(\S\ref{sec:ch5:void-bg}), develop iterated re-interpretation
(\S\ref{sec:ch5:iteration}), study resonance under iteration
(\S\ref{sec:ch5:resonance-iter}), and formalise cultural distance
(\S\ref{sec:ch5:distance}).

% ── §1. The Interpretation Gap ───────────────────────────────────────
\section{The Interpretation Gap}
\label{sec:ch5:gap}

\begin{definition}[Intended Interpretation]
\label{def:intended-interp}
The sender's \emph{intended interpretation} is
$\operatorname{intendedInterp}(r_{\mathrm{intended}}, s)
 = r_{\mathrm{intended}} \compose s$.
\end{definition}

\begin{definition}[Actual Interpretation]
\label{def:actual-interp}
The receiver's \emph{actual interpretation} is
$\operatorname{actualInterp}(r_{\mathrm{actual}}, s)
 = r_{\mathrm{actual}} \compose s$.
\end{definition}

\begin{definition}[Well-Interpreted]
\label{def:well-interpreted}
An interpretation is \emph{well-interpreted} if the actual and intended
interpretations resonate:
\[
  \operatorname{wellInterpreted}(r_{\mathrm{int}}, r_{\mathrm{act}}, s)
  \;\iff\;
  r_{\mathrm{int}} \compose s \;\res\; r_{\mathrm{act}} \compose s.
\]
This is weaker than equality but stronger than ``no relation.''
\end{definition}

\begin{theorem}[Resonant Backgrounds Guarantee Well-Interpretation]
\label{thm:resonant-well-interpreted}
If $r_{\mathrm{int}} \res r_{\mathrm{act}}$, then the signal~$s$ is
well-interpreted.
\end{theorem}

\begin{proof}[Proof sketch]
By resonance-compose-right applied to the resonant backgrounds.
See \texttt{resonant\_backgrounds\_well\_interpreted}.
\end{proof}

\begin{remark}
Members of the same culture (resonant backgrounds) always produce
well-interpreted outcomes for any signal.  Gadamer's ``fusion of
horizons'' is guaranteed when horizons resonate~\cite{gadamer1960}.
\end{remark}

\begin{theorem}[Self-Interpretation Always Succeeds]
\label{thm:self-interpretation}
$\operatorname{wellInterpreted}(r, r, s)$ for all $r, s$.
\end{theorem}

\begin{proof}[Proof sketch]
By reflexivity of resonance on the backgrounds.
See \texttt{self\_interpretation}.
\end{proof}

\begin{remark}
You always understand yourself.  Soliloquy is the degenerate case of
communication.
\end{remark}

\begin{theorem}[Well-Interpretation is Symmetric]
\label{thm:well-interp-symm}
If $\operatorname{wellInterpreted}(r_1, r_2, s)$, then
$\operatorname{wellInterpreted}(r_2, r_1, s)$.
\end{theorem}

\begin{proof}[Proof sketch]
By symmetry of resonance.  See \texttt{well\_interpretation\_symm}.
\end{proof}

\begin{remark}
Mutual intelligibility is symmetric.  If I can understand your ritual,
you can understand mine.  Cultural exchange is inherently
bidirectional---a formalisation of the anthropological observation that
first contact is always reciprocal.
\end{remark}

% ── §2. The Void Background ─────────────────────────────────────────
\section{The Void Background}
\label{sec:ch5:void-bg}

\begin{theorem}[Void Intended Transparency]
\label{thm:void-intended}
$\operatorname{intendedInterp}(\void, s) = s$.
\end{theorem}

\begin{proof}[Proof sketch]
By $\void \compose s = s$.  See \texttt{void\_intended\_transparency}.
\end{proof}

\begin{remark}
A sender who assumes zero shared context intends the signal to be
received ``as is.''  Any deviation from the raw signal is the receiver's
cultural contribution.  This is the formal model of the ``objective
journalist'' fantasy.
\end{remark}

\begin{theorem}[Void Actual Transparency]
\label{thm:void-actual}
$\operatorname{actualInterp}(\void, s) = s$.
\end{theorem}

\begin{proof}[Proof sketch]
By $\void \compose s = s$.  See \texttt{void\_actual\_transparency}.
\end{proof}

\begin{theorem}[Void-Void Perfect Interpretation]
\label{thm:void-void-perfect}
$\operatorname{wellInterpreted}(\void, \void, s)$ for all $s$.
\end{theorem}

\begin{proof}[Proof sketch]
Both interpret $s$ as $s$ itself, so the interpretations are identical
(and hence resonant).
See \texttt{void\_void\_perfect}.
\end{proof}

\begin{remark}
In the (impossible) limit of zero culture, there is zero
misinterpretation.  Misunderstanding is entirely a product of cultural
difference.
\end{remark}

\begin{theorem}[The Hermeneutic Circle]
\label{thm:hermeneutic-circle}
For any $r, s$:
$\operatorname{actualInterp}(r, s) \res \operatorname{actualInterp}(r, s)$.
\end{theorem}

\begin{proof}[Proof sketch]
By reflexivity of resonance.  See \texttt{hermeneutic\_circle}.
\end{proof}

\begin{remark}
Understanding always loops back to self-resonance.  You cannot escape
your own interpretive framework---even when you try to ``step outside,''
the result still resonates with your starting point.  This is the formal
content of Gadamer's hermeneutic circle~\cite{gadamer1960}.
\end{remark}

% ── §3. Iterated Re-Interpretation ───────────────────────────────────
\section{Iterated Re-Interpretation}
\label{sec:ch5:iteration}

\begin{definition}[Iterated Interpretation]
\label{def:iterate-interp}
The $n$-fold re-interpretation of signal~$s$ through background~$r$ is:
\[
  \operatorname{iterateInterp}(r, s, n) =
  \begin{cases}
    s & \text{if } n = 0, \\
    r \compose \operatorname{iterateInterp}(r, s, n-1) & \text{if } n > 0.
  \end{cases}
\]
This models rumination, echo chambers, and iterative cultural processing.
\end{definition}

\begin{lstlisting}[caption={Iterated interpretation in Lean~4},label=lst:iterate]
def iterateInterp (r s : I) : Nat -> I
  | 0 => s
  | n + 1 => IdeaticSpace.compose r (iterateInterp r s n)
\end{lstlisting}

\begin{theorem}[Zero Iteration = Raw Signal]
\label{thm:iterate-zero}
$\operatorname{iterateInterp}(r, s, 0) = s$.
\end{theorem}

\begin{proof}[Proof sketch]
By definition.  See \texttt{iterate\_zero}.
\end{proof}

\begin{theorem}[Single Iteration = Basic Interpretation]
\label{thm:iterate-one}
$\operatorname{iterateInterp}(r, s, 1) = r \compose s$.
\end{theorem}

\begin{proof}[Proof sketch]
By definition.  See \texttt{iterate\_one}.
\end{proof}

\begin{theorem}[Void Iteration Identity]
\label{thm:void-iterate}
$\operatorname{iterateInterp}(\void, s, n) = s$ for all $n$.
\end{theorem}

\begin{proof}[Proof sketch]
By induction on $n$, using $\void \compose x = x$ at the inductive step.
See \texttt{void\_iterate}.
\end{proof}

\begin{remark}
An ``empty mind'' processing a signal repeatedly never changes it.
Without cultural content, there is no interpretive drift.  Echo chambers
require content to echo.
\end{remark}

\begin{theorem}[Iterated Interpretation Depth Bound]
\label{thm:iterate-depth}
\[
  \depth(\operatorname{iterateInterp}(r, s, n))
  \leq n \cdot \depth(r) + \depth(s).
\]
\end{theorem}

\begin{proof}[Proof sketch]
By induction on $n$.  At step $n+1$, the depth-compose bound gives
$\depth(r \compose x) \leq \depth(r) + \depth(x)$; applying the
inductive hypothesis yields the result.
See \texttt{iterate\_depth\_bound}.
\end{proof}

\begin{remark}
Echo chambers have linearly bounded complexity.  Even infinite
rumination through a fixed cultural lens can only add linear depth.
Obsessive rumination is unproductive: each pass adds at most $\depth(r)$
more complexity.
\end{remark}

\begin{theorem}[Void Signal Single Iteration]
\label{thm:void-signal-iterate}
$\operatorname{iterateInterp}(r, \void, 1) = r$.
\end{theorem}

\begin{proof}[Proof sketch]
$r \compose \void = r$ by void-right.
See \texttt{void\_signal\_iterate\_one}.
\end{proof}

\begin{remark}
Processing ``nothing'' through a cultural lens yields the lens itself.
An empty news feed processed through conspiracy-theory schemas yields
the conspiracy schema unchanged.
\end{remark}

% ── §4. Resonance Under Iteration ────────────────────────────────────
\section{Resonance Under Iteration}
\label{sec:ch5:resonance-iter}

\begin{theorem}[Resonant Iteration Step]
\label{thm:resonant-step}
If $r_1 \res r_2$ and $x_1 \res x_2$, then
$r_1 \compose x_1 \;\res\; r_2 \compose x_2$.
\end{theorem}

\begin{proof}[Proof sketch]
By the full resonance-compose lemma.
See \texttt{resonant\_iteration\_step}.
\end{proof}

\begin{theorem}[Full Resonant Iteration]
\label{thm:full-resonant-iteration}
If $r_1 \res r_2$, then for all $n$:
\[
  \operatorname{iterateInterp}(r_1, s, n)
  \;\res\;
  \operatorname{iterateInterp}(r_2, s, n).
\]
\end{theorem}

\begin{proof}[Proof sketch]
By induction on $n$.  Base case: $s \res s$ by reflexivity.  Inductive
step: apply Theorem~\ref{thm:resonant-step} with the inductive
hypothesis.
See \texttt{resonant\_iteration}.
\end{proof}

\begin{remark}
Cultural similarity guarantees interpretive similarity at every stage
of processing.  This is why ``deep'' cultural understanding between
allied groups is possible---resonance propagates through all layers of
interpretation.
\end{remark}

\begin{theorem}[Self-Iteration Self-Resonance]
\label{thm:self-iter-resonance}
$\operatorname{iterateInterp}(r, s, n) \res
 \operatorname{iterateInterp}(r, s, n)$ for all $r, s, n$.
\end{theorem}

\begin{proof}[Proof sketch]
By reflexivity.  See \texttt{self\_iteration\_resonance}.
\end{proof}

\begin{remark}
The echo chamber's output is always ``consistent'' in the sense of
self-resonance---it just may not resonate with \emph{anything else}.
This is the formal trap of intellectual isolation.
\end{remark}

% ── §5. Cultural Distance ────────────────────────────────────────────
\section{Cultural Distance and Interpretation}
\label{sec:ch5:distance}

\begin{theorem}[Composition Preserves Well-Interpretation]
\label{thm:compound-well-interpreted}
If $s_1$ and $s_2$ are both well-interpreted under the same background
pair $(r_1, r_2)$, then
\[
  (r_1 \compose s_1) \compose (r_1 \compose s_2)
  \;\res\;
  (r_2 \compose s_1) \compose (r_2 \compose s_2).
\]
\end{theorem}

\begin{proof}[Proof sketch]
By the full resonance-compose lemma applied to the two well-interpretation
hypotheses.
See \texttt{compound\_well\_interpreted}.
\end{proof}

\begin{remark}
If two cultures understand each other's simple signals, they understand
compound signals too.  Mutual intelligibility at the atomic level
guarantees mutual intelligibility at the molecular level.  This is why
pidgins can bootstrap into creoles.
\end{remark}

\begin{theorem}[Interpretation Depth Gap]
\label{thm:interp-depth-gap}
\[
  \depth(\operatorname{actualInterp}(r_{\mathrm{actual}}, s))
  \leq \depth(r_{\mathrm{actual}}) + \depth(s).
\]
\end{theorem}

\begin{proof}[Proof sketch]
By depth-compose-le.  See \texttt{interpretation\_depth\_gap}.
\end{proof}

\begin{remark}
Misinterpretation does not create unbounded complexity.  Even the most
``wrong'' interpretation is bounded by the receiver's actual depth plus
the signal's depth.
\end{remark}

\begin{theorem}[Iteration Composition / Ritual Accretion]
\label{thm:iteration-composition}
\[
  \operatorname{iterateInterp}(r, s, n + m)
  = \operatorname{iterateInterp}(r,\;
    \operatorname{iterateInterp}(r, s, m),\; n).
\]
\end{theorem}

\begin{proof}[Proof sketch]
By induction on $n$.  See \texttt{iteration\_composition}.
\end{proof}

\begin{remark}
A lifetime of cultural processing ($n + m$ iterations) decomposes into
childhood processing ($m$~iterations) followed by adult processing
($n$~iterations).  Cultural development is \emph{sequentially
decomposable}: we can study ``what did early formation contribute?'' and
``what did later experience add?'' as separable phases.  This is the
formal basis of Bourdieu's distinction between primary and secondary
habitus formation~\cite{bourdieu1984}.
\end{remark}

% ── Summary ──────────────────────────────────────────────────────────
\section*{Chapter Summary}

This chapter developed the formal theory of misinterpretation within
IDT\@.  The \emph{interpretation gap} between intended and actual
interpretation is governed by the resonance between sender and receiver
backgrounds (Theorem~\ref{thm:resonant-well-interpreted}).
Well-interpretation is symmetric
(Theorem~\ref{thm:well-interp-symm}) and guaranteed for identical
backgrounds (Theorem~\ref{thm:self-interpretation}).  \emph{Iterated
re-interpretation} has linearly bounded depth
(Theorem~\ref{thm:iterate-depth}), and \emph{resonant iteration}
(Theorem~\ref{thm:full-resonant-iteration}) shows that culturally
similar backgrounds produce resonant interpretations at every stage.
The \emph{iteration composition} theorem
(Theorem~\ref{thm:iteration-composition}) provides the formal basis for
decomposing cultural development into sequential phases.

This concludes Part~I\@.  We have established the interpretive
framework: repertoires, signalling games, costly signals, resonance
activation, and misinterpretation.  Part~II develops the equilibrium
theory.

All eighteen theorems are verified in
\texttt{FormalAnthropology/IDT\_Signal\_Ch05.lean}.


% ── Part II: Equilibrium and Dynamics ─────────────────────────────────
\part{Equilibrium and Dynamics}
\label{part:equilibrium}

% =============================================================
%  Chapter 6 — Interpretive Equilibrium
%  Lean source: FormalAnthropology/IDT_Signal_Ch06.lean
% =============================================================
\chapter{Interpretive Equilibrium}
\label{ch:equilibrium}

\epigraph{%
  ``A convention is a regularity in behavior, sustained by an interest in
  coordination and an expectation that others will do their part.''}%
  {David Lewis, \textit{Convention} (1969)}

% ----------------------------------------------------------------
\section{Introduction: Why Cultures Don't Change}
\label{sec:eq-intro}

Every anthropologist eventually confronts the same puzzle: if cultures are
``arbitrary'' in Saussure's sense---if there is no deep reason for driving
on the left rather than the right, for shaking hands rather than bowing,
for calling a dog \emph{dog} rather than \emph{chien}---then why do cultures
persist at all?  Why doesn't every generation re-randomize its conventions?

The game-theoretic answer, originating with Thomas Schelling's
\emph{Strategy of Conflict} (1960) and formalized by David Lewis in
\emph{Convention} (1969), is deceptively simple: cultural forms persist
because \emph{no individual can profitably deviate}.  They are Nash equilibria
of an implicit coordination game.  A greeting ritual need not be optimal in
any absolute sense; it need only be \emph{locally stable}---better than any
unilateral alternative.

This chapter transplants that insight into the framework of Ideatic Diffusion
Theory.  We define an \emph{interpretation game}---a triple of a repertoire
(the receiver's cultural background), a sender payoff, and a receiver
payoff---and show that the concept of equilibrium has rich structural
consequences when the action space is an ideatic space $\IS$ equipped with
composition $\compose$, void $\void$, depth $\depth$, and resonance
$\res$.

The central results of the chapter fall into three groups.  First,
\emph{existence}: void (silence) is always an equilibrium when it weakly
dominates, and constant-payoff games admit universal equilibria
(\S\ref{sec:eq-babbling}).  Second, \emph{structure}: outcomes preserve
repertoire length, depth is linearly bounded, and void preserves depth
exactly (\S\ref{sec:eq-structure}--\S\ref{sec:eq-depth}).  Third,
\emph{robustness}: equilibria survive payoff agreement, monotone transforms,
and void extension (\S\ref{sec:eq-resonance}--\S\ref{sec:eq-advanced}).

Throughout, we provide anthropological commentary and Lean~4 code listings
that verify every theorem with zero \texttt{sorry} axioms.


% ================================================================
\section{Core Definitions}
\label{sec:eq-defs}

\begin{definition}[Repertoire]
\label{def:repertoire}
A \emph{repertoire} is a finite list of ideas in $\IS$:
\[
  \mathrm{Rep} \;=\; \mathrm{List}\;\IS.
\]
It represents the receiver's stock of cultural concepts, categories, and
schemas---what they bring to any interpretive encounter.%
\footnote{In Bourdieu's terms, the repertoire is the agent's \emph{habitus}
rendered as a finite enumeration of conceptual dispositions.}
\end{definition}

\begin{definition}[Interpretation]
\label{def:interpret}
Given a repertoire $R = [r_1,\dots,r_n]$ and a signal $s\in\IS$, the
\emph{interpretation} of $s$ through $R$ is
\[
  \mathrm{interpret}(R, s) \;=\; [r_1\compose s,\;\dots,\;r_n\compose s].
\]
Each background idea is composed with the incoming signal, producing a list
of new ideas---one per schema the receiver already holds.
\end{definition}

\begin{lstlisting}[caption={Interpretation in Lean~4.}]
def interpret (rep : Repertoire I) (signal : I) : List I :=
  rep.map (fun r => IdeaticSpace.compose r signal)
\end{lstlisting}

\begin{definition}[Interpretation Game]
\label{def:interp-game}
An \emph{interpretation game} $G = (R, u_S, u_R)$ consists of
\begin{enumerate}
  \item a repertoire $R$ (the receiver's cultural background),
  \item a sender payoff $u_S : \mathrm{List}\;\IS \to \mathbb{N}$,
  \item a receiver payoff $u_R : \mathrm{List}\;\IS \to \mathbb{N}$.
\end{enumerate}
Payoffs depend on the \emph{interpretation} (the list of composed ideas),
not on the signal directly.  This captures the anthropological insight that
meaning is always mediated by the receiver's culture.
\end{definition}

\begin{lstlisting}[caption={The interpretation game structure.}]
structure InterpGame (I : Type*) [IdeaticSpace I] where
  repertoire : Repertoire I
  sender_payoff : List I -> ℕ
  receiver_payoff : List I -> ℕ
\end{lstlisting}

\begin{definition}[Outcome]
\label{def:outcome}
The \emph{outcome} of sending signal $s$ in game $G$ is the interpretation
produced by composing the repertoire with $s$:
\[
  \mathrm{outcome}(G, s) \;=\; \mathrm{interpret}(G.\mathrm{Rep},\; s).
\]
\end{definition}

\begin{definition}[Equilibrium]
\label{def:equilibrium}
A signal $s$ is an \emph{equilibrium} of game $G$ if no alternative signal
$s'$ yields a strictly higher sender payoff:
\[
  \mathrm{IsEquilibrium}(G, s) \;\iff\;
  \forall\, s'.\;
  u_S\!\bigl(\mathrm{outcome}(G, s')\bigr)
  \;\le\;
  u_S\!\bigl(\mathrm{outcome}(G, s)\bigr).
\]
This is the Nash equilibrium condition for the sender, given the receiver's
fixed repertoire.%
\footnote{We suppress the receiver's strategic choice in this formalization,
treating the repertoire as exogenous.  Chapter~\ref{ch:multireceiver}
relaxes this.}
\end{definition}

\begin{lstlisting}[caption={Equilibrium in Lean~4.}]
def IsEquilibrium (G : InterpGame I) (s : I) : Prop :=
  forall (s' : I), G.sender_payoff (outcome G s') <=
    G.sender_payoff (outcome G s)
\end{lstlisting}


% ================================================================
\section{Babbling Equilibrium and Constant-Payoff Games}
\label{sec:eq-babbling}

The deepest result of Lewis's \emph{Convention} is often overlooked: in
\emph{every} signalling game, there exists a ``babbling'' equilibrium in
which signals carry no information.  In IDT, this corresponds to the void
signal.

\begin{theorem}[Outcome of Void]
\label{thm:outcome-void}
$\mathrm{outcome}(G, \void) = G.\mathrm{Rep}.$
Silence leaves meaning unchanged: the interpretation of void through any
repertoire is the repertoire itself.
\end{theorem}

\begin{proof}[Proof sketch]
By induction on the repertoire list.  The base case is trivial.  For the
inductive step, $r\compose\void = r$ by the right-identity axiom of
$\IS$, so the head is preserved and the tail follows from the inductive
hypothesis.
\end{proof}

This result is the linchpin of everything that follows.  It says that
silence is \emph{transparent}: the receiver's world is unchanged.  In
anthropological terms, a ritual that communicates nothing leaves the
participants exactly where they were.

\begin{theorem}[Babbling Equilibrium / Schelling's Silence (6.1)]
\label{thm:babbling-equilibrium}
If no signal produces a strictly better payoff than silence, then $\void$
is an equilibrium:
\[
  \bigl(\forall\,s.\; u_S(\mathrm{outcome}(G,s)) \le
  u_S(\mathrm{outcome}(G,\void))\bigr)
  \;\implies\;
  \mathrm{IsEquilibrium}(G,\void).
\]
\end{theorem}

\begin{proof}[Proof sketch]
Immediate from the definition of \textsf{IsEquilibrium}: the hypothesis
\emph{is} the conclusion.
\end{proof}

Lewis (1969) showed that babbling equilibria exist in every signalling game.
In IDT, the proof is one line---because the algebraic structure of $\IS$
(specifically, the existence of $\void$) guarantees a universal fallback.
This is the formal foundation of Schelling's focal-point theory: when
nothing you say can improve on silence, silence is optimal.

\begin{theorem}[Constant-Payoff Equilibrium (6.2)]
\label{thm:constant-payoff-equilibrium}
If every signal yields the same payoff~$k$, then every signal is an equilibrium:
\[
  \bigl(\forall\,s.\; u_S(\mathrm{outcome}(G,s)) = k\bigr)
  \;\implies\;
  \forall\,s.\;\mathrm{IsEquilibrium}(G,s).
\]
\end{theorem}

\begin{proof}[Proof sketch]
For any candidate deviator $s'$, rewrite both payoffs to $k$ and apply
$k\le k$.
\end{proof}

Anthropologically, constant-payoff games formalize \emph{radical cultural
relativism}: if all rituals are equally valued, none faces selection pressure,
and every cultural form persists indefinitely.  The theorem explains why
highly egalitarian societies display persistent ritual variation---there is
no payoff gradient to drive convergence.

\begin{theorem}[Equilibrium is Reflexive (6.3)]
\label{thm:equilibrium-refl}
Every signal is at least as good as itself:
$u_S(\mathrm{outcome}(G,s)) \le u_S(\mathrm{outcome}(G,s))$.
\end{theorem}

This is the trivial direction of the Nash definition---every strategy weakly
dominates itself.  Though mathematically obvious, it serves as the base case
for inductive proofs of equilibrium preservation in later chapters.


% ================================================================
\section{Outcome Structural Properties}
\label{sec:eq-structure}

Outcomes have a rigid structure: signals transform but never create or
destroy ideas.

\begin{theorem}[Outcome Length Invariance (6.4)]
\label{thm:outcome-length}
$|\mathrm{outcome}(G, s)| = |G.\mathrm{Rep}|$ for all signals $s$.
\end{theorem}

\begin{proof}[Proof sketch]
$\mathrm{outcome}$ is a \texttt{List.map}, which preserves length.
\end{proof}

Cultural contact does not change how many schemas a society holds---it
transforms each one.  The ``size'' of a culture's conceptual inventory is
invariant under external signalling.  This is a non-trivial constraint: it
rules out models in which signals can spontaneously generate new categories
in the receiver's mind.%
\footnote{Category \emph{creation} requires mechanisms beyond simple
composition---see the discussion of ``conceptual blending'' in
Chapter~\ref{ch:lossy}.}

\begin{theorem}[Outcome Length Independence (6.5)]
\label{thm:outcome-length-eq}
Two different signals produce interpretation lists of the same length:
$|\mathrm{outcome}(G, s_1)| = |\mathrm{outcome}(G, s_2)|$.
\end{theorem}

An immediate corollary of Theorem~\ref{thm:outcome-length}: both sides
equal $|G.\mathrm{Rep}|$.  Game-theoretically, the ``action space'' of
interpretations has constant dimensionality across all possible signals.

\begin{theorem}[Empty Repertoire Trivial Outcome (6.6)]
\label{thm:empty-repertoire-outcome}
A receiver with no ideas produces no interpretations:
$\mathrm{outcome}(\langle [], \ldots\rangle, s) = []$.
\end{theorem}

A mind with zero cultural content cannot interpret anything.  There is no
``raw perception''---interpretation requires prior ideas.  This is the
formal counterpart of Kant's dictum: ``Intuitions without concepts are
blind.''

\begin{theorem}[Singleton Repertoire Outcome (6.7)]
\label{thm:singleton-outcome}
A single-idea mind produces exactly one interpretation:
$\mathrm{outcome}(\langle [r], \ldots\rangle, s) = [r\compose s]$.
\end{theorem}

Monocultural societies have exactly one interpretation of every event.
The singleton case is the simplest non-trivial example and will serve as
a running test case in Chapters~\ref{ch:persuasion}--\ref{ch:multireceiver}.


% ================================================================
\section{Depth Bounds on Outcomes}
\label{sec:eq-depth}

The depth function $\depth$ measures interpretive complexity.  How does it
behave under the \texttt{outcome} operation?

\begin{theorem}[Outcome Depth Bound (6.8)]
\label{thm:outcome-depth-bound}
\[
  \sum_{i} \depth\bigl(\mathrm{outcome}(G,s)_i\bigr)
  \;\le\;
  \sum_{i} \depth(r_i) \;+\; |R|\cdot\depth(s).
\]
The total depth of all interpretations is bounded by the repertoire's total
depth plus the repertoire size times the signal depth.
\end{theorem}

\begin{proof}[Proof sketch]
Induction on the repertoire.  At each step, the depth of $r\compose s$ is
at most $\depth(r) + \depth(s)$ by the subadditivity of composition depth.
Summing over $|R|$ elements yields the bound.
\end{proof}

Game-theoretically, the sender's ``persuasive reach'' is linearly bounded
by signal complexity times audience size.  This result foreshadows the
persuasion bounds of Chapter~\ref{ch:persuasion}: even before introducing
formal limits on marginal depth, we see that total impact is $O(|R|\cdot\depth(s))$.

\begin{theorem}[Void Outcome Preserves Depth (6.9)]
\label{thm:void-outcome-depth}
$\mathrm{depthSum}(\mathrm{outcome}(G,\void)) = \mathrm{depthSum}(G.\mathrm{Rep}).$
\end{theorem}

Silence preserves total interpretive complexity exactly---not just
approximately, but with zero loss.  The ``cost'' of silence is exactly zero
in every dimension.

\begin{lstlisting}[caption={Outcome depth bound in Lean~4.}]
theorem outcome_depth_bound
    (rep : Repertoire I) (s : I) (sp rp : List I -> ℕ) :
    depthSum (outcome (<rep, sp, rp> : InterpGame I) s) <=
    depthSum rep + rep.length * IdeaticSpace.depth s
\end{lstlisting}


% ================================================================
\section{Resonance and Equilibrium}
\label{sec:eq-resonance}

Resonance ($\res$) captures structural similarity between ideas.  How does
it interact with equilibrium?

\begin{theorem}[Resonant Signals Produce Resonant Outcomes (6.10)]
\label{thm:resonant-signals-outcomes}
If $s_1\res s_2$, then for each background idea $r$ in the repertoire,
$r\compose s_1 \res r\compose s_2$.
\end{theorem}

Structurally similar rituals---Christmas and Hanukkah, weddings across
cultures---produce structurally similar experiences for each participant.
This is why comparative anthropology is possible at all: resonant signals
yield resonant interpretations, ensuring that cross-cultural comparison is
not an act of pure projection.

\begin{theorem}[Self-Resonance of Outcomes (6.11)]
\label{thm:outcome-self-resonance}
Every outcome resonates with itself at each position:
$r\compose s \res r\compose s$.
\end{theorem}

A trivial consequence of the reflexivity of resonance, but conceptually
important: a single signal always produces self-consistent interpretations.
No signal can produce an internally contradictory outcome.

\begin{theorem}[Equilibrium Under Payoff Agreement (6.12)]
\label{thm:equilibrium-payoff-agreement}
If two games $G_1, G_2$ have the same repertoire and their sender payoffs
agree on all outcomes, then an equilibrium for $G_2$ is an equilibrium
for $G_1$:
\[
  \bigl(G_1.\mathrm{Rep} = G_2.\mathrm{Rep}\bigr)
  \;\wedge\;
  \bigl(\forall\,s.\; u_{S,1}(\mathrm{outcome}(G_1,s)) = u_{S,2}(\mathrm{outcome}(G_2,s))\bigr)
  \;\implies\;
  \bigl(\mathrm{IsEquilibrium}(G_2,s) \implies \mathrm{IsEquilibrium}(G_1,s)\bigr).
\]
\end{theorem}

\begin{proof}[Proof sketch]
Substitute the payoff agreement into the equilibrium condition for $G_2$ and
apply it to $G_1$.
\end{proof}

If a cultural form is stable in one payoff environment, it remains stable in
any environment that produces the same relative rankings.  This is the formal
version of the sociological observation that institutions survive ``rebranding''
as long as the underlying incentive structure is preserved.


% ================================================================
\section{Composition and Sequential Signalling}
\label{sec:eq-composition}

Cultural communication is rarely a single signal.  Rituals have multiple
acts; conversations have turns.  This section studies how equilibrium
interacts with sequential composition.

\begin{theorem}[Sequential Signal Composition (6.13)]
\label{thm:sequential-signal}
$(r\compose s_1)\compose s_2 = r\compose(s_1\compose s_2).$
\end{theorem}

A ritual with two acts can be understood as a single compound act---but only
from the perspective of someone who has already absorbed the first act into
their background.  This is the associativity axiom of $\IS$, applied to the
signalling context.

\begin{theorem}[Equilibrium Stability Under Void Extension (6.14)]
\label{thm:equilibrium-void-extension}
$\mathrm{outcome}(G,\; s\compose\void) = \mathrm{outcome}(G,s).$
\end{theorem}

\begin{proof}[Proof sketch]
For each $r$ in the repertoire, $r\compose(s\compose\void) = r\compose s$
by the right-identity axiom.  Hence the map is pointwise equal.
\end{proof}

Padding a message with silence doesn't change the game.  This is the formal
version of ``cheap talk is irrelevant when it carries no information.''  In
anthropological terms, the elaborate silences that frame ritual speech
(the pause before a toast, the breath before a prayer) are formally inert
with respect to the equilibrium structure.

\begin{theorem}[Dual Repertoire Combination (6.15)]
\label{thm:dual-repertoire}
$\mathrm{outcome}(\langle R_1\!+\!+\!R_2, \ldots\rangle, s)
= \mathrm{outcome}(\langle R_1,\ldots\rangle, s)
\!+\!+\;\mathrm{outcome}(\langle R_2,\ldots\rangle, s).$
\end{theorem}

A bilingual mind's interpretation is the union of each language's
interpretation.  There is no mysterious ``third space'' of bilingual
meaning---just the disjoint combination.%
\footnote{This is a strong claim that will be qualified in
Chapter~\ref{ch:echo}, where iterated self-interpretation can create
genuinely new structures through repeated composition.}

\begin{lstlisting}[caption={Dual repertoire combination in Lean~4.}]
theorem dual_repertoire_outcome
    (rep_1 rep_2 : Repertoire I) (s : I) (sp rp : List I -> ℕ) :
    outcome (<rep_1 ++ rep_2, sp, rp> : InterpGame I) s =
    outcome (<rep_1, sp, rp> : InterpGame I) s ++
    outcome (<rep_2, sp, rp> : InterpGame I) s
\end{lstlisting}


% ================================================================
\section{Advanced Equilibrium Properties}
\label{sec:eq-advanced}

We close with three theorems on the robustness of equilibria under
transformations of the payoff structure.

\begin{theorem}[Maximum Payoff Equilibrium (6.16)]
\label{thm:max-payoff-equilibrium}
If a signal $s$ achieves the maximum possible sender payoff~$M$,
then $s$ is an equilibrium:
\[
  u_S(\mathrm{outcome}(G,s)) = M
  \;\wedge\;
  \forall\,s'.\; u_S(\mathrm{outcome}(G,s')) \le M
  \;\implies\;
  \mathrm{IsEquilibrium}(G,s).
\]
\end{theorem}

The most sacred ritual---the one that maximizes collective effervescence
(Durkheim)---is always in equilibrium because nothing can improve on the
maximum.  This is tautological in isolation, but gains force in combination
with depth bounds: since persuasion is bounded (Chapter~\ref{ch:persuasion}),
the maximum payoff is itself bounded.

\begin{theorem}[Equilibrium Preserved Under Monotone Transform (6.17)]
\label{thm:equilibrium-monotone}
If $f:\mathbb{N}\to\mathbb{N}$ is monotone and $s$ is an equilibrium under
payoff~$u_S$, then $s$ is an equilibrium under $f\circ u_S$:
\[
  (a\le b \implies f(a)\le f(b))
  \;\wedge\;
  \mathrm{IsEquilibrium}(G,s)
  \;\implies\;
  \forall\,s'.\; f(u_S(\mathrm{outcome}(G,s'))) \le f(u_S(\mathrm{outcome}(G,s))).
\]
\end{theorem}

Cultural stability is robust under monotone revaluation.  If a society
decides that all payoffs should be doubled (inflation), squared (increasing
returns), or logarithmically compressed (diminishing marginal utility), the
same cultural forms remain stable.  The theorem identifies the essential
feature: it is the \emph{ordering} of payoffs, not their magnitudes, that
determines equilibrium.

\begin{theorem}[Void Equilibrium in Zero-Payoff Games (6.18)]
\label{thm:void-equilibrium-zero}
In a game where every signal yields zero payoff, void is an equilibrium:
$\mathrm{IsEquilibrium}(\langle R, \lambda\_.\,0, \lambda\_.\,0\rangle,\;\void).$
\end{theorem}

\begin{proof}[Proof sketch]
Every signal gives payoff~$0$, so every signal is an equilibrium
(Theorem~\ref{thm:constant-payoff-equilibrium} with $k=0$).  In particular,
$\void$ is.
\end{proof}

In a society where communication carries no reward (Durkheim's \emph{anomie}),
silence is trivially optimal.  This is the formal model of communicative
breakdown: when nothing matters, nothing is said.

\begin{lstlisting}[caption={Void equilibrium in zero-payoff games.}]
theorem void_equilibrium_zero_game (rep : Repertoire I) :
    IsEquilibrium
      (<rep, fun _ => 0, fun _ => 0> : InterpGame I)
      (IdeaticSpace.void : I)
\end{lstlisting}


% ================================================================
\section{Conclusion}
\label{sec:eq-conclusion}

This chapter established the game-theoretic foundations of interpretive
stability within IDT\@.  We showed that the algebraic structure of ideatic
spaces---composition, void, depth, resonance---yields a rich theory of
equilibrium without any additional axioms.

Three themes will recur in the chapters ahead.  First, \emph{void as
universal fallback}: silence is always available and always stable
(Theorem~\ref{thm:babbling-equilibrium}).  Second, \emph{linear depth
bounds}: the impact of any signal is bounded by its own complexity times
the audience size (Theorem~\ref{thm:outcome-depth-bound}).  Third,
\emph{resonance preservation}: structurally similar inputs yield
structurally similar outputs
(Theorem~\ref{thm:resonant-signals-outcomes}).

Chapter~\ref{ch:persuasion} turns from equilibrium existence to the
\emph{limits} of persuasion: how much can a signal change a receiver's
depth, and why propaganda is subadditive.

% =============================================================
%  Chapter 7 — Persuasion Bounds and Propaganda
%  Lean source: FormalAnthropology/IDT_Signal_Ch07.lean
% =============================================================
\chapter{Persuasion Bounds and Propaganda}
\label{ch:persuasion}

\epigraph{%
  ``The ideas of the ruling class are in every epoch the ruling ideas:
  i.e., the class which is the ruling material force of society is at the
  same time its ruling intellectual force.''}%
  {Karl Marx, \textit{The German Ideology} (1846)}

% ----------------------------------------------------------------
\section{Introduction: The Algebra of ``Hearts and Minds''}
\label{sec:pers-intro}

Every propagandist, preacher, and politician confronts the same constraint,
whether they know it or not: \emph{you cannot add more complexity to someone's
worldview than the complexity of your message allows}.  A bumper sticker
cannot produce a Ph.D.-level understanding; a three-word slogan cannot
install a twenty-variable causal model.  Conversely, a complex treatise
cannot simplify a mind below its current depth---complexity is added, not
subtracted, by composition.

Antonio Gramsci grasped this intuitively when he argued that hegemony requires
not just force but \emph{cultural production}---the sustained manufacture of
complex ideological structures capable of reshaping the depth of popular
consciousness.\footnote{Gramsci, \textit{Prison Notebooks}, Q.~12 (1932).
The concept of ``cultural hegemony'' names exactly the problem that this
chapter formalizes: how much ideological depth can a ruling class actually
install in a population?}  Michel Foucault extended the insight: power
operates not by loud commands but by quiet restructuring of the categories
through which subjects perceive the world---what he called
\emph{power/knowledge}.\footnote{Foucault, \textit{Discipline and Punish}
(1975); \textit{The History of Sexuality}, vol.~1 (1976).}

This chapter provides the mathematical backbone for these intuitions.  We
define \emph{marginal depth}---how much a signal adds to a receiver's
interpretive complexity---and prove that it is bounded above by the signal's
own depth (the \emph{fundamental persuasion bound}).  We then show that
chaining messages is \emph{subadditive}: the total persuasive capacity of
a compound message is at most the sum of its parts.  Repetition helps, but
only linearly---never exponentially.  Even totalitarian propaganda, with
unlimited repetition, faces a linear ceiling.

These results have sharp implications for anthropological theory.  They
explain why ``hearts and minds'' campaigns fail not because of insufficient
repetition, but because of the fundamental \emph{subadditivity} of ideatic
composition.  You cannot shout your way to depth.


% ================================================================
\section{Core Definitions}
\label{sec:pers-defs}

\begin{definition}[Marginal Depth]
\label{def:marginal-depth}
The \emph{marginal depth} of signal $s$ for receiver $r$ is the increase
in depth caused by composition:
\[
  \mathrm{marginalDepth}(r,s) \;=\; \depth(r\compose s) - \depth(r),
\]
where $-$ denotes truncated (natural number) subtraction.
\end{definition}

The marginal depth of a sermon to a parishioner is how much \emph{more}
complex their worldview becomes after hearing it, relative to what they
had before.  It is inherently receiver-dependent: the same lecture may add
five ``layers'' to a novice and zero to an expert.

\begin{definition}[Persuasion Bound]
\label{def:persuasion-bound}
The \emph{persuasion bound} of a signal $s$ is its own depth:
\[
  \mathrm{persuasionBound}(s) \;=\; \depth(s).
\]
This is an \emph{absolute} ceiling, independent of the receiver.
\end{definition}

\begin{lstlisting}[caption={Marginal depth and persuasion bound in Lean~4.}]
def marginalDepth (r s : I) : ℕ :=
  IdeaticSpace.depth (IdeaticSpace.compose r s) -
    IdeaticSpace.depth r

def persuasionBound (s : I) : ℕ := IdeaticSpace.depth s
\end{lstlisting}


% ================================================================
\section{The Fundamental Persuasion Bound}
\label{sec:pers-fundamental}

\begin{theorem}[Marginal Depth Bound (7.1) --- The Fundamental Theorem of Propaganda]
\label{thm:marginal-depth-bound}
For all receivers $r$ and signals $s$:
\[
  \mathrm{marginalDepth}(r,s) \;\le\; \mathrm{persuasionBound}(s)
  \;=\; \depth(s).
\]
\end{theorem}

\begin{proof}[Proof sketch]
By subadditivity of composition depth: $\depth(r\compose s)\le \depth(r)+\depth(s)$.
Subtracting $\depth(r)$ from both sides gives $\mathrm{marginalDepth}(r,s)\le\depth(s)$.
\end{proof}

No matter how receptive the audience, a message cannot add more complexity
than it intrinsically contains.  A bumper sticker cannot produce a
Ph.D.-level understanding.  This is the \textbf{fundamental theorem of
propaganda theory}: persuasive capacity is bounded by signal depth.

The result vindicates Gramsci's insistence that effective hegemony requires
\emph{organic intellectuals}---agents capable of producing signals of
sufficient depth to restructure popular consciousness.  Simple slogans
(low depth) have low persuasion bounds.  Only sustained, complex cultural
production can deeply transform a worldview.

\begin{theorem}[Void Has Zero Marginal Depth (7.2)]
\label{thm:void-marginal-depth}
$\mathrm{marginalDepth}(r, \void) = 0$ for all $r$.
\end{theorem}

\begin{proof}[Proof sketch]
$r\compose\void = r$ by the right-identity axiom, so $\depth(r\compose\void)
- \depth(r) = 0$.
\end{proof}

Silence cannot persuade.  Quaker silence, Buddhist \emph{noble silence}, and
the minute of silence for the dead change nothing about the listener's
depth.  They are powerful \emph{socially}---as markers of solidarity, grief,
or contemplation---but informationally null.%
\footnote{This creates a sharp distinction between the \emph{social} and
\emph{informational} functions of silence, a distinction that Foucault's
analysis of ``incitement to discourse'' (\textit{History of Sexuality},
vol.~1) explores from the opposite direction.}

\begin{theorem}[Void Has Zero Persuasion Bound (7.3)]
\label{thm:void-persuasion-bound}
$\mathrm{persuasionBound}(\void) = 0.$
\end{theorem}

Independent of the receiver, the null message has zero persuasive capacity.
This is a property of the signal itself, not of any particular audience.

\begin{theorem}[Persuasion Bound Nonnegativity (7.4)]
\label{thm:persuasion-nonneg}
$\mathrm{persuasionBound}(s) \ge 0$ for all $s$.
\end{theorem}

Trivially true for natural numbers, but conceptually important: there are
no ``anti-signals'' that reduce the maximum possible effect below zero.
Every signal has non-negative persuasive potential.


% ================================================================
\section{Deeper Signals Are More Persuasive}
\label{sec:pers-monotonicity}

\begin{theorem}[Monotonicity of Persuasion Bound (7.5)]
\label{thm:persuasion-mono}
If $\depth(s_1) \le \depth(s_2)$, then
$\mathrm{persuasionBound}(s_1) \le \mathrm{persuasionBound}(s_2)$.
\end{theorem}

A complex ritual has \emph{more potential} to transform a participant's
worldview than a simple one.  This doesn't mean it \emph{will}---marginal
depth depends on the receiver---but it \emph{can}.  Complexity is a
\emph{necessary} (not sufficient) condition for deep persuasion.

This result gives formal content to Foucault's concept of
\emph{power/knowledge}: the capacity to restructure categories of thought
(persuasion) scales with the complexity of the discursive apparatus
(signal depth).  Simple commands have low persuasion bounds; complex
disciplinary regimes have high ones.

\begin{theorem}[Self-Composition Persuasion Bound (7.6)]
\label{thm:self-compose-persuasion}
$\mathrm{persuasionBound}(s\compose s) \le 2\cdot\mathrm{persuasionBound}(s).$
\end{theorem}

\begin{proof}[Proof sketch]
$\depth(s\compose s)\le \depth(s)+\depth(s) = 2\depth(s)$ by subadditivity.
\end{proof}

Repeating a message doubles the \emph{maximum possible} effect, but the
\emph{actual} effect depends on the receiver.  Propaganda works by
repetition, but with diminishing returns.

\begin{theorem}[ComposeN Persuasion Bound (7.7)]
\label{thm:composeN-persuasion}
$\mathrm{persuasionBound}(\compN{s}{n}) \le n\cdot\mathrm{persuasionBound}(s).$
\end{theorem}

Even with unlimited repetition, persuasive capacity grows at most
linearly---never exponentially.  This is why totalitarian propaganda,
despite its volume, has bounded effectiveness.  The 20th century
demonstrated this repeatedly: neither Nazi Germany nor the Soviet Union
achieved total ideological control despite saturation-level messaging.%
\footnote{Hannah Arendt's \textit{Origins of Totalitarianism} (1951)
documents both the ambition and the failure of totalitarian propaganda
to achieve total depth penetration.}

\begin{lstlisting}[caption={ComposeN persuasion bound in Lean~4.}]
theorem composeN_persuasion_bound (s : I) (n : ℕ) :
    persuasionBound (composeN s n) <=
      n * persuasionBound s
\end{lstlisting}


% ================================================================
\section{Subadditivity: Chaining Doesn't Multiply}
\label{sec:pers-subadditivity}

\begin{theorem}[Persuasion Subadditivity (7.8)]
\label{thm:persuasion-subadditivity}
\[
  \mathrm{persuasionBound}(s_1\compose s_2)
  \;\le\;
  \mathrm{persuasionBound}(s_1) + \mathrm{persuasionBound}(s_2).
\]
\end{theorem}

\begin{proof}[Proof sketch]
Direct from the subadditivity of depth under composition:
$\depth(s_1\compose s_2)\le\depth(s_1)+\depth(s_2)$.
\end{proof}

A two-part message cannot be more persuasive than the sum of its parts.
This rules out ``synergistic propaganda''---the idea that combining
messages could produce super-additive persuasion.  It is the formal
foundation of bounded rationality in cultural transmission.

The subadditivity theorem has immediate policy implications.  If a state
seeks to increase persuasive capacity, it must increase the \emph{depth} of
its messages, not merely their \emph{number}.  Quantity is strictly
dominated by quality as a persuasion strategy.

\begin{theorem}[Three-Message Persuasion Bound (7.9)]
\label{thm:three-message-bound}
\[
  \mathrm{persuasionBound}\bigl((s_1\compose s_2)\compose s_3\bigr)
  \;\le\;
  \mathrm{persuasionBound}(s_1) +
  \mathrm{persuasionBound}(s_2) +
  \mathrm{persuasionBound}(s_3).
\]
\end{theorem}

Multi-act rituals (liturgy with readings, sermon, and communion) have bounded
total impact $\le$ the sum of individual act complexities.  The extension to
$n$~messages follows by induction.

\begin{theorem}[Marginal Depth of Void Receiver (7.10)]
\label{thm:marginal-void-receiver}
$\mathrm{marginalDepth}(\void, s) = \depth(s).$
\end{theorem}

\begin{proof}[Proof sketch]
$\void\compose s = s$ by left-identity, and $\depth(\void)=0$, so
$\depth(s) - 0 = \depth(s)$.
\end{proof}

The ``blank slate'' receiver (\emph{tabula rasa}) absorbs every signal
fully.  This is the formal version of Locke's epistemology---and its
anthropological critique: no real human is a blank slate, so no real human
absorbs a signal fully.  The void receiver is a theoretical limit, never
achieved in practice.


% ================================================================
\section{Persuasion Through Chains}
\label{sec:pers-chains}

What happens when a signal passes through an intermediary?  A teacher
receives a text, processes it through their own background, and transmits
the result to a student.  How much depth can this chain deliver?

\begin{theorem}[Chain Persuasion Bound (7.11)]
\label{thm:chain-persuasion}
$\mathrm{marginalDepth}(r, r_1\compose s) \le \depth(r_1) + \depth(s).$
\end{theorem}

A teacher cannot transmit more depth than their own understanding plus the
source material's complexity.  Great teaching requires both deep knowledge
\emph{and} rich source texts.  This formalizes what every graduate student
knows intuitively: a mediocre advisor with a brilliant text, or a brilliant
advisor with a mediocre text, both produce bounded results.

\begin{theorem}[Marginal Depth of Self-Composition (7.12)]
\label{thm:marginal-self}
$\mathrm{marginalDepth}(r, r) \le \depth(r).$
\end{theorem}

Reflecting on your own ideas (self-composition) cannot add more than your
existing complexity.  Introspection has bounded returns.  This result
anticipates Chapter~\ref{ch:echo}'s analysis of echo chambers: if
self-reflection is bounded, then \emph{iterated} self-reflection is
bounded too.

\begin{theorem}[Marginal Depth Triangle Inequality (7.13)]
\label{thm:marginal-triangle}
$\mathrm{marginalDepth}(r, s_1\compose s_2) \le \depth(s_1) + \depth(s_2).$
\end{theorem}

The marginal impact of a compound message is bounded by the sum of its
parts' complexities.  No synergy, no multiplication---just addition.  This
is the triangle inequality for marginal depth, and it gives the persuasion
space a quasi-metric structure.


% ================================================================
\section{Propaganda Effectiveness Bounds}
\label{sec:pers-propaganda}

We now turn to the specific problem of propaganda: what happens when the
same message is repeated many times?

\begin{theorem}[Bounded Propaganda Effectiveness (7.14)]
\label{thm:bounded-propaganda}
After composing receiver $r$ with $n$ repetitions of signal $s$:
\[
  \depth\bigl(r\compose\compN{s}{n}\bigr)
  \;\le\;
  \depth(r) + n\cdot\depth(s).
\]
\end{theorem}

\begin{proof}[Proof sketch]
Combine the subadditivity of composition depth with the bound
$\depth(\compN{s}{n})\le n\cdot\depth(s)$ from the composeN depth theorem.
\end{proof}

Propaganda saturation---hearing the same message $n$ times---increases
depth at most linearly.  No amount of repetition can produce exponential
radicalization.  This is why ``deprogramming'' is possible: the damage is
bounded, and bounded damage can be reversed.%
\footnote{The bounded-propaganda theorem provides a formal counterargument
to the popular ``brainwashing'' hypothesis.  See Lifton,
\textit{Thought Reform and the Psychology of Totalism} (1961), for the
empirical evidence that even extreme indoctrination has bounded effects.}

\begin{theorem}[Resonant Propaganda (7.15)]
\label{thm:resonant-propaganda}
If $s_1\res s_2$, then $r\compose s_1 \res r\compose s_2$ for all~$r$.
\end{theorem}

Competing narratives that ``resonate'' with the same cultural themes produce
resonant (structurally similar) effects on receivers.  When Fox News and
MSNBC draw on the same patriotic imagery, they produce resonant---not
identical---reactions in their audiences.  The content differs; the
\emph{structure} of the interpretive effect is preserved.

\begin{theorem}[Double Propagandist Bound (7.16)]
\label{thm:double-propagandist}
$\mathrm{marginalDepth}(r, s_1\compose s_2) \le
\mathrm{persuasionBound}(s_1) + \mathrm{persuasionBound}(s_2).$
\end{theorem}

Allied propaganda campaigns---joint military-religious messaging,
corporate-governmental public relations---cannot exceed the sum of their
individual complexities.  Coalition propaganda doesn't synergize; it merely
adds.

\begin{theorem}[Atomic Signal Minimal Persuasion (7.17)]
\label{thm:atomic-persuasion}
If $s$ is atomic, then $\mathrm{persuasionBound}(s) \le 1$.
\end{theorem}

Single words, simple gestures, basic symbols---all have persuasion bound at
most~1.  To deeply transform a worldview, you need \emph{structured}
complexity, not atomic simplicity.  A thumbs-up cannot substitute for a
treatise.%
\footnote{This result formalizes the linguistic insight that single lexemes
have bounded semantic impact.  The power of language lies in
\emph{compositional} complexity, not lexical richness.}

\begin{theorem}[Iterated Propaganda Marginal Decay (7.18)]
\label{thm:iterated-propaganda}
$\mathrm{marginalDepth}(r, \compN{s}{n}) \le n\cdot\mathrm{persuasionBound}(s).$
\end{theorem}

\begin{proof}[Proof sketch]
By the composition depth bound and the composeN depth bound:
$\depth(r\compose\compN{s}{n})\le\depth(r)+n\cdot\depth(s)$,
so $\mathrm{marginalDepth}\le n\cdot\depth(s)$.
\end{proof}

The $k$-th repetition of propaganda is no more effective than the first---the
bound is uniform.  Each repetition has the same ceiling on marginal depth,
but the \emph{actual} marginal depth can only decrease in realistic models.
This formalizes the ``diminishing returns'' of propaganda.

\begin{lstlisting}[caption={Iterated propaganda bound in Lean~4.}]
theorem iterated_propaganda_bound (r s : I) (n : ℕ) :
    marginalDepth r (composeN s n) <=
      n * persuasionBound s
\end{lstlisting}


% ================================================================
\section{Conclusion}
\label{sec:pers-conclusion}

This chapter established the fundamental limits of persuasion within IDT\@.
The central result---the marginal depth bound
(Theorem~\ref{thm:marginal-depth-bound})---states that no signal can add
more complexity than it contains.  The subadditivity theorem
(\ref{thm:persuasion-subadditivity}) rules out synergistic propaganda, and
the bounded propaganda theorem (\ref{thm:bounded-propaganda}) caps the
effect of repetition at linear growth.

Together, these results provide formal foundations for Gramsci's theory of
cultural hegemony (persuasion requires complex cultural production),
Foucault's power/knowledge (persuasive capacity scales with discursive
depth), and the empirical observation that totalitarian propaganda has
bounded effectiveness.

Chapter~\ref{ch:lossy} turns from the \emph{sender's} constraints to the
\emph{receiver's}: what happens when the reception process itself is lossy?

% =============================================================
%  Chapter 8 — Lossy Reception and Oral Tradition
%  Lean source: FormalAnthropology/IDT_Signal_Ch08.lean
% =============================================================
\chapter{Lossy Reception and Oral Tradition}
\label{ch:lossy}

\epigraph{%
  ``The story becomes at each reproduction more concise, more schematized
  \ldots\ Proper names drop out. Vivid details vanish.
  What remains is a general form.''}%
  {Frederic Bartlett, \textit{Remembering} (1932)}

% ----------------------------------------------------------------
\section{Introduction: The Thermodynamics of Storytelling}
\label{sec:lossy-intro}

In 1932, Frederic Bartlett conducted a series of serial reproduction
experiments that became foundational for cognitive anthropology.  He asked a
subject to read a Native American folk tale (``The War of the Ghosts''), then
reproduce it from memory; the next subject read only the reproduction, and
so on through a chain of ten or more retellings.  The results were dramatic:
details dropped out, unfamiliar concepts were rationalized into familiar
categories, and the story simplified with each retelling until it converged
on a stable, culturally assimilated ``skeleton.''\footnote{Bartlett, F.~C.,
\textit{Remembering: A Study in Experimental and Social Psychology}
(Cambridge, 1932), ch.~5.  The experiment has been replicated hundreds of
times with consistent results; see Mesoudi \& Whiten (2008) for a
meta-analysis.}

Oral tradition---the transmission of myths, genealogies, ritual
instructions, and folk knowledge from person to person without writing---is
the primordial form of cultural transmission.  Before the invention of
writing (c.~3200 BCE in Sumer), \emph{all} culture was oral.  Even today,
the vast majority of human cultural knowledge is transmitted orally: recipes,
gossip, jokes, parenting advice, religious instruction.

This chapter formalizes the defining feature of oral transmission:
\emph{lossiness}.  A receiver who processes a signal through an oral
channel---through memory, paraphrase, reinterpretation---can never produce
output more complex than the input.  Depth can only decrease or remain
constant; it cannot increase.  This is the ``second law of cultural
thermodynamics.''\footnote{The analogy to thermodynamic entropy is
imperfect but illuminating.  In thermodynamics, the second law states that
entropy never decreases in an isolated system.  Here, the analog is that
\emph{complexity} (depth) never increases through lossy reception.  The
``entropy'' of cultural transmission is the \emph{loss} of depth.}

We define three levels of receiver: a generic \emph{Receiver} (preserving
void and distributing over composition), a \emph{LossyReceiver} (depth
non-increasing), and a \emph{StrictlyLossy} receiver (depth strictly
decreasing for complex ideas).  We prove that lossy reception composes,
iterates, and ultimately converges: only shallow ideas survive the
``telephone game.''


% ================================================================
\section{Core Structures}
\label{sec:lossy-defs}

\begin{definition}[Receiver]
\label{def:receiver}
A \emph{receiver} is an endomorphism $\varphi:\IS\to\IS$ satisfying:
\begin{enumerate}
  \item \emph{Void preservation:} $\varphi(\void) = \void$.
  \item \emph{Distributivity:} $\varphi(a\compose b) = \varphi(a)\compose\varphi(b)$.
\end{enumerate}
\end{definition}

A receiver is a ``cultural filter''---a formalized version of how a
particular mind processes incoming signals.  The two axioms are natural:
silence in yields silence out, and hearing a compound message is the same
as hearing the parts and composing the results.

\begin{lstlisting}[caption={The Receiver structure in Lean~4.}]
structure Receiver (I : Type*) [IdeaticSpace I] where
  process : I -> I
  process_void :
    process IdeaticSpace.void = IdeaticSpace.void
  process_compose : forall (a b : I),
    process (IdeaticSpace.compose a b) =
    IdeaticSpace.compose (process a) (process b)
\end{lstlisting}

\begin{definition}[Lossy Receiver]
\label{def:lossy-receiver}
A \emph{lossy receiver} is a receiver $\varphi$ such that
\[
  \depth(\varphi(a)) \;\le\; \depth(a) \quad\text{for all } a\in\IS.
\]
Processing never increases depth.
\end{definition}

This is the defining property of oral transmission: you can never come out
of a retelling with \emph{more} complexity than you went in with.  Written
transmission is different (you can re-read, annotate, cross-reference), but
oral transmission is inherently lossy.

\begin{definition}[Strictly Lossy Receiver]
\label{def:strictly-lossy}
A \emph{strictly lossy receiver} is a lossy receiver $\varphi$ such that
\[
  \depth(a) > 1 \;\implies\;
  \depth(\varphi(a)) < \depth(a).
\]
For complex ideas ($\depth > 1$), depth \emph{strictly} decreases.  Ideas of
depth $\le 1$ may be preserved.
\end{definition}

This models the ``telephone game'': every retelling loses something, unless
the message is already maximally simple.  Only atomic ideas survive
perfect-fidelity transmission through a strictly lossy channel.

\begin{lstlisting}[caption={Strictly lossy receiver in Lean~4.}]
structure StrictlyLossy (I : Type*) [IdeaticSpace I]
    extends LossyReceiver I where
  strict_decay : forall (a : I),
    IdeaticSpace.depth a > 1 ->
    IdeaticSpace.depth
      (toLossyReceiver.toReceiver.process a) <
      IdeaticSpace.depth a
\end{lstlisting}


% ================================================================
\section{Basic Properties of Lossy Reception}
\label{sec:lossy-basic}

\begin{theorem}[Lossy Void Preservation (8.1)]
\label{thm:lossy-void-preservation}
All receivers preserve silence: $\varphi(\void) = \void$.
\end{theorem}

No transmission channel creates content from nothing.  If you say nothing,
the listener hears nothing.  This is the zeroth law of communication---prior
to any thermodynamic considerations, it establishes the baseline: the null
signal is universally transparent.

\begin{theorem}[Lossy Depth Bound (8.2)]
\label{thm:lossy-depth-bound}
After lossy processing: $\depth(\varphi(a)) \le \depth(a)$.
\end{theorem}

An oral retelling is never more complex than the original.  This is the
second law of cultural thermodynamics: complexity never increases through
oral transmission.  The theorem is definitional---it \emph{is} the defining
property of a lossy receiver---but stating it as a theorem emphasizes its
central role.%
\footnote{Contrast with \emph{written} transmission, which can \emph{increase}
depth through annotation, commentary, and cross-reference.  The Talmud is
deeper than the Torah; the scholastic summa is deeper than any single
patristic text.  Writing breaks the second law of oral thermodynamics.}

\begin{theorem}[Lossy Composition Bound (8.3)]
\label{thm:lossy-compose-bound}
Processing a composed idea through a lossy receiver:
\[
  \depth(\varphi(a\compose b)) \;\le\; \depth(a) + \depth(b).
\]
\end{theorem}

\begin{proof}[Proof sketch]
By lossiness, $\depth(\varphi(a\compose b))\le\depth(a\compose b)$.  By
subadditivity of composition depth, $\depth(a\compose b)\le\depth(a)+\depth(b)$.
Chaining the inequalities yields the result.
\end{proof}

A story about two events, retold through oral tradition, has depth at most
the sum of the two events' original depths---even before accounting for the
receiver's lossiness.

\begin{theorem}[Lossy Receiver Preserves Void Depth (8.4)]
\label{thm:lossy-void-depth}
$\depth(\varphi(\void)) = 0.$
\end{theorem}

Silence, when transmitted lossily, remains at depth zero.  Combined with
Theorem~\ref{thm:lossy-void-preservation}, this says not just that void
maps to void, but that the image has zero complexity---a stronger
quantitative statement.


% ================================================================
\section{Iterated Lossy Reception}
\label{sec:lossy-iterated}

What happens when a signal passes through a \emph{chain} of lossy receivers?
Bartlett's serial reproduction experiments provide the empirical motivation;
this section provides the formal theory.

\begin{definition}[Iterated Process]
\label{def:iterate-process}
The $n$-fold iteration of a receiver's process function:
\[
  \varphi^{(0)}(a) = a, \qquad
  \varphi^{(n+1)}(a) = \varphi(\varphi^{(n)}(a)).
\]
\end{definition}

\begin{lstlisting}[caption={Iterated process in Lean~4.}]
def iterateProcess (R : Receiver I) : ℕ -> I -> I
  | 0 => fun a => a
  | n + 1 => fun a => R.process (iterateProcess R n a)
\end{lstlisting}

\begin{theorem}[Iterated Lossy Depth Bound (8.5)]
\label{thm:iterated-lossy-depth}
After $n$ applications of a lossy receiver:
$\depth(\varphi^{(n)}(a)) \le \depth(a)$ for all $n$.
\end{theorem}

\begin{proof}[Proof sketch]
By induction on $n$.  The base case ($n=0$) is trivial.  For the inductive
step, apply the lossiness bound to the result of the $(n)$-th iteration,
whose depth is $\le\depth(a)$ by the inductive hypothesis.
\end{proof}

Oral tradition through a chain of $n$ retellings can only lose depth, never
gain it.  Whether the story passes through 3~or 300~people, it can only get
simpler.  This is the formal backbone of Bartlett's experimental finding.

\begin{theorem}[Iterated Void Stability (8.6)]
\label{thm:iterated-void-stable}
$\varphi^{(n)}(\void) = \void$ for all $n$.
\end{theorem}

\begin{proof}[Proof sketch]
Induction on $n$: the base case is identity, and each step applies
$\varphi(\void)=\void$.
\end{proof}

A chain of listeners hearing nothing produces nothing, no matter how long
the chain.  Silence is the unique \emph{absorbing state} of every receiver.

\begin{theorem}[Lossy Reception of Compositions (8.7)]
\label{thm:lossy-distribute}
$\varphi(a\compose b) = \varphi(a)\compose\varphi(b).$
\end{theorem}

Cultural filtering is ``linear'' in the algebraic sense: a receiver's
reaction to a compound stimulus is the compound of their reactions to the
parts.  This is the distributivity axiom of the Receiver structure, stated
as a theorem for emphasis.

In anthropological terms, it means that a listener's interpretation of a
two-part story can be decomposed into their interpretation of part~1 composed
with their interpretation of part~2.  There is no holistic ``Gestalt'' effect
in lossy reception---only part-by-part processing.%
\footnote{This is a strong structural assumption.  Real human cognition
arguably does exhibit Gestalt effects.  The distributivity axiom is a
simplification that enables clean formal results at the cost of ignoring
non-decomposable cognitive processes.}


% ================================================================
\section{Strict Lossiness and Convergence}
\label{sec:lossy-strict}

\begin{theorem}[Strict Decay Bound (8.8)]
\label{thm:strict-decay}
For a strictly lossy receiver and any idea with $\depth(a)>1$:
\[
  \depth(\varphi(a)) < \depth(a).
\]
\end{theorem}

In ``telephone game'' transmission, every complex message becomes strictly
simpler.  This is not a bug---it's a feature.  Cultural evolution works by
simplification.  The myths that survive thousands of years of oral
transmission are the ones that have been \emph{simplified} to the point
where no further simplification occurs.

\begin{theorem}[Shallow Ideas Survive Lossy Reception (8.9)]
\label{thm:shallow-survives}
If $\depth(a)\le 1$, then $\depth(\varphi(a))\le 1$.
\end{theorem}

\begin{proof}[Proof sketch]
By lossiness: $\depth(\varphi(a))\le\depth(a)\le 1$.
\end{proof}

Simple ideas---basic emotions, universal gestures, atomic symbols---survive
any amount of lossy transmission.  This is why certain cultural universals
persist: they're too simple to be degraded further.  The smile, the frown,
the pointing gesture, the head-nod---all have depth $\le 1$ and are thus
impervious to oral degradation.%
\footnote{Ekman's (1971) cross-cultural studies of facial expression provide
empirical support: the six ``basic emotions'' are recognized across cultures,
arguably because they are too simple to be distorted by oral transmission.}

\begin{theorem}[Zero-Depth Ideas Survive Lossy Reception (8.10)]
\label{thm:zero-depth-survives}
If $\depth(a) = 0$, then $\depth(\varphi(a)) = 0$.
\end{theorem}

Void-like ideas (pure silence, empty gestures) are immune to lossy
transmission.  They are fixed points of every lossy receiver.


% ================================================================
\section{Composition of Receivers}
\label{sec:lossy-composition}

\begin{definition}[Composed Receivers]
\label{def:compose-receivers}
Given receivers $\varphi_1,\varphi_2$, their composition is
$(\varphi_2\circ\varphi_1)(a) = \varphi_2(\varphi_1(a))$---the result of
sequential reception.
\end{definition}

\begin{lstlisting}[caption={Receiver composition in Lean~4.}]
def composeReceivers (R_1 R_2 : Receiver I) : Receiver I where
  process := fun a => R_2.process (R_1.process a)
  process_void := by
    rw [R_1.process_void, R_2.process_void]
  process_compose := by
    intro a b
    rw [R_1.process_compose, R_2.process_compose]
\end{lstlisting}

\begin{theorem}[Composition of Lossy Receivers is Lossy (8.11)]
\label{thm:compose-lossy}
If $\varphi_1$ and $\varphi_2$ are lossy, then
$\varphi_2\circ\varphi_1$ is lossy:
\[
  \depth((\varphi_2\circ\varphi_1)(a)) \;\le\; \depth(a)
  \quad\text{for all } a.
\]
\end{theorem}

\begin{proof}[Proof sketch]
$\depth(\varphi_2(\varphi_1(a)))\le\depth(\varphi_1(a))\le\depth(a)$
by applying lossiness twice.
\end{proof}

A chain of imperfect oral transmitters is itself an imperfect transmitter.
Lossiness is closed under composition---there's no way to ``recover'' lost
depth by adding more lossy steps.  This is the impossibility of error
correction in purely oral tradition.  Unlike digital communication (where
redundancy enables error correction), oral channels have no mechanism for
restoring lost depth.

\begin{theorem}[Double Lossy is More Lossy (8.12)]
\label{thm:double-lossy}
$\depth((\varphi_2\circ\varphi_1)(a)) \le \depth(\varphi_1(a)).$
\end{theorem}

Two generations of oral transmission is worse (or equal) than one.  The
second lossy step can only further reduce what the first step already
degraded.  This is the formal version of the folk wisdom that stories ``get
worse with each retelling.''


% ================================================================
\section{Resonance and Lossy Reception}
\label{sec:lossy-resonance}

\begin{theorem}[Receiver Preserves Self-Resonance (8.13)]
\label{thm:receiver-self-resonance}
$\varphi(a) \res \varphi(a)$ for all $a$.
\end{theorem}

Even lossy transmission produces internally consistent output.  A garbled
retelling is garbled, but it still ``makes sense to itself.''  This is
a consequence of the reflexivity of resonance, but it carries anthropological
weight: oral traditions, however degraded, are never \emph{incoherent}.
They simplify, regularize, and schematize, but they always remain
self-consistent.

\begin{theorem}[Lossy Compound Self-Resonance (8.14)]
\label{thm:lossy-compound-resonance}
$\varphi(a\compose b) \res \varphi(a\compose b).$
\end{theorem}

Even through lossy reception, the receiver's output from a compound signal
is internally coherent.  The theorem is trivial (self-resonance is always
true), but it anchors the stronger results of Chapter~\ref{ch:multireceiver},
where we study how \emph{different} receivers' outputs relate to each other.

\begin{theorem}[Void Iterate Always Resonates (8.15)]
\label{thm:void-iterate-resonates}
$\varphi^{(n)}(\void) \res \varphi^{(n)}(\void)$ for all $n$.
\end{theorem}

Chains of silence produce silence, which trivially resonates with itself.
Combined with Theorem~\ref{thm:iterated-void-stable}, this says that
void is not only a fixed point of iteration but a self-resonant one.


% ================================================================
\section{Advanced Properties}
\label{sec:lossy-advanced}

\begin{theorem}[Lossy Iterated Composition Bound (8.16)]
\label{thm:lossy-composeN}
$\depth(\varphi(\compN{a}{n})) \le n\cdot\depth(a).$
\end{theorem}

\begin{proof}[Proof sketch]
By lossiness, $\depth(\varphi(\compN{a}{n}))\le\depth(\compN{a}{n})$.
By the composeN depth bound, $\depth(\compN{a}{n})\le n\cdot\depth(a)$.
\end{proof}

A repeated message (ritual chant, liturgical refrain) loses depth through
lossy reception, bounded by the repetition count times original depth.  Even
before lossiness takes its toll, the starting depth is only linearly bounded
in repetition count.

\begin{theorem}[Receiver Composition is Associative on Output (8.17)]
\label{thm:receiver-assoc}
$\varphi_3(\varphi_2(\varphi_1(a)))
= (\varphi_3\circ\varphi_2\circ\varphi_1)(a).$
\end{theorem}

A three-generation transmission chain (grandparent $\to$ parent $\to$ child)
produces the same result regardless of how we ``bracket'' the generations.
This is the associativity of function composition, applied to the receiver
setting.  It ensures that our model of multi-generational oral tradition is
well-defined: the order of composition doesn't depend on arbitrary grouping.

\begin{theorem}[Depth Monotonicity Through Chains (8.18)]
\label{thm:chain-depth-monotone}
For $m\le n$:
$\depth(\varphi^{(n)}(a)) \le \depth(\varphi^{(m)}(a)).$
\end{theorem}

\begin{proof}[Proof sketch]
By induction on $n$.  If $m=n$, the result is trivial.  If $m<n$, apply
lossiness at step $n$ to reduce to the case $n-1$, then apply the
inductive hypothesis.
\end{proof}

At every point in a chain of oral transmission, the story is no more complex
than it was at any earlier point.  There is no ``regeneration'' through
purely lossy channels.  The depth sequence $\depth(\varphi^{(0)}(a)),
\depth(\varphi^{(1)}(a)), \depth(\varphi^{(2)}(a)), \ldots$ is monotonically
non-increasing.

This is the formal version of a deep anthropological observation: myths
don't get \emph{worse}; they get \emph{stable}.  Oral tradition doesn't
``degrade''---it \emph{converges} to a fixed point determined by the
transmission medium's lossiness.  The stories that survive are the ones
simple enough to be transmitted without further loss.

\begin{lstlisting}[caption={Chain depth monotonicity in Lean~4.}]
theorem chain_depth_monotone
    (R : LossyReceiver I) (a : I) (m n : ℕ)
    (hmn : m <= n) :
    IdeaticSpace.depth
      (iterateProcess R.toReceiver n a) <=
    IdeaticSpace.depth
      (iterateProcess R.toReceiver m a)
\end{lstlisting}


% ================================================================
\section{Conclusion}
\label{sec:lossy-conclusion}

This chapter established the formal theory of lossy reception, providing
mathematical foundations for Bartlett's serial reproduction experiments and
the broader study of oral tradition.

The central results are:
\begin{enumerate}
  \item Lossy reception is closed under composition
    (Theorem~\ref{thm:compose-lossy}): chains of oral transmitters are
    themselves oral transmitters.
  \item Depth is monotonically non-increasing through chains
    (Theorem~\ref{thm:chain-depth-monotone}): stories only get simpler.
  \item Shallow ideas survive (Theorem~\ref{thm:shallow-survives}):
    cultural universals persist because they are too simple to degrade.
\end{enumerate}

The theory explains why oral traditions converge: not because storytellers
are ``bad'' at their job, but because lossiness is a structural feature of
the oral channel.  Written traditions escape this constraint by providing
external memory, but oral traditions are bound by the second law of cultural
thermodynamics.

Chapter~\ref{ch:echo} takes this analysis inward: what happens when a
community applies lossy reception \emph{to its own output}, creating an echo
chamber?

% =============================================================
%  Chapter 9 — Echo Chambers and Iterated Self-Interpretation
%  Lean source: FormalAnthropology/IDT_Signal_Ch09.lean
% =============================================================
\chapter{Echo Chambers and Iterated Self-Interpretation}
\label{ch:echo}

\epigraph{%
  ``Doublethink means the power of holding two contradictory beliefs in
  one's mind simultaneously, and accepting both of them.''}%
  {George Orwell, \textit{Nineteen Eighty-Four} (1949)}

% ----------------------------------------------------------------
\section{Introduction: The Paradox of In-Group Radicalization}
\label{sec:echo-intro}

Popular accounts of echo chambers emphasize \emph{amplification}: a
community that talks only to itself is said to ``radicalize,'' producing
ever-more-extreme versions of its founding ideology.  Social media critics
warn of ``filter bubbles'' (Pariser, 2011) in which algorithmic
recommendation creates spirals of increasing extremism.\footnote{Pariser, E.,
\textit{The Filter Bubble: What the Internet Is Hiding from You} (Penguin,
2011).  For a critical assessment, see Bruns, A., \textit{Are Filter Bubbles
Real?} (Polity, 2019).}

The formal analysis tells a different story.

When a mind with background $r$ repeatedly re-interprets a signal $s$
through its own cultural filter, the result is not an explosion of
complexity but a \emph{collapse}.  Each round of self-interpretation
composes $r$ with the previous result, and the depth of the output is
bounded by $n\cdot\depth(r)+\depth(s)$---\emph{linear} in the number of
rounds.  No exponential radicalization is possible.

The counterintuitive insight of IDT's echo chamber theory is that closed
communities don't \emph{complexify} their ideas---they \emph{simplify}
them.  Radicalization is not the acquisition of deeper ideology; it is the
\emph{erosion of nuance} through iterated self-composition.  The extremist
does not have a more complex worldview than the moderate; they have a
\emph{less} complex one, in which fine distinctions have been compositioned
away by repeated application of the same filter.

This chapter formalizes self-iteration, proves the linear depth bound,
analyzes the role of void backgrounds, studies resonance preservation
through echo, and identifies the fixed-point structure of iterated
self-interpretation.


% ================================================================
\section{Core Definitions}
\label{sec:echo-defs}

\begin{definition}[Self-Iteration (Echo)]
\label{def:self-iterate}
Given background $r\in\IS$ and signal $s\in\IS$, the \emph{$n$-th echo}
is defined recursively:
\[
  \mathrm{selfIterate}(r,s,0) = s,
  \qquad
  \mathrm{selfIterate}(r,s,n\!+\!1) = r\compose\mathrm{selfIterate}(r,s,n).
\]
\end{definition}

Each step represents one ``round'' of in-group discussion: take the current
consensus, filter it through the group's shared worldview $r$, and obtain
the next consensus.

\begin{definition}[Echo Depth]
\label{def:echo-depth}
The \emph{echo depth} at step $n$ is
\[
  \mathrm{echoDepth}(r,s,n) \;=\; \depth(\mathrm{selfIterate}(r,s,n)).
\]
\end{definition}

\begin{lstlisting}[caption={Self-iteration in Lean~4.}]
def selfIterate (r s : I) : ℕ -> I
  | 0 => s
  | n + 1 => IdeaticSpace.compose r (selfIterate r s n)

def echoDepth (r s : I) (n : ℕ) : ℕ :=
  IdeaticSpace.depth (selfIterate r s n)
\end{lstlisting}


% ================================================================
\section{Basic Echo Properties}
\label{sec:echo-basic}

\begin{theorem}[Echo Base Case (9.1)]
\label{thm:echo-base}
$\mathrm{selfIterate}(r, s, 0) = s.$
\end{theorem}

At step~0, the echo is the raw signal, unfiltered.  This is the ``first
hearing''---before enculturation transforms it.  In anthropological terms,
it is the moment of initial contact: the ethnographer's first impression
before theoretical frameworks reshape perception.

\begin{theorem}[Echo Step (9.2)]
\label{thm:echo-step}
$\mathrm{selfIterate}(r, s, n\!+\!1) = r\compose\mathrm{selfIterate}(r,s,n).$
\end{theorem}

Each round of in-group discussion is formally identical: take the current
consensus, filter it through the group's shared worldview, get the next
consensus.  The process is \emph{memoryless}---only the current state
matters, not the history.  This is a Markov property, and it makes the echo
chamber formally analogous to a deterministic dynamical system on $\IS$.

\begin{theorem}[Echo Depth Bound (9.3)]
\label{thm:echo-depth-bound}
For all $n$:
\[
  \mathrm{echoDepth}(r,s,n) \;\le\; n\cdot\depth(r) + \depth(s).
\]
\end{theorem}

\begin{proof}[Proof sketch]
By induction on $n$.  The base case gives $\depth(s)\le 0\cdot\depth(r)+\depth(s)$.
For the inductive step, $\depth(r\compose\mathrm{selfIterate}(r,s,n))
\le\depth(r)+\depth(\mathrm{selfIterate}(r,s,n))
\le\depth(r)+n\cdot\depth(r)+\depth(s)
= (n\!+\!1)\cdot\depth(r)+\depth(s)$.
\end{proof}

This is the \textbf{central theorem of echo chamber theory}: complexity
grows at most \emph{linearly} in the number of rounds.  No exponential
radicalization is possible.  A community hearing the same message 1000~times
cannot produce more complexity than $1000\cdot\depth(r)+\depth(s)$.

The linear bound has immediate implications for the study of radicalization.
If echo chambers could produce exponential depth growth, then small
communities with minimal resources could generate arbitrarily complex
ideologies.  The linear bound rules this out: the complexity ceiling is
determined by the background's depth, and exceeding it requires either a
deeper background or external input.%
\footnote{This does not mean that radicalization is ``harmless.''  Even
linear depth growth can produce dangerous simplifications---see the
discussion of depth \emph{collapse} vs.\ depth \emph{growth} below.  The
point is that the mechanism of harm is simplification, not complexification.}


% ================================================================
\section{Void Background: The Open Mind}
\label{sec:echo-void}

\begin{theorem}[Void Background Preserves Signal (9.4)]
\label{thm:void-background}
If the background is void:
$\mathrm{selfIterate}(\void, s, n) = s$ for all $n$.
\end{theorem}

\begin{proof}[Proof sketch]
Induction on $n$.  Base case: $\mathrm{selfIterate}(\void,s,0)=s$.
Inductive step: $\void\compose\mathrm{selfIterate}(\void,s,n)
= \void\compose s = s$ by the left-identity axiom.
\end{proof}

A perfectly ``open mind''---one with no prior beliefs---cannot transform a
signal through echo.  You need \emph{content} to create an echo; an empty
room doesn't echo.  This is the formal impossibility of the echo chamber
for culturally empty agents.

The result has pedagogical implications: the ideal student (void background)
retains every signal perfectly, with no distortion from prior beliefs.  Of
course, no real student is a blank slate, which is why every teacher
encounters the frustration of students who ``hear'' the lecture through the
filter of their existing misconceptions.

\begin{theorem}[Void Background Echo Depth (9.5)]
\label{thm:void-background-depth}
$\mathrm{echoDepth}(\void, s, n) = \depth(s)$ for all $n$.
\end{theorem}

With void background, echo depth equals signal depth at every step---no
amplification, no degradation.  The signal is perfectly preserved through
arbitrary iterations.

\begin{theorem}[Echo with Void Signal --- Depth Bound (9.6)]
\label{thm:echo-void-signal}
When signal is void:
$\depth(\mathrm{selfIterate}(r, \void, n)) \le n\cdot\depth(r).$
\end{theorem}

An echo chamber with no input signal just repeats its own background---pure
ideological repetition without external stimulus.  The community talks to
itself, producing at most $n$ rounds worth of background complexity.  This
is the formal model of ideological ``navel-gazing'': a group with nothing
new to process, endlessly recycling its own assumptions.


% ================================================================
\section{Resonance in Echo Chambers}
\label{sec:echo-resonance}

\begin{theorem}[Echo Self-Resonance (9.7)]
\label{thm:echo-self-resonance}
$\mathrm{selfIterate}(r,s,n) \res \mathrm{selfIterate}(r,s,n)$ for all $n$.
\end{theorem}

Every stage of an echo chamber is internally consistent.  The community
always ``agrees with itself.''  This is trivially true (self-resonance is
reflexivity) but foundationally important: it means that the output of an
echo chamber, however simplified, is never \emph{incoherent}.  The danger
of the echo chamber is not incoherence but oversimplification.

\begin{theorem}[Resonant Backgrounds Produce Resonant Echoes (9.8)]
\label{thm:resonant-backgrounds}
If $r_1\res r_2$, then for all $n$:
$\mathrm{selfIterate}(r_1, s, n) \res \mathrm{selfIterate}(r_2, s, n).$
\end{theorem}

\begin{proof}[Proof sketch]
Induction on $n$.  Base case: both sides equal $s$, which self-resonates.
Inductive step: $r_1\compose\mathrm{selfIterate}(r_1,s,n)$ resonates with
$r_2\compose\mathrm{selfIterate}(r_2,s,n)$ by the resonance-compose
theorem, using $r_1\res r_2$ and the inductive hypothesis.
\end{proof}

Two similar communities processing the same signal produce similar results
at every stage.  Cultural cousins stay cultural cousins through the echo
process.  This is the formal basis for the observation that, e.g., American
and British media---despite separate echo dynamics---produce structurally
similar political discourses because their cultural backgrounds resonate.

\begin{theorem}[Resonant Signals Produce Resonant Echoes (9.9)]
\label{thm:resonant-signals}
If $s_1\res s_2$, then for all $n$:
$\mathrm{selfIterate}(r, s_1, n) \res \mathrm{selfIterate}(r, s_2, n).$
\end{theorem}

An echo chamber processing two structurally similar inputs produces
structurally similar outputs at every stage.  Two versions of the same news
story, fed into the same community, produce resonant echo trajectories.
The \emph{structural similarity} of the inputs is preserved through
arbitrary iterations.

\begin{lstlisting}[caption={Resonant backgrounds theorem in Lean~4.}]
theorem resonant_backgrounds_resonant_echoes
    {r_1 r_2 : I} (s : I)
    (hr : IdeaticSpace.resonates r_1 r_2) :
    forall (n : ℕ), IdeaticSpace.resonates
      (selfIterate r_1 s n) (selfIterate r_2 s n)
\end{lstlisting}


% ================================================================
\section{Composition and Iteration}
\label{sec:echo-composition}

\begin{theorem}[Two-Step Echo = Single Compound Step (9.10)]
\label{thm:two-step-echo}
$\mathrm{selfIterate}(r, s, 2) = r\compose(r\compose s).$
\end{theorem}

Two rounds of echo chamber discussion is equivalent to the background
filtering the background-filtered signal.  This is the formal version of
``groupthink'': the community's consensus at step~2 is just the community's
reaction to its own reaction.  The formula makes the recursive nature of
groupthink explicit: it is composition with the same element, applied twice.

\begin{theorem}[Three-Step Echo = Triple Composition (9.11)]
\label{thm:three-step-echo}
$\mathrm{selfIterate}(r, s, 3) = r\compose(r\compose(r\compose s)).$
\end{theorem}

Three rounds of self-interpretation.  The result is a three-layer
composition---formally identical to a three-generation oral transmission
chain where all generations share the same culture.  This connection to
Chapter~\ref{ch:lossy} is deep: an echo chamber is a \emph{homogeneous}
transmission chain, where every link is the same lossy receiver.

\begin{theorem}[Echo Depth at Step~1 (9.12)]
\label{thm:echo-step1}
$\mathrm{echoDepth}(r, s, 1) \le \depth(r) + \depth(s).$
\end{theorem}

The first echo adds at most the background's complexity.  This is the
``first filter''---the initial cultural processing of a raw signal.  In
media studies terms, it is the ``framing'' stage, where raw events are
filtered through editorial perspective.

\begin{theorem}[Echo Depth at Step~2 (9.13)]
\label{thm:echo-step2}
$\mathrm{echoDepth}(r, s, 2) \le 2\cdot\depth(r) + \depth(s).$
\end{theorem}

The second echo adds at most another round of background complexity.  Linear
growth, not exponential.  By step~2, the community has ``discussed the
discussion''---and the total complexity is at most twice the background
depth plus the original signal depth.


% ================================================================
\section{Fixed Points and Convergence}
\label{sec:echo-fixed}

\begin{theorem}[Void is Echo Fixed Point (9.14)]
\label{thm:void-fixed-point}
$\mathrm{selfIterate}(\void, \void, n) = \void$ for all $n$.
\end{theorem}

Complete silence in an empty echo chamber remains silence forever.  This is
the ``nothing'' fixed point---the unique state from which no echo can
emerge.  It is both a mathematical triviality and a philosophical limit
case: absolute emptiness is absolutely stable.

\begin{theorem}[Echo Preserves Void Signal Depth (9.15)]
\label{thm:echo-void-depth}
$\mathrm{echoDepth}(r, \void, n) \le n\cdot\depth(r).$
\end{theorem}

An echo chamber with no external input produces at most $n$~rounds' worth of
background repetition---pure self-reinforcement without new information.
This models the ``stale discourse'' phenomenon: a community that receives
no new input can only recirculate its existing depth, with the total bounded
by how many times the background has been applied.


% ================================================================
\section{Advanced Echo Properties}
\label{sec:echo-advanced}

\begin{theorem}[Echo is Left-Nested Composition (9.16)]
\label{thm:echo-left-nested}
$\mathrm{selfIterate}(r, s, n\!+\!1) = r\compose\mathrm{selfIterate}(r,s,n).$
\end{theorem}

This is definitional, but worth stating explicitly.  The echo process is
\emph{always} ``background first, then previous result.''  The background
is the dominant frame---it comes first in every composition.  This is why
echo chambers are \emph{conservative} in the formal sense: the background
always takes precedence over the incoming signal.

In Bourdieu's terms, the \emph{habitus} (background) structures every
perception (composition) so that the field's doxa is reproduced at each
step.  The echo chamber is the formal model of \emph{symbolic reproduction}
---the process by which dominant cultural categories perpetuate themselves
through iterated interpretation.

\begin{theorem}[Parallel Echo Chambers with Same Signal (9.17)]
\label{thm:parallel-echo}
If $r_1\res r_2$, then
$\mathrm{selfIterate}(r_1, s, 1) \res \mathrm{selfIterate}(r_2, s, 1).$
\end{theorem}

\begin{proof}[Proof sketch]
At step~1: $r_1\compose s \res r_2\compose s$ by resonance-compose,
using $r_1\res r_2$ and $s\res s$ (reflexivity).
\end{proof}

Two similar communities reacting to the same news story produce similar
``hot takes'' (step-1 interpretations).  This is the formal basis of
polling: structurally similar demographics respond similarly to the same
stimulus.  The theorem guarantees that resonant backgrounds produce resonant
first reactions---not identical ones, but structurally related ones.

\begin{theorem}[Echo Depth Monotonicity Relative to Signal Depth (9.18)]
\label{thm:echo-depth-lower}
$\depth(s) \le \mathrm{echoDepth}(\void, s, n)$ for all $n$.
\end{theorem}

\begin{proof}[Proof sketch]
By Theorem~\ref{thm:void-background},
$\mathrm{selfIterate}(\void,s,n)=s$, so
$\mathrm{echoDepth}(\void,s,n)=\depth(s)$.
\end{proof}

An ``open'' echo chamber (void background) never reduces the signal.
Complexity is preserved when there is no cultural filter.  This is the ideal
of academic freedom: an institution with no ideological filter preserves the
complexity of every input.  The theorem formalizes what universities aspire
to be (and rarely are): void backgrounds that transmit without distortion.%
\footnote{The reality, of course, is that universities have strong
backgrounds---disciplinary paradigms, methodological commitments,
institutional cultures---that filter incoming signals just as powerfully as
any other echo chamber.  Kuhn's \textit{Structure of Scientific Revolutions}
(1962) is precisely about the non-void background of scientific communities.}

\begin{lstlisting}[caption={Echo depth lower bound for void background.}]
theorem echo_depth_lower_void_background
    (s : I) (n : ℕ) :
    IdeaticSpace.depth s <=
      echoDepth (IdeaticSpace.void : I) s n
\end{lstlisting}


% ================================================================
\section{Conclusion}
\label{sec:echo-conclusion}

This chapter overturned the popular narrative about echo chambers.  The
formal analysis shows that iterated self-interpretation produces not
\emph{amplification} but \emph{bounded linear growth} of complexity
(Theorem~\ref{thm:echo-depth-bound}).  Echo chambers are depth
\emph{sinks}, not depth \emph{sources}.

Three results stand out:
\begin{enumerate}
  \item \emph{Linear depth bound} (Theorem~\ref{thm:echo-depth-bound}):
    no exponential radicalization through echo.
  \item \emph{Resonance preservation}
    (Theorems~\ref{thm:resonant-backgrounds}--\ref{thm:resonant-signals}):
    similar communities produce similar echoes.
  \item \emph{Void preservation} (Theorem~\ref{thm:void-background}):
    open minds don't echo.
\end{enumerate}

The analysis suggests that the real danger of echo chambers is not
complexity explosion but \emph{nuance erosion}: the repeated application of
a fixed filter smooths away fine distinctions, leaving only the coarse
structure of the background.  Radicalization, in IDT's framework, is a loss
of depth, not a gain.

Chapter~\ref{ch:multireceiver} moves from the single-receiver echo chamber
to the multi-receiver setting: what happens when the same signal is sent to
many different minds simultaneously?

% =============================================================
%  Chapter 10 — Multi-Receiver Signalling and Public Ritual
%  Lean source: FormalAnthropology/IDT_Signal_Ch10.lean
% =============================================================
\chapter{Multi-Receiver Signalling and Public Ritual}
\label{ch:multireceiver}

\epigraph{%
  ``The culture of a people is an ensemble of texts, themselves ensembles,
  which the anthropologist strains to read over the shoulders of those to
  whom they properly belong.''}%
  {Clifford Geertz, \textit{The Interpretation of Cultures} (1973)}

% ----------------------------------------------------------------
\section{Introduction: One Ritual, Many Minds}
\label{sec:mr-intro}

The defining feature of public ritual is that the \emph{same} signal is
sent to \emph{multiple} receivers with \emph{different} cultural
backgrounds.  A Catholic mass is witnessed simultaneously by the devout
grandmother who has attended every Sunday for sixty years, the skeptical
teenager dragged along by family obligation, and the visiting anthropologist
scribbling field notes.  Each receives the same sensory input; each produces
a radically different interpretation.

Clifford Geertz captured this multiplicity in his concept of ``thick
description'': the same cockfight in Bali means ``gambling'' to one
observer, ``status contest'' to another, ``art'' to a third, and ``data''
to the ethnographer.\footnote{Geertz, C., ``Deep Play: Notes on the
Balinese Cockfight,'' in \textit{The Interpretation of Cultures} (Basic
Books, 1973), pp.~412--453.}  Victor Turner extended the insight to ritual
proper: the same ceremony of \emph{Nkang'a} among the Ndembu produces
different (but structurally related) experiences in initiands, elders, and
spectators.\footnote{Turner, V., \textit{The Forest of Symbols: Aspects of
Ndembu Ritual} (Cornell, 1967).}

The previous chapters studied signalling to a \emph{single} receiver.  This
chapter generalizes to the multi-receiver case.  We define an
\emph{audience}---a list of repertoires, each representing a different
mind---and study what happens when a single signal is broadcast to the
entire audience simultaneously.

The central anthropological insight is that public ritual \emph{works} not
because everyone gets the same meaning, but because everyone gets
\emph{resonant} meanings.  The grandmother, the teenager, and the
anthropologist all experience something structurally related---though not
identical---because resonance is preserved under composition with a common
signal.  This is the formal foundation of Durkheim's \emph{conscience
collective}: shared ritual produces shared structure, not shared content.


% ================================================================
\section{Core Definitions}
\label{sec:mr-defs}

\begin{definition}[Audience]
\label{def:audience}
An \emph{audience} is a list of repertoires:
\[
  \mathrm{Audience}(\IS) \;=\; \mathrm{List}(\mathrm{Repertoire}(\IS))
  \;=\; \mathrm{List}(\mathrm{List}\;\IS).
\]
Each element represents one mind's stock of ideas.  The audience is
heterogeneous: different minds may have different numbers of ideas and
different depths.
\end{definition}

\begin{definition}[Audience Outcome]
\label{def:audience-outcome}
The \emph{audience outcome} of sending signal $s$ to audience
$A = [R_1,\ldots,R_m]$ is the list of interpretation lists:
\[
  \mathrm{audienceOutcome}(A, s) \;=\;
  [\mathrm{interpret}(R_1,s),\;\ldots,\;\mathrm{interpret}(R_m,s)].
\]
\end{definition}

\begin{lstlisting}[caption={Audience and audience outcome in Lean~4.}]
abbrev Audience (I : Type*) := List (Repertoire I)

def audienceOutcome
    (audience : Audience I) (s : I) : List (List I) :=
  audience.map (fun rep => interpret rep s)
\end{lstlisting}

\begin{definition}[Total Audience Depth]
\label{def:total-audience-depth}
The \emph{total audience depth} of an outcome is the sum of all depth sums
across all minds:
\[
  \mathrm{totalAudienceDepth}(\mathit{outcomes}) \;=\;
  \sum_{i}\;\mathrm{depthSum}(\mathit{outcomes}_i).
\]
\end{definition}


% ================================================================
\section{Structural Invariants}
\label{sec:mr-structure}

The first group of theorems establishes that the multi-receiver operation
preserves the fundamental structural properties of signalling.

\begin{theorem}[Audience Size Preservation (10.1)]
\label{thm:audience-size}
$|\mathrm{audienceOutcome}(A, s)| = |A|.$
\end{theorem}

\begin{proof}[Proof sketch]
$\mathrm{audienceOutcome}$ is a \texttt{List.map}, which preserves length.
\end{proof}

A public ritual has exactly as many ``experiences'' as there are
participants.  No ghost observers, no merged consciousnesses---each mind
is counted exactly once.  This is a non-trivial modeling choice: it rules
out ``collective consciousness'' as a separate entity above and beyond the
individual minds.%
\footnote{Durkheim's \emph{conscience collective} is often misread as
positing a supra-individual mind.  In IDT, collective consciousness is
emergent: it is the structural similarity (resonance) among individual
outcomes, not a separate entity.}

\begin{theorem}[Per-Receiver Interpretation Count (10.2)]
\label{thm:per-receiver-count}
$|\mathrm{interpret}(R, s)| = |R|.$
\end{theorem}

Each mind produces exactly as many interpretations as it has ideas in its
repertoire.  A five-category mind produces five interpretations; a
fifty-category mind produces fifty.  The signal doesn't create or destroy
ideas---it transforms each one.

\begin{theorem}[Same Signal Different Minds --- Length Preserved (10.3)]
\label{thm:same-signal-different}
For any audience member $R\in A$:
$|\mathrm{interpret}(R, s)| = |R|.$
\end{theorem}

The ``same ceremony, different meanings'' principle: each participant brings
their own framework and gets their own number of meanings, determined by
\emph{their} background, not the signal.  A child at a symphony may have
three categories (``loud,'' ``soft,'' ``boring''); a musicologist may have
three hundred.  The symphony produces three interpretations for the child
and three hundred for the musicologist.


% ================================================================
\section{Void Signal: Universal Transparency}
\label{sec:mr-void}

\begin{theorem}[Void Signal Transparent For Single Receiver (10.4)]
\label{thm:void-transparent-single}
$\mathrm{interpret}(R, \void) = R.$
\end{theorem}

\begin{proof}[Proof sketch]
By induction on $R$.  For each idea $r$, $r\compose\void = r$ by the
right-identity axiom.
\end{proof}

Silence leaves every mind unchanged.  The minute of silence, the Quaker
meeting, the Buddhist \emph{noble silence}---all exploit this: by saying
nothing, every participant's existing worldview remains intact.  Silence
is the only signal that is universally transparent.

\begin{theorem}[Void Signal Transparent For All Audiences (10.5)]
\label{thm:void-transparent-audience}
$\mathrm{audienceOutcome}(A, \void) = A.$
\end{theorem}

\begin{proof}[Proof sketch]
By induction on the audience list.  Each repertoire maps to itself
by Theorem~\ref{thm:void-transparent-single}.
\end{proof}

Universal transparency of silence is the formal foundation of the ``null
ritual''---a ceremony that changes nothing.  Every culture has such
ceremonies: they serve social functions (marking solidarity, observing
convention) without altering anyone's ideas.  The formal result shows that
void is the \emph{unique} signal with this property (assuming the repertoire
is non-trivially structured).

\begin{lstlisting}[caption={Void transparency for audiences in Lean~4.}]
theorem void_transparent_audience (audience : Audience I) :
    audienceOutcome audience (IdeaticSpace.void : I)
      = audience
\end{lstlisting}


% ================================================================
\section{Depth Bounds Across Audiences}
\label{sec:mr-depth}

\begin{theorem}[Single Receiver Depth Bound (10.6)]
\label{thm:single-receiver-depth}
\[
  \mathrm{depthSum}(\mathrm{interpret}(R, s))
  \;\le\;
  \mathrm{depthSum}(R) + |R|\cdot\depth(s).
\]
\end{theorem}

\begin{proof}[Proof sketch]
Induction on $R$.  At each step, $\depth(r\compose s)\le\depth(r)+\depth(s)$
by subadditivity.  Summing over all $|R|$ elements gives the bound.
\end{proof}

Each audience member's ``persuasion exposure'' is linearly bounded.  This
extends the persuasion bounds of Chapter~\ref{ch:persuasion} to the
per-receiver setting: even in a public ritual, each participant's total
interpretive transformation is bounded by their own cultural depth plus the
ceremony's complexity.

\begin{theorem}[Audience Depth Bound --- Per Member (10.7)]
\label{thm:audience-member-depth}
For each $R\in A$:
$\mathrm{depthSum}(\mathrm{interpret}(R,s)) \le \mathrm{depthSum}(R)+|R|\cdot\depth(s).$
\end{theorem}

A PhD and a child at the same concert are both bounded, but differently.
The PhD's larger repertoire means a larger absolute bound, but the \emph{per-idea}
bound ($\depth(r)+\depth(s)$ for each schema) is the same.  The signal's
persuasive impact is democratic in structure but heterogeneous in magnitude.

\begin{theorem}[Empty Audience Has Empty Outcome (10.8)]
\label{thm:empty-audience}
$\mathrm{audienceOutcome}([], s) = [].$
\end{theorem}

A ritual with no participants produces no effects.  Performance requires an
audience.  This is the formal version of the old question: if a tree falls
in a forest and no one is there to hear it, does it make a sound?  In IDT:
if a signal is sent to an empty audience, no interpretations are produced.

\begin{theorem}[Singleton Audience Outcome (10.9)]
\label{thm:singleton-audience}
$\mathrm{audienceOutcome}([R], s) = [\mathrm{interpret}(R,s)].$
\end{theorem}

An audience of one produces exactly one interpretation list.  This is the
formal model of individual meditation or private prayer: a single mind
processing a signal in isolation.  It reduces the multi-receiver framework
to the single-receiver case of Chapters~\ref{ch:equilibrium}--\ref{ch:echo}.


% ================================================================
\section{Resonance Across Audiences}
\label{sec:mr-resonance}

This section contains the deepest results of the chapter.  Resonance---the
structural similarity between ideas---propagates through multi-receiver
signalling in three ways: through resonant signals, through resonant
receivers, and through the combination of both.

\begin{theorem}[Resonant Signals Same Mind (10.10)]
\label{thm:resonant-signals-same}
If $s_1\res s_2$, then for each background idea $r$:
$r\compose s_1 \res r\compose s_2.$
\end{theorem}

Christmas and Hanukkah are ``resonant signals''---structurally similar
festivals involving light, family gathering, and gift-giving.  Each
participant's experience of one resonates with their experience of the
other.  This is why comparative religion is intellectually productive: the
structural similarity of the signals guarantees structural similarity of the
interpretations, regardless of the receiver's background.

\begin{theorem}[Same Signal Resonant Receivers (10.11)]
\label{thm:same-signal-resonant}
If $r_1\res r_2$, then for any signal $s$:
$r_1\compose s \res r_2\compose s.$
\end{theorem}

Two members of the same culture (resonant backgrounds) always produce
resonant interpretations of the same event.  This is Durkheim's
\emph{conscience collective}: cultural consensus is a structural consequence
of shared backgrounds.  The theorem gives it precise mathematical content:
resonant backgrounds + common signal = resonant experiences.%
\footnote{Durkheim, \'E., \textit{Les Formes \'el\'ementaires de la vie
religieuse} (1912).  The concept of \emph{conscience collective} has been
criticized as vague; the resonance theorem provides a formal sharpening.}

\begin{theorem}[Full Resonance: Resonant Receivers + Resonant Signals (10.12)]
\label{thm:full-resonance}
If $r_1\res r_2$ and $s_1\res s_2$, then
$r_1\compose s_1 \res r_2\compose s_2.$
\end{theorem}

\begin{proof}[Proof sketch]
Direct from the resonance-compose axiom: resonance is preserved under
composition when both components resonate.
\end{proof}

Members of allied cultures witnessing similar rituals produce resonant
experiences.  This is Turner's \emph{communitas}: spontaneous solidarity
arising when similar people undergo similar experiences.  The theorem shows
that communitas is not mystical---it is a structural consequence of
resonance preservation under composition.

\begin{theorem}[Self-Resonance of All Audience Outcomes (10.13)]
\label{thm:audience-self-resonance}
For each audience member with $r\in R$:
$r\compose s \res r\compose s.$
\end{theorem}

Each participant's experience of a ritual is internally consistent,
regardless of how different participants' experiences are from each other.
The grandmother's experience is coherent; the teenager's experience is
coherent; the anthropologist's experience is coherent.  What may not be
coherent is the \emph{comparison} between them---but that is a question
about \emph{inter-resonance}, not \emph{self-resonance}.

\begin{lstlisting}[caption={Full resonance theorem in Lean~4.}]
theorem full_resonance
    {r_1 r_2 s_1 s_2 : I}
    (hr : IdeaticSpace.resonates r_1 r_2)
    (hs : IdeaticSpace.resonates s_1 s_2) :
    IdeaticSpace.resonates
      (IdeaticSpace.compose r_1 s_1)
      (IdeaticSpace.compose r_2 s_2) :=
  IdeaticSpace.resonance_compose r_1 r_2 s_1 s_2 hr hs
\end{lstlisting}


% ================================================================
\section{Audience Composition and Merging}
\label{sec:mr-merging}

\begin{theorem}[Audience Concatenation (10.14)]
\label{thm:audience-concat}
\[
  \mathrm{audienceOutcome}(A_1\!+\!+\!A_2,\; s)
  \;=\;
  \mathrm{audienceOutcome}(A_1, s)\!+\!+\;\mathrm{audienceOutcome}(A_2, s).
\]
\end{theorem}

Merging two audiences then signalling is the same as signalling each
audience separately then merging the results.  A joint ceremony for two
communities (e.g., an interfaith service) produces results that are just the
concatenation of what each community would have experienced separately.

There is no mysterious ``synergy'' in public ritual---just the union of
individual experiences.  This is a strong anti-holistic result: it says that
the collective effect of a public ritual is \emph{decomposable} into
per-community effects.  Turner's communitas, whatever its phenomenological
reality, is structurally a concatenation, not a fusion.

\begin{theorem}[Replicated Audience (10.15)]
\label{thm:replicated-audience}
An audience of $n$ identical minds produces $n$ identical interpretation lists:
\[
  \mathrm{audienceOutcome}(\mathrm{replicate}(n, R),\; s)
  \;=\;
  \mathrm{replicate}(n, \mathrm{interpret}(R, s)).
\]
\end{theorem}

\begin{proof}[Proof sketch]
Induction on $n$.  The base case is trivial.  For the inductive step, the
head of the list is $\mathrm{interpret}(R,s)$, and the tail follows by the
inductive hypothesis.
\end{proof}

A perfectly homogeneous community produces perfectly homogeneous responses
to any signal.  This is the theoretical limit of cultural uniformity---the
formal ``ideal type'' of a totalitarian society in which every citizen has
been given the same repertoire.  In practice, no society achieves this: even
the most totalitarian regimes harbor underground variation.%
\footnote{On the limits of cultural homogenization in totalitarian states,
see Fitzpatrick, S., \textit{Everyday Stalinism} (Oxford, 1999);
Gellately, R., \textit{Backing Hitler} (Oxford, 2001).}

\begin{lstlisting}[caption={Replicated audience in Lean~4.}]
theorem replicated_audience
    (rep : Repertoire I) (s : I) (n : ℕ) :
    audienceOutcome (List.replicate n rep) s =
    List.replicate n (interpret rep s)
\end{lstlisting}


% ================================================================
\section{Advanced Multi-Receiver Properties}
\label{sec:mr-advanced}

\begin{theorem}[Sequential Signal to Audience (10.16)]
\label{thm:sequential-audience}
For each audience member:
$r\compose(s_1\compose s_2) = (r\compose s_1)\compose s_2.$
\end{theorem}

A two-act ritual (overture + main act) is equivalent to the main act
performed for an audience that has already absorbed the overture.  The first
act transforms the audience; the second act addresses the transformed
audience.  This is the associativity of composition applied to the ritual
setting, and it justifies the common anthropological practice of analyzing
multi-phase rituals as sequential transformations of the participants'
interpretive state.

\begin{theorem}[Void Audience Member is Transparent (10.17)]
\label{thm:void-audience-member}
$\mathrm{interpret}([\void], s) = [s].$
\end{theorem}

\begin{proof}[Proof sketch]
$\void\compose s = s$ by the left-identity axiom.
\end{proof}

The ``naive observer'' with no cultural baggage sees the signal as-is.  This
formalizes the (impossible) ideal of Husserl's \emph{epoch\'e}---bracketing
all presuppositions to encounter the phenomenon ``as it really is.''  In
IDT, the void audience member is the formal model of phenomenological
reduction: a mind with no background that processes signals transparently.

Of course, no real observer achieves this.  Every ethnographer brings a
repertoire; every scientist brings a paradigm.  The void audience member is
a limiting case, useful for theoretical analysis but unattainable in
practice.%
\footnote{Gadamer's critique of Husserl in \textit{Truth and Method} (1960)
makes precisely this point: understanding always involves
\emph{Vorurteile} (pre-judgments), and the attempt to eliminate them is
itself a pre-judgment.}

\begin{theorem}[Depth Bound for Void Outcome Across Audience (10.18)]
\label{thm:void-outcome-preserved}
$\mathrm{depthSum}(\mathrm{interpret}(R, \void)) = \mathrm{depthSum}(R).$
\end{theorem}

The ``information cost'' of silence is exactly zero for every audience member
simultaneously.  This is the formal foundation of the \emph{status quo}:
doing nothing changes nothing for anyone, regardless of their background.

Combined with the audience concatenation theorem
(Theorem~\ref{thm:audience-concat}), this means that the void signal
preserves the total audience depth exactly:
$\mathrm{totalAudienceDepth}(\mathrm{audienceOutcome}(A,\void))
= \mathrm{totalAudienceDepth}(A)$.
Silence is cost-free at the individual level \emph{and} at the collective
level.

\begin{lstlisting}[caption={Void outcome depth preservation in Lean~4.}]
theorem void_outcome_depth_preserved
    (rep : Repertoire I) :
    depthSum (interpret rep (IdeaticSpace.void : I))
      = depthSum rep
\end{lstlisting}


% ================================================================
\section{Conclusion}
\label{sec:mr-conclusion}

This chapter generalized the signalling framework from single receivers to
heterogeneous audiences, providing formal foundations for the anthropological
study of public ritual.

The central results are:
\begin{enumerate}
  \item \emph{Structural invariance}
    (Theorems~\ref{thm:audience-size}--\ref{thm:same-signal-different}):
    audience size and per-receiver interpretation count are preserved.
  \item \emph{Universal transparency of void}
    (Theorems~\ref{thm:void-transparent-single}--\ref{thm:void-transparent-audience}):
    silence changes no one.
  \item \emph{Resonance propagation}
    (Theorems~\ref{thm:resonant-signals-same}--\ref{thm:full-resonance}):
    resonant signals and resonant receivers produce resonant experiences.
  \item \emph{Decomposability} (Theorem~\ref{thm:audience-concat}):
    multi-community rituals decompose into per-community effects.
\end{enumerate}

The theory vindicates Geertz's ``thick description'' program: different
observers \emph{do} produce different interpretations, but these
interpretations are structurally related through the resonance relation.
It formalizes Durkheim's \emph{conscience collective} as resonance among
individual outcomes, and Turner's \emph{communitas} as the full resonance
theorem (Theorem~\ref{thm:full-resonance}).

Together, Chapters~\ref{ch:equilibrium}--\ref{ch:multireceiver} provide a
complete game-theoretic framework for the study of signalling in ideatic
spaces: from equilibrium existence (Chapter~\ref{ch:equilibrium}) through
persuasion bounds (Chapter~\ref{ch:persuasion}), lossy reception
(Chapter~\ref{ch:lossy}), echo chambers (Chapter~\ref{ch:echo}), to the
multi-receiver case studied here.  The framework is fully formalized in
Lean~4, with every theorem machine-verified and zero \texttt{sorry} axioms.


% ── Part III: Populations and Markets ─────────────────────────────────
\part{Populations and Markets}
\label{part:populations}

% ================================================================
% Chapter 11: Competing Senders and the Marketplace of Ideas
% ================================================================

\chapter{Competing Senders and the Marketplace of Ideas}
\label{ch:marketplace}

\epigraph{The ideas of the ruling class are in every epoch the ruling ideas,
i.e.\ the class which is the ruling \emph{material} force of society is at the
same time its ruling \emph{intellectual} force.}{Karl Marx, \textit{The German Ideology} (1846)}

% ----------------------------------------------------------------
\section{Introduction}
% ----------------------------------------------------------------

In every human ecology, multiple idea-systems compete for the same audience.
Missionaries contend with shamans for the interpretive loyalty of
newly contacted peoples; Enlightenment rationalism and Romantic nationalism
vie for Europe's educated classes; Instagram influencers and grandparents
offer competing visions of the good life to the same adolescent.
The \emph{marketplace of ideas}---a metaphor usually left at the level of
metaphor---is, in Ideatic Diffusion Theory, a structure amenable to rigorous
analysis.

This chapter formalizes the marketplace.  We introduce the \texttt{IdeaMarket}
structure, which collects a list of competing \emph{senders} (signals), a
receiver's background \emph{repertoire}, and a payoff function defined in
terms of interpretive depth.  A sender $s_1$ \emph{dominates} sender $s_2$
when $s_1$ produces deeper interpretations for every element of the
receiver's repertoire.  We then prove eighteen theorems that expose the
structural constraints on ideological competition: void never dominates,
self-dominance is trivial, depth chains bound competitive advantage, and
coalition of senders is subadditive---connecting to Gramsci's cultural
hegemony, Berger's sacred canopy, and rational-choice theory of religion.

The key insight is that competitive signalling in ideatic space is not a
free-for-all: the subadditivity of depth imposes hard ceilings on how much
any sender---or any coalition of senders---can achieve.  Cultural hegemony,
in the formal sense, is always bounded above by the audience's accumulated
cultural capital.  This structural result places precise limits on the
efficacy of propaganda, evangelism, and ideological persuasion.

Throughout this chapter, we work in an ideatic space $(\IS, \compose, \void, \depth, \res)$
satisfying the IDT axioms.  Signals and repertoire elements are ideas
$s, r \in \IS$; interpretation is composition; payoff is depth.\footnote{The
interpretation model here is identical to that introduced in Chapter~1:
a receiver with background $r$ encountering signal $s$ produces the
interpretation $r \compose s$.}

\subsection{Historical Foundations of the Marketplace Metaphor}

The ``marketplace of ideas'' is among the most influential metaphors in Western
political and intellectual history, yet its genealogy is rarely traced with
care.  The phrase itself is modern, but the concept it encodes---that
truth emerges from the competitive interplay of rival doctrines---has deep
roots in both liberal political theory and the sociology of religion.

John Stuart Mill's \textit{On Liberty} (1859) provides the canonical
philosophical argument.  Mill's defense of free expression rests not on
any natural right but on an \emph{epistemic} claim: suppression of opinion
is harmful because it deprives society of the opportunity to exchange
error for truth, or to sharpen truth by collision with error.  ``The
peculiar evil of silencing the expression of an opinion is, that it is
robbing the human race,'' Mill writes---a statement that, in our
formalism, asserts that removing a sender from $\mathbf{S}$ can only
reduce the maximum achievable payoff across the repertoire $\mathbf{R}$.
Mill's argument implicitly assumes a kind of monotonicity: more senders,
weakly more total payoff.  Our framework makes this assumption precise
and identifies the conditions under which it holds.

The phrase ``marketplace of ideas'' itself is typically attributed to
Justice Oliver Wendell Holmes Jr.'s celebrated dissent in
\textit{Abrams v.\ United States} (1919), in which he wrote: ``the best
test of truth is the power of the thought to get itself accepted in the
competition of the market.''  Holmes's market metaphor imports the
apparatus of classical economics---competition, selection, equilibrium---into
the domain of belief.  In IDT, we take this importation seriously: the
\texttt{IdeaMarket} structure is a formal market, and dominance is a formal
notion of competitive advantage.  But unlike the neoclassical market, the
ideatic market operates under subadditivity rather than constant returns to
scale, and payoff is bounded by receiver capacity rather than by supply.

The sociology of religion provides a parallel tradition.  Rodney Stark and
William Sims Bainbridge, in \textit{A Theory of Religion} (1987), proposed
a rational-choice model of religious competition: churches are ``firms''
producing supernatural compensators, believers are ``consumers'' selecting
among competing suppliers, and the vigor of religious life is a function of
market structure.  Stark's later work, \textit{The Rise of Christianity}
(1996), argued that Christianity's triumph in the late Roman Empire was not
a miracle but a consequence of superior ``market positioning''---a signal
that resonated across a wider range of receiver backgrounds than its
competitors.  Our formal notion of dominance
(Definition~\ref{def:dominance}) captures precisely this competitive
dynamic: a sender dominates when it achieves higher depth on \emph{every}
element of the repertoire.

Richard Dawkins's concept of the \emph{meme}, introduced in \textit{The
Selfish Gene} (1976), offers another precursor.  Dawkins proposed that
cultural units---memes---replicate, mutate, and compete for survival in a
manner analogous to genes.  Daniel Dennett, in \textit{Darwin's Dangerous
Idea} (1995), extended the memetic framework, arguing that cultural
evolution operates through the differential reproduction of ideas in
minds construed as ecological niches.  In our framework, memes are senders,
minds are receivers with repertoires, and ``survival'' corresponds to
achieving non-zero payoff: a meme ``survives'' in a mind when
$\depth(r \compose s) > 0$ for at least one $r \in \mathbf{R}$.

Antonio Gramsci's theory of cultural hegemony, developed in the
\textit{Quaderni del carcere} (1929--1935), provides the critical
counterpoint to liberal marketplace optimism.  For Gramsci, the marketplace
is not a level playing field but a terrain shaped by power.  The ruling
class achieves \emph{egemonia} not merely by producing superior signals but
by structuring the receivers' repertoires---through education, media, and
institutional socialization---so that its signals \emph{necessarily}
dominate.  In IDT terms, hegemony is achieved not by maximizing
$\depth(s)$ but by shaping $\mathbf{R}$ so that
$\depth(r \compose s_{\text{ruling}}) \ge \depth(r \compose s_{\text{rival}})$
for all $r \in \mathbf{R}$.  The theorems of this chapter apply regardless
of how $\mathbf{R}$ was formed; the critical question---\emph{who shapes
the repertoire?}---lies outside the formal system but is illuminated by it.

Edward Herman and Noam Chomsky's ``propaganda model,'' articulated in
\textit{Manufacturing Consent} (1988), extends Gramsci's insight to mass
media.  In their account, the marketplace of ideas is systematically
distorted by ownership concentration, advertising dependency, and
institutional sourcing.  Our payoff bound
(Theorem~\ref{thm:payoff-bound}) provides a structural complement: even
in a perfectly competitive market, payoff is bounded by audience capacity.
Propaganda's efficacy is limited not only by institutional filters but by
the mathematical ceiling that subadditivity imposes.

The modern attention economy, theorized by Herbert Simon (``a wealth of
information creates a poverty of attention''), Tim Wu (\textit{The Attention
Merchants}, 2016), and others, represents the latest iteration of
marketplace competition.  Social media platforms---Facebook's News Feed,
TikTok's For You page, YouTube's recommendation algorithm---are literal
idea markets, with competing senders (creators, advertisers, political
actors) vying for a finite attention budget.  The receiver's repertoire
$\mathbf{R}$ is finite, and each sender's signal competes for interpretive
engagement with each element of $\mathbf{R}$.  The total payoff across all
senders is bounded by the receiver's total interpretive capacity.  The
formal results of this chapter thus apply directly to the digital
marketplace of the twenty-first century: algorithmic curation selects
senders, but the structural constraints---subadditivity, bounded payoff,
resonance preservation---remain invariant.

% ----------------------------------------------------------------
\section{Core Definitions}
\label{sec:ch11-defs}
% ----------------------------------------------------------------

\begin{definition}[Idea Market]
\label{def:idea-market}
An \emph{idea market} is a tuple $\mathcal{M} = (\mathbf{S}, \mathbf{R})$
where $\mathbf{S} = [s_1, \dots, s_m]$ is a list of \emph{competing senders}
(signals) and $\mathbf{R} = [r_1, \dots, r_n]$ is the receiver's
\emph{background repertoire}.  All elements belong to an ideatic space~$\IS$.
\end{definition}

\begin{lstlisting}[caption={IdeaMarket structure in Lean~4}]
structure IdeaMarket (I : Type*) [IdeaticSpace I] where
  senders    : List I
  repertoire : List I
\end{lstlisting}

\begin{definition}[Interpretation]
\label{def:interpret-ch11}
The \emph{interpretation} of a signal $s$ against a repertoire
$\mathbf{R} = [r_1, \dots, r_n]$ is the list
\[
  \operatorname{interpret}(\mathbf{R}, s)
  \;=\; [r_1 \compose s,\; r_2 \compose s,\; \dots,\; r_n \compose s].
\]
Each entry $r_i \compose s$ is the idea produced when background $r_i$
receives signal $s$.
\end{definition}

\begin{definition}[Total Payoff]
\label{def:total-payoff}
The \emph{total interpretive payoff} of signal $s$ on repertoire $\mathbf{R}$
is
\[
  \operatorname{totalPayoff}(\mathbf{R}, s)
  \;=\; \sum_{i=1}^{n} \depth(r_i \compose s).
\]
This aggregates the depth of every interpretation the signal induces.
\end{definition}

\begin{definition}[Dominance]
\label{def:dominance}
Sender $s_1$ \emph{dominates} sender $s_2$ on repertoire $\mathbf{R}$ when
\[
  \forall\, r \in \mathbf{R},\quad
  \depth(r \compose s_1) \;\ge\; \depth(r \compose s_2).
\]
In anthropological terms, $s_1$ is \emph{more culturally productive} than
$s_2$ for this audience.
\end{definition}

\begin{lstlisting}[caption={Dominance in Lean~4}]
def dominates (s_1 s_2 : I) (rep : List I) : Prop :=
  forall r in rep, depth (compose r s_1) >= depth (compose r s_2)
\end{lstlisting}

The notion of dominance is a pointwise ordering on the payoff function.  It
is deliberately strong: $s_1$ must beat $s_2$ on \emph{every} member of the
audience, not merely on average.  This models \emph{Gramscian} hegemony, in
which the ruling ideology must be accepted across all strata of civil
society, not merely by a majority.\footnote{Gramsci's concept of
\emph{egemonia} in the \textit{Quaderni del carcere} requires precisely
this universality: hegemony is not statistical dominance but categorical
dominance across the entire \emph{blocco storico}.}

% ----------------------------------------------------------------
\section{Void and Dominance}
\label{sec:ch11-void}
% ----------------------------------------------------------------

\begin{theorem}[Void Dominates is Identity (11.1)]
\label{thm:void-dom-identity}
The void signal, when received, simply returns the repertoire unchanged:
\[
  \operatorname{interpret}(\mathbf{R}, \void) \;=\; \mathbf{R}.
\]
\end{theorem}

\begin{proof}[Proof sketch]
By induction on the list $\mathbf{R}$.  For each element $r$, the void
axiom gives $r \compose \void = r$, so
$\operatorname{map}(\lambda\, r.\; r \compose \void, \mathbf{R}) = \mathbf{R}$.
\end{proof}

Silence, the null message, is never the \emph{winning} strategy in a
competitive marketplace---it merely preserves the status quo.  Gramsci's
insight is precisely this: hegemony requires \emph{active} cultural
production, not mere silence.  The void signal is the ideational equivalent
of ideological abdication.

\begin{example}[The Roman Marketplace of Cults]
In the late Roman Republic, the traditional Roman pantheon competed with
imported mystery cults---Isis from Egypt, Mithras from Persia, Cybele from
Phrygia---for the devotion of urban populations. Each cult offered a
different ``signal'' to the same audience. The depth of interpretation
varied with the receiver's background: a Stoic philosopher produced a
deeper interpretation of Mithraic ritual than a dock worker, not because
the ritual differed but because $\depth(r_{\text{philosopher}}) >
\depth(r_{\text{worker}})$. The eventual dominance of Christianity can be
modeled as the emergence of a sender whose signal resonated across the
widest range of receiver backgrounds.\footnote{Rodney Stark,
\textit{The Rise of Christianity} (1996), argues that Christianity's
competitive advantage lay precisely in its broad resonance across social
strata.}

The formal apparatus of this chapter illuminates the mechanism.  Each cult
$c_i \in \mathbf{S}$ had its characteristic payoff profile across the
urban repertoire $\mathbf{R}_{\text{Rome}}$.  Mithraism achieved high depth
among soldiers ($\depth(r_{\text{miles}} \compose s_{\text{Mithra}})$ was
large) but near-zero depth among women, whom it excluded.  The cult of Isis
appealed to merchants and sailors but had limited reach among the
senatorial class.  Christianity's competitive strategy---universal
soteriology, egalitarian ecclesiology, portable ritual---produced moderate
but \emph{uniform} depth across the entire repertoire.  It did not dominate
any single cult on its home turf, but it dominated in the aggregate:
$\operatorname{totalPayoff}(\mathbf{R}_{\text{Rome}}, s_{\text{Christ}}) >
\operatorname{totalPayoff}(\mathbf{R}_{\text{Rome}}, s_{c_i})$ for each
rival $c_i$.  The marketplace selected for breadth, not peak depth---a
structural prediction of our payoff framework.
\end{example}

The Roman example illustrates a general principle that recurs throughout
this chapter: in a diverse marketplace, the sender that achieves moderate
depth across \emph{all} receivers often outperforms the sender that achieves
extreme depth in a narrow niche.  This is a direct consequence of the
structure of the total payoff function, which sums depth across the entire
repertoire.  A signal tuned to a single receiver type can achieve high
local depth but low total payoff; a signal tuned to the entire population
achieves lower local depth but higher total payoff.  The tension between
specialization and universality is thus a structural feature of the
marketplace, not merely a strategic choice.

\begin{theorem}[Self-Dominance is Trivial (11.2)]
\label{thm:self-dom}
Every signal dominates itself:
\[
  \forall\, s \in \IS,\quad \operatorname{dominates}(s, s, \mathbf{R}).
\]
\end{theorem}

\begin{proof}[Proof sketch]
For each $r \in \mathbf{R}$, $\depth(r \compose s) \ge \depth(r \compose s)$
holds by reflexivity of~$\le$.
\end{proof}

Self-consistency is the minimal requirement of any ideological system.
That a religion ``dominates itself'' says nothing about competitive
advantage---it is the zero-information baseline.  Philosophically, this is
Hegel's \emph{Sichselbstgleichheit}: self-identity is the formal precondition
of all determination, never the determination itself.

\begin{theorem}[Void Signal Payoff Equals Repertoire Depth (11.3)]
\label{thm:void-payoff}
The total payoff from the void signal equals the total depth of the repertoire:
\[
  \operatorname{totalPayoff}(\mathbf{R}, \void) \;=\; \operatorname{depthSum}(\mathbf{R}).
\]
\end{theorem}

\begin{proof}[Proof sketch]
By Theorem~\ref{thm:void-dom-identity},
$\operatorname{interpret}(\mathbf{R}, \void) = \mathbf{R}$, so
$\operatorname{depthSum}(\operatorname{interpret}(\mathbf{R}, \void))
= \operatorname{depthSum}(\mathbf{R})$.
\end{proof}

A society's interpretive capacity under silence is exactly its accumulated
cultural depth.  This is the \emph{baseline} against which all competing
signals are measured.  In Bourdieu's terms, the void signal reveals nothing
but the audience's pre-existing \emph{capital culturel}.\footnote{Bourdieu,
\textit{La Distinction} (1979), defines cultural capital as the accumulated
stock of cultural knowledge, skills, and dispositions.  Our
$\operatorname{depthSum}(\mathbf{R})$ is its formal analogue.}

\begin{theorem}[Dominance Reflexivity (11.4)]
\label{thm:dom-refl}
Dominance is reflexive: $\operatorname{dominates}(s, s, \mathbf{R})$ for all $s$
and $\mathbf{R}$.
\end{theorem}

\begin{proof}[Proof sketch]
Immediate from Theorem~\ref{thm:self-dom}.
\end{proof}

In the language of partial orders, dominance is a \emph{preorder} (reflexive
and---as we shall see in later chapters---transitive in an appropriate sense).
This is the formal structure underlying marketplace competition: a preorder on
signals, indexed by audiences.

\begin{remark}[Connection to Attention Economics]
The ideatic marketplace anticipates Herbert Simon's observation that
``a wealth of information creates a poverty of attention.'' In our
framework, the receiver's repertoire $\mathbf{R}$ is finite, and each
sender competes for interpretive engagement with each element of
$\mathbf{R}$. The total payoff across all senders is bounded by the
receiver's total interpretive capacity---a formal version of the
``attention budget'' that dominates contemporary media economics.

Simon's insight, formulated in 1971, has acquired new urgency in the
age of algorithmic content curation.  When a platform's recommendation
engine selects which senders to present to a receiver, it is performing
a constrained optimization over the idea market $\mathcal{M}$: choosing
a subset $\mathbf{S}' \subseteq \mathbf{S}$ to maximize (some proxy for)
total payoff subject to the attention constraint $|\mathbf{S}'| \le k$.
The structural results of this chapter---boundedness, subadditivity,
resonance preservation---apply to this optimization problem and constrain
its solutions.  No algorithm, however sophisticated, can circumvent the
subadditivity ceiling.
\end{remark}

\begin{example}[The Protestant Reformation as Marketplace Competition]
The Protestant Reformation, inaugurated by Luther's \textit{Ninety-Five
Theses} in 1517, is perhaps history's most consequential example of a new
sender entering an established ideatic marketplace.  For a millennium,
the Catholic Church had enjoyed something approaching monopoly status in
Western European Christendom: $\mathbf{S} \approx \{s_{\text{Catholic}}\}$.
Luther, Calvin, Zwingli, and the radical reformers introduced competing
signals $s_{\text{Lutheran}}, s_{\text{Calvinist}}, s_{\text{Anabaptist}}$,
each calibrated to different segments of the receiver repertoire.

Luther's signal was optimized for the German-speaking urban middle class:
vernacular scripture, congregational hymns, the doctrine of the
priesthood of all believers.  For a Wittenberg burgher $r_{\text{burgher}}$,
$\depth(r_{\text{burgher}} \compose s_{\text{Luther}}) \gg
\depth(r_{\text{burgher}} \compose s_{\text{Catholic}})$---the Lutheran
signal resonated more deeply with his existing cultural repertoire
(vernacular literacy, anti-clerical sentiment, commercial independence).
But for a Bavarian peasant embedded in a deeply Catholic milieu,
$r_{\text{peasant}}$, the Catholic signal still dominated.  The
Reformation thus illustrates the \emph{audience-dependence} of dominance:
$s_{\text{Luther}}$ dominated $s_{\text{Catholic}}$ on some repertoires
but not others.  The eventual confessional geography of Europe---the
principle of \textit{cuius regio, eius religio}---reflects the fact that
dominance was regional, not universal.

Calvin's Geneva, by contrast, represents a case of \emph{repertoire
reshaping}: the Consistory systematically restructured receivers'
backgrounds (through catechism, discipline, sumptuary laws) so that
$\depth(r_{\text{Genevan}} \compose s_{\text{Calvin}})$ was maximized.
This is Gramscian hegemony avant la lettre---competitive advantage
achieved by shaping the audience, not just the signal.
\end{example}

% ----------------------------------------------------------------
\section{Depth Bounds on Competition}
\label{sec:ch11-bounds}
% ----------------------------------------------------------------

\begin{theorem}[Hermeneutic Payoff Bound (11.5)]
\label{thm:payoff-bound}
The total payoff of any signal $s$ on repertoire $\mathbf{R}$ is bounded:
\[
  \operatorname{totalPayoff}(\mathbf{R}, s)
  \;\le\; \operatorname{depthSum}(\mathbf{R}) + |\mathbf{R}| \cdot \depth(s).
\]
\end{theorem}

\begin{proof}[Proof sketch]
By induction on $\mathbf{R}$.  At each step, the subadditivity axiom gives
$\depth(r \compose s) \le \depth(r) + \depth(s)$, and the bound accumulates
linearly.
\end{proof}

Even the most profound revelation cannot exceed the audience's capacity plus a
linear contribution from the signal itself.  Mass media reaches many, but depth
per receiver is bounded by pre-existing cultural capital.  This theorem is the
formal backbone of what communications scholars call the
\emph{knowledge gap hypothesis}: the information-rich get richer because their
higher $\operatorname{depthSum}(\mathbf{R})$ allows them to extract more from
every signal.\footnote{Tichenor, Donohue, and Olien, ``Mass Media Flow and
Differential Growth in Knowledge,'' \textit{Public Opinion Quarterly}
34 (1970), 159--170.}

\begin{example}[Algorithmic Competition on Social Media Platforms]
Contemporary social media platforms instantiate the idea market structure
with unprecedented literalness.  Consider TikTok's ``For You'' page.  The
platform maintains an enormous sender pool
$\mathbf{S} = \{s_1, s_2, \dots, s_N\}$ (all available videos), and each
user brings a repertoire $\mathbf{R}_u$ shaped by prior viewing history,
demographic background, and expressed preferences.  The recommendation
algorithm selects a subset $\mathbf{S}' \subset \mathbf{S}$ to present to
the user, implicitly attempting to maximize engagement---a proxy for our
total payoff $\operatorname{totalPayoff}(\mathbf{R}_u, s)$.

The subadditivity bound (Theorem~\ref{thm:payoff-bound}) has an immediate
practical consequence: no matter how sophisticated the algorithm, the total
interpretive value extracted is bounded by the user's existing repertoire
depth plus a linear contribution from the content.  This places a hard
ceiling on ``algorithmic radicalization'' theories: the depth of the
interpretation produced by an extremist signal $s_{\text{radical}}$ is
bounded by $\depth(r) + \depth(s_{\text{radical}})$ for each repertoire
element $r$.  A receiver with a shallow repertoire
($\operatorname{depthSum}(\mathbf{R}_u)$ is small) can be moved by the
signal, but only to a bounded degree.  The algorithm cannot manufacture
depth that neither sender nor receiver possesses.

Furthermore, the resonance preservation theorem
(Theorem~\ref{thm:homogeneous-audience}) explains the ``filter bubble''
phenomenon: users with resonant backgrounds
($r_i \res r_j$ for $r_i, r_j \in \mathbf{R}_u$) produce resonant
interpretations of the same content, reinforcing homogeneity.  The platform
does not create the bubble; it amplifies a structural property of
resonance-preserving composition.
\end{example}

The social media example illustrates how the abstract formal results of this
chapter apply to concrete twenty-first-century institutions.  The marketplace
metaphor, born in Mill's liberalism, refined in Holmes's jurisprudence, and
challenged by Gramsci's critical theory, finds its most literal instantiation
in the algorithmic content platforms of the digital age.  Yet the structural
constraints remain invariant: subadditivity, bounded payoff, resonance
preservation.  These are not historical contingencies but mathematical
necessities inherent in the structure of ideatic space.

\begin{theorem}[Single-Element Repertoire Payoff (11.6)]
\label{thm:singleton-payoff}
For a singleton repertoire $[r]$, the payoff of signal $s$ is exactly:
\[
  \operatorname{totalPayoff}([r], s) = \depth(r \compose s).
\]
\end{theorem}

\begin{proof}[Proof sketch]
Direct unfolding of definitions: $\operatorname{interpret}([r], s) =
[r \compose s]$, and $\operatorname{depthSum}([r \compose s]) =
\depth(r \compose s)$.
\end{proof}

The simplest ``audience''---one mind, one background idea---has its payoff
fully determined by the composition of background and signal.  This is the
irreducible unit of cultural reception, the phenomenological atom of
\emph{Verstehen}.

\begin{theorem}[Empty Repertoire Yields Zero (11.7)]
\label{thm:empty-rep-zero}
If no one is listening, every signal has zero payoff:
$\operatorname{totalPayoff}([], s) = 0$.
\end{theorem}

\begin{proof}[Proof sketch]
The empty list maps to the empty list, whose depth sum is zero.
\end{proof}

A prophet without followers, a book without readers---ideational production
without reception is culturally null.  The tree that falls in an empty forest
produces no interpretive depth.  This is not mysticism but arithmetic: zero
receivers yield zero payoff.

% ----------------------------------------------------------------
\section{Competitive Composition}
\label{sec:ch11-comp}
% ----------------------------------------------------------------

\begin{theorem}[Signal Composition Depth Bound (11.8)]
\label{thm:composed-signal-bound}
Composing two signals $s_1, s_2$ yields a compound signal whose depth on
any repertoire element $r$ is triply bounded:
\[
  \depth(r \compose (s_1 \compose s_2))
  \;\le\; \depth(r) + \depth(s_1) + \depth(s_2).
\]
\end{theorem}

\begin{proof}[Proof sketch]
Apply subadditivity twice: first to $s_1 \compose s_2$, then to
$r \compose (s_1 \compose s_2)$, and combine via transitivity.
\end{proof}

Syncretism---composing competing mythologies into a single system---cannot
create more depth than the sum of parts.  The Hellenistic synthesis of Greek
philosophy and Egyptian religion, the syncretic Christianity of colonial Latin
America, the New Age blending of Eastern and Western spiritualities: all are
structurally constrained by this triple bound.  The ``best of both worlds'' is
an arithmetic ceiling, not an aspirational metaphor.

\begin{theorem}[Void Competitor Adds Nothing, Right (11.9)]
\label{thm:void-comp-right}
Adding void as a second sender to compose with $s$ does not change $s$:
$s \compose \void = s$.
\end{theorem}

\begin{proof}[Proof sketch]
This is the right-identity axiom of the monoid.
\end{proof}

\begin{theorem}[Void Competitor Adds Nothing, Left (11.10)]
\label{thm:void-comp-left}
Void composed on the left with any signal yields the signal unchanged:
$\void \compose s = s$.
\end{theorem}

\begin{proof}[Proof sketch]
This is the left-identity axiom of the monoid.
\end{proof}

A ``silent partner'' in cultural production is no partner at all.  Coalitions
with non-contributors are structurally inert.  This is why purely nominal
alliances---the Holy Roman Empire, which Voltaire quipped was ``neither holy,
nor Roman, nor an empire''---add no interpretive depth.  The void contributor
is formally invisible.

The results of this section have deep implications for understanding
cultural syncretism.  The triple bound on composed signals
(Theorem~\ref{thm:composed-signal-bound}) formalizes a constraint that
anthropologists of religion have long observed informally: syncretic
religious systems---Vodou combining West African \emph{orisha} with
Catholic saints, Sikhism blending Hindu devotionalism with Islamic
monotheism, Bahá'í synthesizing Abrahamic and Zoroastrian themes---never
achieve a depth exceeding the sum of their components.  This is not a
failure of synthesis but a structural inevitability.  The ``sum of parts''
ceiling applies universally, regardless of the creativity or profundity
of the synthesizer.

The void-composition results (Theorems~\ref{thm:void-comp-right}
and~\ref{thm:void-comp-left}) complement the syncretism analysis by
establishing the base case: adding nothing to a cultural system changes
nothing.  This seemingly trivial result has a non-trivial implication for
theories of cultural decline.  When a society's cultural production
approaches void---when its intellectuals, artists, and religious
leaders produce signals of negligible depth---the marketplace degenerates.
The total payoff reduces to the existing repertoire's depth, and no
competitive dynamics operate.  This is the formal analogue of Spengler's
``Zivilisation'' phase, in which exhausted cultures merely repeat existing
forms without generating new depth.\footnote{Oswald Spengler,
\textit{Der Untergang des Abendlandes} (1918--1922).  Spengler's
distinction between \emph{Kultur} (creative cultural production) and
\emph{Zivilisation} (exhausted repetition) maps onto the distinction
between non-void and near-void senders in our framework.}

% ----------------------------------------------------------------
\section{Resonance and Competitive Advantage}
\label{sec:ch11-resonance}
% ----------------------------------------------------------------

\begin{theorem}[Resonant Signals Produce Resonant Interpretations (11.11)]
\label{thm:resonant-competitors}
If two competing senders produce resonant signals, $s_1 \res s_2$, then any
receiver with background $r$ produces resonant interpretations:
\[
  s_1 \res s_2
  \quad\Longrightarrow\quad
  (r \compose s_1) \res (r \compose s_2).
\]
\end{theorem}

\begin{proof}[Proof sketch]
By the left-composition preservation of resonance
(\texttt{resonance\_compose\_left}).
\end{proof}

Competing denominations of the same religion---Catholic vs.\ Orthodox, Sunni
vs.\ Shia---produce \emph{resonant} interpretations in the laity.  The
faithful perceive both as ``basically the same religion'' precisely because the
signals resonate.  This theorem is the formal basis of ecumenism: if the
signals resonate, no receiver can distinguish them at the level of
resonance.\footnote{The ecumenical movement, from the 1910 Edinburgh Conference
onward, rests on exactly this structural claim: beneath doctrinal differences,
the core Christian signals resonate.}

\begin{theorem}[Homogeneous Audience (11.12)]
\label{thm:homogeneous-audience}
If two receivers have resonant backgrounds, $r_1 \res r_2$, and receive the
same signal $s$, their interpretations are resonant:
\[
  r_1 \res r_2
  \quad\Longrightarrow\quad
  (r_1 \compose s) \res (r_2 \compose s).
\]
\end{theorem}

\begin{proof}[Proof sketch]
By right-composition preservation of resonance
(\texttt{resonance\_compose\_right}).
\end{proof}

Members of the same cultural milieu---resonant backgrounds---confronted with
the same competing signal will interpret it similarly.  This is why propaganda
works best within culturally homogeneous populations: the more resonant the
audience, the more uniform the reception.  Conversely, propaganda aimed at a
culturally heterogeneous audience (non-resonant backgrounds) produces
divergent interpretations---the structural limit of mass persuasion in diverse
societies.

\begin{theorem}[Full Cultural Resonance (11.13)]
\label{thm:full-resonance}
Resonant backgrounds plus resonant signals imply resonant interpretations:
\[
  r_1 \res r_2 \;\wedge\; s_1 \res s_2
  \quad\Longrightarrow\quad
  (r_1 \compose s_1) \res (r_2 \compose s_2).
\]
\end{theorem}

\begin{proof}[Proof sketch]
By the full composition-preservation of resonance
(\texttt{resonance\_compose}).
\end{proof}

Allied cultures (resonant backgrounds) exposed to similar myths (resonant
signals) will always arrive at compatible worldviews.  This is the formal
mechanism behind what Redfield called the ``Great Tradition / Little
Tradition'' unity: literate urban elites and illiterate peasants, as long as
their backgrounds resonate and the signals they receive resonate, will
inhabit structurally compatible interpretive universes.\footnote{Robert
Redfield, \textit{Peasant Society and Culture} (1956).  Redfield's
bi-partition of tradition into ``great'' and ``little'' variants assumes
precisely this kind of resonance preservation across social strata.}

\begin{remark}[Open Questions in Competitive Signalling]
Several natural questions arise from the framework developed in this
chapter that remain open for future investigation:

\begin{enumerate}
\item \textbf{Dynamic competition.}  Our model is static: senders and
  repertoire are fixed.  In reality, senders \emph{adapt} their signals in
  response to competitive pressure, and receivers' repertoires evolve
  through exposure.  A dynamic extension would model the co-evolution of
  $\mathbf{S}$ and $\mathbf{R}$ over time, potentially connecting to
  evolutionary game theory and replicator dynamics.

\item \textbf{Coalitional signalling.}  When multiple senders form a
  coalition (an alliance of churches, a media conglomerate, a political
  coalition), the coalition's composite signal is the composition
  $s_1 \compose s_2 \compose \cdots \compose s_k$.  The subadditivity
  bound applies, but optimal coalition formation---which subset of senders
  should ally?---is a combinatorial optimization problem whose complexity
  is not yet characterized in our framework.

\item \textbf{Receiver strategic behavior.}  Our model treats receivers as
  passive: they interpret signals with their existing repertoire but do not
  strategically choose which signals to attend to.  Incorporating receiver
  choice would connect the framework to mechanism design and the economics
  of attention.

\item \textbf{Stochastic payoff.}  In practice, depth may be noisy or
  probabilistic---a receiver may or may not ``get'' a given signal on any
  particular exposure.  A stochastic extension of the payoff function would
  connect to information-theoretic models of cultural transmission.
\end{enumerate}

These questions point toward a richer theory of cultural competition that
integrates the static structural results of this chapter with dynamic,
strategic, and stochastic considerations.
\end{remark}

% ----------------------------------------------------------------
\section{Marketplace Structural Theorems}
\label{sec:ch11-structural}
% ----------------------------------------------------------------

\begin{theorem}[Interpretation Preserves Length (11.14)]
\label{thm:interp-length}
The number of interpretations equals the repertoire size:
$|\operatorname{interpret}(\mathbf{R}, s)| = |\mathbf{R}|$.
\end{theorem}

\begin{proof}[Proof sketch]
$\operatorname{interpret}$ is a \texttt{map}, which preserves list length.
\end{proof}

Cultural contact does not change the \emph{number} of categories a society
has---it transforms each one.  The French anthropological term
\emph{bricolage} (Lévi-Strauss) describes precisely this: the bricoleur
rearranges existing schema rather than adding new ones.  Our theorem
formalizes this conservation law of cultural schema.

The length-preservation result has a subtle but important anthropological
implication.  Cultural contact is often described in terms of
\emph{accretion}---new categories being added to a society's conceptual
inventory.  Our theorem shows that, at the formal level of ideatic
composition, interpretation is \emph{transformative} rather than
\emph{additive}: each element of the repertoire is transformed by the
incoming signal, but no new elements are created.  The apparent addition
of new categories in empirical cultural contact (the introduction of
``democracy'' into a society that lacked the concept, for instance)
must be modeled not as a new element appearing in the interpretation list
but as a pre-existing schema being transformed beyond recognition---the
bricoleur's old tool repurposed for a new function.  This is a strong
claim, and it reflects the limitations of our model as much as the
structure of cultural reality.  A richer model might allow the
interpretation function to produce lists of different length, but the
simplicity of the present framework---and its formal tractability---is
purchased at the price of this conservation law.

\begin{theorem}[Additivity of Repertoires (11.15)]
\label{thm:rep-additivity}
Interpreting a signal with a merged repertoire equals concatenating
interpretations of the sub-repertoires:
\[
  \operatorname{interpret}(\mathbf{R}_1 \mathbin{+\!\!\!+} \mathbf{R}_2, s)
  \;=\;
  \operatorname{interpret}(\mathbf{R}_1, s)
  \mathbin{+\!\!\!+}
  \operatorname{interpret}(\mathbf{R}_2, s).
\]
\end{theorem}

\begin{proof}[Proof sketch]
By the distributivity of \texttt{map} over list concatenation.
\end{proof}

When two groups merge---through marriage alliance, conquest, or
migration---their joint interpretation of any signal equals the union of their
separate interpretations.  Cultural integration is structurally additive.
This additivity result is the formal basis of the ``mosaic'' model of
multicultural societies: each group's interpretive contribution is preserved
intact in the combined repertoire.

\begin{theorem}[Composed Signal Associativity (11.16)]
\label{thm:composed-assoc}
Interpreting a composition of signals through background $r$ is the same
as sequential interpretation:
\[
  r \compose (s_1 \compose s_2)
  \;=\;
  (r \compose s_1) \compose s_2.
\]
\end{theorem}

\begin{proof}[Proof sketch]
By associativity of the monoid operation, applied symmetrically.
\end{proof}

Receiving two messages in sequence---first $s_1$, then $s_2$---is
structurally identical to receiving their composition.  The order of cultural
exposure \emph{telescopes} via associativity: a semester's lectures, received
one by one, produce the same interpretation as receiving their composition
all at once.  This is both the strength and the limitation of the model: it
captures the structural content of sequential exposure but not the
psychological dynamics of surprise, anticipation, or fatigue.

\begin{theorem}[Marginal Depth Contribution Bound (11.17)]
\label{thm:marginal-bound}
Adding a new element $r$ to the repertoire increases total payoff by at most
$\depth(r) + \depth(s)$:
\[
  \operatorname{totalPayoff}(r :: \mathbf{R}, s)
  \;\le\;
  \depth(r) + \depth(s) + \operatorname{totalPayoff}(\mathbf{R}, s).
\]
\end{theorem}

\begin{proof}[Proof sketch]
The new payoff contribution is $\depth(r \compose s) \le \depth(r) + \depth(s)$
by subadditivity.
\end{proof}

Educating one more person---adding one more background to the audience---increases
total societal understanding by at most that person's depth plus the signal's.
Diminishing returns in mass education are structurally guaranteed.  This
theorem resonates with Illich's critique of institutional education: past a
certain point, expanding the student body adds progressively less interpretive
value per capita.\footnote{Ivan Illich, \textit{Deschooling Society} (1971).
Illich's argument that schooling yields diminishing returns finds formal
support in our bound.}

\begin{theorem}[Iterated Signal Depth Bound (11.18)]
\label{thm:iterated-signal}
Composing a signal with itself $n$ times yields depth at most $n \cdot \depth(s)$:
\[
  \depth(\compN{s}{n}) \;\le\; n \cdot \depth(s).
\]
\end{theorem}

\begin{proof}[Proof sketch]
By the depth bound on $\operatorname{composeN}$, which follows from induction
and subadditivity.
\end{proof}

Repetition in ritual---chanting a mantra $n$ times, reciting a creed at every
service, replaying a national anthem daily---has bounded interpretive effect.
The subadditivity of depth means repetition is inherently
\emph{anti-inflationary}: you cannot create unbounded meaning through mere
repetition.  This constrains theories of ritual efficacy that posit cumulative
spiritual deepening through sheer repetition.  The effect of the $n$-th
repetition is bounded by $\depth(s)$, not by $n \cdot \depth(s)$.  The mantra
may be powerful, but it is not infinitely so.

\begin{lstlisting}[caption={Iterated signal depth bound in Lean~4}]
theorem iterated_signal_bound (s : I) (n : ℕ) :
    depth (composeN s n) <= n * depth s :=
  depth_composeN s n
\end{lstlisting}

% ----------------------------------------------------------------
\section{Implications for Media Ecology}
\label{sec:ch11-media}
% ----------------------------------------------------------------

The formal results of this chapter, while derived in the abstract setting of
ideatic space, have direct and substantive implications for the interdisciplinary
field of media ecology---the study of how media environments shape human
perception, cognition, and social organization.  We conclude the main body of
the chapter by drawing out four such implications.

\subsection{The Structural Limits of Propaganda}

The payoff bound (Theorem~\ref{thm:payoff-bound}) and the marginal depth
contribution bound (Theorem~\ref{thm:marginal-bound}) together establish a
hard ceiling on the efficacy of propaganda.  Consider a state actor emitting
a propaganda signal $s_{\text{prop}}$ to a population with repertoire
$\mathbf{R}$.  The total interpretive impact is at most
$\operatorname{depthSum}(\mathbf{R}) + |\mathbf{R}| \cdot
\depth(s_{\text{prop}})$.  This bound is independent of the propagandist's
skill, resources, or institutional power.  It depends only on the audience's
pre-existing depth and the signal's intrinsic complexity.

This result is both sobering and reassuring.  It is sobering because it shows
that audiences with deep repertoires---educated, culturally rich
populations---are capable of extracting \emph{more} depth from propaganda, not
less.  The bound increases with $\operatorname{depthSum}(\mathbf{R})$.  It is
reassuring because it shows that propaganda cannot create depth \emph{ex
nihilo}: the ``big lie,'' repeated indefinitely, is bounded by
$n \cdot \depth(s_{\text{lie}})$ (Theorem~\ref{thm:iterated-signal}), not
by $n^2$ or $e^n$.  Repetition is subadditive.

Jacques Ellul's distinction between ``propaganda of agitation'' and
``propaganda of integration'' finds formal expression here.  Agitation
propaganda ($\depth(s)$ is low, targeting a large $|\mathbf{R}|$) achieves
breadth at the expense of depth.  Integration propaganda ($\depth(s)$ is
high, targeting a resonant $\mathbf{R}$) achieves depth at the expense of
breadth.  Neither can escape the payoff bound.\footnote{Jacques Ellul,
\textit{Propagandes} (1962).}

\subsection{Filter Bubbles and Echo Chambers}

The homogeneous audience theorem (Theorem~\ref{thm:homogeneous-audience})
provides a formal account of the ``filter bubble'' phenomenon identified by
Eli Pariser.  When an algorithmic recommender system selects content for a user
whose repertoire elements are mutually resonant ($r_i \res r_j$ for all
$r_i, r_j \in \mathbf{R}_u$), any signal $s$ produces mutually resonant
interpretations: $(r_i \compose s) \res (r_j \compose s)$.  The user's
interpretive world becomes \emph{more} resonant with each exposure, not less.

The formal mechanism is clear: resonance is preserved under composition
(right context), so a resonant repertoire exposed to any signal remains
resonant.  The ``bubble'' is not an artifact of the algorithm but a
structural consequence of resonance preservation.  The algorithm merely
\emph{accelerates} a dynamics that is inherent in the algebra of
ideatic space.  To ``burst'' a filter bubble, one would need to introduce
repertoire elements that are \emph{not} resonant with the existing
repertoire---a structural intervention, not merely an algorithmic one.

\subsection{The Attention Economy as Bounded Competition}

The modern attention economy, in which content creators compete for a finite
resource (human attention), maps directly onto our formal framework.  The
receiver's attention budget corresponds to the finiteness of $\mathbf{R}$:
a receiver can attend to only finitely many signals, and each interpretation
$r \compose s$ consumes interpretive resources.  The total payoff across all
attended signals is bounded by $\operatorname{depthSum}(\mathbf{R}) +
|\mathbf{R}| \cdot \max_s \depth(s)$.

This bound has implications for platform design.  A platform that presents
$k$ signals to a receiver with repertoire of size $n$ achieves total payoff
at most $n \cdot k$ times the sum of average depths.  Increasing $k$ (showing
more content) does not increase payoff per item; it merely spreads the
receiver's interpretive capacity more thinly.  This is the formal basis of
the ``attention scarcity'' that drives the economics of digital platforms:
the scarce resource is not information (signals are abundant) but
interpretive capacity (repertoire is bounded).

\subsection{Cultural Biodiversity and Marketplace Health}

A marketplace with many diverse senders ($|\mathbf{S}|$ is large, senders are
pairwise non-resonant) offers a richer competitive environment than a
marketplace dominated by a few resonant senders.  The full cultural resonance
theorem (Theorem~\ref{thm:full-resonance}) shows that resonant senders produce
resonant interpretations: if $s_1 \res s_2$, then their interpretations in any
receiver are also resonant.  A marketplace dominated by resonant senders thus
produces a \emph{homogeneous} interpretive landscape---many signals, but all
``saying the same thing'' in the resonance sense.

Cultural biodiversity---the analogue of biological diversity in ecological
theory---requires non-resonant senders.  A healthy marketplace, in our
formalism, is one in which the sender pool $\mathbf{S}$ spans a large portion
of ideatic space, ensuring that different receivers with different repertoires
find signals that resonate with their particular backgrounds.  The marketplace
metaphor thus acquires an ecological dimension: just as a healthy ecosystem
requires species diversity, a healthy idea marketplace requires sender
diversity.  Monopolistic control over the sender pool---whether by a state, a
church, or a platform---reduces cultural biodiversity and, with it, the
potential for deep interpretation across the full range of receiver backgrounds.

% ----------------------------------------------------------------
\section{Conclusion}
\label{sec:ch11-conclusion}
% ----------------------------------------------------------------

The marketplace of ideas, when formalized in IDT, reveals a surprising
structural rigidity.  Competition between senders is not anarchic but
constrained: void never wins, dominance is reflexive, payoff is bounded by
audience capacity, composition is subadditive, and resonant signals produce
resonant interpretations regardless of competitive context.

These results formalize intuitions that social theorists have long articulated
in qualitative terms.  Gramsci's hegemony requires active cultural production
(Theorem~\ref{thm:void-dom-identity}); Bourdieu's cultural capital constrains
reception (Theorem~\ref{thm:payoff-bound}); Redfield's Great/Little Tradition
unity rests on resonance preservation (Theorem~\ref{thm:full-resonance});
Illich's diminishing returns in education follow from marginal bounds
(Theorem~\ref{thm:marginal-bound}).

In the next chapter, we turn from competition to \emph{cooperation}: stable
matching of senders and receivers based on cultural compatibility, formalizing
the two-sided matching structures that underlie kinship, marriage exchange,
and interfaith dialogue.

% ================================================================
% Chapter 12: Stable Matching and Cultural Compatibility
% ================================================================

\chapter{Stable Matching and Cultural Compatibility}
\label{ch:matching}

\epigraph{The prohibition of incest is less a rule prohibiting marriage with
the mother, sister, or daughter, than a rule obliging the mother, sister, or
daughter to be \emph{given} to others.  It is the supreme rule of the gift.}%
{Claude Lévi-Strauss, \textit{Les Structures élémentaires de la parenté} (1949)}

% ----------------------------------------------------------------
\section{Introduction}
% ----------------------------------------------------------------

Two-sided matching is ubiquitous in human social organization.  Marriage
exchange---the foundational institution of kinship---pairs lineages on the
basis of compatibility.  Student--mentor pairing in academic apprenticeship,
patron--client relationships in Mediterranean and Mesoamerican societies,
and interfaith dialogue all involve matching ``senders'' with ``receivers''
on the basis of \emph{cultural compatibility}.

In Ideatic Diffusion Theory, compatibility reduces to resonance.  Two ideas
are compatible when they evoke each other; a match is fruitful when the
composition of the paired ideas achieves significant depth.  This chapter
formalizes these notions and proves eighteen theorems that connect matching
theory to Lévi-Strauss's kinship structures, Malinowski's Kula ring,
Mauss's gift economy, and Bourdieu's marriage strategies.

The central contribution is a suite of theorems showing that compatibility is
reflexive, symmetric, and preserved under composition---precisely the
algebraic properties that make stable matching \emph{possible}.  Without
these properties, no stable assignment of senders to receivers could exist.
Lévi-Strauss's intuition that kinship is a ``total social fact'' finds its
formal expression in the closure of compatibility under the monoid
operation.\footnote{Lévi-Strauss's notion of kinship as \emph{fait social
total} (following Mauss) implies that every aspect of social life---economic,
ritual, juridical---is implicated in the matching structure.  Our algebraic
closure results formalize this totality.}

\subsection{Historical Foundations of Matching Theory}

The mathematical theory of stable matching originates with David Gale and
Lloyd Shapley's seminal 1962 paper, ``College Admissions and the Stability
of Marriage.''  Gale and Shapley proved that for any set of men and women
with complete and transitive preference orderings, a stable matching
exists---one in which no unmatched pair mutually prefers each other to their
assigned partners.  Their deferred acceptance algorithm, elegant in its
simplicity, showed that stability is not merely a theoretical possibility but
a constructively achievable outcome.

The practical significance of stable matching was recognized by the Nobel
Committee in 2012, when Lloyd Shapley (jointly with Alvin Roth) received the
Sveriges Riksbank Prize in Economic Sciences.  Roth's contribution was to
demonstrate that stable matching theory could be applied to real-world
institutional design: the National Resident Matching Program (NRMP) for
medical residencies, school choice mechanisms in New York and Boston, and
kidney exchange programs.  Roth's work showed that abstract mathematical
properties---stability, strategy-proofness, efficiency---have direct
consequences for human welfare.\footnote{Alvin E.\ Roth, ``The Economist as
Engineer: Game Theory, Experimentation, and Computation as Tools for Design
Economics,'' \textit{Econometrica} 70(4), 2002, pp.\ 1341--1378.}

The anthropological study of matching long predates its mathematical
formalization.  Claude Lévi-Strauss's \textit{Les Structures élémentaires
de la parenté} (1949) is, in essence, a theory of marriage matching.
Lévi-Strauss distinguished two fundamental types of exchange:
\emph{restricted exchange} (échange restreint), in which two groups
directly swap marriage partners, and \emph{generalized exchange} (échange
généralisé), in which three or more groups form a directed cycle of
partner transfer.  In restricted exchange, if clan A gives women to clan B,
then clan B gives women to clan A---a bilateral matching.  In generalized
exchange, A gives to B, B gives to C, C gives to A---a cyclic matching
that cannot be decomposed into bilateral pairs.

Lévi-Strauss's analysis was deeply influenced by Marcel Mauss's theory of
the gift (\textit{Essai sur le don}, 1925).  Mauss argued that exchange is
the foundation of social solidarity: the obligation to give, to receive,
and to reciprocate constitutes the ``total social fact'' (fait social total)
that binds individuals into society.  For Lévi-Strauss, marriage exchange is
the paradigmatic instance of Mauss's triple obligation: lineages are obliged
to give women (the prohibition of incest), to receive women (exogamy), and
to reciprocate (the alliance system).  Our formal notion of compatibility---
defined as resonance between sender and receiver---captures the structural
precondition for this exchange: two lineages can exchange only if their
totemic symbols resonate.

The cross-cousin marriage systems that preoccupied Lévi-Strauss provide
particularly elegant illustrations of matching theory.  In \emph{Dravidian
kinship} (exemplified by many South Indian and Australian Aboriginal
societies), the marriageable category is defined by a binary partition of
kin into ``parallel'' (same-sex sibling links) and ``cross'' (opposite-sex
sibling links) relatives.  Ego must marry a cross-cousin---specifically,
in the most restricted systems, a bilateral cross-cousin (mother's
brother's daughter, who is also father's sister's daughter).  This system
generates a stable matching between two exogamous moieties, each moiety
constituting a ``sender'' pool for the other.

The Kariera system of Western Australia, analyzed by A.\ R.\
Radcliffe-Brown in 1913, provides the simplest four-section system:
society is divided into four named sections, and marriage rules specify
exactly which section may marry which.  The Kariera system is equivalent
to a bipartite graph with two moieties and a deterministic matching rule
within each generation.  In our framework, the four sections correspond
to four ideas $a_1, a_2, a_3, a_4 \in \IS$, with compatibility
$a_i \res a_j$ holding if and only if $i$ and $j$ belong to
marriageable sections.  The reflexivity, symmetry, and composition-preservation
of compatibility (Theorems~\ref{thm:self-compat}--\ref{thm:comp-preserves})
ensure that the Kariera system is algebraically well-defined across
generations.

Modern matching markets extend these ancient patterns into new domains.
Dating applications (Tinder, Hinge, Bumble) are literal matching markets
in which individuals present signals (profiles) and express compatibility
judgments (swipes).  The underlying algorithms---while proprietary---are
descendants of the Gale-Shapley framework, adapted for incomplete
information and stochastic preferences.  In our formalism, a dating app
user's profile is a signal $s \in \IS$, their preferences constitute a
repertoire $\mathbf{R}$, and a ``match'' occurs when mutual compatibility
$s_A \res s_B$ holds.  The reflexivity of compatibility
(Theorem~\ref{thm:self-compat}) guarantees that every user is trivially
compatible with themselves---the formal reason that self-matching must be
explicitly excluded by platform design.

Organ donation matching, formalized by Roth, Sönmez, and Ünver (2004),
provides another modern instance.  Here, the ``ideas'' are biological: a
donor kidney's immunological profile is a signal, a recipient's antibody
profile is a repertoire element, and compatibility is immunological
crossmatch compatibility.  The algebraic properties we prove in this
chapter---reflexivity, symmetry, composition preservation---do \emph{not}
generally hold for immunological compatibility (crossmatch can be
asymmetric), which is precisely why kidney exchange requires more
sophisticated algorithms than simple stable matching.  Our framework thus
identifies the structural properties that distinguish culturally-mediated
matching from biologically-constrained matching.

Assortative mating---the tendency for individuals to mate with similar
partners---has been extensively documented in sociology.  Educational
homogamy (the tendency to marry partners of similar educational attainment)
has increased dramatically in the United States and Western Europe since
the 1960s, with profound consequences for income inequality.\footnote{Robert
D.\ Mare, ``Five Decades of Educational Assortative Mating,''
\textit{American Sociological Review} 56(1), 1991, pp.\ 15--32.}
In our framework, assortative mating corresponds to a preference for
matches in which sender and receiver have similar depth:
$|\depth(s) - \depth(r)| \le \epsilon$ for some threshold $\epsilon$.
The composition quality bound (Theorem~\ref{thm:quality-bound}) implies
that assortatively matched pairs achieve quality close to
$2 \cdot \depth(a)$ (when $\depth(s) \approx \depth(r) \approx \depth(a)$),
while disassortatively matched pairs (one deep, one shallow) achieve
quality bounded by $\depth(\text{deep}) + \depth(\text{shallow})$---a
larger bound but with diminishing marginal returns from the shallow
partner's contribution.

% ----------------------------------------------------------------
\section{Core Definitions}
\label{sec:ch12-defs}
% ----------------------------------------------------------------

\begin{definition}[Compatibility]
\label{def:compatible}
Two ideas $a, b \in \IS$ are \emph{compatible} when they resonate:
\[
  \operatorname{compatible}(a, b)
  \;\;\stackrel{\text{def}}{=}\;\;
  a \res b.
\]
In kinship theory, this is the analogue of \emph{marriageability}: two
lineages are compatible when their totemic symbols evoke each other.
\end{definition}

\begin{definition}[Match Pair]
\label{def:match-pair}
A \emph{match pair} is a structure $\langle \text{sender}, \text{receiver} \rangle$
pairing a sending idea (gift, signal, myth) with a receiving idea (background,
habitus, totem).  A match pair is \emph{compatible} when sender and receiver
resonate.
\end{definition}

\begin{lstlisting}[caption={MatchPair structure in Lean~4}]
structure MatchPair (I : Type*) [IdeaticSpace I] where
  sender   : I
  receiver : I

def MatchPair.isCompatible (p : MatchPair I) : Prop :=
  compatible p.sender p.receiver
\end{lstlisting}

\begin{definition}[Composition Quality]
\label{def:comp-quality}
The \emph{composition quality} of a sender--receiver pair is the depth of
their composition:
\[
  \operatorname{quality}(s, r) \;=\; \depth(s \compose r).
\]
Higher quality means deeper cultural synthesis.
\end{definition}

% ----------------------------------------------------------------
\section{Compatibility Fundamentals}
\label{sec:ch12-compat}
% ----------------------------------------------------------------

\begin{theorem}[Self-Compatibility (12.1)]
\label{thm:self-compat}
Every idea is compatible with itself: $\forall\, a \in \IS,\; a \res a$.
\end{theorem}

\begin{proof}[Proof sketch]
By reflexivity of resonance.
\end{proof}

Endogamy---marrying within one's own group---is always ``compatible'' in the
formal sense.  Lévi-Strauss argued that exogamy is culturally
\emph{mandated} precisely because endogamy is always the structural default.
The prohibition of incest is not the prohibition of something impossible but
of something \emph{too easy}: the self-match is always available, and culture
must intervene to prevent it.

\begin{remark}[Connection to Graph Theory and Bipartite Matching]
The compatibility relation $\res$, being reflexive
(Theorem~\ref{thm:self-compat}) and symmetric
(Theorem~\ref{thm:compat-symm}), defines an undirected graph on the
ideatic space $\IS$: the \emph{compatibility graph} $G_{\res} = (\IS, E)$
where $\{a, b\} \in E$ iff $a \res b$.  Reflexivity means every vertex
has a self-loop; symmetry means every edge is undirected.  A stable
matching in the sense of Gale and Shapley is then a maximum matching in
this graph, subject to preference orderings.

When the compatibility graph is bipartite---as in moiety-based marriage
systems, where society is divided into two exogamous halves---the matching
problem reduces to bipartite matching, for which efficient polynomial-time
algorithms exist (e.g., the Hopcroft-Karp algorithm).  The four-section
systems (Kariera, Aranda) correspond to bipartite graphs with additional
structure.  The eight-section systems analyzed by Lévi-Strauss correspond
to more complex graph structures whose matching properties depend on the
specific compatibility rules encoded in the kinship terminology.

This graph-theoretic perspective connects our algebraic framework to the
extensive literature on matching in combinatorial optimization, including
Hall's marriage theorem, König's theorem, and the Birkhoff-von Neumann
theorem on doubly stochastic matrices.  Each of these classical results
has an anthropological interpretation when the vertices of the graph are
cultural ideas and the edges are resonance relations.
\end{remark}

\begin{theorem}[Compatibility is Symmetric (12.2)]
\label{thm:compat-symm}
If $a$ is compatible with $b$, then $b$ is compatible with $a$:
\[
  a \res b \;\Longrightarrow\; b \res a.
\]
\end{theorem}

\begin{proof}[Proof sketch]
By symmetry of resonance.
\end{proof}

In Lévi-Strauss's \emph{restricted exchange}, if clan A can marry into clan B,
then clan B can marry into clan A.  The symmetry of marriage exchange---the
\emph{connubium}---is a structural consequence of resonance symmetry.  This
rules out asymmetric kinship systems in which marriageability is one-directional,
a constraint that Lévi-Strauss himself noted distinguishes restricted from
generalized exchange.\footnote{In \emph{generalized exchange}, asymmetry
\emph{is} allowed---clan A gives to B, B gives to C, C gives to A---but
the pairwise compatibility remains symmetric; only the \emph{flow} of
gifts is directed.}

\begin{example}[The Kula Ring: Circular Exchange and Compatibility]
Bronisław Malinowski's study of the Kula ring in the Trobriand Islands
(\textit{Argonauts of the Western Pacific}, 1922) provides a classic
illustration of matching based on compatibility.  In the Kula system,
ceremonial objects circulate among island communities in a fixed ring:
\emph{soulava} (necklaces) travel clockwise, \emph{mwali} (armbands)
travel counterclockwise.  Each exchange is between two partners---a
sender and a receiver---who are bound by a long-term Kula relationship.

In our framework, each Kula partner brings a cultural repertoire shaped
by their island's traditions, their rank within the chieftainship
hierarchy, and their personal exchange history.  A successful Kula
exchange requires compatibility: the sender's object (signal) must
resonate with the receiver's expectations (repertoire element).  A
\emph{soulava} of great prestige resonates differently in the hands of a
paramount chief ($\depth(r_{\text{chief}} \compose s_{\text{soulava}})$
is high) than in the hands of a commoner
($\depth(r_{\text{commoner}} \compose s_{\text{soulava}})$ is lower).

The symmetry of compatibility (Theorem~\ref{thm:compat-symm}) is
instantiated in the Kula by the norm of reciprocity: if chief A finds
chief B a worthy Kula partner (compatible), then chief B finds chief A
equally worthy.  The exchange relationship is necessarily bilateral, even
though the objects flow in a single direction at any given moment.
Malinowski emphasized that the Kula partnership, once established, is
``one of the permanent and overriding features of the natives'
life''---a stable match in both the colloquial and the technical sense.
\end{example}

\begin{theorem}[Void Self-Compatibility (12.3)]
\label{thm:void-self-compat}
Void is compatible with itself: $\void \res \void$.
\end{theorem}

\begin{proof}[Proof sketch]
An instance of self-compatibility (Theorem~\ref{thm:self-compat}).
\end{proof}

The null exchange---two parties exchanging nothing---is trivially compatible.
This is the structural zero-point of all reciprocity, the ``degree zero'' of
Mauss's gift economy.  It is the baseline from which all actual exchange
departs.

\begin{theorem}[Absorbed Void Preserves Compatibility (12.4)]
\label{thm:absorbed-void}
If $a \res b$, then $(a \compose \void) \res b$.
\end{theorem}

\begin{proof}[Proof sketch]
Since $a \compose \void = a$ by the void axiom, this is immediate.
\end{proof}

Appending silence to a cultural signal does not change its compatibility.
A ritual followed by silence is the same ritual---the ``afterglow'' of
ceremony is structurally inert.  The Durkheimian \emph{effervescence
collective} dissipates, but the signal's compatibility endures.

% ----------------------------------------------------------------
\section{Composition and Compatibility}
\label{sec:ch12-comp}
% ----------------------------------------------------------------

\begin{theorem}[Composition Preserves Compatibility (12.5)]
\label{thm:comp-preserves}
If $(a, a')$ and $(b, b')$ are compatible pairs, then their compositions
are also compatible:
\[
  a \res a' \;\wedge\; b \res b'
  \quad\Longrightarrow\quad
  (a \compose b) \res (a' \compose b').
\]
\end{theorem}

\begin{proof}[Proof sketch]
By the full composition-preservation of resonance
(\texttt{resonance\_compose}).
\end{proof}

If two couples are each internally compatible, their ``children'' (compositions)
are also compatible with each other.  Cultural compatibility is
\emph{heritable} through ideatic composition.  This is Lévi-Strauss's
\emph{atom of kinship}: the minimal structure of kinship---two
allied families producing compatible offspring---is formally guaranteed by
this theorem.\footnote{Lévi-Strauss introduced the ``atom of kinship'' in
his 1945 article ``Structural Analysis in Linguistics and in Anthropology.''
The atom comprises a man, his wife, their child, and the wife's brother: four
positions related by three relations (descent, marriage, avunculate).  Our
theorem captures the compositional closure that makes this atom stable.}

\begin{example}[Indian Caste Endogamy and Jati Compatibility]
The Indian caste system, particularly the institution of
\emph{jati} (birth-based sub-caste), provides one of the most elaborate
and long-lived matching systems in human history.  Each jati constitutes a
closed endogamous unit: marriage is permitted only within one's own jati,
or, in some reform movements, within a limited set of ``compatible'' jatis.

In our framework, each jati $j$ corresponds to an idea in $\IS$, encoding
the ritual status, occupational identity, dietary practices, and kinship
norms of the group.  Endogamy is the requirement that matching occurs only
between compatible ideas: $j_1 \res j_2$ holds only when $j_1$ and $j_2$
belong to marriageable categories.  Traditionally, the strictest form of
endogamy requires $j_1 = j_2$---self-matching---which is always guaranteed
by Theorem~\ref{thm:self-compat}.

The composition-preservation theorem (Theorem~\ref{thm:comp-preserves})
has a specific implication for the Indian system: if two jatis are
compatible and another pair of jatis is compatible, then the composed
cultural identities of their offspring are also compatible.  This ensures
the \emph{heritability} of caste compatibility across generations---precisely
the structural property that has allowed the jati system to persist for
millennia.  The system is self-reinforcing because compatibility, once
established, is algebraically guaranteed to reproduce itself through
composition.

B.\ R.\ Ambedkar's critique of caste endogamy---articulated in his 1916
paper ``Castes in India: Their Mechanism, Genesis and Development''---can be
read as a critique of this very algebraic closure.  Ambedkar argued that
caste is maintained by ``endogamy superimposed on exogamy'': the
prohibition of inter-caste marriage creates closed compatibility classes,
while the mandate of exogamy within the caste ensures that the compatibility
class reproduces itself.  Our formal framework captures exactly this
mechanism: self-compatibility guarantees the endogamous baseline, and
composition preservation guarantees hereditary closure.
\end{example}

\begin{theorem}[Right-Context Compatibility (12.6)]
\label{thm:right-context}
Compatible senders remain compatible when placed in the same right context:
\[
  a \res b
  \quad\Longrightarrow\quad
  (a \compose c) \res (b \compose c).
\]
\end{theorem}

\begin{proof}[Proof sketch]
By right-composition preservation of resonance.
\end{proof}

If two myths resonate, adding the same ritual context to both preserves their
resonance.  Shared ritual frames maintain cultural compatibility---which is
why liturgical uniformity (the Latin Mass, the Pali Canon) serves as a
\emph{resonance stabilizer} across geographically dispersed communities.

\begin{theorem}[Left-Context Compatibility (12.7)]
\label{thm:left-context}
Compatible ideas remain compatible when preceded by the same context:
\[
  a \res b
  \quad\Longrightarrow\quad
  (c \compose a) \res (c \compose b).
\]
\end{theorem}

\begin{proof}[Proof sketch]
By left-composition preservation of resonance.
\end{proof}

A shared cultural preamble---educational background, initiation rite---preserves
the resonance between ideas.  Preliminary socialization ensures continued
compatibility, which is why initiation rites (Masonic degrees, academic
comprehensive exams, military basic training) serve as \emph{compatibility
guarantors}: by imposing a common left-context $c$, they ensure that
subsequent ideas $a, b$ remain resonant in the initiates' minds.

\begin{example}[Academic Peer Review as Cultural Matching]
The peer review system in academic publishing instantiates a matching
problem with direct parallels to our formal framework.  When a journal
editor receives a manuscript (sender signal $s$), they must match it with
reviewers whose expertise (receiver repertoire) is compatible.  A paper
on Bayesian phylogenetics must be matched with a reviewer who can produce
a deep interpretation---$\depth(r_{\text{expert}} \compose s_{\text{paper}})$
is high---rather than a reviewer whose interpretation is shallow or null.

The left-context compatibility result (Theorem~\ref{thm:left-context})
applies directly: if two papers $a, b$ on related topics resonate
($a \res b$), then a reviewer with any background $c$ produces resonant
interpretations of both: $(c \compose a) \res (c \compose b)$.  This is
why reviewers are typically assigned clusters of related papers---the
resonance between papers ensures consistent evaluation.

The right-context result (Theorem~\ref{thm:right-context}) applies to
reviewer training: if two reviewers $c_1, c_2$ share a common disciplinary
background (are both ``preceded by'' the same graduate training $c$), then
their compatibility with any given paper is preserved.  The
standardization of doctoral training within a discipline serves as the
``shared left-context'' that ensures reviewer compatibility with the
discipline's literature.
\end{example}

% ----------------------------------------------------------------
\section{Quality Bounds}
\label{sec:ch12-quality}
% ----------------------------------------------------------------

\begin{theorem}[Composition Quality Bound (12.8)]
\label{thm:quality-bound}
The quality of any match is bounded by the sum of depths:
\[
  \operatorname{quality}(s, r) \;\le\; \depth(s) + \depth(r).
\]
\end{theorem}

\begin{proof}[Proof sketch]
Direct application of the subadditivity axiom.
\end{proof}

The depth of cultural synthesis---marriage, trade, dialogue---cannot exceed
the combined depth of the participants.  Two shallow cultures produce a
shallow synthesis.  This constrains utopian visions of ``fusion'': the
Esperanto of cultures is bounded above by the complexity of the cultures
being fused.

\begin{theorem}[Void Match Quality, Right (12.9)]
\label{thm:void-quality-right}
Matching with void yields one's own depth: $\operatorname{quality}(a, \void) = \depth(a)$.
\end{theorem}

\begin{proof}[Proof sketch]
$a \compose \void = a$ by the void axiom.
\end{proof}

The ``null exchange''---giving silence, receiving nothing---reveals one's own
depth.  This is the formal basis of introspection: in the absence of external
input, one sees only one's own complexity.  The contemplative
traditions---monastic silence, Zen \emph{zazen}, Quaker worship---formalize
this structural fact as spiritual practice.

\begin{theorem}[Void Match Quality, Left (12.10)]
\label{thm:void-quality-left}
The void sender reveals the receiver's depth: $\operatorname{quality}(\void, a) = \depth(a)$.
\end{theorem}

\begin{proof}[Proof sketch]
$\void \compose a = a$ by the void axiom.
\end{proof}

A culture that sends nothing (isolationism) reveals only what the receiver
already has.  The ``mirror effect'' of cultural isolation: the isolated
society, receiving no external signals, can only see its own reflection.

\begin{theorem}[Void-Void Match (12.11)]
\label{thm:void-void}
The quality of matching void with void is zero: $\operatorname{quality}(\void, \void) = 0$.
\end{theorem}

\begin{proof}[Proof sketch]
$\void \compose \void = \void$, and $\depth(\void) = 0$.
\end{proof}

When neither party brings anything, nothing is produced.  The ``empty
exchange'' has zero cultural value---the structural null point of all social
interaction.  This is Hobbes's state of nature, formalized: before the
social contract introduces non-void ideas, cultural quality is identically
zero.

% ----------------------------------------------------------------
\section{Match Pair Properties}
\label{sec:ch12-pairs}
% ----------------------------------------------------------------

\begin{theorem}[Compatible Pairs are Self-Evident (12.12)]
\label{thm:self-match}
The match pair $(a, a)$ is always compatible.
\end{theorem}

\begin{proof}[Proof sketch]
By self-compatibility (Theorem~\ref{thm:self-compat}).
\end{proof}

Self-matching (narcissistic culture) is always structurally valid, even if
socially prohibited.  The incest taboo in Lévi-Strauss's theory exists
precisely to prevent the structurally trivial but socially destructive
self-match.  Without the taboo, every lineage would choose itself---the
path of least structural resistance.

\begin{theorem}[Void Pair is Compatible (12.13)]
\label{thm:void-pair}
The match pair $(\void, \void)$ is compatible.
\end{theorem}

\begin{proof}[Proof sketch]
By self-compatibility of void.
\end{proof}

The ``null marriage''---two culturally empty entities exchanging nothing---is
trivially compatible.  This is the structural zero of kinship: before culture
introduces non-void ideas, all matches are vacuously valid.

\begin{theorem}[Composition Quality Associativity (12.14)]
\label{thm:quality-assoc}
Quality is associative under composition:
\[
  \operatorname{quality}((a \compose b), c)
  \;=\;
  \operatorname{quality}(a, (b \compose c)).
\]
\end{theorem}

\begin{proof}[Proof sketch]
Both sides unfold to $\depth((a \compose b) \compose c) = \depth(a \compose (b \compose c))$
by associativity.
\end{proof}

In a three-party exchange (mediator between two clans), the total quality is
invariant under regrouping.  The ``middle-man'' can be absorbed into either
side without changing the structural outcome---a formal justification for the
role of the \emph{go-between} in cross-cultural negotiation.

\begin{remark}[Open Questions on Dynamic Compatibility]
The matching theory developed in this chapter is fundamentally
\emph{static}: compatibility is a fixed relation on $\IS$, and match
quality is a fixed function of depth.  In practice, cultural compatibility
is dynamic---it evolves over time as matched pairs interact, adapt, and
transform each other.  Several open questions arise:

\begin{enumerate}
\item \textbf{Compatibility drift.}  Can two initially compatible ideas
  become incompatible through repeated composition?  Our framework precludes
  this (composition preserves compatibility by Theorem~\ref{thm:comp-preserves}),
  but this may reflect a limitation of the model rather than a feature of
  cultural reality.  Marriages do fail; alliances do dissolve.  A richer
  model might introduce a dynamic compatibility relation
  $\res_t$ that evolves with time.

\item \textbf{Matching under incomplete information.}  In the
  Gale-Shapley framework, agents know their complete preference orderings.
  In cultural matching, compatibility is often discovered only through
  interaction---one does not know whether two ideas resonate until one
  composes them.  A model incorporating costly compatibility testing
  (``courtship'') would connect to the economics of search and inspection
  goods.

\item \textbf{Many-to-many matching.}  Our match pairs are dyadic, but
  many cultural institutions involve polyadic matching: a student matched
  with multiple mentors, a polycultural society integrating many traditions
  simultaneously.  Extending the framework to $k$-tuples rather than pairs
  would require generalizing the compatibility and quality notions.

\item \textbf{Stability under perturbation.}  Is a stable matching robust
  to small changes in the compatibility relation?  If one pair's resonance
  is disrupted (e.g., by cultural change), does the entire matching unravel,
  or can it be locally repaired?  This connects to the literature on robust
  matching in mechanism design.
\end{enumerate}
\end{remark}

% ----------------------------------------------------------------
\section{Structural Matching Theorems}
\label{sec:ch12-structural}
% ----------------------------------------------------------------

\begin{theorem}[Double Composition Quality Bound (12.15)]
\label{thm:double-bound}
The quality of matching two compound ideas is quadruply bounded:
\[
  \operatorname{quality}((a \compose b), (c \compose d))
  \;\le\;
  \depth(a) + \depth(b) + \depth(c) + \depth(d).
\]
\end{theorem}

\begin{proof}[Proof sketch]
Apply subadditivity to the outer composition, then to each inner
composition, and combine via transitivity.
\end{proof}

The synthesis of two compound ideas---the alliance of two already-complex
cultural systems---cannot exceed the total depth of all components.  Colonial
encounters between complex civilizations (Aztec--Spanish, Chinese--British)
produce syntheses bounded by the sum of all four sub-cultural depths involved.

\begin{theorem}[Iterated Matching Depth (12.16)]
\label{thm:iterated-match}
Matching a signal with itself $n$ times via $\compN{a}{n}$ yields depth at
most $n \cdot \depth(a)$.
\end{theorem}

\begin{proof}[Proof sketch]
By the iterated composition depth bound (Theorem~\ref{thm:iterated-signal}
from Chapter~\ref{ch:marketplace}).
\end{proof}

Repeated cultural exchange---annual rituals, recurring trade---has bounded
cumulative depth.  The Kula ring's repeated exchanges are subadditive: each
cycle adds at most $\depth(\text{gift})$ to the cultural synthesis.
Malinowski's observation that the Kula values lie not in the objects
themselves but in the circulation is consistent with our bound: the \emph{same}
gift, circulated $n$ times, yields at most $n \cdot \depth(\text{gift})$,
not exponentially growing prestige.\footnote{Bronisław Malinowski,
\textit{Argonauts of the Western Pacific} (1922).}

\begin{theorem}[Synthesis Self-Compatibility (12.17)]
\label{thm:synthesis-self}
If $a$ and $b$ are compatible, then their composition $a \compose b$ is
compatible with itself.
\end{theorem}

\begin{proof}[Proof sketch]
Every idea is self-compatible (Theorem~\ref{thm:self-compat}).
\end{proof}

A culturally valid marriage (compatible pair) produces offspring (composition)
that is internally consistent.  The synthesis is never self-contradictory
in the resonance sense---a formal guarantee of the coherence of
cross-cultural offspring.

\begin{theorem}[Exchange Symmetry of Quality (12.18)]
\label{thm:exchange-symm}
If $a \res b$, then both $(a \compose b)$ and $(b \compose a)$ are
self-compatible:
\[
  a \res b
  \;\Longrightarrow\;
  (a \compose b) \res (a \compose b)
  \;\wedge\;
  (b \compose a) \res (b \compose a).
\]
\end{theorem}

\begin{proof}[Proof sketch]
Both directions follow from self-compatibility.
\end{proof}

In Mauss's \textit{Essai sur le don}, the gift and counter-gift are
structurally parallel but not identical.  We have $a \compose b \ne b \compose a$
in general---composition is not commutative---formalizing the asymmetry inherent
in reciprocal exchange.  The gift carries the giver's stamp; the counter-gift
carries the recipient's.  Yet both compositions are self-consistent: neither the
gift nor the counter-gift contradicts itself, even if they differ from each
other.  Reciprocity does not require identity---only mutual
self-consistency.\footnote{Marcel Mauss, \textit{Essai sur le don} (1925).
The \emph{hau}---the spirit of the gift that compels return---need not be
identical to the original gift, only resonant with it.}

\begin{lstlisting}[caption={Exchange symmetry in Lean~4}]
theorem exchange_compositions_self_compatible {a b : I}
    (h : compatible a b) :
    compatible (compose a b) (compose a b) /\
    compatible (compose b a) (compose b a) :=
  <self_compatible _, self_compatible _>
\end{lstlisting}

\begin{remark}[Connection to Chapter~\ref{ch:marketplace}: From Competition to Cooperation]
Chapters~\ref{ch:marketplace} and~\ref{ch:matching} are structural
complements.  In Chapter~\ref{ch:marketplace}, multiple senders compete
for a single receiver's interpretive engagement; the central question is
\emph{which signal dominates?}  In this chapter, senders and receivers are
matched pairwise; the central question is \emph{which pairings are stable?}

The formal connection is precise.  The dominance relation of
Chapter~\ref{ch:marketplace}---$s_1$ dominates $s_2$ on repertoire
$\mathbf{R}$ when
$\depth(r \compose s_1) \ge \depth(r \compose s_2)$ for all
$r \in \mathbf{R}$---can be interpreted as a preference ordering: the
receiver ``prefers'' $s_1$ to $s_2$.  When this preference ordering is
used to construct a Gale-Shapley instance, the resulting stable matching
is the \emph{equilibrium of the marketplace}.  Competition
(Chapter~\ref{ch:marketplace}) determines preferences; matching
(this chapter) determines allocation.

This connection illuminates a fundamental tension in cultural systems.  The
marketplace rewards \emph{breadth}---signals that achieve non-zero depth
across many receivers.  But matching rewards \emph{depth}---pairings that
achieve maximal composition quality for specific pairs.  A society that
optimizes for marketplace competition (broadcasting uniform signals) may
sacrifice matching quality (deep dyadic synthesis), and vice versa.  The
tension between broadcast media (marketplace logic) and intimate
relationships (matching logic) is a structural feature of ideatic space,
not merely a sociological observation.
\end{remark}

\subsection{Matching, Solidarity, and Cultural Reproduction}

The eighteen theorems of this chapter collectively establish that the
algebraic structure of ideatic space is \emph{hospitable to matching}: the
properties of reflexivity, symmetry, composition-preservation, and bounded
quality ensure that stable pairings exist, persist across generations, and
produce bounded but non-trivial cultural synthesis.  This hospitality is
not accidental.  The IDT axioms were chosen, in part, because they
guarantee exactly these matching-theoretic properties.  The axioms encode,
at the most abstract level, the structural preconditions for the
reciprocal exchange that Mauss identified as the foundation of social life.

Durkheim's distinction between \emph{mechanical solidarity} (solidarity
based on similarity) and \emph{organic solidarity} (solidarity based on
complementarity) maps onto two extreme regimes of our matching framework.
In mechanical solidarity, all ideas are mutually resonant
($a \res b$ for all $a, b \in \IS$), and every match is trivially
compatible.  The compatibility graph is complete, and any matching is stable.
In organic solidarity, compatibility is selective---$a \res b$ holds only
for specific pairs---and matching requires careful allocation.  The
transition from mechanical to organic solidarity, which Durkheim associated
with the division of labor, corresponds formally to the sparsification of
the compatibility graph: fewer edges, more constrained matching, more
complex equilibria.

The cultural reproduction of matching systems---their persistence across
generations---is guaranteed by composition-preservation
(Theorem~\ref{thm:comp-preserves}).  If this generation's matches are
compatible, then the ``offspring'' (compositions) of those matches are also
compatible with each other.  The matching system reproduces itself through
algebraic closure.  This is the formal mechanism underlying what Bourdieu
called \emph{reproduction}: the tendency of social structures to perpetuate
themselves through the very practices they engender.\footnote{Pierre
Bourdieu, \textit{La Reproduction} (1970), co-authored with Jean-Claude
Passeron.}

% ----------------------------------------------------------------
\section{Conclusion}
\label{sec:ch12-conclusion}
% ----------------------------------------------------------------

The matching theory developed in this chapter reveals that cultural
compatibility---defined as resonance---possesses exactly the algebraic
properties needed for stable pairing.  Compatibility is reflexive
(Theorem~\ref{thm:self-compat}), symmetric
(Theorem~\ref{thm:compat-symm}), and preserved under composition
(Theorem~\ref{thm:comp-preserves}).  These are not accidental features but
structural necessities: without them, the kinship systems catalogued by
Lévi-Strauss, the gift economies described by Mauss, and the patron--client
networks analyzed by Bourdieu would lack the algebraic foundation to persist
across generations.

Quality bounds (Theorem~\ref{thm:quality-bound}) ensure that no match
produces unbounded cultural synthesis, and void-match results
(Theorems~\ref{thm:void-quality-right}--\ref{thm:void-void}) establish the
baseline against which all exchange is measured.  In the next chapter, we
move from dyadic matching to \emph{coalitional} composition: the
cooperative game theory of cultural groups, formalized through Durkheim's
concept of solidarity.

% ================================================================
% Chapter 13: Coalition Value and Totemic Solidarity
% ================================================================

\chapter{Coalition Value and Totemic Solidarity}
\label{ch:coalition}

\epigraph{The totality of beliefs and sentiments common to the average members
of a society forms a determinate system with a life of its own; one may call it
the \emph{collective} or \emph{common consciousness.}}{Émile Durkheim,
\textit{De la division du travail social} (1893)}

% ----------------------------------------------------------------
\section{Introduction}
% ----------------------------------------------------------------

Human groups are not mere collections of individuals---they are
\emph{coalitions} that compose their ideas jointly.  Durkheim's
\emph{mechanical solidarity} arises when members share a
\emph{conscience collective}: a jointly composed idea-system deeper than
any individual's contribution.  The totem, in Durkheim's analysis of
Australian religions, is precisely this: the symbol that expresses the
group's collective composition, the sign under which individual ideas are
subsumed into a shared cultural complex.\footnote{Durkheim,
\textit{Les Formes élémentaires de la vie religieuse} (1912).  The totem
is not merely a label but the material expression of the clan's
\emph{conscience collective}---what we formalize as the coalition composition.}

This chapter formalizes cooperative game theory in ideatic space.  We
introduce \emph{coalition composition}---the fold of a list of ideas via
the monoid operation, starting from void---and \emph{coalition value},
the depth of the joint composition.  We define \emph{marginal contribution},
the ideatic analogue of the Shapley marginal value, and prove eighteen
theorems showing that coalition formation is subadditive, void members
contribute nothing, associativity makes coalition merging well-defined, and
resonant coalitions produce resonant joint ideas.

The structural lesson is austere: there is no free lunch in cultural
synthesis.  The combined depth of a coalition never exceeds the sum of
individual depths.  The genius raises the group by at most their own
complexity; the free rider contributes nothing; and a coalition of voids
has zero cultural value regardless of its size.

% ----------------------------------------------------------------
\subsection{Historical Context: From Game Theory to Solidarity}
% ----------------------------------------------------------------

The mathematical foundations of coalition theory originate in John von~Neumann
and Oskar Morgenstern's \textit{Theory of Games and Economic Behavior}
(1944), which introduced cooperative game theory as a framework for analyzing
situations where binding agreements among players are possible.  Von~Neumann
and Morgenstern's key insight was that when agents can form coalitions, the
relevant strategic unit is no longer the individual but the \emph{group}:
the coalition's aggregate payoff determines each member's bargaining position.
The \emph{characteristic function} $v: 2^N \to \mathbb{R}$ assigns to each
subset~$S$ of the player set~$N$ a value $v(S)$ representing the total payoff
that coalition~$S$ can guarantee itself.  Our coalition value function is
the ideatic analogue: instead of monetary payoffs, it measures the cultural
depth that a group of idea-bearers can collectively achieve.\footnote{Von
Neumann and Morgenstern, \textit{Theory of Games and Economic Behavior}
(Princeton, 1944).  The cooperative game framework they developed has been
applied across economics, political science, and operations research; our
application to cultural composition is, to our knowledge, the first
systematic formalization in this direction.}

Lloyd Shapley's celebrated \emph{value} (1953) provides a unique allocation
of the coalition's total payoff to individual members, satisfying axioms of
efficiency, symmetry, linearity, and the null-player property.  The
\emph{Shapley value} of player~$i$ in a cooperative game with characteristic
function~$v$ is:
\[
  \phi_i(v) = \sum_{S \subseteq N \setminus \{i\}}
  \frac{|S|!\;(|N| - |S| - 1)!}{|N|!}
  \bigl[v(S \cup \{i\}) - v(S)\bigr].
\]
This formula averages the marginal contribution of player~$i$ across all
possible orderings in which the coalition might form.  Our marginal
contribution (Definition~\ref{def:marginal}) captures the same intuition in
the ideatic setting: the depth that an individual adds to a coalition depends
on when they join and what is already composed.  Shapley's null-player axiom
is our Theorem~\ref{thm:void-marginal}: a player who contributes nothing to
every coalition receives zero allocation.\footnote{Lloyd~S.\ Shapley, ``A
Value for $n$-Person Games,'' \textit{Annals of Mathematics Studies}~28
(1953), pp.\ 307--317.  For the connection between marginal contributions
and fair allocation, see also Alvin~E.\ Roth, ed., \textit{The Shapley
Value} (Cambridge, 1988).}

The \emph{core} of a cooperative game---the set of allocations under which
no coalition has an incentive to deviate---provides a complementary notion of
stability.  An allocation is in the core if and only if every coalition
receives at least its characteristic function value.  In ideatic terms,
a cultural arrangement is stable when every possible sub-group feels
adequately represented in the collective composition.  When the core is
empty, no stable arrangement exists, and the culture is prone to fission:
discontented sub-coalitions will break away to form their own collective
representations.  The subadditivity of our coalition value
(Theorem~\ref{thm:pair-subadd}) guarantees that the ``grand coalition''---the
entire society---always has a non-empty core, since no sub-coalition can
achieve more than the sum of its members' individual depths.

These game-theoretic tools find their anthropological counterpart in
Émile Durkheim's theory of solidarity.  In \textit{De la division du travail
social} (1893), Durkheim distinguished two ideal types of social cohesion.
\emph{Mechanical solidarity} characterizes segmentary societies---those
organized into structurally identical, self-sufficient units (clans,
lineages, moieties).  Here, solidarity arises from \emph{resemblance}:
members share a strong conscience collective, a common body of beliefs and
sentiments that integrates them through their very similarity.  The
collective representation is the fold of many similar contributions---our
coalition composition of ideas that largely resonate with one another.
\emph{Organic solidarity}, by contrast, characterizes complex societies
organized through the division of labor.  Here, solidarity arises from
\emph{complementarity}: members contribute different, specialized ideas
whose composition yields a collective representation richer than any
individual could produce alone.  The depth of the coalition exceeds any
singleton's depth precisely because the ideas are diverse.\footnote{Durkheim,
\textit{De la division du travail social} (Paris, 1893).  The distinction
between mechanical and organic solidarity maps onto our framework as follows:
mechanical solidarity corresponds to coalitions of resonant (similar) ideas,
while organic solidarity corresponds to coalitions of complementary
(non-resonant, high-depth) ideas.  Both are subadditive, but organic
solidarity tends to approach the subadditivity bound more closely.}

Durkheim's analysis of totemism in \textit{Les Formes élémentaires de la
vie religieuse} (1912) provides the ethnographic grounding for our
formalization.  Among the Australian Aboriginal peoples---particularly the
Aranda (Arrernte) of Central Australia and the Kariera of Western
Australia---social organization is structured by elaborate section systems
that partition the population into named groups, each associated with a
totemic emblem drawn from the natural world (kangaroo, emu, witchetty grub,
rain, etc.).  The Aranda system employs eight sections, arranged in a
pattern of prescribed marriage classes and moiety divisions.  The Kariera
system uses four sections.  In both cases, the totem functions as the
\emph{symbol} of the group's coalition composition: it is the material sign
under which individual contributions are subsumed into a collective
representation.  Durkheim argued that the totem is not merely a label or
taxonomic convenience but the very \emph{expression} of the group's social
force---the conscience collective made visible.\footnote{On Australian
section systems, see A.~P.~Elkin, \textit{The Australian Aborigines} (1938);
W.~Lloyd Warner, \textit{A Black Civilization} (1937); and more recently,
Patrick McConvell and Barry Alpher, ``The Formal Analysis of Kin Avoidance
Language in Australia,'' in \textit{Kinship Systems} (2002).}

The formal framework developed in this chapter also connects to contemporary
applications that Durkheim could not have anticipated.  Wikipedia, the
collaborative encyclopedia, exemplifies coalition composition at scale:
each editor contributes an idea (an edit, a paragraph, a citation) that is
composed with the existing article to produce a collective representation.
The article's ``depth''---its comprehensiveness, accuracy, and
sophistication---is the coalition value of the editor community.
Open-source software development follows a similar logic: each contributor's
code is composed with the existing codebase through version control, and the
project's complexity is the coalition value of the development team.  In
both cases, subadditivity holds: the project's depth never exceeds the sum
of all individual contributions, though it may fall far short of that
bound due to redundancy, conflicts, and the overhead of coordination.

Social capital theory, as developed by James Coleman, Robert Putnam, and
Pierre Bourdieu, provides a further bridge between cooperative game theory
and anthropological analysis.\footnote{James~S.\ Coleman, ``Social Capital
in the Creation of Human Capital,'' \textit{American Journal of Sociology}
94 (1988), pp.\ S95--S120; Robert~D.\ Putnam, \textit{Bowling Alone}
(2000); Pierre Bourdieu, ``The Forms of Capital,'' in J.~Richardson, ed.,
\textit{Handbook of Theory and Research for the Sociology of Education}
(1986).}  For Coleman, social capital inheres in the structure of
relations among persons---precisely the coalition structure we formalize.
For Putnam, social capital manifests as ``bridging'' (between diverse
groups) and ``bonding'' (within homogeneous groups), a distinction that maps
onto our framework: bonding capital corresponds to coalitions of resonant
ideas (mechanical solidarity), while bridging capital corresponds to
coalitions of diverse ideas (organic solidarity).  For Bourdieu, capital
exists in multiple forms---economic, cultural, social, symbolic---and is
convertible across forms through the habitus.  Our depth metric captures
one dimension of Bourdieu's cultural capital; the composition operation
captures the habitus's role in mediating between individual and collective
representations.

% ----------------------------------------------------------------
\section{Core Definitions}
\label{sec:ch13-defs}
% ----------------------------------------------------------------

\begin{definition}[Coalition Composition]
\label{def:coalition-compose}
The \emph{coalition composition} of a list of ideas $[a_1, a_2, \dots, a_n]$
is defined recursively:
\[
  \operatorname{coalitionCompose}([]) = \void, \qquad
  \operatorname{coalitionCompose}(a :: \text{rest}) =
    \operatorname{coalitionCompose}(\text{rest}) \compose a.
\]
This models a group's \emph{collective idea}---the joint composition of all
members' contributions.  In Durkheim's terms, this \emph{is} the conscience
collective.
\end{definition}

\begin{lstlisting}[caption={Coalition composition in Lean~4}]
def coalitionCompose : List I -> I
  | []         => void
  | a :: rest  => compose (coalitionCompose rest) a
\end{lstlisting}

Note the right-to-left fold: the earliest member in the list contributes
last, layered atop the previous members' composition.  This models the
anthropological fact that latecomers to a tradition interpret through---and
are interpreted against---the existing collective representation.

\begin{definition}[Coalition Value]
\label{def:coalition-value}
The \emph{coalition value} of a group is the depth of its joint composition:
\[
  \operatorname{coalitionValue}(\text{coalition}) =
  \depth(\operatorname{coalitionCompose}(\text{coalition})).
\]
\end{definition}

\begin{definition}[Marginal Contribution]
\label{def:marginal}
The \emph{marginal contribution} of idea $a$ to a coalition is the
difference in coalition value when $a$ is added:
\[
  \operatorname{marginal}(a, C) =
  \operatorname{coalitionValue}(a :: C) - \operatorname{coalitionValue}(C).
\]
This is the ideatic analogue of the Shapley marginal contribution in
cooperative game theory.
\end{definition}

% ----------------------------------------------------------------
\section{Empty and Singleton Coalitions}
\label{sec:ch13-empty}
% ----------------------------------------------------------------

\begin{theorem}[Empty Coalition is Void (13.1)]
\label{thm:empty-void}
The empty coalition composes to void: $\operatorname{coalitionCompose}([]) = \void$.
\end{theorem}

\begin{proof}[Proof sketch]
By definition of the base case.
\end{proof}

A group with no members has no culture.  This is the structural zero of
social organization---the ``state of nature'' in Hobbes's thought
experiment, but formalized as the void of ideatic composition.  Before the
first idea is contributed, the collective representation is empty.

\begin{theorem}[Empty Coalition Has Zero Value (13.2)]
\label{thm:empty-zero}
The cultural complexity of the empty group is zero:
$\operatorname{coalitionValue}([]) = 0$.
\end{theorem}

\begin{proof}[Proof sketch]
$\operatorname{coalitionCompose}([]) = \void$, and $\depth(\void) = 0$.
\end{proof}

\emph{Ex nihilo nihil fit}---from nothing, nothing comes.  A society begins
with zero cultural depth and must build it through composition.  The
Hobbesian ``war of all against all'' is not merely unpleasant; it is
culturally sterile.

\begin{theorem}[Singleton Coalition (13.3)]
\label{thm:singleton-coal}
A one-member coalition's composition is $\void \compose a$:
$\operatorname{coalitionCompose}([a]) = \void \compose a$.
\end{theorem}

\begin{proof}[Proof sketch]
By unfolding the recursive definition once.
\end{proof}

\begin{theorem}[Singleton Coalition Value (13.4)]
\label{thm:singleton-value}
A one-member coalition has value equal to that member's depth:
$\operatorname{coalitionValue}([a]) = \depth(a)$.
\end{theorem}

\begin{proof}[Proof sketch]
$\void \compose a = a$ by the void axiom, so
$\depth(\void \compose a) = \depth(a)$.
\end{proof}

The isolated individual contributes exactly their own depth to culture---no
more, no less.  Robinson Crusoe's cultural footprint is precisely his personal
complexity.  This theorem anchors the individual as the unit of cultural
measurement: before composition with others, each person's contribution is
exactly what they carry.\footnote{This should be contrasted with Durkheim's
anti-individualist position, which holds that the conscience collective is
\emph{emergent} and irreducible to individual contributions.  Our theorem
shows that at the singleton level, there is no emergence---emergence
requires at least two members.}

% ----------------------------------------------------------------
\section{Void Members and Inert Contributions}
\label{sec:ch13-void}
% ----------------------------------------------------------------

\begin{theorem}[Adding Void Preserves Coalition Composition (13.5)]
\label{thm:void-member}
Adding void to the front of a coalition does not change its composition:
\[
  \operatorname{coalitionCompose}(\void :: C) = \operatorname{coalitionCompose}(C).
\]
\end{theorem}

\begin{proof}[Proof sketch]
$\operatorname{coalitionCompose}(C) \compose \void =
\operatorname{coalitionCompose}(C)$ by the right-identity axiom.
\end{proof}

A ``member'' who brings nothing---the free rider, the passive observer---does
not alter the group's collective idea.  Durkheim's distinction between
mechanical and organic solidarity relies on each member \emph{contributing};
void members are structurally invisible.  This is the formal basis of
Olson's collective action problem: individuals who contribute nothing receive
the benefits of the collective good without increasing it.

\begin{theorem}[Void Member Has Zero Marginal Contribution (13.6)]
\label{thm:void-marginal}
Adding void to a coalition changes nothing about coalition value:
$\operatorname{marginal}(\void, C) = 0$.
\end{theorem}

\begin{proof}[Proof sketch]
Since $\operatorname{coalitionCompose}(\void :: C) =
\operatorname{coalitionCompose}(C)$ by Theorem~\ref{thm:void-member}, the
difference is zero.
\end{proof}

The free rider contributes nothing.  This formalizes Olson's central theorem
in \textit{The Logic of Collective Action}: non-contributors do not enhance
the collective good.\footnote{Mancur Olson, \textit{The Logic of Collective
Action} (1965).  Olson's argument that rational individuals will not
contribute to public goods unless compelled is captured by our result that
void contributions have zero marginal value.}

% ----------------------------------------------------------------
\section{Two-Member Coalitions}
\label{sec:ch13-dyad}
% ----------------------------------------------------------------

\begin{theorem}[Two-Member Coalition (13.7)]
\label{thm:two-member}
A two-member coalition $[a, b]$ has composition
$(\void \compose b) \compose a$.
\end{theorem}

\begin{proof}[Proof sketch]
By unfolding the definition twice.
\end{proof}

\begin{theorem}[Two-Member Coalition Simplified (13.8)]
\label{thm:two-member-simp}
The two-member coalition $[a, b]$ simplifies to $b \compose a$.
\end{theorem}

\begin{proof}[Proof sketch]
$\void \compose b = b$ by the void axiom, so
$(\void \compose b) \compose a = b \compose a$.
\end{proof}

The dyad---pair bond, marriage, master--apprentice---is the simplest
non-trivial coalition.  Its collective idea is simply the composition of
both members' ideas.  Simmel's insight that the dyad is qualitatively
different from larger groups is confirmed algebraically: the dyad is the
only coalition whose value reduces to a single composition, without nested
structure.\footnote{Georg Simmel, ``The Number of Members as Determining the
Sociological Form of the Group'' (1902).  Simmel argued that the dyad lacks
the super-individual structure that emerges in triads and larger groups.
Our simplification theorem confirms this: $[a, b]$ reduces to $b \compose a$,
a simple composition with no nesting.}

\begin{lstlisting}[caption={Two-member coalition simplified in Lean~4}]
theorem two_member_simplified (a b : I) :
    coalitionCompose [a, b] = compose b a := by
  simp [coalitionCompose, void_left]
\end{lstlisting}

% ----------------------------------------------------------------
\section{Subadditivity and Value Bounds}
\label{sec:ch13-bounds}
% ----------------------------------------------------------------

\begin{theorem}[Coalition Value Subadditivity for Pairs (13.9)]
\label{thm:pair-subadd}
The value of a two-member coalition is bounded by the sum of individual
depths:
\[
  \operatorname{coalitionValue}([a, b]) \;\le\; \depth(a) + \depth(b).
\]
\end{theorem}

\begin{proof}[Proof sketch]
By Theorem~\ref{thm:two-member-simp},
$\operatorname{coalitionValue}([a, b]) = \depth(b \compose a)
\le \depth(b) + \depth(a)$ by subadditivity.
\end{proof}

Two cultures merging cannot produce a joint culture deeper than the sum of
their parts.  Cultural synthesis is bounded---there is no ``$1 + 1 = 3$'' in
ideatic composition.  This constrains utopian visions of cultural fusion and
places a hard ceiling on the ``synergy'' of intercultural dialogue.  The
deepest possible synthesis of Shakespeare and Bashō is at most
$\depth(\text{Shakespeare}) + \depth(\text{Bashō})$.

\begin{example}[The Amish Barn Raising]
\label{ex:barn-raising}
The Amish barn raising exemplifies coalition composition in practice.
Each participant brings a specific skill---carpentry, roofing, masonry---and
the collective output is the composed structure.  The coalition value
(the completed barn) exceeds what any individual could achieve alone, yet
Theorem~\ref{thm:pair-subadd} ensures it cannot exceed the sum of all
individual depths.  The barn raising also illustrates
Theorem~\ref{thm:void-marginal}: community members who attend but
contribute nothing (void members) do not increase the coalition's cultural
output, though they may serve important social functions not captured by
the depth metric.

More precisely, let $a_1, a_2, \dots, a_n$ denote the ideas (skills,
techniques, knowledge) brought by the $n$ participants.  The barn as a
completed structure is $\operatorname{coalitionCompose}([a_1, \dots, a_n])$,
and its complexity---the coalition value---is the depth of this composition.
Subadditivity gives $\operatorname{coalitionValue}([a_1, \dots, a_n])
\le \sum_{i=1}^n \depth(a_i)$: the barn is at most as complex as the sum
of all individual skills.  In practice, the inequality is strict because
skills overlap (multiple carpenters), and the coordination overhead
(scheduling, communication) effectively reduces the realized depth below
the theoretical maximum.\footnote{Donald Kraybill, \textit{The Riddle of Amish
Culture} (2001), documents how barn raisings serve simultaneously as
economic cooperation and social bonding rituals.  The event reinforces
Durkheim's mechanical solidarity: participants share a common value system
(Ordnung) that makes their individual contributions resonant.}
\end{example}

\begin{remark}[Connection to Olson's Collective Action Problem]
\label{rem:olson}
Mancur Olson's \textit{The Logic of Collective Action} (1965) identifies
a fundamental tension in group behavior: rational self-interested agents
will under-contribute to public goods because they can free-ride on
others' contributions.  Theorem~\ref{thm:void-marginal} provides the
formal skeleton of this insight.  A free rider whose contribution is void
adds nothing to the coalition value, yet benefits from the collective
representation.  However, our framework also reveals a limit to Olson's
pessimism.  In ideatic space, the free rider is \emph{structurally
invisible}: they neither help nor harm the coalition's depth.  The
collective good---the coalition composition---is robust to the presence
of void members.  The problem is not that free riders destroy value but
that they fail to create it, leaving the coalition at a depth below its
potential maximum.  Olson's solution---selective incentives to encourage
participation---corresponds in our framework to mechanisms that convert
void members into non-void contributors, thereby increasing the realized
coalition value toward its theoretical bound.
\end{remark}

\begin{theorem}[Marginal Contribution Bounded by Depth (13.10)]
\label{thm:marginal-bounded}
Adding member $a$ to any coalition increases value by at most $\depth(a)$:
\[
  \operatorname{coalitionValue}(a :: C) \;\le\;
  \operatorname{coalitionValue}(C) + \depth(a).
\]
\end{theorem}

\begin{proof}[Proof sketch]
$\operatorname{coalitionCompose}(a :: C) =
\operatorname{coalitionCompose}(C) \compose a$, and
$\depth(\operatorname{coalitionCompose}(C) \compose a)
\le \depth(\operatorname{coalitionCompose}(C)) + \depth(a)$
by subadditivity.
\end{proof}

No individual can contribute more than their own depth to a group.  The
genius raises the group by at most their own complexity---the \emph{bounded
marginal returns} of individual talent in collective settings.  This is the
formal counterpart to the sociological observation that even extraordinary
individuals (Marx, Darwin, Freud) shift the collective culture by a
bounded amount proportional to their personal depth.

\begin{theorem}[Marginal Contribution at Most Depth (13.11)]
\label{thm:marginal-le-depth}
The marginal contribution of $a$ is at most $\depth(a)$:
\[
  \operatorname{marginal}(a, C) \;\le\; \depth(a).
\]
\end{theorem}

\begin{proof}[Proof sketch]
Since $\operatorname{coalitionValue}(a :: C) \le
\operatorname{coalitionValue}(C) + \depth(a)$ by
Theorem~\ref{thm:marginal-bounded}, the difference
$\operatorname{coalitionValue}(a :: C) - \operatorname{coalitionValue}(C)
\le \depth(a)$.
\end{proof}

An individual's marginal contribution to culture is bounded by their personal
depth.  No initiate brings more to a secret society than they carry.  This
constrains the ``great man'' theory of history: even the greatest individual's
marginal impact is bounded above.

\begin{example}[Wikipedia as Coalition Composition]
\label{ex:wikipedia}
Wikipedia's collaborative editing process provides a striking modern
illustration of coalition composition.  Each editor contributes an idea---a
paragraph of text, a citation, a structural reorganization, a factual
correction---that is composed with the existing article.  The article at
any given moment is $\operatorname{coalitionCompose}([e_1, e_2, \dots, e_n])$
where $e_i$ is the $i$-th edit in the article's history.  The article's
quality and comprehensiveness---its coalition value---is the depth of
this composition.

Several of our theorems find direct expression in Wikipedia's dynamics.
Theorem~\ref{thm:singleton-value} shows that a single editor's article
has exactly their own depth: one-person articles reflect one person's
knowledge.  Theorem~\ref{thm:pair-subadd} shows that merging two stub
articles produces an article no deeper than the sum of the stubs' depths.
Theorem~\ref{thm:void-marginal} explains why many minor edits (fixing
typos, adjusting formatting) do not substantially increase article depth:
these are near-void contributions whose marginal value is negligible.
And Theorem~\ref{thm:void-list} explains why articles created by bots
that contribute no substantive content remain at zero depth regardless
of their edit count.

The Wikipedia model also highlights a phenomenon not fully captured by
our current theorems: \emph{edit wars}, in which competing editors
repeatedly overwrite each other's contributions.  In our framework, an
edit war is a sequence of compositions and de-compositions (reversions)
that leaves the coalition value oscillating rather than monotonically
increasing.  The depth metric is indifferent to the social drama; it
tracks only the net compositional state.\footnote{For a sociological
analysis of Wikipedia's collaborative dynamics, see Yochai Benkler,
\textit{The Wealth of Networks} (2006), ch.~3; and Aaron Halfaker et al.,
``The Rise and Decline of an Open Collaboration System,''
\textit{American Behavioral Scientist} 57:5 (2013), pp.\ 664--688.}
\end{example}

\begin{remark}[Emergent Properties and the Limits of Subadditivity]
\label{rem:emergence}
Subadditivity (Theorem~\ref{thm:pair-subadd}) states that the coalition
value cannot \emph{exceed} the sum of individual depths.  But it says nothing
about how close the coalition value comes to that bound.  In practice,
coalition compositions exhibit a wide range of behaviors: some approach
the additive bound closely (when contributions are orthogonal and
non-redundant), while others fall far short (when contributions overlap or
interfere).

Durkheim's concept of \emph{emergence}---the idea that the collective
consciousness is qualitatively different from, and irreducible to, the sum
of individual consciousnesses---finds a precise location in our framework.
Emergence, in our terms, is the phenomenon whereby
$\operatorname{coalitionValue}(C) \ne \sum_{a \in C} \depth(a)$.  The
\emph{gap} between the coalition value and the additive bound is the
``emergence deficit''---the amount by which composition falls short of
mere aggregation.  When this deficit is small, the coalition is ``nearly
additive'' and emergence is weak.  When it is large, interactions between
members substantially reshape the collective idea, and emergence is strong.
The subadditivity theorem guarantees that emergence is always
\emph{non-positive}: coalitions never produce more than the sum of their
parts, though they may produce qualitatively different compositions that
are, in a structural sense, ``more than'' the mere concatenation of
individual ideas.
\end{remark}

% ----------------------------------------------------------------
\section{Associativity and Merging}
\label{sec:ch13-assoc}
% ----------------------------------------------------------------

\begin{theorem}[Coalition Composition Unfold (13.12)]
\label{thm:coal-unfold}
$\operatorname{coalitionCompose}(a :: \text{rest}) =
\operatorname{coalitionCompose}(\text{rest}) \compose a$.
\end{theorem}

\begin{proof}[Proof sketch]
This is definitional---the recursive case of \texttt{coalitionCompose}.
\end{proof}

Adding a new member to a coalition means composing their idea with the
existing collective idea.  Initiation rites formalize this structural
operation: the initiate's individual idea is composed with the group's
collective representation to produce a new, slightly altered collective.

\begin{theorem}[Three-Member Coalition (13.13)]
\label{thm:three-member}
The three-member coalition $[a, b, c]$ has composition
$((\void \compose c) \compose b) \compose a$.
\end{theorem}

\begin{proof}[Proof sketch]
By three applications of the recursive case.
\end{proof}

\begin{theorem}[Three-Member Simplified (13.14)]
\label{thm:three-member-simp}
The three-member coalition simplifies to $(c \compose b) \compose a$.
\end{theorem}

\begin{proof}[Proof sketch]
$\void \compose c = c$ by the void axiom, then unfold.
\end{proof}

The triad---Simmel's ``third''---introduces mediation.  The three-member
coalition's collective idea involves nested composition: the mediator
(middle member) connects the other two.  This is why Simmel considered the
triad qualitatively different from the dyad: the triad has internal
structure (nested composition) that the dyad lacks.  Formally, $[a, b]$
is a flat composition $b \compose a$, while $[a, b, c]$ is a nested
composition $(c \compose b) \compose a$---the nesting is the formal
signature of mediation.

\begin{example}[Jazz Improvisation as Coalition Composition]
\label{ex:jazz}
A jazz ensemble---say, the Miles Davis Quintet circa 1965, with Davis
(trumpet), Wayne Shorter (saxophone), Herbie Hancock (piano), Ron Carter
(bass), and Tony Williams (drums)---provides a vivid illustration of
coalition composition in real time.  Each musician contributes an idea
(a melodic line, a harmonic voicing, a rhythmic pattern) that is
composed with the ongoing collective sound.  The ensemble's music at
any moment is $\operatorname{coalitionCompose}([m_1, m_2, m_3, m_4, m_5])$
where $m_i$ is the $i$-th musician's contribution.

The triad theorem (Theorem~\ref{thm:three-member-simp}) is directly
relevant: any three musicians in the quintet form a trio whose collective
idea is a nested composition $(c \compose b) \compose a$.  The nesting
matters: the rhythm section (bass and drums) provides the compositional
foundation upon which the harmonic instrument (piano) layers, and the
soloists (trumpet, saxophone) compose atop this pre-existing structure.
The order of composition is not arbitrary---it reflects the functional
hierarchy of the ensemble.

Subadditivity (Theorem~\ref{thm:pair-subadd}) explains why even the
greatest ensemble cannot exceed the sum of its members' individual
virtuosity.  But the jazz ensemble also illustrates the \emph{emergence
deficit} discussed in Remark~\ref{rem:emergence}: the ensemble's music is
qualitatively different from a mere superposition of five solo
performances.  Interaction reshapes each contribution, producing a
coalition composition whose character cannot be predicted from individual
depths alone.  The depth of ``Kind of Blue'' is bounded by but not equal to
$\depth(\text{Davis}) + \depth(\text{Coltrane}) + \depth(\text{Evans}) +
\depth(\text{Chambers}) + \depth(\text{Cobb})$.\footnote{For a
musicological analysis of jazz ensemble interaction, see Ingrid Monson,
\textit{Saying Something: Jazz Improvisation and Interaction} (1996);
and Paul Berliner, \textit{Thinking in Jazz} (1994).}
\end{example}

% ----------------------------------------------------------------
\section{Resonance in Coalitions}
\label{sec:ch13-resonance}
% ----------------------------------------------------------------

\begin{theorem}[Self-Resonance of Coalition Value (13.15)]
\label{thm:coal-self-res}
Every coalition's collective idea resonates with itself:
\[
  \operatorname{coalitionCompose}(C) \res \operatorname{coalitionCompose}(C).
\]
\end{theorem}

\begin{proof}[Proof sketch]
By reflexivity of resonance.
\end{proof}

Every group's collective representation is self-consistent in the resonance
sense.  A culture never fails to ``recognize itself.''  Durkheim's
conscience collective is always internally coherent: the collective
representation may be shallow or deep, simple or complex, but it never
contradicts itself at the level of resonance.\footnote{This should not be
confused with logical consistency.  Resonance is weaker than logical
entailment: two ideas can resonate without being logically compatible.
The point is that the collective idea always \emph{recognizes} itself,
even if it contains contradictions by other metrics.}

\begin{theorem}[Resonant Members Yield Resonant Coalitions (13.16)]
\label{thm:resonant-singletons}
If $a \res a'$, then $\operatorname{coalitionCompose}([a]) \res
\operatorname{coalitionCompose}([a'])$.
\end{theorem}

\begin{proof}[Proof sketch]
Singleton coalitions simplify to the member itself (after void composition),
so the result follows from the hypothesis.
\end{proof}

If two individuals' ideas resonate---they ``get'' each other---their
one-person ``cultures'' are compatible.  This is the micro-foundation of
cultural affinity: interpersonal resonance at the individual level
translates directly to inter-cultural resonance at the singleton-coalition
level.

\begin{theorem}[Coalition Value of Void List (13.17)]
\label{thm:void-list}
A coalition consisting entirely of $n$ voids has zero value:
$\operatorname{coalitionValue}(\operatorname{replicate}(n, \void)) = 0$.
\end{theorem}

\begin{proof}[Proof sketch]
By induction on $n$: each void member preserves the existing composition
(Theorem~\ref{thm:void-member}), and the base case is the empty coalition
with value zero.
\end{proof}

A group of non-contributors produces no cultural depth---``a meeting of empty
minds.''  Bureaucracies without substance have zero cultural value regardless
of how many members they have.  The formal futility of scaling a null
organization: ten thousand voids compose to void.

\begin{lstlisting}[caption={Void coalition zero value in Lean~4}]
theorem void_coalition_zero_value : forall (n : ℕ),
    coalitionValue (List.replicate n void) = 0
  | 0     => by simp [coalitionValue, coalitionCompose, depth_void]
  | n + 1 => by
    simp [coalitionValue, coalitionCompose, List.replicate_succ, void_right]
    exact void_coalition_zero_value n
\end{lstlisting}

\begin{example}[Monastic Communities as Coalitions]
\label{ex:monastery}
The Benedictine monastery offers a counterpoint to the jazz ensemble:
a coalition designed to maximize collective depth through disciplined,
rule-governed composition.  The Rule of St.~Benedict prescribes a daily
schedule of prayer (\textit{opus Dei}), manual labor (\textit{labora}),
and sacred reading (\textit{lectio divina}) that structures each monk's
contribution to the communal life.  In our terms, each monk contributes
an idea shaped by the Rule, and the monastery's collective
representation---its spiritual culture---is the coalition composition
of these disciplined contributions.

Theorem~\ref{thm:void-list} is particularly illuminating here.  The
monastic tradition has always recognized the danger of \emph{acedia}---the
``noonday demon'' of spiritual torpor---which, in our framework,
corresponds to a monk whose contribution has become void.  A monastery
full of acedic monks has zero coalition value, regardless of its
population.  The elaborate monastic disciplines of self-examination,
confession, and spiritual direction are precisely mechanisms for detecting
and correcting void contributions before they accumulate.

The Benedictine emphasis on \textit{stabilitas loci}---the vow to remain
in one monastery for life---also finds formal expression.  By keeping the
coalition membership stable, the monastery ensures that each member's
contribution is composed with a familiar collective representation,
reducing the coordination overhead that plagues more fluid
coalitions.\footnote{On monastic organization as a form of social
engineering, see C.~H.~Lawrence, \textit{Medieval Monasticism} (3rd ed.,
2001); and Terrence Kardong, \textit{Benedict's Rule: A Translation and
Commentary} (1996).}
\end{example}

\begin{remark}[Connections to Chapter~\ref{ch:matching} and Matching Theory]
\label{rem:matching-connection}
The coalition theory developed in this chapter connects directly to the
matching theory of Chapter~\ref{ch:matching}.  There, we analyzed how
individual ideas are paired through the matching operation; here, we
study how those matched individuals compose into groups.  The connection
is precise: a coalition composition is a \emph{sequence} of pairwise
compositions, folded from right to left.  Every property of pairwise
composition (subadditivity, void identity, associativity) lifts to the
coalition level through induction on the list structure.

The matching-theoretic perspective suggests a natural refinement of
coalition value: instead of measuring only the depth of the grand
coalition's composition, one might also consider the depth profile of all
possible sub-coalitions---the \emph{characteristic function} of the
cooperative game.  Whether this richer structure yields additional
anthropological insights (e.g., about the stability of kinship systems
or the fission-fusion dynamics of band societies) remains an open question
for future work.
\end{remark}

\begin{theorem}[Coalition Value Monotonicity Bound (13.18)]
\label{thm:value-mono}
Adding any member to a coalition changes the value by at most the member's
depth:
\[
  \operatorname{coalitionValue}(a :: C) \;\le\;
  \operatorname{coalitionValue}(C) + \depth(a).
\]
\end{theorem}

\begin{proof}[Proof sketch]
This is Theorem~\ref{thm:marginal-bounded} restated.
\end{proof}

Cultural change is incremental.  Adding one person to a society can shift its
collective depth by at most that person's individual depth.  Revolutions
require deep individuals---shallow newcomers are structurally bounded in
their effect.  This is the formal counterpart to Kuhn's observation that
scientific revolutions require individuals of exceptional depth (Copernicus,
Newton, Einstein) to shift the collective paradigm.\footnote{Thomas Kuhn,
\textit{The Structure of Scientific Revolutions} (1962).  Kuhn's
``revolutionary scientist'' is, in our terms, an individual with
exceptionally high $\depth(a)$, whose marginal contribution is correspondingly
large---but still bounded.}

% ----------------------------------------------------------------
\section{The Anthropology of Collective Action}
\label{sec:ch13-collective-action}
% ----------------------------------------------------------------

The theorems of this chapter, taken together, provide a formal grammar for
the anthropological study of collective action.  We briefly survey how
this grammar illuminates several major themes in the anthropological
literature, drawing connections between the abstract algebra of coalition
composition and the concrete ethnography of human cooperation.

\subsection{Segmentary Lineage Systems}

Evans-Pritchard's classic analysis of the Nuer of southern Sudan (1940)
describes a \emph{segmentary lineage system} in which political coalitions
form and dissolve along genealogical lines.\footnote{E.~E.~Evans-Pritchard,
\textit{The Nuer} (Oxford, 1940).}  In the event of conflict, lineage
segments unite at successively higher levels of genealogical inclusion:
brothers unite against cousins, cousins against more distant kin, and
so on, following the principle ``I against my brother, my brother and I
against our cousin, our lineage against the world.''  In our framework,
each level of segmentary inclusion corresponds to a coalition composition
at a different scale.  The coalition value at each level is bounded by
subadditivity (Theorem~\ref{thm:pair-subadd}), and the marginal
contribution of each segment is bounded by its depth
(Theorem~\ref{thm:marginal-le-depth}).  The segmentary system is thus
a hierarchically nested sequence of coalition compositions, each
satisfying the structural constraints we have proven.

\subsection{Gift Exchange and Reciprocity}

Malinowski's analysis of the Kula ring (1922) and Mauss's comparative
theory of the gift (1925) both describe systems of reciprocal exchange
in which the act of giving creates social bonds---coalitions formed
through the exchange of culturally valued objects.\footnote{Bronisław
Malinowski, \textit{Argonauts of the Western Pacific} (1922); Marcel
Mauss, \textit{Essai sur le don} (1925).}  In Mauss's formulation, the
gift carries a portion of the giver's \emph{hau} (spiritual force),
which compels reciprocation.  In our terms, the gift is an idea with
positive depth, and giving it is an act of composition: the receiver's
cultural state is composed with the gift to produce a new state.  The
obligation to reciprocate ensures that the exchange is bidirectional,
producing a dyadic coalition (Theorem~\ref{thm:two-member-simp}) whose
value is at most the sum of both parties' depths.

\subsection{Collective Effervescence}

Durkheim's concept of \emph{collective effervescence}---the heightened
emotional energy generated by group ritual---maps onto our framework
as follows.  During collective ritual (the corroboree, the pilgrimage,
the political rally), individuals compose their ideas in rapid succession,
producing a coalition composition of exceptional density.  The emotional
intensity of the event reflects the rate of composition, not its ultimate
depth: subadditivity (Theorem~\ref{thm:pair-subadd}) still constrains
the final depth, but the subjective experience of rapid composition
creates the impression of transcendence.  Durkheim's observation that
participants feel themselves in the presence of a force ``greater than
themselves'' corresponds to the perception that the coalition value
exceeds any individual's depth---which it can, though it cannot exceed
the sum.

\subsection{Free Riding and Institutional Design}

The formal analysis of void members (Theorems~\ref{thm:void-member}
and~\ref{thm:void-marginal}) connects to a vast literature in
institutional economics on the design of mechanisms to prevent free
riding.  Elinor Ostrom's \textit{Governing the Commons} (1990) documents
how communities around the world have developed institutional rules---
monitoring, graduated sanctions, conflict-resolution mechanisms---to
ensure that members contribute to collective goods rather than free-riding
on others' contributions.\footnote{Elinor Ostrom, \textit{Governing the
Commons} (Cambridge, 1990).  Ostrom received the Nobel Prize in Economics
in 2009 for her work demonstrating that communities can self-organize to
manage common-pool resources without either privatization or state
regulation.}  In our framework, Ostrom's institutions are mechanisms for
converting void members into non-void contributors, thereby increasing
the coalition value.  The effectiveness of these institutions can be
measured by the gap between the \emph{actual} coalition value and the
\emph{counterfactual} value if all members contributed their full depth.

\subsection{Digital Communities and Online Collaboration}

The rise of digital platforms has created new forms of coalition
composition at unprecedented scale.  Open-source software projects
(Linux, Apache, Wikipedia) involve thousands of contributors whose
individual ideas are composed through version control systems and
collaborative editing interfaces.  The coalition value of these
projects---measured by code complexity, article comprehensiveness, or
feature richness---is the depth of the grand coalition's composition.
Our theorems apply directly: the project's depth is bounded by the sum
of all contributors' depths (subadditivity), non-contributing members
(lurkers, inactive accounts) add nothing (void marginal contribution),
and the project's depth grows at most linearly with the number of
substantive contributors (monotonicity bound).

These digital coalitions also exhibit phenomena that stretch our
framework in productive directions.  The ``1\% rule'' of online
communities---the empirical observation that roughly 1\% of users
create content, 9\% edit or curate, and 90\% passively consume---is a
natural consequence of Theorem~\ref{thm:void-list}: the 90\% of
passive consumers are void members whose presence does not increase the
coalition value.  Yet their presence may matter for other reasons
(audience size, network effects, advertising revenue) not captured by
the depth metric.\footnote{For the 1\% rule, see Ben McConnell and Jackie
Huba, ``The 1\% Rule,'' \textit{Church of the Customer Blog} (2006);
and Jakob Nielsen, ``Participation Inequality: Lurkers vs.\
Contributors in Internet Communities'' (2006).}

% ----------------------------------------------------------------
\section{Conclusion}
\label{sec:ch13-conclusion}
% ----------------------------------------------------------------

The cooperative game theory of ideatic space reveals a structural discipline
in coalition formation.  The empty coalition has zero value
(Theorem~\ref{thm:empty-zero}); singleton value equals individual depth
(Theorem~\ref{thm:singleton-value}); dyadic value is subadditive
(Theorem~\ref{thm:pair-subadd}); void members contribute nothing
(Theorem~\ref{thm:void-marginal}); and coalitions of voids are culturally
sterile (Theorem~\ref{thm:void-list}).

These results formalize Durkheim's insight that solidarity is not mere
aggregation but \emph{composition}---and composition is structurally
constrained.  The next chapter applies these coalition-theoretic tools to a
different institutional setting: the \emph{auction}, where ideas compete for
interpretive priority through depth-based bidding, connecting to the
anthropology of potlatch, prestige goods, and ritual economy.

% ================================================================
% Chapter 14: Ideatic Auctions and Ritual Economy
% ================================================================

\chapter{Ideatic Auctions and Ritual Economy}
\label{ch:auction}

\epigraph{Conspicuous consumption of valuable goods is a means of reputability
to the gentleman of leisure.}{Thorstein Veblen,
\textit{The Theory of the Leisure Class} (1899)}

% ----------------------------------------------------------------
\section{Introduction}
% ----------------------------------------------------------------

In many cultural contexts, ideas compete for interpretive priority through
\emph{depth-based bidding}.  The potlatch of the Pacific Northwest Coast, in
which rival chiefs destroy or give away enormous quantities of wealth to
establish precedence, is the paradigmatic example: the deepest cultural
contribution wins prestige.  But the pattern is far more widespread.
The competitive almsgiving of medieval Christendom, the monumental
architecture of Egyptian pharaohs, the conspicuous sacrifice in ancient
temple economies, the academic citation game---all are forms of
``auction'' where the most complex (deepest) offering wins.

Veblen's theory of conspicuous consumption captures the demand side of this
dynamic: prestige accrues to the display of costly (deep) goods.
Bourdieu's concept of symbolic capital captures the supply side: cultural
producers bid for status by offering complex cultural products to
audiences whose capacity to appreciate them (their cultural capital)
constrains the impact of even the most extravagant bid.\footnote{Bourdieu,
\textit{La Distinction} (1979) and \textit{Les Règles de l'art} (1992).
The ``field of cultural production'' is precisely our auction, with
depth as the currency and resonance as the condition of successful
reception.}

This chapter formalizes ideatic auctions.  We introduce \emph{bid value}
(depth), a \emph{maxDepthIdea} selector, and \emph{auction interpretation}
(the winning bid composed with the receiver's background).  We prove
eighteen theorems showing that void always loses, deeper bids win,
composing bids is subadditive, and the winner's interpretation is bounded
by audience capacity.

% ----------------------------------------------------------------
\section{Core Definitions}
\label{sec:ch14-defs}
% ----------------------------------------------------------------

\begin{definition}[Bid Value]
\label{def:bid-value}
The \emph{bid value} of an idea $a$ is its depth:
$\operatorname{bidValue}(a) = \depth(a)$.
In the ritual economy, depth equals prestige equals bid amount.
\end{definition}

\begin{definition}[Maximum Depth Idea]
\label{def:max-depth}
Given a fallback idea and a list of bids, $\operatorname{maxDepthIdea}$
selects the idea with the greatest depth via a left fold that keeps the
deeper of each pair:
\[
  \operatorname{maxDepthIdea}(\text{best}, []) = \text{best}, \qquad
  \operatorname{maxDepthIdea}(\text{best}, x :: \text{rest}) =
  \begin{cases}
    \operatorname{maxDepthIdea}(x, \text{rest}) & \text{if } \depth(x) \ge \depth(\text{best}), \\
    \operatorname{maxDepthIdea}(\text{best}, \text{rest}) & \text{otherwise}.
  \end{cases}
\]
\end{definition}

\begin{lstlisting}[caption={Maximum depth idea selector in Lean~4}]
def maxDepthIdea : I -> List I -> I
  | best, []       => best
  | best, x :: rest =>
    if depth x >= depth best then maxDepthIdea x rest
    else maxDepthIdea best rest
\end{lstlisting}

\begin{definition}[Auction Interpretation]
\label{def:auction-interp}
The \emph{auction interpretation} composes the receiver's background with
the winning (deepest) bid:
\[
  \operatorname{auctionInterpret}(\text{bids}, \text{fallback}, \text{receiver})
  = \text{receiver} \compose \operatorname{maxDepthIdea}(\text{fallback}, \text{bids}).
\]
The receiver's background mediates the winning bid: the audience does not
absorb the winning idea raw but interprets it through their own cultural
lens.
\end{definition}

% ----------------------------------------------------------------
\section{Void as Non-Competitor}
\label{sec:ch14-void}
% ----------------------------------------------------------------

\begin{theorem}[Void Bid Value is Zero (14.1)]
\label{thm:void-bid}
The void idea has zero bid value: $\operatorname{bidValue}(\void) = 0$.
\end{theorem}

\begin{proof}[Proof sketch]
$\depth(\void) = 0$ by the void depth axiom.
\end{proof}

In the potlatch, offering nothing is the ultimate humiliation.  Zero depth
equals zero cultural capital equals zero bidding power.  The formal basis of
``losing face'': the chief who brings no gifts has bid $\void$ and receives
zero prestige.  Mauss's \emph{obligation to give} is, at bottom, an
obligation to bid non-void.\footnote{Mauss, \textit{Essai sur le don}
(1925): ``To refuse to give, to fail to invite, just as to refuse to
accept, is tantamount to declaring war; it is to reject the bond of
alliance and commonality.''}

\begin{theorem}[Void Loses to Any Bid (14.2)]
\label{thm:void-loses}
Given any bid $a$ and void as fallback, the maximum has non-negative depth:
$\operatorname{bidValue}(\operatorname{maxDepthIdea}(\void, [a])) \ge 0$.
\end{theorem}

\begin{proof}[Proof sketch]
Depth is always non-negative.
\end{proof}

Any cultural offering, no matter how shallow, beats silence.  This is the
formal foundation of ``showing up'': participation is always at least as
good as absence in competitive ritual.  The anthropological wisdom that
``the worst thing you can do is not come'' is a theorem.

\begin{theorem}[Deeper Beats Void (14.3)]
\label{thm:deeper-void}
If $\depth(a) \ge 1$, then $a$ beats void:
$\depth(\operatorname{maxDepthIdea}(\void, [a])) \ge 1$.
\end{theorem}

\begin{proof}[Proof sketch]
When $\depth(a) \ge 1 \ge 0 = \depth(\void)$, the selector chooses $a$.
\end{proof}

In competitive gift-giving, the party who brings \emph{something} always
outranks the party who brings nothing.  Mauss's obligation to give is
structurally enforced: non-givers are structurally dominated.  Even the
most modest gift---a single shell necklace, a handshake, a word of
greeting---outbids silence.

% ----------------------------------------------------------------
\section{Selection Properties}
\label{sec:ch14-selection}
% ----------------------------------------------------------------

\begin{theorem}[Single Bidder Wins (14.4)]
\label{thm:single-bidder}
With only one bid, the winner is determined by comparison with the fallback:
\[
  \operatorname{maxDepthIdea}(\text{fallback}, [a]) =
  \begin{cases}
    a & \text{if } \depth(a) \ge \depth(\text{fallback}), \\
    \text{fallback} & \text{otherwise}.
  \end{cases}
\]
\end{theorem}

\begin{proof}[Proof sketch]
Direct unfolding of the definition with an empty remainder.
\end{proof}

A monopoly on cultural production---single priesthood, state media---means
the sole producer always ``wins'' the auction (assuming their depth exceeds
the fallback tradition).  Competition is absent; interpretation is
determined unilaterally.  This is the formal structure of theocracy:
one bidder, guaranteed victory.

\begin{theorem}[Fallback Survives Empty Auction (14.5)]
\label{thm:empty-auction}
With no bids, the fallback wins:
$\operatorname{maxDepthIdea}(\text{fallback}, []) = \text{fallback}$.
\end{theorem}

\begin{proof}[Proof sketch]
The base case of the recursion.
\end{proof}

When no one competes, the status quo persists.  The ``default culture''
(tradition, orthodoxy) wins by forfeit when no challenger appears.  This
is the formal mechanism of cultural conservatism: in the absence of
competing bids, the existing cultural framework is preserved intact.

\begin{theorem}[Bid Value is Non-Negative (14.6)]
\label{thm:bid-nonneg}
All bids have non-negative value: $\operatorname{bidValue}(a) \ge 0$.
\end{theorem}

\begin{proof}[Proof sketch]
Depth is a natural number.
\end{proof}

There are no ``negative contributions'' in ideatic space.  Every idea has
non-negative complexity.  Even the most trivial cultural offering---a
monosyllable, a blank canvas---has at least zero depth.  The floor of
cultural value is zero, never negative.  This is a modeling choice with
significant anthropological implications: we assume that ideas cannot
\emph{destroy} depth, only fail to add it.\footnote{One could extend the
model with signed depth to capture ``anti-ideas'' or cultural
destruction, but the current axiomatization does not permit this.  The
non-negativity of depth is one of IDT's stronger structural commitments.}

% ----------------------------------------------------------------
\section{Composition and Bidding}
\label{sec:ch14-composition}
% ----------------------------------------------------------------

\begin{theorem}[Composed Bid Subadditivity (14.7)]
\label{thm:composed-bid}
The bid value of composing two ideas is at most the sum of their bid values:
\[
  \operatorname{bidValue}(a \compose b)
  \;\le\;
  \operatorname{bidValue}(a) + \operatorname{bidValue}(b).
\]
\end{theorem}

\begin{proof}[Proof sketch]
By subadditivity of depth.
\end{proof}

Combining two gifts into one package does not create more prestige than
giving them separately.  The potlatch chief who presents a composite gift
of blankets and copper shields gains at most the sum of their separate
values.  Subadditivity prevents ``prestige inflation'' through bundling.
This is the formal refutation of the marketing claim that ``the whole is
greater than the sum of its parts''---in ideatic auctions, the whole is
at most equal to the sum.

\begin{example}[The Contemporary Art Market]
\label{ex:art-market}
The contemporary art market, epitomized by the auction houses Sotheby's
and Christie's, provides a nearly literal instantiation of the ideatic
auction.  When a painting by, say, Jean-Michel Basquiat is offered at
auction, each bidder places a monetary bid that functions as a proxy for
their valuation of the work's cultural depth.  The winning bidder---the
one who values the work most highly (or who has the most resources to
signal status)---acquires the right to ``interpret'' the work: to display
it, to incorporate it into their collection narrative, to claim the
prestige associated with ownership.

Theorem~\ref{thm:composed-bid} applies directly: the prestige of
displaying two works together (a Basquiat alongside a Warhol) is at most
the sum of their individual prestige values.  Bundling artworks does not
create prestige beyond the additive bound.  Theorem~\ref{thm:void-bid}
explains why the artist who produces nothing (the rumored but never
delivered masterpiece) commands zero prestige in the market, no matter
how great their reputation.  And Theorem~\ref{thm:winner-bound} captures
a well-known phenomenon in art criticism: the depth of an artwork's
reception depends as much on the audience's sophistication ($\depth(r)$)
as on the work itself.  A Rothko color field painting that elicits
profound meditative experience in a viewer steeped in Abstract
Expressionism may produce only bewilderment in a viewer without that
background.\footnote{For the sociology of the art market, see Don Thompson,
\textit{The \$12 Million Stuffed Shark} (2008); and Olav Velthuis,
\textit{Talking Prices: Symbolic Meanings of Prices on the Market for
Contemporary Art} (2005).}
\end{example}

\begin{remark}[Connection to Mechanism Design]
\label{rem:mechanism-design}
The ideatic auction framework connects to the broader theory of
\emph{mechanism design}---the ``reverse engineering'' of game theory, in
which the designer chooses the rules of the game to achieve a desired
outcome.  In our setting, the ``mechanism'' is the cultural institution
that determines how bids are solicited, evaluated, and rewarded.  The
potlatch is one mechanism; the academic peer review system is another;
the art auction is a third.  Each implements a different set of rules
for depth-based competition, but all are subject to the same structural
constraints: void bids lose, composition is subadditive, and reception
is bounded.

The mechanism design perspective suggests a normative question: what
cultural institutions \emph{should} we design?  If the goal is to maximize
the depth of the winning bid (cultural excellence), then the incentive
compatibility of the Vickrey auction suggests that institutions should
encourage participants to reveal their full depth---to ``bid'' their
best work without strategic withholding.  Academic peer review, at its
best, approximates this ideal: researchers submit their deepest work for
evaluation by experts.  But peer review also suffers from well-known
failures (bias, conservatism, the Matthew effect) that distort the
depth-based selection process.  The mechanism design framework provides
tools for diagnosing and correcting these failures.\footnote{For mechanism
design applied to scientific institutions, see Partha Dasgupta and
Paul David, ``Toward a New Economics of Science,'' \textit{Research
Policy} 23:5 (1994), pp.\ 487--521.}
\end{remark}

\begin{theorem}[Void Composition Preserves Bid, Right (14.8)]
\label{thm:void-bid-right}
$\operatorname{bidValue}(a \compose \void) = \operatorname{bidValue}(a)$.
\end{theorem}

\begin{proof}[Proof sketch]
$a \compose \void = a$ by the void axiom.
\end{proof}

\begin{theorem}[Void Composition Preserves Bid, Left (14.9)]
\label{thm:void-bid-left}
$\operatorname{bidValue}(\void \compose a) = \operatorname{bidValue}(a)$.
\end{theorem}

\begin{proof}[Proof sketch]
$\void \compose a = a$ by the void axiom.
\end{proof}

Adding ``nothing'' to a gift does not increase its prestige.  Padding a
potlatch offering with empty gestures is structurally transparent---the
audience sees through it.  Prefacing a gift with silence (the dramatic pause
before the big reveal) adds nothing to the gift's structural depth.
Ceremony may have psychological effects, but it is orthogonal to substance
in the depth metric.

\begin{example}[Medieval Relic Trade as Prestige Auction]
\label{ex:relics}
The medieval European trade in saints' relics provides a vivid historical
example of the ideatic auction.  From the early Middle Ages through the
Reformation, relics---bones, clothing, instruments of martyrdom, even
phials of saints' blood---were objects of intense competition among
churches, monasteries, and secular rulers.  A relic's ``bid value'' was
its depth: the proximity of the relic to the original saint, the
reliability of its provenance chain (\textit{translatio}), the number
and quality of miracles attributed to it, and the liturgical traditions
that had accumulated around it.

The most prestigious relics---the Crown of Thorns acquired by Louis~IX
for the Sainte-Chapelle, the relics of St.~James at Compostela, the
Holy Blood of Bruges---were ``deep'' in precisely our technical sense:
they were embedded in rich networks of narrative, ritual, and institutional
memory.  Theorem~\ref{thm:void-bid} explains why a church without relics
had zero prestige in the medieval devotional economy: it had bid void.
Theorem~\ref{thm:composed-bid} explains why combining multiple minor
relics (a practice known as \textit{relic bundling}) produced prestige at
most equal to the sum of the individual relics' values---never more.

The relic trade also illustrates Theorem~\ref{thm:winner-bound}: the
prestige impact of even the most extraordinary relic was bounded by the
receiving community's capacity to appreciate it.  A relic of the True
Cross placed in a remote village chapel produced less cultural impact
than the same relic placed in a great cathedral with an educated clergy
capable of composing elaborate liturgical celebrations around it.  The
relic's depth was constant, but $\depth(r)$---the receiver's depth---varied
enormously.\footnote{On the relic trade, see Patrick Geary, \textit{Furta
Sacra} (1978); and Charles Freeman, \textit{Holy Bones, Holy Dust: How
Relics Shaped the History of Medieval Europe} (2011).}
\end{example}

% ----------------------------------------------------------------
\section{Interpretation of Winning Bids}
\label{sec:ch14-interp}
% ----------------------------------------------------------------

\begin{theorem}[Winner's Interpretation Depth Bound (14.10)]
\label{thm:winner-bound}
The depth of the auction interpretation is bounded:
\[
  \depth(\operatorname{auctionInterpret}(\text{bids}, \text{fallback}, r))
  \;\le\;
  \depth(r) + \depth(\operatorname{maxDepthIdea}(\text{fallback}, \text{bids})).
\]
\end{theorem}

\begin{proof}[Proof sketch]
By subadditivity applied to $r \compose \operatorname{maxDepthIdea}(\cdots)$.
\end{proof}

The cultural impact of the winning bid is bounded by the receiver's
background plus the bid's depth.  Even the most extravagant potlatch is
limited by the audience's capacity to appreciate it.  Bourdieu's insight
that cultural capital constrains reception is formalized: a peasant
watching a court ballet and a courtier watching the same performance
produce interpretations of vastly different depth, not because the ballet
differs but because $\depth(r)$ differs.\footnote{Bourdieu's analysis of
aesthetic judgment in \textit{La Distinction} shows that reception of
cultural goods is stratified by the receiver's accumulated cultural
capital---our $\depth(r)$.}

\begin{example}[Academic Citation as Depth-Based Bidding]
\label{ex:citation}
The academic publication system provides a modern instance of depth-based
auction dynamics.  Each published paper is a ``bid'' in the marketplace
of ideas, with its depth determined by the sophistication of its
arguments, the breadth of its evidence, and the novelty of its
contributions.  The ``auction'' takes place through peer review (the
selection mechanism) and citation (the post-auction valuation).

A paper's bid value ($\operatorname{bidValue}(a) = \depth(a)$) reflects
its intellectual complexity.  Papers with greater depth tend to win the
``auction'' of attention: they are published in more prestigious venues,
cited more frequently, and exert greater influence on subsequent research.
Theorem~\ref{thm:void-bid} captures a basic truth of academic life: a
researcher who publishes nothing has zero academic prestige.
Theorem~\ref{thm:composed-bid} applies to collaborative papers: the
depth of a co-authored paper is at most the sum of the co-authors'
individual depths, though it may be less due to coordination costs and
redundancy.

Theorem~\ref{thm:winner-bound} is particularly illuminating: the impact
of a paper on a reader is bounded by the reader's own depth.  A
groundbreaking paper in algebraic geometry will have enormous impact on a
reader with deep mathematical background ($\depth(r)$ large) but minimal
impact on a reader without the prerequisites ($\depth(r)$ small).  This
is the formal basis of the ``prerequisites problem'' in education: students
cannot absorb material that exceeds their background depth, regardless of
the material's intrinsic quality.\footnote{On the sociology of academic
citation, see Robert K.\ Merton, ``The Matthew Effect in Science,''
\textit{Science} 159:3810 (1968), pp.\ 56--63; and Pierre Bourdieu,
\textit{Homo Academicus} (1984).}
\end{example}

\begin{theorem}[Empty Auction Interpretation (14.11)]
\label{thm:empty-interp}
With no bids, the interpretation is the receiver composed with the fallback:
$\operatorname{auctionInterpret}([], \text{fallback}, r) = r \compose \text{fallback}$.
\end{theorem}

\begin{proof}[Proof sketch]
$\operatorname{maxDepthIdea}(\text{fallback}, []) = \text{fallback}$.
\end{proof}

When no external cultural force competes, the receiver interprets only
through their default framework.  Cultural isolation means self-referential
interpretation: the audience, receiving no new bids, composes its background
with tradition alone.

\begin{theorem}[Void Receiver Gets Raw Winner (14.12)]
\label{thm:void-receiver}
If the receiver's background is void, the interpretation is just the winning
bid:
$\operatorname{auctionInterpret}(\text{bids}, \text{fallback}, \void) =
\operatorname{maxDepthIdea}(\text{fallback}, \text{bids})$.
\end{theorem}

\begin{proof}[Proof sketch]
$\void \compose w = w$ by the void axiom.
\end{proof}

The ``blank slate'' receiver---the newborn, the cultural outsider---simply
absorbs the winning idea without transformation.  Cultural naivety means
unmediated reception.  This is the anthropological equivalent of ``going
native'': the observer without pre-existing cultural categories receives the
dominant cultural signal raw, unfiltered by prior interpretive
frameworks.\footnote{The trope of the ethnographer ``going native'' (from
Malinowski onward) captures precisely this: the field worker who sheds their
own cultural background ($r \to \void$) receives the host culture's signals
without the mediation of home categories.}

\begin{example}[Social Media ``Likes'' as Micro-Auctions]
\label{ex:social-media}
Social media platforms implement a continuous, decentralized version of
the ideatic auction.  Each post is a bid for attention, and the
platform's algorithmic feed serves as the auctioneer, selecting the
``deepest'' (most engaging) content for display.  The ``likes,''
``shares,'' and ``retweets'' that follow are a form of post-auction
valuation: the audience registers its interpretation of the winning bids.

The depth-as-prestige identification
(Definition~\ref{def:bid-value}) maps directly onto the platform's
engagement metrics: content with greater depth (complexity, novelty,
emotional resonance) tends to attract more engagement, which the
platform interprets as higher value.  Theorem~\ref{thm:void-bid}
explains why users who post nothing---``lurkers''---accumulate zero social
media prestige.  Theorem~\ref{thm:void-loses} explains why any post, no
matter how trivial, outperforms silence in the attention economy.

However, the social media auction also reveals tensions in our framework.
Platform algorithms optimize for \emph{engagement} rather than
\emph{depth}, and the two do not always coincide: shallow, emotionally
provocative content may generate more engagement than deep, nuanced
analysis.  This suggests that real-world auctions may use metrics other
than depth for bid evaluation---a limitation of the pure depth-as-prestige
identification that our formalization makes explicit.\footnote{For the
attention economy, see Tim Wu, \textit{The Attention Merchants} (2016);
and Zeynep Tufekci, ``YouTube, the Great Radicalizer,'' \textit{The New
York Times} (March 10, 2018).}
\end{example}

\begin{remark}[The Limits of Depth-as-Prestige Identification]
\label{rem:depth-prestige}
Definition~\ref{def:bid-value} identifies bid value with depth:
$\operatorname{bidValue}(a) = \depth(a)$.  This identification is the
central modeling choice of this chapter, and it deserves scrutiny.  In
many cultural contexts, the identification is apt: prestige accrues to
complexity, sophistication, and elaboration.  The Kwakiutl potlatch,
the medieval cathedral, the academic monograph---all derive their
prestige from depth.

But there are important counterexamples.  Zen Buddhism values simplicity
over complexity; the haiku achieves prestige through compression rather
than elaboration; minimalist art (Judd, Andre, Flavin) stakes its claim
on the refusal of depth.  In these contexts, prestige may be inversely
related to depth, or orthogonal to it.  Our formalization captures only
the ``maximalist'' regime in which deeper is more prestigious.  A more
general framework might parameterize the bid-value function, allowing
$\operatorname{bidValue}(a) = f(\depth(a))$ for some monotone (or even
non-monotone) function~$f$, thereby accommodating both the potlatch
(where $f$ is increasing) and the haiku (where $f$ might be decreasing
above a threshold).
\end{remark}

% ----------------------------------------------------------------
\section{Structural Auction Theorems}
\label{sec:ch14-structural}
% ----------------------------------------------------------------

\begin{theorem}[Iterated Bid Depth Bound (14.13)]
\label{thm:iterated-bid}
Repeating a bid $n$ times via composition yields depth at most
$n \cdot \operatorname{bidValue}(a)$:
\[
  \operatorname{bidValue}(\compN{a}{n}) \;\le\; n \cdot \operatorname{bidValue}(a).
\]
\end{theorem}

\begin{proof}[Proof sketch]
By the iterated composition depth bound.
\end{proof}

Repetitive gift-giving---annual potlatch cycles, recurring tribute---has
bounded cumulative prestige.  You cannot achieve infinite status through mere
repetition.  The subadditivity of depth imposes diminishing returns:
the $n$-th iteration of a potlatch adds at most $\depth(a)$, and the
total is bounded by $n \cdot \depth(a)$.

\begin{theorem}[Atomic Bid Bound (14.14)]
\label{thm:atomic-bid}
Atomic ideas have bid value at most one: if $a$ is atomic, then
$\operatorname{bidValue}(a) \le 1$.
\end{theorem}

\begin{proof}[Proof sketch]
By the depth bound on atomic ideas.
\end{proof}

The simplest, most elemental cultural offerings---a single word, a basic
gesture, a minimal ritual act---have minimal prestige value.  Prestige
comes from complexity: the elaborate ceremony outranks the simple nod.
This formalizes the anthropological observation that prestige goods are
invariably \emph{complex} goods: the jade celt, the knotted quipu, the
illuminated manuscript---all derive their prestige from depth, not from
mere existence.

\begin{theorem}[Resonant Bids Produce Resonant Interpretations (14.15)]
\label{thm:resonant-bids}
If two bids resonate, $b_1 \res b_2$, any common receiver produces resonant
interpretations:
\[
  b_1 \res b_2
  \;\Longrightarrow\;
  (r \compose b_1) \res (r \compose b_2).
\]
\end{theorem}

\begin{proof}[Proof sketch]
By left-composition preservation of resonance.
\end{proof}

Competing ritual systems that resonate---e.g., Catholic and Orthodox
Christianity---produce resonant experiences in the same audience.  The laity
perceives structural similarity even across competing institutions.  This
is the formal basis of religious tolerance: if the competing bids resonate,
the audience's interpretations will also resonate, reducing the perceived
distance between rival systems.

\begin{theorem}[Resonant Receivers Agree on Winner (14.16)]
\label{thm:resonant-receivers}
If two receivers resonate, $r_1 \res r_2$, and evaluate the same bid, their
interpretations resonate:
\[
  r_1 \res r_2
  \;\Longrightarrow\;
  (r_1 \compose b) \res (r_2 \compose b).
\]
\end{theorem}

\begin{proof}[Proof sketch]
By right-composition preservation of resonance.
\end{proof}

Members of the same culture watching the same potlatch produce resonant
evaluations.  Cultural consensus on prestige is a structural consequence of
shared background.  This is why prestige hierarchies are culture-specific:
a Kwakiutl potlatch and a Parisian art auction assign prestige to different
bids, but within each culture, members with resonant backgrounds agree on
what is prestigious.

\begin{theorem}[Bid Associativity (14.17)]
\label{thm:bid-assoc}
Composing bids is associative:
$(a \compose b) \compose c = a \compose (b \compose c)$.
\end{theorem}

\begin{proof}[Proof sketch]
By associativity of the monoid.
\end{proof}

In a multi-stage potlatch (gifts given over successive days), the grouping of
days does not affect the total composition.  The ``grand total'' of a
festival is invariant under scheduling---structure trumps sequence.  Whether
the chief gives blankets on Monday and copper shields on Tuesday, or reverses
the order, the total compositional depth is the same (though the individual
days' compositions may differ).

\begin{remark}[Dynamic Auctions and Repeated Games]
\label{rem:dynamic-auctions}
The theorems in this chapter analyze \emph{static} auctions---single-round
competitions in which bids are evaluated once.  But many cultural
competitions are \emph{dynamic}: they unfold over time, with bids placed
in successive rounds and outcomes depending on the history of previous
bids.  The annual potlatch cycle, the multi-year academic tenure process,
the lifelong artistic career---all are repeated auctions in which current
bids are evaluated against the background of past performance.

In a repeated auction, Theorem~\ref{thm:iterated-bid} provides the
fundamental constraint: repeated bidding yields depth at most
$n \cdot \depth(a)$, where $n$ is the number of rounds and $a$ is the
bid.  This linear bound on cumulative prestige has important implications.
It means that an artist who produces works of constant depth~$d$ in each
of $n$ periods accumulates prestige at most $nd$---a linear growth rate
that precludes the exponential prestige accumulation sometimes imagined
in popular accounts of artistic genius.

The theory of repeated games (Aumann, 1959; Friedman, 1971) shows that
repeated interaction can sustain cooperative outcomes not possible in
one-shot games.  In the ideatic setting, this suggests that repeated
cultural competition can sustain \emph{norms} of depth: participants
bid deeply not merely to win the current round but to maintain a reputation
for depth that benefits them in future rounds.  The potlatch chief who
underbids in one year risks being excluded from future exchanges---a
dynamic incentive not captured by our static analysis but consistent with
its structural constraints.\footnote{For repeated game theory, see Drew
Fudenberg and Jean Tirole, \textit{Game Theory} (1991), ch.~5; and for
applications to cultural institutions, see Avner Greif, \textit{Institutions
and the Path to the Modern Economy} (2006).}
\end{remark}

\begin{theorem}[Depth-Zero Bid Composure (14.18)]
\label{thm:zero-depth-comp}
If both bidders have zero depth, their composition has zero depth:
\[
  \depth(a) = 0 \;\wedge\; \depth(b) = 0
  \;\Longrightarrow\;
  \depth(a \compose b) = 0.
\]
\end{theorem}

\begin{proof}[Proof sketch]
By subadditivity, $\depth(a \compose b) \le \depth(a) + \depth(b) = 0$,
and depth is non-negative, so $\depth(a \compose b) = 0$.
\end{proof}

Two culturally empty offerings combined produce nothing.  ``Nothing plus
nothing is nothing''---the formal futility of combining shallow gestures.
Two empty speeches, two blank canvases, two silent prayers: their
composition is as void as each alone.  The potlatch chief who bundles
two worthless gifts has still bid zero.

\begin{lstlisting}[caption={Depth-zero composition in Lean~4}]
theorem zero_depth_composition {a b : I}
    (ha : depth a = 0) (hb : depth b = 0) :
    depth (compose a b) = 0 := by
  have h := depth_compose_le a b
  omega
\end{lstlisting}

% ----------------------------------------------------------------
\section{From Potlatch to Platform: Auctions Across Scales}
\label{sec:ch14-across-scales}
% ----------------------------------------------------------------

The ideatic auction is a structural pattern that recurs across an
extraordinary range of human institutions, from the smallest face-to-face
community to the largest digital platform.  In this section, we trace the
auction pattern across three scales---local, institutional, and
global---showing how the theorems of this chapter apply at each level.

\subsection{Local Auctions: Face-to-Face Prestige Competition}

At the smallest scale, ideatic auctions occur in face-to-face interactions
where individuals compete for interpretive priority through displays of
cultural depth.  The dinner-party conversation, the academic seminar, the
religious disputation---all are local auctions in which participants bid
their ideas and the ``winner'' (the deepest contribution) sets the terms
of subsequent discussion.

Erving Goffman's interaction ritual theory (1967) provides the
micro-sociological grounding for local auctions.\footnote{Erving Goffman,
\textit{Interaction Ritual} (1967); and Randall Collins, \textit{Interaction
Ritual Chains} (2004), which extends Goffman's framework to explain the
accumulation of ``emotional energy'' through successful interaction
rituals.}  In Goffman's framework, each participant in a social encounter
presents a ``face''---a claimed social value---that is evaluated by others.
The participant whose face is most successfully maintained wins the
interaction's ``pot'' of emotional energy and social approval.  In our
terms, each participant's ``face'' is their bid, and the depth of the
bid determines its success.

Theorem~\ref{thm:single-bidder} captures the dynamics of monopolized
conversation: when one participant dominates (a single bidder), they
win the auction by default, provided their depth exceeds the fallback
(the conventional expectation).  Theorem~\ref{thm:void-bid} explains
why silence (void bidding) is socially costly in most interaction
contexts: the participant who says nothing forfeits the competition and
risks ``losing face'' in Goffman's sense.

\subsection{Institutional Auctions: Organized Cultural Competition}

At the institutional scale, auctions are formalized through rules,
procedures, and evaluation criteria.  The academic peer review system,
the art exhibition jury, the literary prize committee, the religious
ordination process---all are institutional auctions in which bids
(papers, artworks, manuscripts, candidates) are evaluated by designated
judges according to established criteria.

Institutional auctions differ from local auctions in several respects
that our framework captures.  First, the evaluation is \emph{delegated}:
the ``receiver'' in Theorem~\ref{thm:winner-bound} is not the general
public but a specialized panel with (typically) high $\depth(r)$.  This
delegation increases the interpretive depth of the evaluation but also
introduces potential biases (the panel's particular composition shapes
which bids are recognized as ``deep'').

Second, institutional auctions are often \emph{iterated}: the Nobel
Prize is awarded annually, tenure reviews occur on a fixed schedule,
art exhibitions recur seasonally.  Theorem~\ref{thm:iterated-bid}
constrains the cumulative prestige of iterated bidding, while the
repeated-game dynamics discussed in Remark~\ref{rem:dynamic-auctions}
explain why institutional participants invest in reputation-building
over time.

Third, institutional auctions impose \emph{format constraints}: papers
must follow certain conventions, artworks must fit exhibition spaces,
candidates must meet eligibility criteria.  These constraints function
as filters on the space of permissible bids, potentially excluding
high-depth contributions that do not conform to institutional norms.
The tension between format constraints and depth maximization is a
perennial source of institutional controversy---from the Salon des
Refusés (1863) to the replication crisis in psychology.

\subsection{Global Auctions: Platforms and Markets}

At the global scale, digital platforms have created auction mechanisms
of unprecedented scope and speed.  Google's search algorithm ranks
web pages by a depth-like metric (PageRank, which measures the
structural centrality of a page in the link graph).  Amazon's
recommendation system selects products through a form of collaborative
filtering that implicitly auctions items for user attention.  TikTok's
algorithmic feed is a continuous auction in which short video ``bids''
compete for screen time based on engagement metrics that proxy for depth.

These global auctions share the structural properties formalized in our
theorems.  Void content receives zero engagement
(Theorem~\ref{thm:void-bid}).  Deeper content tends to outperform
shallower content (Theorem~\ref{thm:deeper-void}).  The impact of
content on any given user is bounded by that user's background
(Theorem~\ref{thm:winner-bound}).  And the composition of content
(playlists, reading lists, curated feeds) is subadditive
(Theorem~\ref{thm:composed-bid}).

The global auction scale also introduces phenomena that push against our
framework's assumptions.  Network effects can create ``viral'' content
whose apparent depth is inflated by social contagion rather than
intrinsic complexity.  Algorithmic curation may optimize for engagement
rather than depth, creating incentives for ``clickbait''---content
designed to exploit cognitive biases rather than demonstrate genuine
cultural depth.  These phenomena suggest extensions of the framework
in which the bid-value function incorporates social and algorithmic
amplification factors beyond the intrinsic depth of the idea itself.

\subsection{Cross-Scale Connections}

The three scales of auction---local, institutional, and global---are not
independent but mutually reinforcing.  An academic's success in local
auctions (seminars, conferences) builds the reputation that enables
success in institutional auctions (publications, grants, prizes), which
in turn amplifies their presence in global auctions (media coverage,
public intellectual status).  This cross-scale amplification is consistent
with our framework: each level of auction is subject to the same
structural constraints, and success at one level provides the depth
(accumulated cultural capital) that enables competitive bidding at the
next level.

The cross-scale perspective also illuminates the historical transition
from potlatch to platform.  The potlatch was a local auction embedded
in a specific community; its prestige effects were geographically
bounded.  The modern art market is an institutional auction with global
reach; its prestige effects span continents.  The attention economy of
social media is a global auction operating in real time; its prestige
effects are instantaneous and worldwide.  At each stage, the auction's
scope has expanded, but the structural constraints---subadditivity,
void loss, bounded reception---remain invariant.  The potlatch chief
and the TikTok influencer are subject to the same ideatic algebra.

% ----------------------------------------------------------------
\section{Conclusion}
\label{sec:ch14-conclusion}
% ----------------------------------------------------------------

The ideatic auction formalizes a pervasive pattern in human cultural
life: competition for interpretive priority through depth.  The
structural results of this chapter---void always loses
(Theorem~\ref{thm:void-bid}), composition is subadditive
(Theorem~\ref{thm:composed-bid}), reception is bounded by audience
capacity (Theorem~\ref{thm:winner-bound}), and zero-depth bids compose to
zero (Theorem~\ref{thm:zero-depth-comp})---impose hard constraints on
the ritual economy.

These constraints connect to three major anthropological traditions.
Mauss's gift theory is formalized through the obligation to bid non-void
(Theorem~\ref{thm:void-bid}); Veblen's conspicuous consumption is
captured by the depth-as-prestige identification
(Definition~\ref{def:bid-value}); and Bourdieu's symbolic capital is
reflected in the receiver-bounded interpretation
(Theorem~\ref{thm:winner-bound}).

In the final chapter of this sequence, we turn from competitive bidding
to \emph{curation}: the receiver's active management of their own
repertoire, connecting to museum theory, education policy, and the
political economy of knowledge.

% ================================================================
% Chapter 15: Repertoire Curation and Museum Theory
% ================================================================

\chapter{Repertoire Curation and Museum Theory}
\label{ch:curation}

\epigraph{The cultural arbitrary which a system of education transmits,
reproduces, and inculcates is \emph{always} the cultural arbitrary of the
dominant classes.}{Pierre Bourdieu and Jean-Claude Passeron,
\textit{La Reproduction} (1970)}

% ----------------------------------------------------------------
\section{Introduction}
% ----------------------------------------------------------------

The preceding chapters have treated the receiver as a fixed entity---a
repertoire passively awaiting signals.  But receivers are not passive.
Every mind \emph{curates} its repertoire: selecting which ideas to retain,
which to discard, and which new ideas to admit.  The museum curator
choosing which artifacts to display, the teacher designing a syllabus, the
censor deciding which books to ban, the autodidact selecting their next
reading---all are engaged in the same structural operation: filtering and
expanding a repertoire according to some criterion.

Bourdieu's concept of \emph{cultural capital}---the accumulated, curated
stock of ideas that determines one's interpretive capacity---is the
sociological correlate of the repertoire.  Cultural capital is not merely
\emph{possessed}; it is actively \emph{managed}.  The bourgeois family
curates its children's education; the academy curates its canon; the state
curates the national curriculum.  Each act of curation is simultaneously
an act of inclusion and exclusion, of memory and forgetting.\footnote{On the
politics of canon formation, see John Guillory, \textit{Cultural Capital:
The Problem of Literary Canon Formation} (1993).  Guillory argues that the
literary canon is a form of cultural capital whose curation serves class
reproduction---a thesis that our formal framework substantiates.}

This chapter formalizes repertoire curation and expansion.  We introduce
$\operatorname{curate}$---filtering a repertoire by a predicate---and
$\operatorname{expand}$---adding a new idea to the front.  We then prove
eighteen theorems that connect these operations to museum theory, education
policy, censorship dynamics, and the political economy of knowledge.

\subsection*{The Museum Without Walls}

Andr\'{e} Malraux's \textit{Le Mus\'{e}e imaginaire} (1947) introduced the
revolutionary idea of a ``museum without walls'': through photographic
reproduction, every artwork becomes available to every viewer regardless
of geography.  Malraux argued that this technological shift fundamentally
altered the nature of curation---the curator was no longer the gatekeeper
of a physical collection but an \emph{editor} assembling a virtual
repertoire from the totality of world art.  In our formal framework,
Malraux's insight translates directly: the advent of cheap reproduction
means that the set $\IS$ available for curation expands dramatically, but
the curatorial predicate $p$ remains the structural bottleneck.  The
museum without walls does not eliminate curation; it makes curation
\emph{more necessary}, because the uncurated set has grown beyond any
individual's capacity to survey.\footnote{Malraux's argument anticipates
the information-overload problems of the digital age.  When every image
is available online, the curatorial function becomes \emph{more} valuable,
not less.  The ``museum without walls'' is, in our terms, a repertoire
of effectively infinite cardinality---making the choice of predicate $p$
the decisive act.}

Tony Bennett's \textit{The Birth of the Museum} (1995) offered a
complementary perspective, analyzing museums as \emph{disciplinary
institutions} in the Foucauldian sense.  For Bennett, the museum does not
merely curate knowledge; it produces docile subjects by training visitors
to see the world through the museum's organizing categories.  The
repertoire $\mathbf{R}$ displayed by a museum is not a neutral selection
but a \emph{technology of citizenship}, shaping the visitor's own
repertoire through repeated exposure.  Bennett's analysis maps onto our
formal framework as follows: the museum's displayed repertoire
$\operatorname{curate}(\mathbf{R}_{\text{total}}, p_{\text{institutional}})$
becomes the basis for the visitor's expansion operations.  The visitor
expands their personal repertoire $\mathbf{R}_{\text{visitor}}$ by
elements drawn from the museum's curated display, so that institutional
curation upstream constrains individual expansion downstream.

\subsection*{Education as Repertoire Building}

The connection between curation and education runs deeper than analogy.
John Dewey's progressive pedagogy, articulated in \textit{Democracy and
Education} (1916), conceived of education as the guided expansion of
experience---a process we formalize as iterated application of
$\operatorname{expand}$.  Dewey rejected the ``banking model'' in which
knowledge is deposited into passive minds; instead, he argued that
education must be an active process of inquiry in which the learner
selects (curates) their experiences.  The Deweyan student is simultaneously
expanding and curating: each new experience is an expansion, and the
reflective evaluation of that experience is a curation step that
determines which elements persist in the repertoire.

Paulo Freire radicalized Dewey's insight in \textit{Pedagogy of the
Oppressed} (1968), arguing that the ``banking model'' of education is
itself a form of oppressive curation.  When the teacher controls all
expansion operations---deciding which ideas enter the student's
repertoire---and the student has no voice in the curatorial predicate $p$,
education becomes domination.  Freire's alternative, ``problem-posing
education,'' gives the student co-authorship over both operations:
the student participates in deciding what to learn
($\operatorname{expand}$) and how to evaluate what has been learned
($\operatorname{curate}$).  In our formal terms, Freirean liberation
consists in democratizing the predicate $p$---transferring curatorial
authority from the oppressor to the community of learners.

\subsection*{Censorship as Negative Curation}

If education is positive curation---selecting \emph{for} inclusion---then
censorship is negative curation: selecting \emph{for exclusion}.  The
\textit{Index Librorum Prohibitorum}, maintained by the Catholic Church
from 1559 to 1966, was among the most systematic curatorial operations
in history.  The Index defined a predicate
$p_{\text{Index}}(r) = \texttt{false}$ for thousands of works by
Copernicus, Galileo, Descartes, Hume, Voltaire, and others, effectively
curating the Catholic intellectual repertoire by negation.  The structural
consequence, formalized in our Theorem~\ref{thm:curate-length}, is
clear: $|\operatorname{curate}(\mathbf{R}_{\text{Catholic}},
p_{\text{Index}})| < |\mathbf{R}_{\text{Catholic}}|$.  The censored
repertoire is strictly smaller---and therefore, by
Theorem~\ref{thm:curated-interp}, generates strictly fewer
interpretations of any signal.

Contemporary book banning---from school library challenges in the United
States to internet censorship in authoritarian regimes---follows the same
structural logic.  Each banned book is an element removed from the
communal repertoire; each act of censorship narrows the society's
interpretive capacity.  The formal framework makes precise what civil
libertarians assert intuitively: censorship makes us collectively dumber,
in the strict sense of reducing the number of available interpretations.

\subsection*{The Canon Wars}

The ``canon wars'' of the 1980s and 1990s---the debate over whether
Western universities should center their curricula on the ``Great Books''
of the Western tradition or expand to include non-Western and minority
voices---are, in IDT terms, a dispute about the appropriate predicate
$p$.  Defenders of the traditional canon (Allan Bloom's \textit{The
Closing of the American Mind}, 1987; Harold Bloom's \textit{The Western
Canon}, 1994) argued for a predicate that privileged aesthetic excellence
as judged by established critical traditions.  Advocates of multicultural
expansion (Henry Louis Gates Jr., bell hooks, Edward Said) argued that
the traditional predicate systematically excluded the cultural
productions of women, people of color, and non-Western societies.

Our formal framework illuminates both sides of the debate.  The
traditionalists are correct that some curatorial predicate is necessary
(Theorem~\ref{thm:universal-curate}: universal acceptance is not
curation at all, and the resulting repertoire is unnavigable).  The
multiculturalists are correct that a narrow predicate reduces
interpretive capacity (Theorem~\ref{thm:curated-interp}: fewer retained
ideas means fewer interpretations of any signal).  The resolution, if
there is one, lies in designing predicates that balance navigability
against breadth---a problem our theorems frame precisely but do not
solve, because the choice of $p$ is irreducibly political.

\subsection*{Digital Curation and the Paradox of Choice}

The digital age has introduced a new species of curation: algorithmic
recommendation.  Spotify's Discover Weekly, YouTube's recommendation
engine, and TikTok's For You page all implement curatorial predicates
$p_{\text{algorithm}}$ that filter the vast ideatic space of available
content into a personalized repertoire for each user.  These algorithms
curate at a scale and speed impossible for any human curator, processing
millions of items to produce a feed of dozens.  The structural properties
are identical to those of the museum curator: the algorithm filters
($\operatorname{curate}$), the user absorbs ($\operatorname{expand}$),
and the resulting repertoire determines interpretive capacity.

Barry Schwartz's \textit{The Paradox of Choice} (2004) identified the
psychological cost of excessive options: when the uncurated set $\IS$ is
too large, individuals become paralyzed, anxious, and dissatisfied.
Schwartz's paradox is, in our terms, a consequence of the gap between
$|\IS|$ and human cognitive capacity to evaluate curatorial predicates.
The algorithm resolves this paradox by imposing a predicate on the user's
behalf---but at the cost of surrendering curatorial autonomy to an opaque
system whose predicate $p_{\text{algorithm}}$ is optimized for engagement
metrics rather than interpretive depth.  The formal question becomes:
does algorithmic curation increase or decrease
$\operatorname{depthSum}(\mathbf{R}_{\text{user}})$?  Our theorems
(Theorem~\ref{thm:curated-depthsum}) show that curation can only
\emph{reduce} total depth; whether the algorithm's curation reduces
depth less than the user's own imperfect curation is an empirical
question that the formal framework poses but cannot answer
\textit{a priori}.

% ----------------------------------------------------------------
\section{Core Definitions}
\label{sec:ch15-defs}
% ----------------------------------------------------------------

\begin{definition}[Curation]
\label{def:curate}
To \emph{curate} a repertoire $\mathbf{R}$ is to filter it by a Boolean
predicate $p$:
\[
  \operatorname{curate}(\mathbf{R}, p)
  \;=\; [r \in \mathbf{R} \mid p(r) = \texttt{true}].
\]
This models selective retention: the museum curator choosing artifacts, the
teacher choosing concepts, the censor choosing permitted ideas.
\end{definition}

\begin{definition}[Expansion]
\label{def:expand}
To \emph{expand} a repertoire $\mathbf{R}$ with a new idea $a$ is to
prepend it:
\[
  \operatorname{expand}(\mathbf{R}, a) \;=\; a :: \mathbf{R}.
\]
This models learning: acquiring a new concept, encountering a new artifact,
absorbing a new cultural input.
\end{definition}

\begin{lstlisting}[caption={Curation and expansion in Lean~4}]
def curate (rep : List I) (p : I -> Bool) : List I :=
  rep.filter p

def expand (rep : List I) (new_idea : I) : List I :=
  new_idea :: rep
\end{lstlisting}

\begin{definition}[Interpretation]
\label{def:interpret-ch15}
When a receiver with repertoire $\mathbf{R}$ encounters signal $s$, they
produce interpretations:
\[
  \operatorname{interpret}(\mathbf{R}, s)
  \;=\; [r_1 \compose s,\; r_2 \compose s,\; \dots,\; r_n \compose s].
\]
\end{definition}

The interplay of curation, expansion, and interpretation is the central
drama of this chapter.  Curation \emph{reduces} the repertoire (and hence
the number and depth of possible interpretations); expansion
\emph{increases} it.  The structural question is: how do these operations
compose, and what invariants do they preserve?

\begin{example}[The Alexandrian Library as Curation Catastrophe]
\label{ex:alexandria}
The Library of Alexandria, founded in the third century BCE under the
Ptolemaic dynasty, represented the ancient world's most ambitious attempt
at \emph{anti-curation}: the aspiration to collect every written work in
existence.  In our formal terms, the Alexandrian project sought to
approximate $\mathbf{R}_{\text{Alexandria}} \approx \IS_{\text{written}}$
---a repertoire equal to the entire space of written ideas.  By
Theorem~\ref{thm:universal-curate}, this amounts to the identity
curation $\operatorname{curate}(\IS, \lambda x.\, \texttt{true}) = \IS$,
which is structurally equivalent to no curation at all.

The destruction of the Library---whether by fire, neglect, or
conquest---constitutes the most famous curation catastrophe in history.
The curatorial predicate shifted from $\lambda x.\, \texttt{true}$
(preserve everything) to something far more restrictive, dictated by the
contingencies of survival: which scrolls happened to be copied, which
survived the fires, which were translated into Arabic or Latin.  The
resulting repertoire $\mathbf{R}_{\text{surviving}}$ is a strict subset
of $\mathbf{R}_{\text{Alexandria}}$, and by
Theorem~\ref{thm:curated-depthsum},
$\operatorname{depthSum}(\mathbf{R}_{\text{surviving}}) \le
\operatorname{depthSum}(\mathbf{R}_{\text{Alexandria}})$.  The lost
depth---the works of Sappho, Aristotle's second book of the
\textit{Poetics}, countless treatises on mathematics, astronomy, and
philosophy---represents an irrecoverable reduction in humanity's
interpretive capacity.
\end{example}

\begin{example}[The Chinese Imperial Examination as Repertoire Filter]
\label{ex:keju}
The Chinese imperial examination system (\textit{k\={e}j\v{u}}),
operational from 605 CE to 1905, was among the longest-running and most
consequential curatorial operations in human history.  The examinations
required mastery of a fixed canon: the Four Books and Five Classics of
Confucian literature.  In formal terms, the examination defined a
curatorial predicate $p_{\text{k\={e}j\v{u}}}$ that accepted ideas
consonant with the Confucian classics and rejected all others.  Every
aspiring official curated their personal repertoire according to this
predicate: $\mathbf{R}_{\text{scholar}} = \operatorname{curate}(
\mathbf{R}_{\text{available}}, p_{\text{k\={e}j\v{u}}})$.

The system's strength was its meritocratic inclusiveness: any man
(and, in rare cases, any woman) could in principle expand their
repertoire sufficiently to pass.  Its weakness was the narrowness of the
predicate: by the late imperial period, the ``eight-legged essay''
format had ossified the curatorial criterion to the point where
$p_{\text{k\={e}j\v{u}}}$ rejected virtually all modern scientific and
technological ideas.  When Western ideas arrived in force in the
nineteenth century, the Qing dynasty possessed a governing class whose
repertoires had been curated for literary elegance and moral philosophy
but systematically stripped of empirical science.  The result was a
catastrophic mismatch between the curated repertoire and the signals
arriving from the industrializing West.
\end{example}

% ----------------------------------------------------------------
\section{Curation Fundamentals}
\label{sec:ch15-curation}
% ----------------------------------------------------------------

\begin{theorem}[Curation Preserves Membership (15.1)]
\label{thm:curate-subset}
Every element of the curated repertoire was in the original:
\[
  \forall\, x \in \operatorname{curate}(\mathbf{R}, p),\quad x \in \mathbf{R}.
\]
\end{theorem}

\begin{proof}[Proof sketch]
By the property of list filtering: every element of
$\operatorname{filter}(p, \mathbf{R})$ belongs to $\mathbf{R}$.
\end{proof}

A museum can only display what it owns.  Curation is selection, not
creation---the formal basis of the archivist's constraint.  You cannot
curate what you do not have.  This seemingly obvious fact has deep
implications for cultural policy: any canon-formation exercise is bounded
by the existing cultural stock.  A national curriculum can only draw from
the ideas already present in the culture; it cannot conjure new ones through
the act of selection.\footnote{This distinguishes curation from
\emph{creation}.  In our model, creation is the introduction of a genuinely
new element $a \in \IS$ not already in the repertoire; curation is the
selection among existing elements.  The two operations have fundamentally
different structural properties.}

\begin{theorem}[Curation Reduces Length (15.2)]
\label{thm:curate-length}
The curated repertoire is never longer than the original:
\[
  |\operatorname{curate}(\mathbf{R}, p)| \;\le\; |\mathbf{R}|.
\]
\end{theorem}

\begin{proof}[Proof sketch]
Filtering can only remove elements, never add them.
\end{proof}

Curation only removes---it never inflates.  A censored library is always
smaller than the uncensored original.  This formalizes the ``information
loss'' inherent in any act of selection.  Every curatorial decision is,
structurally, a decision to \emph{forget}.

\begin{theorem}[Vacuous Curation is Empty (15.3)]
\label{thm:vacuous-curate}
Filtering by a predicate that rejects everything yields the empty list:
$\operatorname{curate}(\mathbf{R}, \lambda x.\, \texttt{false}) = []$.
\end{theorem}

\begin{proof}[Proof sketch]
No element passes the filter.
\end{proof}

Total censorship---rejecting every idea---produces cultural void.  This is
the formal limit of suppression: complete censorship equals complete cultural
erasure.  The totalitarian ideal of the \emph{tabula rasa}---erasing all
prior culture to build anew---is structurally equivalent to reducing the
repertoire to the empty list.  Pol Pot's ``Year Zero,'' Mao's Cultural
Revolution, the \emph{damnatio memoriae} of the Roman Senate: all are
attempts at vacuous curation.\footnote{The historical record suggests that
vacuous curation is never perfectly achieved: underground literature,
oral tradition, and material culture survive even the most thorough
purges.  The theorem describes the \emph{structural limit}, not the
empirical reality.}

\begin{theorem}[Universal Curation is Identity (15.4)]
\label{thm:universal-curate}
Filtering by a predicate that accepts everything yields the original:
$\operatorname{curate}(\mathbf{R}, \lambda x.\, \texttt{true}) = \mathbf{R}$.
\end{theorem}

\begin{proof}[Proof sketch]
Every element passes the filter.
\end{proof}

A museum that displays everything is just a warehouse.  ``Curating'' with no
selection criterion is structurally equivalent to no curation at all.  This
is the formal basis of the distinction between the \emph{archive}
(everything preserved) and the \emph{exhibition} (a curated selection).
The uncurated repertoire is the Borgesian library: complete, but
unnavigable.\footnote{Jorge Luis Borges, ``The Library of Babel'' (1941).
Borges's infinite library contains every possible book but is useless
precisely because it lacks curation.}

\begin{remark}[Connection to Information Theory: Channel Capacity]
\label{rem:channel-capacity}
The duality between curation (reducing the repertoire) and universal
acceptance (identity) has a precise analogue in Shannon's information
theory.  A communication channel with capacity $C$ can transmit at most
$C$ bits per second reliably; attempting to transmit more introduces
errors.  Similarly, a cognitive agent with finite processing capacity
can only maintain and effectively deploy a repertoire of bounded size.
The curatorial predicate $p$ functions as a \emph{channel code},
selecting which elements of the vast ideatic space $\IS$ to retain
within the agent's finite capacity.

Shannon's noisy channel coding theorem states that optimal codes exist
which achieve the capacity bound---but finding them is computationally
hard.  The analogous statement for curation is: for any cognitive
capacity $C$, there exists an optimal predicate $p^*$ that maximizes
$\operatorname{depthSum}(\operatorname{curate}(\mathbf{R}, p^*))$
subject to $|\operatorname{curate}(\mathbf{R}, p^*)| \le C$.  The
curator's dilemma is that this optimal predicate is, in general,
intractable to compute---hence the irreducibly judgmental character of
curatorial work.  No algorithm can replace the curator's taste, for
the same reason that no algorithm can efficiently find optimal channel
codes in the general case.
\end{remark}

\begin{example}[Spotify's Recommendation Algorithm as Automated Curation]
\label{ex:spotify}
Spotify's recommendation engine curates from a catalog of over 100
million tracks.  Its predicate $p_{\text{Spotify}}$ is a learned function
of the user's listening history, demographic data, audio features, and
collaborative filtering across millions of users.  In our framework,
each weekly ``Discover Weekly'' playlist is a curated repertoire:
$\mathbf{R}_{\text{playlist}} = \operatorname{curate}(
\mathbf{R}_{\text{catalog}}, p_{\text{Spotify}})$.

By Theorem~\ref{thm:curate-subset}, every recommended track was already
in the catalog; the algorithm creates nothing new.  By
Theorem~\ref{thm:curate-length}, the playlist is vastly shorter than the
catalog---typically 30 tracks from 100 million.  The crucial question is
whether the algorithmic predicate selects for depth or for familiarity.
Empirical evidence suggests that recommendation algorithms optimize for
engagement (clicks, completions, saves), which correlates imperfectly
with interpretive depth.  The result is a ``filter bubble'' in which the
user's curated repertoire converges to a narrow region of ideatic space,
reducing the diversity of interpretive possibilities available when
encountering novel signals.
\end{example}

% ----------------------------------------------------------------
\section{Expansion Fundamentals}
\label{sec:ch15-expansion}
% ----------------------------------------------------------------

\begin{theorem}[Expansion Increases Length (15.5)]
\label{thm:expand-length}
Adding a new idea increases the repertoire by exactly one:
\[
  |\operatorname{expand}(\mathbf{R}, a)| = |\mathbf{R}| + 1.
\]
\end{theorem}

\begin{proof}[Proof sketch]
Prepending one element increases list length by one.
\end{proof}

Each learning event adds exactly one new schema.  Education is
\emph{incremental}---you learn one thing at a time.  This constrains
theories of ``epiphany'' or ``gestalt shift'' to be single-idea additions
in the formal model.  The Humboldt-ian ideal of \emph{Bildung}---the
gradual, systematic expansion of the mind---is structurally captured by
iterated application of $\operatorname{expand}$.

\begin{theorem}[New Idea is in Expanded Repertoire (15.6)]
\label{thm:new-idea-mem}
The newly added idea belongs to the expanded repertoire:
$a \in \operatorname{expand}(\mathbf{R}, a)$.
\end{theorem}

\begin{proof}[Proof sketch]
The new idea is the head of the list.
\end{proof}

What you learn, you possess.  The student who has been taught a concept
\emph{has} that concept in their repertoire.  This is the structural basis
of education as repertoire expansion: successful pedagogy is precisely the
operation $\operatorname{expand}(\mathbf{R}_{\text{student}}, a_{\text{lesson}})$.

\begin{theorem}[Old Ideas Survive Expansion (15.7)]
\label{thm:old-survive}
All ideas in the original repertoire remain in the expanded version:
\[
  \forall\, x \in \mathbf{R},\quad x \in \operatorname{expand}(\mathbf{R}, a).
\]
\end{theorem}

\begin{proof}[Proof sketch]
The original list is the tail of the expanded list.
\end{proof}

Learning does not erase prior knowledge.  The student who learns calculus
still remembers arithmetic.  Expansion is \emph{monotone}---repertoires only
grow.  This is an important modeling assumption: in our framework, forgetting
requires an explicit curation step (filtering out the unwanted idea), not
mere neglect.  The distinction between passive forgetting and active
censorship is, in IDT, the distinction between non-operation and
curation.\footnote{In cognitive science, the debate between ``decay'' and
``interference'' theories of forgetting maps onto our distinction between
curation (active removal) and repertoire stagnation (non-expansion).}

% ----------------------------------------------------------------
\section{Curation and Interpretation}
\label{sec:ch15-curate-interp}
% ----------------------------------------------------------------

\begin{theorem}[Curated Interpretation is Sublist (15.8)]
\label{thm:curated-interp}
Interpreting a curated repertoire gives at most as many interpretations as
the full repertoire:
\[
  |\operatorname{interpret}(\operatorname{curate}(\mathbf{R}, p), s)|
  \;\le\;
  |\operatorname{interpret}(\mathbf{R}, s)|.
\]
\end{theorem}

\begin{proof}[Proof sketch]
$|\operatorname{interpret}(\cdot, s)| = |\cdot|$ (interpretation preserves
length), and $|\operatorname{curate}(\mathbf{R}, p)| \le |\mathbf{R}|$.
\end{proof}

A censored mind produces fewer interpretations.  The number of possible
readings of a text is bounded by the reader's (curated) repertoire size.
A narrower education yields fewer perspectives---the formal basis of the
humanist argument for breadth in the curriculum.  The liberally educated
mind, with its expansive repertoire, generates more interpretations than
the narrowly trained specialist.

\begin{remark}[Censorship as Curation with Political Consequences]
\label{rem:censorship-political}
The connection between curation and censorship deserves sustained
attention, because the formal equivalence conceals a crucial normative
distinction.  The museum curator and the government censor both apply
predicates to repertoires, but they differ in three respects: (1) the
scope of application (one collection vs.\ an entire society), (2) the
coercive power backing the predicate (aesthetic judgment vs.\ legal
penalty), and (3) the reversibility of the operation (the deaccessioned
painting can be reacquired; the burned book is gone forever).

Our formal framework captures the \emph{structural} identity of these
operations while remaining silent on their normative valence.  This is
a feature, not a bug: the theorems show that the same mathematical
constraints govern benign curation and malign censorship, which means
that arguments about the \emph{structural} effects of censorship need
not appeal to contested moral premises.
Theorem~\ref{thm:curated-interp}---that curation reduces interpretive
capacity---is value-neutral as a mathematical fact, but devastating as
an argument against censorship: it shows that \emph{any} censorship,
regardless of its motivation, reduces the society's collective capacity
to interpret signals.
\end{remark}

\begin{theorem}[Empty Curation Yields No Interpretation (15.9)]
\label{thm:empty-curate-interp}
If curation removes everything, there are no interpretations:
\[
  \operatorname{interpret}(\operatorname{curate}(\mathbf{R},
  \lambda x.\, \texttt{false}), s) = [].
\]
\end{theorem}

\begin{proof}[Proof sketch]
By Theorem~\ref{thm:vacuous-curate}, the curated repertoire is empty,
and mapping over the empty list yields the empty list.
\end{proof}

The culturally emptied mind cannot interpret anything.  Total censorship
followed by signal reception equals zero interpretive output.  This is the
``blank slate after brainwashing'' theorem: the totalitarian regime that
succeeds in erasing all prior culture leaves its subjects incapable of
interpreting \emph{any} signal, including its own propaganda.  The self-
defeating nature of complete censorship is a structural theorem.

% ----------------------------------------------------------------
\section{Expansion and Interpretation}
\label{sec:ch15-expand-interp}
% ----------------------------------------------------------------

\begin{theorem}[Expanded Interpretation Adds One (15.10)]
\label{thm:expanded-interp}
Interpreting an expanded repertoire produces one more interpretation than
the original:
\[
  |\operatorname{interpret}(\operatorname{expand}(\mathbf{R}, a), s)|
  \;=\;
  |\operatorname{interpret}(\mathbf{R}, s)| + 1.
\]
\end{theorem}

\begin{proof}[Proof sketch]
Expansion adds one element; interpretation preserves length.
\end{proof}

Each new idea in the repertoire adds exactly one new interpretive
possibility.  Education's return is precisely one new perspective per concept
learned---the fundamental unit of \emph{Bildung}.  The marginal return on
education, measured in interpretive possibilities, is constant: one concept
in, one new interpretation out.

\begin{theorem}[Expansion Interpretation Structure (15.11)]
\label{thm:expand-structure}
Interpreting an expanded repertoire prepends the new interpretation:
\[
  \operatorname{interpret}(\operatorname{expand}(\mathbf{R}, a), s)
  \;=\;
  (a \compose s) :: \operatorname{interpret}(\mathbf{R}, s).
\]
\end{theorem}

\begin{proof}[Proof sketch]
By definition: $\operatorname{expand}(\mathbf{R}, a) = a :: \mathbf{R}$,
and mapping over a cons is $f(a) :: \operatorname{map}(f, \mathbf{R})$.
\end{proof}

The new concept's interpretation appears first, followed by all prior
interpretations unchanged.  Learning adds a new lens without distorting
existing ones---the ``additive'' model of education.  This is structurally
optimistic: new knowledge does not corrupt old knowledge.  The Baconian
vision of knowledge as cumulative finds formal support here.

\begin{theorem}[Void Expansion Doesn't Change Interpretations (15.12)]
\label{thm:void-expansion}
Adding void to the repertoire adds the raw signal $s$ as a new interpretation:
\[
  \operatorname{interpret}(\operatorname{expand}(\mathbf{R}, \void), s)
  \;=\;
  s :: \operatorname{interpret}(\mathbf{R}, s).
\]
\end{theorem}

\begin{proof}[Proof sketch]
$\void \compose s = s$ by the void axiom, so the new interpretation is $s$
itself.
\end{proof}

Adding ``empty receptivity''---the Keatsian \emph{negative capability}---to
one's repertoire adds the raw signal as a new interpretation.  The
educationally valuable ``empty slot'' lets you see the signal unmediated.
This is the formal justification for intellectual humility: maintaining a
void element in one's repertoire ensures that one always has access to the
unfiltered signal, free from the bias of pre-existing categories.\footnote{Keats's
concept of ``negative capability''---``being in uncertainties, mysteries,
doubts, without any irritable reaching after fact and reason''---is
precisely the act of maintaining $\void$ in the repertoire: a position of
receptivity that produces the raw signal as interpretation.}

\begin{lstlisting}[caption={Void expansion raw signal in Lean~4}]
theorem void_expansion_raw_signal (rep : List I) (s : I) :
    interpret (expand rep void) s = s :: interpret rep s := by
  simp [expanded_interp_structure, void_left]
\end{lstlisting}

% ----------------------------------------------------------------
\section{DepthSum Bounds}
\label{sec:ch15-depthsum}
% ----------------------------------------------------------------

\begin{theorem}[Curated DepthSum Bound (15.13)]
\label{thm:curated-depthsum}
The depth sum of a curated repertoire is at most the depth sum of the
original:
\[
  \operatorname{depthSum}(\operatorname{curate}(\mathbf{R}, p))
  \;\le\;
  \operatorname{depthSum}(\mathbf{R}).
\]
\end{theorem}

\begin{proof}[Proof sketch]
By induction on $\mathbf{R}$.  At each step, if the element passes the
filter, both sides include $\depth(r)$; if it fails, only the right side
includes it, maintaining the inequality.
\end{proof}

Curation can only \emph{reduce} total cultural depth.  A museum that hides
artifacts loses depth; a censored library is shallower than the uncensored
one.  Information destruction is irreversible in the depth metric.  This
theorem has stark implications for cultural policy: every act of censorship
is a permanent reduction in the society's total interpretive capacity.  The
censored ideas' depth is \emph{lost}, not merely hidden.

\begin{theorem}[Expanded DepthSum (15.14)]
\label{thm:expanded-depthsum}
Expanding the repertoire adds exactly the new idea's depth:
\[
  \operatorname{depthSum}(\operatorname{expand}(\mathbf{R}, a))
  \;=\;
  \depth(a) + \operatorname{depthSum}(\mathbf{R}).
\]
\end{theorem}

\begin{proof}[Proof sketch]
$\operatorname{expand}(\mathbf{R}, a) = a :: \mathbf{R}$, and
$\operatorname{depthSum}(a :: \mathbf{R}) = \depth(a) +
\operatorname{depthSum}(\mathbf{R})$.
\end{proof}

Each new idea adds exactly its depth to the total cultural capital.
Education is a \emph{linear} accumulation of depth---Bourdieu's cultural
capital grows additively with each new acquisition.  The student who learns
a concept of depth $d$ increases their total cultural capital by exactly $d$.
This linearity is both empowering (every lesson counts) and sobering (there
are no shortcuts---depth must be earned one idea at a time).

\begin{theorem}[Void Expansion DepthSum (15.15)]
\label{thm:void-expansion-depth}
Adding void to the repertoire does not change depthSum:
\[
  \operatorname{depthSum}(\operatorname{expand}(\mathbf{R}, \void))
  \;=\;
  \operatorname{depthSum}(\mathbf{R}).
\]
\end{theorem}

\begin{proof}[Proof sketch]
$\depth(\void) = 0$, so $0 + \operatorname{depthSum}(\mathbf{R}) =
\operatorname{depthSum}(\mathbf{R})$.
\end{proof}

Adding ``nothing'' to one's cultural stock adds zero depth.  The void is
costless to maintain in the repertoire---the formal basis of ``keeping an
open mind'' being free in the depth metric.  Intellectual humility (carrying
$\void$ in one's repertoire) costs nothing and gains the benefit of raw
signal access (Theorem~\ref{thm:void-expansion}).

\begin{example}[Museum Deaccessioning Controversies]
\label{ex:deaccessioning}
Museum deaccessioning---the practice of removing works from a permanent
collection, typically through sale---illustrates the irreversibility of
curation in practice.  When the Berkshire Museum in 2017 proposed selling
Norman Rockwell paintings and Hudson River School landscapes to fund
operational costs, the art world erupted.  The controversy turns on a
formal point: deaccessioning applies a new predicate
$p_{\text{financial}}$ (``does this work justify its storage and
insurance costs?'') to a repertoire that was originally curated by a
different predicate $p_{\text{aesthetic}}$ (``does this work have
artistic merit?'').

By Theorem~\ref{thm:curate-idempotent}, re-curating by the same
predicate has no effect.  But curating by a \emph{different} predicate
compounds the restriction:
$\operatorname{curate}(\operatorname{curate}(\mathbf{R},
p_{\text{aesthetic}}), p_{\text{financial}})$ yields a repertoire that
satisfies \emph{both} predicates, and is therefore (in general) strictly
smaller than either single-predicate curation.  The lost works satisfy
$p_{\text{aesthetic}}$ but not $p_{\text{financial}}$; they are
aesthetically valuable but financially burdensome.  The museum's
cultural depth (Theorem~\ref{thm:curated-depthsum}) shrinks by the sum
of the deaccessioned works' depths.

The deaccessioning debate thus reduces to a question about predicate
ordering: should financial predicates ever override aesthetic ones in
cultural institutions?  Our framework does not answer this normative
question, but it precisely quantifies the cost: each deaccessioned work
reduces $\operatorname{depthSum}$ by exactly $\depth(r_{\text{sold}})$,
and this reduction is permanent once the work enters a private collection.
\end{example}

% ----------------------------------------------------------------
\section{Structural Curation Theorems}
\label{sec:ch15-structural}
% ----------------------------------------------------------------

\begin{theorem}[Double Curation is Single Curation (15.16)]
\label{thm:curate-idempotent}
Curating a curated repertoire by the same predicate is idempotent:
\[
  \operatorname{curate}(\operatorname{curate}(\mathbf{R}, p), p)
  \;=\;
  \operatorname{curate}(\mathbf{R}, p).
\]
\end{theorem}

\begin{proof}[Proof sketch]
$\operatorname{filter}(p, \operatorname{filter}(p, \mathbf{R}))
= \operatorname{filter}(p \wedge p, \mathbf{R})
= \operatorname{filter}(p, \mathbf{R})$
since $p \wedge p = p$.
\end{proof}

Censoring an already-censored library by the same criterion has no further
effect.  The ``diminishing returns'' of redundant censorship---once you have
removed all ``heretical'' books, removing them again changes nothing.  This
idempotency result has practical implications for bureaucratic censorship:
the second censor applying the same criteria as the first is structurally
redundant, a waste of institutional resources.  Totalitarian regimes that
layer censorship committees atop each other (with identical mandates) achieve
nothing beyond the first pass.\footnote{Historical examples include the
multiple overlapping censorship bodies of the Soviet Union (\textit{Glavlit},
\textit{Goskomizdat}, Party committees), each applying essentially the same
ideological filter.  Our theorem shows this redundancy is structural, not
merely bureaucratic.}

\begin{theorem}[Expand Then Curate Accepting (15.17)]
\label{thm:expand-curate-accept}
If the new idea passes the predicate, expanding then curating keeps it:
\[
  p(a) = \texttt{true}
  \;\Longrightarrow\;
  \operatorname{curate}(\operatorname{expand}(\mathbf{R}, a), p)
  = a :: \operatorname{curate}(\mathbf{R}, p).
\]
\end{theorem}

\begin{proof}[Proof sketch]
Filtering $a :: \mathbf{R}$ with $p$: since $p(a) = \texttt{true}$,
$a$ is kept, and the tail is filtered as usual.
\end{proof}

If a newly acquired idea passes the censorship test, it survives curation.
``Approved learning'' is permanent---the accepted idea stays in the
repertoire through subsequent curation rounds.  This is the formal basis
of canonical inclusion: once an idea is admitted to the canon (passes the
predicate), it persists through all future curation by the same criterion.
Dante, admitted to the Western canon, survives every subsequent round of
canon revision that uses the same evaluative framework.

\begin{theorem}[Expand Then Curate Rejecting (15.18)]
\label{thm:expand-curate-reject}
If the new idea fails the predicate, expanding then curating removes it:
\[
  p(a) = \texttt{false}
  \;\Longrightarrow\;
  \operatorname{curate}(\operatorname{expand}(\mathbf{R}, a), p)
  = \operatorname{curate}(\mathbf{R}, p).
\]
\end{theorem}

\begin{proof}[Proof sketch]
Filtering $a :: \mathbf{R}$ with $p$: since $p(a) = \texttt{false}$,
$a$ is discarded, and the tail is filtered as before.
\end{proof}

Ideas that fail the censorship test are immediately purged.  The
``forbidden knowledge'' theorem: learning something that the regime rejects
is structurally equivalent to never having learned it, after curation.  This
formalizes the futility of teaching banned concepts in totalitarian
educational systems---the student may \emph{encounter} the forbidden idea
(via $\operatorname{expand}$), but the subsequent curation step erases it
from the repertoire as if the encounter never happened.

The theorem is disturbing in its implications.  If $p$ represents the
ideological filter of a totalitarian state, then
$\operatorname{curate}(\operatorname{expand}(\mathbf{R}, a_{\text{forbidden}}), p)
= \operatorname{curate}(\mathbf{R}, p)$: the forbidden idea leaves no trace.
The dissident's effort to teach subversive ideas is, under perfect
censorship, structurally null.  Of course, perfect censorship is an
idealization---but the theorem shows what perfect censorship would
\emph{mean}: structural erasure of dissent.

\paragraph{Extended proof discussion.}
The expand-then-curate-rejecting theorem (15.18) merits careful
reflection on its proof structure, because it reveals a deep asymmetry
between expansion and curation.  The proof proceeds by case analysis on
the head of the list $a :: \mathbf{R}$.  When $p(a) = \texttt{false}$,
the filter function skips the head and recurses on the tail, producing
exactly $\operatorname{filter}(p, \mathbf{R})$.  The key insight is that
the filter is \emph{memoryless}: it does not record that an element was
tested and rejected.  There is no ``trace'' of the rejected idea in the
output.  This memorylessness is what makes censorship structurally
complete in the formal model: the curated repertoire contains no evidence
that anything was removed.

Contrast this with real-world censorship, where the act of banning a
book often draws attention to it (the ``Streisand effect'').  The
discrepancy arises because real censorship operates in a social context
where the \emph{act} of curation is itself observable as a signal.  In
our model, the curatorial operation is transparent to downstream
consumers of the repertoire: they see only
$\operatorname{curate}(\mathbf{R}, p)$, not the pair
$(\mathbf{R}, p)$.  Modeling censorship-awareness would require a
second-order theory in which agents can observe and interpret curatorial
predicates as signals---a direction we leave for future work.

\begin{remark}[The Filter Bubble Problem]
\label{rem:filter-bubble}
Eli Pariser's concept of the ``filter bubble'' (\textit{The Filter
Bubble}, 2011) describes the phenomenon in which algorithmic curation
progressively narrows an individual's information environment.  In our
formal terms, a filter bubble arises when the curatorial predicate is
\emph{adaptive}: $p_{t+1}$ is a function of
$\operatorname{curate}(\mathbf{R}, p_t)$.  If the algorithm curates
based on what the user has previously consumed, and the user consumes
only what is curated, a feedback loop emerges.

Formally, let $p_0$ be the initial predicate, and define
$\mathbf{R}_{t+1} = \operatorname{curate}(\mathbf{R}_t, p_t)$ where
$p_t$ depends on $\mathbf{R}_{t-1}$.  By Theorem~\ref{thm:curate-length},
$|\mathbf{R}_{t+1}| \le |\mathbf{R}_t|$, so the sequence
$\{|\mathbf{R}_t|\}$ is non-increasing.  If $p_t$ ever strictly
rejects an element of $\mathbf{R}_t$, the repertoire shrinks
\emph{irreversibly} (assuming no external expansion).  The filter
bubble is thus a monotone descent in repertoire size: once an idea is
filtered out, the adaptive predicate---which learns from the filtered
repertoire---has no way to ``rediscover'' it.  This is the formal
diagnosis of the filter bubble: adaptive curation without external
expansion converges to a fixed point of minimal repertoire that merely
confirms the user's existing ideas.

Pariser's solution---deliberate exposure to diverse viewpoints---is
formally an $\operatorname{expand}$ operation that injects ideas outside
the current curated repertoire, breaking the feedback loop by
introducing elements that the adaptive predicate has not yet evaluated.
The interplay between Theorems~\ref{thm:expand-curate-accept}
and~\ref{thm:expand-curate-reject} then determines whether these
injected ideas survive subsequent curation rounds.
\end{remark}

\begin{lstlisting}[caption={Expand-then-curate in Lean~4}]
theorem expand_curate_removes {new_idea : I} {p : I -> Bool}
    (hp : p new_idea = false) (rep : List I) :
    curate (expand rep new_idea) p = curate rep p := by
  simp [expand, curate, List.filter, hp]
\end{lstlisting}

% ----------------------------------------------------------------
\section{The Politics of Repertoire}
\label{sec:ch15-politics}
% ----------------------------------------------------------------

The theorems of this chapter, though stated in the neutral language of
predicates and lists, have inescapably political implications.  Every act
of curation is an exercise of power: the power to determine what ideas
are available for interpretation, and therefore what interpretations are
possible.  This section draws out the political consequences of the
formal framework.

\subsection{Who Holds the Predicate?}

The central political question of repertoire curation is: \emph{who
chooses the predicate $p$?}  In a democratic society, the answer is
ideally ``the community of interpreters themselves,'' through mechanisms
such as democratic curriculum design, public library governance, and
editorial independence.  In authoritarian regimes, the predicate is
imposed by the state: a single party, a censorship bureau, or a
supreme leader determines what ideas are admitted to the national
repertoire.

The formal framework reveals that this difference is not merely
quantitative (how much censorship) but qualitative (who controls the
structural operation).  The predicate $p$ is a \emph{function}, and the
choice of function space matters: a community that can choose from a
rich space of predicates (diverse curricula, multiple publishers,
independent media) has a structurally different curatorial capacity than
one restricted to a single state-imposed predicate.

\subsection{Curation and Inequality}

Bourdieu's analysis of cultural capital implies that curation is
unevenly distributed.  The children of the bourgeoisie enter school
with repertoires already curated by family exposure to art, literature,
and cosmopolitan experience.  The children of the working class enter
with repertoires curated by material necessity.  The school then applies
a uniform predicate---the curriculum---to both groups, but the
\emph{expansion} operations are different: the bourgeois child
encounters new ideas that resonate with an already-curated repertoire,
while the working-class child encounters ideas for which no resonating
elements exist in their current repertoire.

Formally, if $\mathbf{R}_{\text{bourgeois}}$ already contains elements
that resonate with the curriculum's new ideas, then expansion is
``sticky''---the new ideas compose meaningfully with existing ones,
producing high-depth interpretations.  If
$\mathbf{R}_{\text{working}}$ lacks such resonating elements, expansion
still adds the new idea (Theorem~\ref{thm:new-idea-mem}), but the
interpretive compositions it enables may have lower depth due to the
absence of resonant partners.  Educational inequality is thus a
\emph{compositional} phenomenon: it arises not from differential access
to ideas per se, but from differential resonance between new ideas and
existing repertoires.

\subsection{Repertoire Diversity as Public Good}

A society's collective repertoire---the union of all individual
repertoires---determines its collective interpretive capacity.  If every
citizen curates by the same predicate (ideological conformity), the
collective repertoire is no larger than a single individual's.  If
citizens curate by diverse predicates (ideological pluralism), the
collective repertoire is the union of many differently curated subsets,
and is therefore larger.

This suggests a formal argument for intellectual diversity as a public
good: a society with diverse curatorial predicates has a strictly larger
collective repertoire, and therefore a strictly greater collective
capacity to interpret novel signals.  When an unprecedented crisis
arrives---a pandemic, an economic shock, a technological disruption---the
society with a diverse repertoire has more interpretive resources to draw
upon than the homogeneous society.  Censorship, which imposes a uniform
predicate, is thus a structural risk factor for societal fragility:
it reduces the diversity of interpretive resources available when they
are most needed.

\paragraph{Extended proof discussion: Depth bounds and curatorial
composition.}
The interaction between depth bounds and iterated curation
deserves more detailed analysis.  Consider a sequence of curations
$\operatorname{curate}(\operatorname{curate}(\mathbf{R}, p_1), p_2)$.
By Theorem~\ref{thm:curated-depthsum} applied twice:
\[
  \operatorname{depthSum}(\operatorname{curate}(
    \operatorname{curate}(\mathbf{R}, p_1), p_2))
  \;\le\;
  \operatorname{depthSum}(\operatorname{curate}(\mathbf{R}, p_1))
  \;\le\;
  \operatorname{depthSum}(\mathbf{R}).
\]
Each curation pass can only decrease depth, so iterated curation by
different predicates produces a monotone-decreasing sequence of depth
sums.  This means that the order of application does not matter for the
\emph{direction} of the inequality (both orderings lose depth), but it
may matter for the \emph{amount} of depth lost.

Specifically, if $p_1$ removes deep elements and $p_2$ removes shallow
elements, the total depth loss is the sum of the depths of all removed
elements.  But if $p_1$ removes some elements that $p_2$ would also have
removed, the compound curation $\operatorname{curate}(\cdot, p_1 \wedge
p_2)$---applying both predicates simultaneously---may yield a different
depth sum than either sequential application.  The key structural
insight is that $\operatorname{curate}(\operatorname{curate}(\mathbf{R},
p_1), p_2) = \operatorname{curate}(\mathbf{R}, p_1 \wedge p_2)$
when both predicates are applied to the same elements, which follows
from the fact that $\operatorname{filter}(p_2,
\operatorname{filter}(p_1, \mathbf{R})) =
\operatorname{filter}(p_1 \wedge p_2, \mathbf{R})$ for any list
$\mathbf{R}$.  This filter-conjunction lemma is the formal basis of the
observation that layered censorship (first political, then moral, then
aesthetic) is equivalent to a single pass with a conjunctive predicate.

% ----------------------------------------------------------------
\section{Conclusion}
\label{sec:ch15-conclusion}
% ----------------------------------------------------------------

Repertoire curation, formalized in IDT, reveals the structural anatomy of
cultural management.  Curation is subtractive
(Theorem~\ref{thm:curate-length}), idempotent
(Theorem~\ref{thm:curate-idempotent}), and depth-reducing
(Theorem~\ref{thm:curated-depthsum}).  Expansion is additive
(Theorem~\ref{thm:expand-length}), monotone
(Theorem~\ref{thm:old-survive}), and depth-increasing
(Theorem~\ref{thm:expanded-depthsum}).  The interaction of the two---expand
then curate---is determined entirely by whether the new idea passes the
filter (Theorems~\ref{thm:expand-curate-accept}
and~\ref{thm:expand-curate-reject}).

These results connect to four domains of cultural practice.  In
\emph{museum theory}, curation is the fundamental operation of exhibition
design---and our theorems show that it is structurally subtractive.  In
\emph{education policy}, expansion is the mechanism of \emph{Bildung},
and each lesson adds exactly one interpretive possibility.  In
\emph{censorship dynamics}, vacuous curation produces cultural void
(Theorem~\ref{thm:vacuous-curate}), and the expand-then-curate-rejecting
result (Theorem~\ref{thm:expand-curate-reject}) formalizes the structural
erasure of dissent.  In \emph{Bourdieu's sociology}, cultural capital grows
linearly with expansion (Theorem~\ref{thm:expanded-depthsum}) and can only
shrink through curation (Theorem~\ref{thm:curated-depthsum}).

The five chapters of this sequence---marketplace competition
(Chapter~\ref{ch:marketplace}), stable matching
(Chapter~\ref{ch:matching}), coalition value
(Chapter~\ref{ch:coalition}), ideatic auction
(Chapter~\ref{ch:auction}), and repertoire curation---together provide a
comprehensive game-theoretic framework for cultural dynamics within
Ideatic Diffusion Theory.  Each chapter isolates a different mode of
cultural interaction---competition, pairing, cooperation, bidding,
management---and derives structural constraints that hold across all
instantiations of the IDT axioms.  The resulting picture is one of
\emph{bounded cultural possibility}: every mode of interaction is
constrained by subadditivity, identity, and resonance, the three pillars
on which the theory rests.


% ── Part IV: Cooperation, Networks, and Synthesis ─────────────────────
\part{Cooperation, Networks, and Synthesis}
\label{part:synthesis}

% ============================================================
%  Chapter 16 — Cooperative Composition and Gift Exchange
% ============================================================
\chapter{Cooperative Composition and Gift Exchange}
\label{ch:cooperative}

\begin{quote}
\textit{``The gift not yet repaid debases the man who accepted it.''}\\
--- Marcel Mauss, \textit{The Gift} (1925)
\end{quote}

\vspace{1em}

\section{Introduction}

Marcel Mauss's celebrated \textit{Essai sur le don} (1925) demonstrated that
gift exchange is never merely economic: to give is to create obligation, to
receive is to assume debt, and to reciprocate is to weave a social bond.
In societies from the Trobriand Islands to the Pacific Northwest, the
circulation of gifts constitutes a total social fact---simultaneously
economic, juridical, moral, and aesthetic.\footnote{Mauss distinguished three
obligations of the gift: to give, to receive, and to reciprocate. Each is
indispensable; failure at any stage ruptures the social fabric.}

In Ideatic Diffusion Theory, cooperation is \emph{composition}: when two agents
contribute their ideas to a joint project, the result is $\compose(a, b)$.
The order matters---gifts are not symmetric!---and the void $\void$
contributes nothing.  This chapter formalizes cooperative game theory within
ideatic space, establishing twenty theorems that characterize joint
composition, multi-party projects, depth bounds on cooperative output, and
resonance preservation under collaboration.

The analogy between $\compose$ and gift-giving is precise.  Agent $a$
``gives first,'' setting the frame of the exchange; agent $b$ reciprocates
within that frame.  The result $\compose(a,b)$ is a compound idea whose
depth is bounded by the sum of the contributors' depths---one cannot create
profundity from banality---and whose resonance properties are inherited from
the inputs.  When we extend to multi-party projects via iterated
composition, we obtain a formal model of Malinowski's \textit{Kula ring}, of
potlatch sequences, and of any process in which meaning accrues through
sequential cooperative acts.\footnote{Bronislaw Malinowski's
\textit{Argonauts of the Western Pacific} (1922) documented how shell
valuables circulate through a ring of trading partners, each exchange adding
a layer of social meaning.}

We proceed as follows.  Section~\ref{sec:ch16-defs} introduces the core
definitions---\texttt{jointCompose}, \texttt{symmetricJoint}, and
\texttt{JointProject}---together with their Lean formalizations.
Section~\ref{sec:ch16-void} establishes the role of void as the
non-contributor.  Section~\ref{sec:ch16-assoc} proves associativity of
multi-party cooperation.  Section~\ref{sec:ch16-resonance} explores
resonance preservation.  Section~\ref{sec:ch16-depth} derives depth
bounds on cooperative output.  Section~\ref{sec:ch16-multi} extends the
analysis to multi-party and iterated cooperation.  Section~\ref{sec:ch16-concl}
concludes.

\subsection*{Mauss's Three Obligations and the Structure of Reciprocity}

Mauss's analysis of gift exchange identified three fundamental
obligations: the obligation \emph{to give}, the obligation \emph{to
receive}, and the obligation \emph{to reciprocate}.  Each obligation is
indispensable; failure at any stage ruptures the social bond.  In our
formal framework, these three obligations map onto structural properties
of composition.  The obligation to give corresponds to the act of
contributing $a$ to a joint composition $\compose(a, b)$---the initiator
must provide an input.  The obligation to receive corresponds to the
requirement that the second argument $b$ accept the compositional frame
established by $a$---compositionally, $b$ operates within the context
set by $a$, which it cannot refuse without leaving the exchange entirely.
The obligation to reciprocate corresponds to the non-triviality of $b$:
reciprocation with $\void$ (Theorem~\ref{thm:void_joint_right}) leaves
the initiator's gift unchanged, which is structurally equivalent to
non-reciprocation.  A genuine return gift requires $b \neq \void$, so
that $\compose(a, b) \neq a$.

This three-fold structure reveals why gift exchange creates lasting
social bonds in a way that market transactions do not.  A market
transaction terminates at the moment of exchange: money for goods, and
the obligation is discharged.  A gift, by contrast, creates an
asymmetric temporal structure: the obligation to reciprocate persists
until a counter-gift is offered, and the counter-gift itself creates a
new obligation, generating an indefinitely extended chain of mutual
obligation.  In compositional terms, the gift exchange generates an
iterated sequence $\compose(a_1, \compose(b_1, \compose(a_2,
\compose(b_2, \ldots))))$ that is structurally open-ended---there is no
natural termination point, because each reciprocation is itself a gift
that demands reciprocation.\footnote{This open-endedness distinguishes
gift exchange from contract: a contract specifies termination conditions,
while a gift relationship is, in principle, interminable.  See Marshall
Sahlins, \textit{Stone Age Economics} (1972), Chapter 5.}

\subsection*{Sahlins's Typology of Reciprocity}

Marshall Sahlins, in \textit{Stone Age Economics} (1972), proposed a
typology of reciprocity that maps naturally onto our compositional
framework.  \emph{Generalized reciprocity} is the altruistic extreme:
$a$ gives without expectation of return, corresponding to
$\compose(a, \void)$---the gift is offered and the return is void.
\emph{Balanced reciprocity} involves equivalent counter-gifts:
$\compose(a, b)$ where $\depth(b) \approx \depth(a)$, so that the
exchange is approximately symmetric in depth.  \emph{Negative
reciprocity} is the exploitative extreme: one party attempts to extract
maximum value while giving minimum in return, corresponding to
$\compose(a, b)$ where $\depth(a) \gg \depth(b)$ or $\depth(b) \gg
\depth(a)$, and the asymmetry is intentional rather than freely chosen.

Sahlins arranged these types on a continuum correlated with social
distance: generalized reciprocity prevails among close kin (parents
give to children without expecting repayment), balanced reciprocity
governs relations among trading partners and allies, and negative
reciprocity characterizes dealings with strangers and enemies.  Our
framework captures this continuum through the depth ratio
$\depth(a)/\depth(b)$: values near 1 indicate balanced exchange,
values much greater or less than 1 indicate asymmetric exchange, and
the direction of asymmetry (which party contributes more depth)
determines whether the exchange is generous or exploitative.

\subsection*{Gregory, Weiner, and the Inalienability Problem}

Chris Gregory's \textit{Gifts and Commodities} (1982) drew a sharp
distinction between gift economies and commodity economies.  In a
commodity economy, the exchange transfers \emph{alienable} objects
between \emph{independent} parties.  In a gift economy, the exchange
transfers \emph{inalienable} possessions between \emph{interdependent}
parties.  The key difference is that the gift retains a connection to
its giver: the \textit{hau} (spirit) of the gift, in Maori terms,
compels reciprocation.

Annette Weiner, in \textit{Inalienable Possessions} (1992), extended
this analysis by identifying a paradoxical practice she called
``keeping-while-giving.''  Certain treasured possessions---royal
regalia, sacred heirlooms, clan emblems---circulate as gifts but are
never fully alienated from their original owners.  The giver retains
a claim on the object even after giving it away.  In our compositional
framework, this corresponds to an idea $a$ that participates in
composition $\compose(a, b)$ without being ``consumed'': $a$ remains
available for future compositions even after its contribution to the
current joint project.  This is structurally guaranteed by our model,
since composition does not destroy its inputs---$a$ and $b$ continue
to exist as elements of $\IS$ after $\compose(a, b)$ is computed.

The inalienability of the gift maps onto the persistence of ideas in
ideatic space.  When a scholar shares an idea through publication, the
idea enters the public repertoire without leaving the scholar's
repertoire---this is ``keeping-while-giving'' in the academic context.
The formal framework captures this naturally: expansion adds an idea
to the recipient's repertoire without removing it from the giver's.

\subsection*{Modern Gift Exchange: From Blood to Code}

The logic of gift exchange extends far beyond pre-modern societies.
Richard Titmuss's \textit{The Gift Relationship} (1970) analyzed blood
donation as a form of generalized reciprocity: donors give blood to
strangers, expecting no direct return, motivated by solidarity with an
abstract community.  In our terms, each blood donation is a contribution
$a_{\text{blood}}$ to a joint project
$\texttt{output}(\langle [a_1, a_2, \ldots, a_n] \rangle)$ whose
output is the functioning of the healthcare system.  By
Theorem~\ref{thm:empty_project_void}, the healthcare system without
donors produces $\void$---nothing.  Each donor's contribution is
indispensable in the aggregate.

Open-source software represents another form of cooperative
composition.  When a programmer contributes code to a public
repository, they compose their work with that of other contributors:
$\compose(\text{existing\_codebase}, \text{new\_contribution})$.  The
Linux kernel, with its thousands of contributors over three decades,
is a monumental joint project whose output is the iterated composition
of individual contributions.  By
Theorem~\ref{thm:iterated_self_cooperation_depth}, the depth of this
project is bounded by the sum of individual depths---a bound that, for
Linux, runs into the millions of person-years of intellectual investment.

Jacques Derrida's \textit{Given Time} (1991) posed a radical
philosophical challenge to the very concept of the gift.  Derrida
argued that a true gift---one given without any expectation of return,
without even the consciousness that one is giving---is structurally
impossible within an economy of exchange.  The moment the gift is
recognized as a gift, it enters the economy of obligation and ceases
to be gratuitous.  In our framework, Derrida's argument translates to
the observation that $\compose(a, b)$ is always structurally
\emph{responsive}: the second argument $b$ is shaped by $a$, and the
result $\compose(a, b)$ bears the trace of both contributions.  A
``pure gift'' would require $\compose(a, b) = b$---that the giver's
contribution leave no trace---which is precisely the void condition
$a = \void$.  Derrida's paradox is thus that the only true gift is
no gift at all.\footnote{Lewis Hyde's \textit{The Gift} (1983) offers
a less paradoxical account, arguing that creative gifts---artworks,
stories, scientific discoveries---circulate through communities in ways
that create bonds of gratitude and obligation.  Hyde's ``gift economy''
is closer to Mauss's balanced reciprocity than to Derrida's impossible
ideal.}

% ============================================================
\section{Core Definitions}
\label{sec:ch16-defs}

We begin with the fundamental structures of cooperative composition.

\begin{definition}[Joint Composition]
\label{def:jointCompose}
Let $a, b \in \IS$.  The \emph{joint composition} of $a$ and $b$ is
\[
  \texttt{jointCompose}(a, b) \;:=\; \compose(a, b).
\]
In Mauss's terms, agent $a$ gives first and agent $b$ reciprocates.
The result is a shared compound idea.
\end{definition}

\begin{lstlisting}[caption={Joint composition in Lean}]
def jointCompose (a b : I) : I := IdeaticSpace.compose a b
\end{lstlisting}

Joint composition is simply composition itself, but the naming emphasizes the
cooperative framing: two agents are contributing sequentially to a shared
outcome.  The asymmetry is essential.  In Mauss's analysis, the initiator of
an exchange determines its moral character; the recipient responds within a
framework already established.

\begin{definition}[Symmetric Joint]
\label{def:symmetricJoint}
Two ideas $a, b$ are \emph{symmetrically joint} when the order of
composition does not matter:
\[
  \texttt{symmetricJoint}(a, b) \;:\Longleftrightarrow\;
  \compose(a, b) = \compose(b, a).
\]
\end{definition}

\begin{lstlisting}[caption={Symmetric joint in Lean}]
def symmetricJoint (a b : I) : Prop :=
  IdeaticSpace.compose a b = IdeaticSpace.compose b a
\end{lstlisting}

Symmetric joints are rare.  In most real exchanges---ritual, economic,
discursive---who gives first shapes the meaning of the entire exchange.
The symmetric case corresponds to perfectly reciprocal partnerships where
neither party holds agenda-setting power.

\begin{definition}[Joint Project]
\label{def:JointProject}
A \emph{joint project} consists of a list of participants, each contributing
an idea.  The project's \emph{output} is the left-fold of all contributions
via composition, starting from $\void$:
\[
  \texttt{output}(p) \;:=\; \text{foldl}(\compose, \void, p.\text{participants}).
\]
The \emph{total depth} of a joint project is $\sum_{a \in p} \depth(a)$.
\end{definition}

\begin{lstlisting}[caption={Joint project structure in Lean}]
structure JointProject (I : Type*) [IdeaticSpace I] where
  participants : List I

def JointProject.output (p : JointProject I) : I :=
  p.participants.foldl IdeaticSpace.compose IdeaticSpace.void

def JointProject.totalDepth (p : JointProject I) : N :=
  depthSum p.participants
\end{lstlisting}

The left-fold semantics mean that the first contributor sets the initial
frame (after void), and each subsequent contributor composes onto the
accumulating result.  This models the sequential accretion of meaning in
collaborative enterprises---from barn-raisings to co-authored papers.

\begin{example}[The Kula Ring as Iterated Cooperative Composition]
\label{ex:kula-ring}
The Kula Ring, documented by Bronislaw Malinowski in \textit{Argonauts
of the Western Pacific} (1922) and further analyzed by Nancy Munn in
\textit{The Fame of Ganda} (1986), is among the best-studied systems of
ceremonial gift exchange.  In the Kula, shell valuables (\textit{vaga},
arm-shells, and \textit{mwali}, necklaces) circulate in opposite
directions through a ring of trading partners spanning the islands of
eastern Papua New Guinea.

In our formal framework, the Kula Ring is an iterated joint project.
Let $a_1, a_2, \ldots, a_n$ represent the traders arranged in a ring,
and let $g_i$ represent the gift offered by trader $a_i$.  The
composition $\compose(g_1, \compose(g_2, \ldots, \compose(g_{n-1},
g_n)\ldots))$ represents the accumulated social meaning of the exchange
cycle.  By Theorem~\ref{thm:gift_triangle_depth} (generalized to $n$
parties), the depth of this composed meaning is bounded by
$\sum_{i=1}^n \depth(g_i)$.

What makes the Kula remarkable, from our formal perspective, is that
the valuables themselves gain depth through circulation.  Munn showed
that a shell's \emph{fame}---its social value---increases with each
exchange, because the history of its previous owners becomes part of
its meaning.  In IDT terms, the gift $g_i$ at step $i$ is not the
same as the ``raw'' shell but $\compose(g_{i-1}, g_i)$: the shell
\emph{composes with} each new exchange context.  The bounded depth of
each step (by subadditivity) means that the shell's accumulated meaning
grows at most linearly with the number of exchanges---but it does grow,
and this growth is precisely what gives old, much-traveled shells their
extraordinary prestige.
\end{example}

\begin{example}[Blood Donation as Generalized Reciprocity]
\label{ex:titmuss}
Richard Titmuss's \textit{The Gift Relationship} (1970) compared
voluntary (British) and commercial (American) blood donation systems,
arguing that the voluntary system produced higher-quality blood,
greater supply reliability, and stronger social cohesion.  In IDT
terms, the voluntary system models \emph{generalized reciprocity}
(Sahlins's typology): each donor $a_i$ contributes to a joint project
$\texttt{output}(\langle [a_1, \ldots, a_n] \rangle)$ without knowing
or caring about the identity of the recipient.

The key formal property is that generalized reciprocity requires
$b = \void$ in the donor's local exchange: the donor gives
($\compose(a_{\text{donor}}, \void) = a_{\text{donor}}$ by
Theorem~\ref{thm:void_joint_right}) without receiving any direct
counter-gift.  The ``return'' is not personal reciprocation but the
existence of the system itself---the assurance that if the donor ever
needs blood, others will give in the same spirit.  This is reciprocity
displaced in time and social space: the counter-gift comes from the
\emph{system}, not from any individual partner.

Titmuss's finding that voluntary systems outperform commercial ones
suggests that generalized reciprocity can be more efficient than
balanced exchange when the ``commodity'' (blood) is difficult to verify
for quality.  The formal insight is that when $\depth(a)$ is
observable only imperfectly, the gift logic (which appeals to social
bonds and moral obligation) may elicit higher-depth contributions than
the market logic (which offers monetary compensation for a commodity
whose depth---quality---is hard to price).
\end{example}

% ============================================================
\section{The Void as Non-Contributor}
\label{sec:ch16-void}

Void is the identity element of composition, and in the cooperative context
it represents non-participation: silence before or after a gift.

\begin{theorem}[Void Contributes Nothing---Left (16.1)]
\label{thm:void_joint_left}
For all $b \in \IS$,
\[
  \texttt{jointCompose}(\void, b) = b.
\]
\end{theorem}

\noindent\textit{Proof sketch.}\quad
Direct from the left-identity law of $\compose$: $\compose(\void, b) = b$.
\hfill$\square$

\medskip

\noindent\textbf{Commentary.}
Silence before a gift changes nothing about the gift.  The Maussian
interpretation is that non-participation at the outset imposes no frame on
the exchange; the recipient's contribution stands alone.

\begin{theorem}[Void Contributes Nothing---Right (16.2)]
\label{thm:void_joint_right}
For all $a \in \IS$,
\[
  \texttt{jointCompose}(a, \void) = a.
\]
\end{theorem}

\noindent\textit{Proof sketch.}\quad
From the right-identity law: $\compose(a, \void) = a$.
\hfill$\square$

\medskip

\noindent\textbf{Commentary.}
A gift followed by silence is just the gift.  Reciprocation with void---an
empty counter-gift---leaves the initiator's contribution unaltered.  This
formalizes the asymmetry that Mauss identified: the absence of reciprocation
does not annihilate the gift, but it leaves the social bond incomplete.

\begin{theorem}[Void Symmetry (16.7)]
\label{thm:void_symmetric_left}
For all $a \in \IS$, $\texttt{symmetricJoint}(\void, a)$.
\end{theorem}

\noindent\textit{Proof sketch.}\quad
Both $\compose(\void, a)$ and $\compose(a, \void)$ equal $a$, so
the equation holds trivially.
\hfill$\square$

\medskip

\noindent\textbf{Commentary.}
Void always commutes---silence is the universal symmetric partner.
Non-participation is the one ``gift'' whose order never matters.
This is both mathematically trivial and anthropologically significant:
the absence of contribution creates no ordering conflict.

\begin{theorem}[Void Symmetry---Self (16.8)]
\label{thm:void_symmetric_self}
$\texttt{symmetricJoint}(\void, \void)$.
\end{theorem}

\noindent\textit{Proof sketch.}\quad
$\compose(\void, \void) = \void = \compose(\void, \void)$.
\hfill$\square$

\medskip

\noindent\textbf{Commentary.}
As Wittgenstein remarked, ``Whereof one cannot speak, thereof one must be
silent''---and the order of two silences is immaterial.  This theorem, while
trivial, anchors the degenerate case of the theory.

% ============================================================
\section{Associativity and Multi-Party Cooperation}
\label{sec:ch16-assoc}

The associativity of composition ensures that multi-party cooperation is
well-defined regardless of how sub-coalitions are bracketed.

\begin{theorem}[Associativity of Multi-Party Cooperation (16.3)]
\label{thm:joint_assoc}
For all $a, b, c \in \IS$,
\[
  \texttt{jointCompose}(\texttt{jointCompose}(a, b), c) \;=\;
  \texttt{jointCompose}(a, \texttt{jointCompose}(b, c)).
\]
\end{theorem}

\noindent\textit{Proof sketch.}\quad
Unfolds to $\compose(\compose(a,b), c) = \compose(a, \compose(b,c))$,
which is the associativity axiom of $\compose$.
\hfill$\square$

\medskip

\noindent\textbf{Commentary.}
In a ritual with three participants, it doesn't matter whether $A$ and $B$
cooperate first and then engage $C$, or $A$ engages a pre-formed $B$--$C$
coalition.  The overall composite meaning is the same.  This is a powerful
structural result: it means that multi-party cooperation is intrinsically
well-defined, independent of the contingent order in which sub-coalitions
form.\footnote{In cooperative game theory, associativity of the
characteristic function is often assumed but rarely proved from first
principles.  Here it follows from the algebraic structure of ideatic
space.}

\begin{theorem}[Four-Party Cooperation Associativity (16.18)]
\label{thm:four_party_assoc}
For all $a, b, c, d \in \IS$,
\[
  \texttt{jointCompose}(\texttt{jointCompose}(\texttt{jointCompose}(a, b), c), d) \;=\;
  \texttt{jointCompose}(a, \texttt{jointCompose}(b, \texttt{jointCompose}(c, d))).
\]
\end{theorem}

\noindent\textit{Proof sketch.}\quad
Iterated application of the associativity axiom of $\compose$.
\hfill$\square$

\medskip

\noindent\textbf{Commentary.}
In organizational theory, four-person teams can be sub-divided many ways:
$(AB)(CD)$, $A(B(CD))$, $((AB)C)D$, and so on.  Associativity guarantees
that all bracketings produce the same final product.  The insight extends to
$n$-party cooperation by induction, establishing that the output of a
\texttt{JointProject} is invariant under any re-bracketing of its participant
list.

\begin{theorem}[Empty Project Produces Void (16.9)]
\label{thm:empty_project_void}
\[
  \texttt{output}(\langle [\ ] \rangle) = \void.
\]
\end{theorem}

\noindent\textit{Proof sketch.}\quad
The fold over an empty list returns the initial accumulator $\void$.
\hfill$\square$

\medskip

\noindent\textbf{Commentary.}
An unfunded, unstaffed project yields silence.  This is the organizational
analogue of the void identity: without participants, no meaning is produced.

\begin{theorem}[Singleton Project Preserves Contribution (16.10)]
\label{thm:singleton_project}
For all $a \in \IS$,
\[
  \texttt{output}(\langle [a] \rangle) = a.
\]
\end{theorem}

\noindent\textit{Proof sketch.}\quad
The fold computes $\compose(\void, a) = a$ by left-identity.
\hfill$\square$

\medskip

\noindent\textbf{Commentary.}
A solo artisan's work is their own idea, unmediated by any collaborator.
The project structure does not distort a lone contributor's output.

\begin{example}[Open-Source Software as Cooperative Composition]
\label{ex:open-source}
Open-source software development provides a paradigmatic modern example
of cooperative composition.  Consider the Linux kernel, initiated by
Linus Torvalds in 1991 and developed collaboratively by thousands of
contributors.  Each contribution---a patch, a driver, a subsystem---is
composed with the existing codebase: the project output at time $t$ is
$\texttt{output}(\langle [c_1, c_2, \ldots, c_t] \rangle)$, the
left-fold composition of all contributions in chronological order.

By Theorem~\ref{thm:empty_project_void}, the kernel before any
contribution is $\void$: an empty project produces no output.  By
Theorem~\ref{thm:singleton_project}, Torvalds's initial contribution
$c_1$ is the project in its first form.  Each subsequent contribution
$c_i$ composes with the accumulated result, and by
Theorem~\ref{thm:joint_depth_bound}, the depth of the new output is
bounded by the depth of the existing codebase plus the depth of the new
contribution.

The open-source model is structurally a gift economy: contributors
``give'' their code to the project without receiving direct monetary
compensation.  The counter-gift is the use of the project itself (a
form of generalized reciprocity), the prestige associated with
contribution (a form of symbolic capital in Bourdieu's sense), and the
learning that comes from composing one's ideas with those of expert
collaborators.  Eric Raymond's distinction between the ``cathedral''
(centrally designed software) and the ``bazaar'' (decentralized
open-source development) maps onto our framework as the difference
between a single-composer project and a multi-party joint project with
heterogeneous contributors and emergent compositional structure.
\end{example}

\begin{example}[Academic Co-authorship as Joint Composition]
\label{ex:coauthorship}
Academic co-authorship is a particularly transparent instance of joint
composition.  When two scholars write a paper together, each brings
their individual repertoire---their knowledge, methods, and
theoretical commitments---and the resulting paper is $\compose(a, b)$
where $a$ and $b$ represent the respective contributions.

The order of composition (who is ``first author'') has both social and
structural significance.  In many fields, first authorship signals
primary intellectual contribution---the first author ``sets the frame''
within which the second author's contribution is interpreted.  This
maps precisely onto the asymmetry of $\compose(a, b) \neq \compose(b,
a)$ in the general case (non-symmetric joints, per
Definition~\ref{def:symmetricJoint}).

By Theorem~\ref{thm:joint_depth_bound}, the depth of the co-authored
paper is bounded by $\depth(a) + \depth(b)$: two scholars cannot
produce a paper deeper than the sum of their individual intellectual
capacities.  By Theorem~\ref{thm:trivial_cooperation}, if both
contributions have depth zero, the joint paper has depth zero---
shallow thinking combined with shallow thinking produces only shallow
output.  This formalizes the common academic complaint that
collaborations between mediocre researchers produce mediocre results:
the depth bound is real, and no amount of collaboration can overcome
the limitation of the inputs.
\end{example}

% ============================================================
\section{Resonance Under Cooperation}
\label{sec:ch16-resonance}

Resonance---the fundamental relation of ideatic compatibility---behaves
well under cooperative composition.  This section establishes that
cooperation preserves, propagates, and in some cases strengthens resonance.

\begin{theorem}[Resonant Partners Produce Resonant Outcomes (16.4)]
\label{thm:resonant_partners_resonant_output}
If $\text{resonates}(a_1, a_2)$ and $\text{resonates}(b_1, b_2)$, then
\[
  \text{resonates}\bigl(\texttt{jointCompose}(a_1, b_1),\;
  \texttt{jointCompose}(a_2, b_2)\bigr).
\]
\end{theorem}

\noindent\textit{Proof sketch.}\quad
Follows from the \texttt{resonance\_compose} lemma, which establishes
that resonance is preserved under componentwise composition.
\hfill$\square$

\medskip

\noindent\textbf{Commentary.}
When gift-givers are spiritually aligned---in Mauss's language, when the
\textit{hau} (spirit of the gift) circulates freely---their combined gifts
carry a resonance that mismatched givers cannot achieve.  This theorem
formalizes the intuition that cultural consonance between cooperating parties
produces consonant outcomes.

\begin{theorem}[Reciprocal Exchange Resonance (16.15)]
\label{thm:reciprocal_resonance}
If $\text{resonates}(a, b)$, then
\[
  \text{resonates}\bigl(\texttt{jointCompose}(a, b),\;
  \texttt{jointCompose}(b, a)\bigr).
\]
\end{theorem}

\noindent\textit{Proof sketch.}\quad
Since $\text{resonates}(a, b)$ implies $\text{resonates}(b, a)$ by
symmetry, applying \texttt{resonance\_compose} to the pairs
$(a, b)$ and $(b, a)$ yields the result.
\hfill$\square$

\medskip

\noindent\textbf{Commentary.}
The counter-gift resonates with the gift when giver and receiver are in
spiritual accord.  This is Mauss's central insight formalized: reciprocal
exchange between resonant partners creates a \emph{cycle} of resonant
meaning, regardless of which direction the gift flows.

\begin{theorem}[Self-Resonance of Joint Output (16.11)]
\label{thm:joint_self_resonance}
For all $a, b \in \IS$,
\[
  \text{resonates}\bigl(\texttt{jointCompose}(a, b),\;
  \texttt{jointCompose}(a, b)\bigr).
\]
\end{theorem}

\noindent\textit{Proof sketch.}\quad
Self-resonance is a general property: every idea resonates with itself.
\hfill$\square$

\medskip

\noindent\textbf{Commentary.}
Any completed exchange, however imperfect, achieves internal coherence.
This is the minimal condition for a meaningful cooperative outcome.

\begin{theorem}[Right-Resonance Preservation (16.12)]
\label{thm:joint_resonance_right}
If $\text{resonates}(a, b)$, then for all $c$,
$\text{resonates}\bigl(\texttt{jointCompose}(a, c),\; \texttt{jointCompose}(b, c)\bigr)$.
\end{theorem}

\noindent\textit{Proof sketch.}\quad
Follows from \texttt{resonance\_compose\_right}: resonant agents cooperating
with the same partner produce resonant results.
\hfill$\square$

\medskip

\noindent\textbf{Commentary.}
Durkheim's insight that agents sharing collective consciousness produce
similar ritual outcomes is formalized here: if $a$ and $b$ are culturally
resonant, their joint works with any fixed third party $c$ will resonate.

\begin{theorem}[Left-Resonance Preservation (16.13)]
\label{thm:joint_resonance_left}
If $\text{resonates}(a, b)$, then for all $c$,
$\text{resonates}\bigl(\texttt{jointCompose}(c, a),\; \texttt{jointCompose}(c, b)\bigr)$.
\end{theorem}

\noindent\textit{Proof sketch.}\quad
From \texttt{resonance\_compose\_left}: the same frame applied to resonant
inputs produces resonant outputs.
\hfill$\square$

\medskip

\noindent\textbf{Commentary.}
Bourdieu's concept of habitus finds formal expression here: the same habitus
(frame $c$) applied to resonant inputs produces resonant interpretations.
The frame preserves the resonance of its inputs.

\begin{remark}[The Asymmetry of Composition and Power in Gift Exchange]
\label{rem:asymmetry-power}
The non-commutativity of composition---$\compose(a, b) \neq \compose(b, a)$
in general---has deep implications for the politics of gift exchange.
In Mauss's analysis, the initiator of an exchange exercises a form of
power: by giving first, the initiator establishes the ``frame'' within
which the recipient must respond.  The Kwakiutl potlatch illustrates
this dynamic: the chief who gives the most lavish feast establishes a
standard that rivals must match or exceed, structuring the entire field
of reciprocal obligation.

In compositional terms, the asymmetry means that $\compose(a, b)$ bears
the ``signature'' of $a$ more prominently than that of $b$: the first
argument establishes the context, and the second argument operates
within it.  When the depth of $a$ far exceeds the depth of $b$
($\depth(a) \gg \depth(b)$), the joint output is dominated by $a$'s
contribution---a formal model of the asymmetric gift that creates
subordination rather than solidarity.  Pierre Bourdieu recognized this
dynamic in \textit{The Logic of Practice} (1980): ``the gift that
cannot be reciprocated creates a lasting asymmetry of obligation, which
is another name for power.''

The formal diagnosis of exploitative gift exchange is thus:
$\depth(a) \gg \depth(b)$ combined with social pressure to accept the
composition $\compose(a, b)$ as binding.  Development aid, with its
attendant conditionalities, often takes this form: the donor's
contribution sets a frame that the recipient cannot escape without
rejecting the gift entirely.
\end{remark}

\begin{remark}[Connection to Chapter~14's Auction Theory]
\label{rem:auction-connection}
The cooperative composition framework of this chapter is the
\emph{complement} of the competitive auction framework developed in
Chapter~\ref{ch:auction}.  In an auction, agents compete to have their
ideas selected; in cooperative composition, agents collaborate to
produce a joint idea.  The two frameworks meet in the concept of the
\emph{coalition}: a group of agents who cooperate internally while
competing externally.

Formally, the joint project $\texttt{output}(\langle [a_1, \ldots,
a_k] \rangle)$ produces a composite idea that the coalition then
``enters'' into the marketplace competition of
Chapter~\ref{ch:marketplace} or ``bids'' in the auction of
Chapter~\ref{ch:auction}.  The depth bound of
Theorem~\ref{thm:joint_depth_bound} constrains the coalition's
competitive capacity: no coalition can produce an output deeper than
the sum of its members' depths, which limits the competitive advantage
of collaboration.  This connection between cooperative and competitive
frameworks mirrors the real-world dynamics of firms, political parties,
and cultural movements: internal cooperation generates a collective
product that then competes in the external marketplace of ideas.
\end{remark}

% ============================================================
\section{Depth Bounds on Cooperative Output}
\label{sec:ch16-depth}

Depth measures interpretive complexity.  This section establishes that
cooperative composition respects subadditivity: the complexity of a joint
venture cannot exceed the total intellectual capital invested.

\begin{theorem}[Depth Bound on Joint Output (16.5)]
\label{thm:joint_depth_bound}
For all $a, b \in \IS$,
\[
  \depth\bigl(\texttt{jointCompose}(a, b)\bigr) \;\leq\;
  \depth(a) + \depth(b).
\]
\end{theorem}

\noindent\textit{Proof sketch.}\quad
Direct from the subadditivity of depth under composition:
$\depth(\compose(a,b)) \leq \depth(a) + \depth(b)$.
\hfill$\square$

\medskip

\noindent\textbf{Commentary.}
The complexity of a joint venture cannot exceed the total intellectual
capital invested by both partners.  This is a fundamental resource constraint
on cooperation: two partners of bounded depth can only produce output of
bounded depth.  The analogy to economic production functions is precise---the
output is bounded by the sum of inputs.

\begin{theorem}[Joint With Self Equals ComposeN 2 (16.6)]
\label{thm:joint_self_is_composeN2}
For all $a \in \IS$,
\[
  \texttt{jointCompose}(a, a) = \texttt{composeN}(a, 2).
\]
\end{theorem}

\noindent\textit{Proof sketch.}\quad
By definition, $\texttt{composeN}(a, 2) = \compose(a, \compose(a, \void))
= \compose(a, a)$.
\hfill$\square$

\medskip

\noindent\textbf{Commentary.}
Self-reinforcement through repeated practice is formally identical to
composing an idea with itself.  In ritual theory, the repetition of a
ceremony deepens its effect; here, the depth of the repeated composition is
at most $2 \cdot \depth(a)$ by subadditivity.

\begin{theorem}[Depth Zero Closure Under Cooperation (16.14)]
\label{thm:trivial_cooperation}
If $\depth(a) = 0$ and $\depth(b) = 0$, then
$\depth\bigl(\texttt{jointCompose}(a, b)\bigr) = 0$.
\end{theorem}

\noindent\textit{Proof sketch.}\quad
From the depth-zero closure property of composition.
\hfill$\square$

\medskip

\noindent\textbf{Commentary.}
Shallow discourse combined with shallow discourse can never produce depth.
This is the critical-theory insight that Adorno and Horkheimer might have
appreciated: the culture industry's endless recombination of trivial content
produces only trivial content.\footnote{Theodor Adorno and Max Horkheimer,
\textit{Dialectic of Enlightenment} (1944), argued that mass culture
recombines familiar elements without producing genuine novelty.}

\begin{theorem}[Cooperation Monotonicity for Depth (16.17)]
\label{thm:void_cooperation_depth}
For all $a \in \IS$,
$\depth\bigl(\texttt{jointCompose}(a, \void)\bigr) = \depth(a)$.
\end{theorem}

\noindent\textit{Proof sketch.}\quad
Since $\compose(a, \void) = a$, the depth is preserved.
\hfill$\square$

\medskip

\noindent\textbf{Commentary.}
Free-riders (void contributors) don't add complexity to the joint output,
but they don't reduce it either.  The depth of the project is unchanged
by the presence or absence of non-contributing participants---a result with
clear implications for team management and organizational design.

\begin{remark}[The Spirit of the Gift (\textit{Hau}) as Resonance]
\label{rem:hau-resonance}
Mauss famously invoked the Maori concept of \textit{hau}---the ``spirit
of the gift''---to explain why gifts must be reciprocated.  The
\textit{hau} of an object compels its recipient to make a return; failure
to reciprocate invites spiritual danger.  While Mauss's ethnographic
interpretation has been debated (Marshall Sahlins offered an influential
reinterpretation in \textit{Stone Age Economics}), the structural insight
remains: something in the gift \emph{persists} beyond the physical
transfer, creating an ongoing connection between giver and receiver.

In our framework, the \textit{hau} maps naturally onto \emph{resonance}.
When $a$ gives a gift $g$ to $b$, and $g$ resonates with $b$'s
repertoire ($\text{resonates}(g, r_b)$ for some $r_b \in
\mathbf{R}_b$), the gift creates a compositional connection:
$\compose(g, r_b)$ is a new idea in $b$'s interpretive space that bears
the trace of both the gift and the recipient's prior understanding.
This resonance-trace is the formal analogue of the \textit{hau}: it is
the structural mechanism by which the gift ``stays with'' the giver
even after being transferred.

Theorem~\ref{thm:resonant_partners_resonant_output} formalizes the
positive dynamics of the \textit{hau}: when giver and receiver are in
resonance, their joint compositions produce resonant outcomes.  The
\textit{hau} circulates because resonance is preserved under
composition: the spirit of the gift propagates through the exchange
network as long as the partners remain culturally aligned.  When
resonance breaks down---when the gift fails to resonate with the
recipient's repertoire---the \textit{hau} is lost, and the gift becomes
a mere commodity: an alienable object stripped of its spiritual
connection to the giver.
\end{remark}

% ============================================================
\section{Multi-Party Projects and Iterated Cooperation}
\label{sec:ch16-multi}

We now extend the analysis to multi-party projects, iterated
self-cooperation, and the celebrated gift-exchange cycles of
anthropological theory.

\begin{theorem}[Joint Project Depth Sum Bound (16.16)]
\label{thm:project_depth_bound}
If every participant $a$ in project $p$ satisfies $\depth(a) \leq d$, then
\[
  \texttt{totalDepth}(p) \;\leq\; |p.\text{participants}| \cdot d.
\]
\end{theorem}

\noindent\textit{Proof sketch.}\quad
The total depth is a sum of individual depths, each bounded by $d$.
The sum of $n$ terms each at most $d$ is at most $n \cdot d$.
\hfill$\square$

\medskip

\noindent\textbf{Commentary.}
Total intellectual capital invested constrains what the project can achieve.
This is the project-management analogue of the depth bound: a team of $n$
members, each of bounded capability $d$, produces a project whose total
complexity is at most $n \cdot d$.

\begin{theorem}[Iterated Self-Cooperation Depth Bound (16.19)]
\label{thm:iterated_self_cooperation_depth}
For all $a \in \IS$ and $n \in \mathbb{N}$,
\[
  \depth\bigl(\texttt{composeN}(a, n)\bigr) \;\leq\; n \cdot \depth(a).
\]
\end{theorem}

\noindent\textit{Proof sketch.}\quad
By induction on $n$, using the subadditivity of depth at each step.
\hfill$\square$

\medskip

\noindent\textbf{Commentary.}
Repeating a ritual $n$ times has bounded cumulative depth, growing at most
linearly with repetitions.  This formalizes the intuition that practice
deepens understanding, but at a rate bounded by the practitioner's own
depth.  The logarithmic ``diminishing returns'' often observed in ritual
practice are consistent with this linear upper bound.

\begin{theorem}[Gift Exchange Triangle (16.20)]
\label{thm:gift_triangle_depth}
For all $a, b, c \in \IS$,
\[
  \depth\bigl(\texttt{jointCompose}(\texttt{jointCompose}(a, b), c)\bigr)
  \;\leq\; \depth(a) + \depth(b) + \depth(c).
\]
\end{theorem}

\noindent\textit{Proof sketch.}\quad
Apply the depth bound twice: first to $\compose(a, b)$ and $c$, obtaining
$\depth(\compose(\compose(a,b), c)) \leq \depth(\compose(a,b)) + \depth(c)$,
then to $a$ and $b$, obtaining $\depth(\compose(a,b)) \leq \depth(a) +
\depth(b)$.  Combining by transitivity yields the three-term bound.
\hfill$\square$

\medskip

\noindent\textbf{Commentary.}
This theorem formalizes Malinowski's Kula ring: gifts circulating through a
ring of three partners have bounded total complexity, ensuring the exchange
remains manageable.  The three-term depth bound generalizes to $n$-term
bounds by induction (see Theorem~\ref{thm:iterated_self_cooperation_depth}),
but the triangular case is particularly significant because it corresponds
to the minimal non-trivial cycle of exchange.

% ============================================================
\section{From Mauss to the Digital Commons}
\label{sec:ch16-digital-commons}

The transition from Mauss's analysis of Melanesian shell exchange to
contemporary digital gift economies is not as great a leap as it might
appear.  The structural properties formalized in this chapter---identity,
associativity, resonance preservation, depth subadditivity---hold
regardless of whether the ``gifts'' are shell necklaces, blood, lines
of code, or Wikipedia edits.  This section traces the formal continuities
and examines what changes when gift exchange goes digital.

\subsection{The Creative Commons as Compositional Infrastructure}

The Creative Commons licensing system, founded in 2001, provides a
legal framework for ``keeping-while-giving'' in the digital age.  When
an author releases a work under a Creative Commons license, they give
the work to the public while retaining certain rights (attribution,
non-commercial use, share-alike).  In our framework, this is precisely
Weiner's ``keeping-while-giving'': the idea $a$ enters the public
repertoire via $\operatorname{expand}(\mathbf{R}_{\text{public}}, a)$
without leaving the author's repertoire
$\mathbf{R}_{\text{author}}$.

The ``share-alike'' provision (CC BY-SA) has a particularly elegant
formal interpretation.  It requires that any derivative work---any
composition $\compose(a, b)$ that builds on the original---must itself
be released under the same license.  This is a \emph{compositional
constraint}: the predicate ``is licensed CC BY-SA'' is preserved under
composition.  If $a$ is CC BY-SA and $b$ is composed with it, then
$\compose(a, b)$ must also be CC BY-SA.  The share-alike condition
thus propagates through the compositional network, ensuring that the
gift economy of creative works is self-sustaining: each act of giving
(releasing under CC BY-SA) constrains all future compositions to
remain within the gift economy.

\subsection{Wikipedia as a Multi-Party Joint Project}

Wikipedia, with its hundreds of thousands of active editors, is perhaps
the largest joint project in human history.  Each Wikipedia article is a
joint project in our formal sense: $\texttt{output}(\langle [e_1, e_2,
\ldots, e_n] \rangle)$, where $e_i$ is the $i$-th edit.  By
Theorem~\ref{thm:empty_project_void}, an article before any edits is
$\void$---the empty page.  By Theorem~\ref{thm:singleton_project}, the
first edit $e_1$ establishes the article.  Each subsequent edit composes
with the accumulated content, and the associativity of composition
(Theorem~\ref{thm:joint_assoc}) ensures that the article's meaning is
independent of how we bracket the editorial history.

The ``edit war'' phenomenon---in which opposing editors repeatedly
overwrite each other's contributions---reveals the limits of our
compositional model.  In a strict edit war, the composition oscillates
between $\compose(a, b)$ and $\compose(b, a)$, which are generically
different (non-symmetric joints).  The resolution of edit wars through
consensus mechanisms (talk pages, mediation, administrative fiat)
corresponds to the social process of selecting a canonical composition
order---a choice that our formal framework shows is structurally
consequential, since the output depends on the order of contributions.

\subsection{Cryptocurrency and the Gift-Commodity Boundary}

The emergence of non-fungible tokens (NFTs) and digital collectibles
blurs the boundary between gift and commodity economies in ways that
our framework illuminates.  An artist who ``mints'' an NFT creates a
digital artifact that functions simultaneously as a commodity (it has a
market price) and a gift (it carries the ``aura'' of the creator,
creating a personal connection between artist and collector).  The
\textit{hau} of the NFT is its provenance: the blockchain record of
its creation and subsequent transfers, which is formally analogous to
the compositional history $\compose(a_1, \compose(a_2, \ldots))$ that
accumulates social meaning through circulation.

\paragraph{Extended proof discussion: The gift exchange triangle.}
Theorem~\ref{thm:gift_triangle_depth}, the gift exchange triangle,
merits a more detailed proof analysis.  The theorem states that for
three participants $a, b, c$, the depth of $\compose(\compose(a, b),
c)$ is at most $\depth(a) + \depth(b) + \depth(c)$.  The proof
proceeds by two applications of subadditivity:
\begin{enumerate}
  \item First, $\depth(\compose(\compose(a, b), c)) \le
    \depth(\compose(a, b)) + \depth(c)$ by direct subadditivity.
  \item Second, $\depth(\compose(a, b)) \le \depth(a) + \depth(b)$
    by subadditivity again.
  \item Combining by transitivity of $\le$:
    $\depth(\compose(\compose(a, b), c)) \le
    (\depth(a) + \depth(b)) + \depth(c) =
    \depth(a) + \depth(b) + \depth(c)$.
\end{enumerate}
This telescoping argument generalizes to $n$ participants by induction:
for any joint project $\langle [a_1, \ldots, a_n] \rangle$, the depth
of the output is at most $\sum_{i=1}^n \depth(a_i)$.  The inductive
step uses exactly the same two-step argument, applying subadditivity to
peel off the last participant and invoking the inductive hypothesis for
the remaining $n-1$.

The bound is tight when composition is ``depth-preserving''---i.e., when
$\depth(\compose(a, b)) = \depth(a) + \depth(b)$.  Whether any
non-trivial ideatic space achieves this equality is an open question.
The bound is loose when composition is ``depth-absorbing''---i.e., when
$\depth(\compose(a, b)) \ll \depth(a) + \depth(b)$, as when combining
two redundant ideas produces nothing new.  The gap between the bound and
the actual depth measures the degree of ``compositional redundancy''
between the contributors---a quantity related to mutual information in
information-theoretic terms.

\paragraph{Extended proof discussion: Associativity and its
computational significance.}
The associativity theorem (Theorem~\ref{thm:joint_assoc}) is
algebraically elementary---it follows directly from the associativity
axiom of composition---but its computational and organizational
significance is profound.  In a distributed computing context,
associativity means that a multi-party computation can be parallelized:
instead of computing $\compose(a_1, \compose(a_2, \compose(a_3,
a_4)))$ sequentially (depth $O(n)$), we can compute
$\compose(\compose(a_1, a_2), \compose(a_3, a_4))$ in parallel
(depth $O(\log n)$).  The outputs are identical by associativity.

In organizational terms, this means that hierarchical and flat team
structures produce the same joint output.  A hierarchical organization
computes $\compose(\compose(a_1, a_2), \compose(a_3, a_4))$---pairs
cooperate and then their outputs are composed at the next level.  A
flat organization computes $\compose(a_1, \compose(a_2, \compose(a_3,
a_4)))$---each member contributes sequentially to an accumulating
whole.  Associativity guarantees identical results, which means that
the choice between hierarchical and flat structures is driven by
considerations of communication cost and coordination overhead, not
by the quality of the output.

% ============================================================
\section{Conclusion}
\label{sec:ch16-concl}

This chapter has established the formal foundations of cooperative
composition in ideatic space.  Twenty theorems characterize the algebraic
and metric properties of joint composition, multi-party projects, and
iterated cooperation.

The central results are:
\begin{itemize}
  \item \textbf{Identity:} Void contributes nothing
    (Theorems~\ref{thm:void_joint_left}--\ref{thm:void_joint_right}).
  \item \textbf{Associativity:} Multi-party cooperation is bracket-independent
    (Theorems~\ref{thm:joint_assoc}--\ref{thm:four_party_assoc}).
  \item \textbf{Resonance preservation:} Resonant partners produce resonant
    outcomes (Theorems~\ref{thm:resonant_partners_resonant_output}--\ref{thm:joint_resonance_left}).
  \item \textbf{Depth subadditivity:} Cooperative complexity cannot exceed
    the sum of individual complexities (Theorems~\ref{thm:joint_depth_bound}--\ref{thm:gift_triangle_depth}).
\end{itemize}

Mauss's \textit{Essai sur le don} revealed that gift exchange is a total
social fact.  Our formalization shows that it is also a total \emph{algebraic}
fact: the structure of ideatic space---its identity, associativity, and
subadditivity---completely determines the bounds and resonance properties
of cooperative meaning-making.  In the next chapter, we turn to what happens
when cooperation breaks down: the problem of bargaining over composition order.

% ============================================================
%  Chapter 17 — Bargaining Over Composition Order
% ============================================================
\chapter{Bargaining Over Composition Order}
\label{ch:bargaining}

\begin{quote}
\textit{``The definition of the situation is itself the first move in the game.''}\\
--- Victor Turner, \textit{The Ritual Process} (1969)
\end{quote}

\vspace{1em}

\section{Introduction}

Composition in ideatic space is associative but \emph{not} commutative.
In general, $\compose(a, b) \neq \compose(b, a)$: who speaks first shapes
the meaning of the entire exchange.  This fundamental asymmetry is the
engine of power in signalling games.  The ability to set the agenda---to
determine what is discussed first, which frame is laid down before the other
party responds---is, as Victor Turner recognized, a primary form of social
control.\footnote{Turner's analysis of ritual process emphasizes that the
\emph{order} of ritual stages determines the transformation undergone by
participants.  The liminal phase must precede reaggregation; reversing the
order produces a different---often incoherent---ritual.}

In political discourse, this is the problem of \emph{agenda-setting power}.
In legal reasoning, it is the question of which precedent is applied first.
In Hegelian dialectics, it is the difference between thesis-antithesis and
antithesis-thesis: the reversal produces a different synthesis.

This chapter formalizes the bargaining problem that arises from
non-commutativity.  We introduce the \texttt{compositionOrder} pair, the
\texttt{orderInvariant} predicate, and the \texttt{OrderPair} structure,
then prove eighteen theorems characterizing when order matters, when it
doesn't, and what structural properties hold regardless of ordering.

The key insight is that while the \emph{content} of composition depends on
order, certain \emph{structural} properties---depth bounds, self-resonance,
and the behavior of void---are order-invariant.  This creates a formal
analogy to fair bargaining: the pie may be divided differently depending on
who moves first, but the total size of the pie is bounded the same way
regardless.

\subsection{Axiomatic Bargaining and the Composition Problem}

The mathematical study of bargaining begins with John Nash's landmark
1950 paper, which showed that under four axioms---Pareto efficiency,
symmetry, invariance to affine transformations, and independence of
irrelevant alternatives---there exists a unique bargaining solution, the
\emph{Nash bargaining solution}.\footnote{John F.\ Nash, ``The Bargaining
Problem,'' \textit{Econometrica} 18, no.\ 2 (1950): 155--162.}  Nash's
framework assumes that the bargaining set is fixed and symmetric.  But in
ideatic composition, the bargaining problem is more fundamental: the parties
must agree not merely on a division of surplus but on the \emph{order in
which their contributions are composed}.  Because $\compose(a,b) \neq
\compose(b,a)$ in general, the choice of ordering is itself a strategic
variable, not a parameter.

Ariel Rubinstein's alternating-offers model (1982) extended Nash's static
framework into a dynamic setting, showing that when parties take turns
making proposals, the equilibrium outcome converges to the Nash solution as
the discount factor approaches unity.\footnote{Ariel Rubinstein, ``Perfect
Equilibrium in a Bargaining Model,'' \textit{Econometrica} 50, no.\ 1
(1982): 97--109.}  In our framework, Rubinstein's alternating offers
correspond directly to alternating composition orders: if $A$ proposes
$\compose(a,b)$ and $B$ proposes $\compose(b,a)$, the bargaining process is
a negotiation over which composition is realized.  The key distinction is
that in Rubinstein's model, the proposals converge to a common outcome; in
non-commutative composition, there may be no convergence---the two orderings
may yield genuinely different ideas.

\subsection{Agenda-Setting, Power, and Framing}

Political scientists have long recognized that controlling the
\emph{agenda}---the sequence in which issues are considered---is a primary
form of political power.  William Riker's theory of
\emph{heresthetics}\footnote{William H.\ Riker, \textit{The Art of
Political Manipulation} (1986).} identifies agenda manipulation as one of
three fundamental political strategies: you win not by persuading voters
to change their preferences, but by changing which alternatives they
consider and in what order.  John Kingdon's theory of ``policy
windows''\footnote{John W.\ Kingdon, \textit{Agendas, Alternatives, and
Public Policies}, 2nd ed.\ (2003).} similarly emphasizes that the timing
and sequencing of political proposals determines which policies are adopted.

Steven Lukes's three-dimensional theory of power provides a deeper
framework.\footnote{Steven Lukes, \textit{Power: A Radical View}, 2nd ed.\
(2005).}  The \emph{first dimension} is overt decision-making power; the
\emph{second} is the power to set the agenda, determining which issues are
even considered; the \emph{third} is the power to shape desires and
perceptions, so that dominated parties do not even recognize the
alternatives foreclosed by the agenda.  In IDT terms, the first dimension
corresponds to choosing $a$ versus $b$; the second to choosing
$\compose(a,b)$ versus $\compose(b,a)$; and the third to the ability to
make agents unaware that $\compose(b,a)$ is even a possibility.

Kahneman and Tversky's prospect theory\footnote{Daniel Kahneman and Amos
Tversky, ``Prospect Theory: An Analysis of Decision Under Risk,''
\textit{Econometrica} 47, no.\ 2 (1979): 263--291.} demonstrated that the
\emph{framing} of a decision problem---whether outcomes are presented as
gains or losses relative to a reference point---systematically alters
choices.  Framing is a form of composition ordering: to frame a policy as
``preventing losses'' rather than ``forgoing gains'' is to compose the
evaluative frame \emph{before} the policy content, producing
$\compose(\text{loss-frame}, \text{policy})$ rather than
$\compose(\text{gain-frame}, \text{policy})$.  Our formalism shows that
these two compositions are, in general, genuinely different ideas---not
merely psychologically different presentations of the same idea.

The concept of the ``Overton window''---the range of policies considered
politically acceptable at a given time---can also be understood as a
constraint on composition order.  The window determines which ideas can
serve as the \emph{first} component in a policy composition: only ideas
within the window are available as opening frames.  Shifting the Overton
window is thus equivalent to expanding the set of permissible first-movers
in the bargaining over composition order.

\subsection{Historical Parallels}

The Treaty of Westphalia (1648) offers a paradigmatic illustration.  The
Catholic and Protestant powers negotiated not merely over territorial
boundaries but over which \emph{framework} would serve as the basis for
the emerging state system.  The principle \textit{cuius regio, eius
religio} was itself a meta-bargain: the sovereign's religious identity
would be composed first, and the subjects' beliefs would follow.  In modern
terms, the Westphalian settlement resolved a composition-order dispute by
establishing a fixed ordering rule.

Similarly, the Camp David Accords (1978) required Anwar Sadat and Menachem
Begin to negotiate not just territorial concessions but the \emph{sequence}
in which issues would be addressed.  Jimmy Carter's genius as mediator lay
in restructuring the agenda: by separating the Sinai question from the
Palestinian question and addressing them in a specific order, he changed
the composition of the negotiation itself.\footnote{William B.\ Quandt,
\textit{Camp David: Peacemaking and Politics} (1986).}

In legislative settings, the power to determine the order of business---the
``rules of the game''---is often more consequential than any individual
vote.  The Speaker of the House, the Senate Majority Leader, and committee
chairs all derive power primarily from their ability to set the
\emph{sequence} in which proposals are composed.  Our formalization gives
this intuition mathematical content: the composition order determines the
outcome, and the agenda-setter controls the composition order.

\section{Core Definitions}
\label{sec:ch17-defs}

\begin{definition}[Composition Order Pair]
\label{def:compositionOrder}
The \emph{composition order pair} of $a, b \in \IS$ is
\[
  \texttt{compositionOrder}(a, b) \;:=\;
  \bigl(\compose(a, b),\; \compose(b, a)\bigr) \;\in\; \IS \times \IS.
\]
This captures both possible sequencings of a bilateral exchange: ``I go
first'' versus ``you go first.''
\end{definition}

\begin{lstlisting}[caption={Composition order pair in Lean}]
def compositionOrder (a b : I) : I * I :=
  (IdeaticSpace.compose a b, IdeaticSpace.compose b a)
\end{lstlisting}

\begin{definition}[Order Invariance]
\label{def:orderInvariant}
Two ideas $a, b$ are \emph{order-invariant} when composition commutes:
\[
  \texttt{orderInvariant}(a, b) \;:\Longleftrightarrow\;
  \compose(a, b) = \compose(b, a).
\]
This is a strong condition: most idea-compositions are \emph{not}
order-invariant.  When it holds, there is no agenda-setting advantage.
\end{definition}

\begin{lstlisting}[caption={Order invariance in Lean}]
def orderInvariant (a b : I) : Prop :=
  IdeaticSpace.compose a b = IdeaticSpace.compose b a
\end{lstlisting}

\begin{definition}[Order Pair Structure]
\label{def:OrderPair}
An \emph{order pair} packages two ideas with both their compositions
and proofs that the compositions are correctly formed:
\[
  \texttt{OrderPair} = \bigl(l, r, f, s \;\mid\;
    f = \compose(l, r),\; s = \compose(r, l)\bigr).
\]
\end{definition}

\begin{lstlisting}[caption={Order pair structure in Lean}]
structure OrderPair (I : Type*) [IdeaticSpace I] where
  left : I
  right : I
  first : I
  second : I
  ab_eq : first = IdeaticSpace.compose left right
  ba_eq : second = IdeaticSpace.compose right left
\end{lstlisting}

% ============================================================
\section{Illustrative Examples}
\label{sec:ch17-examples}

Before proceeding to the theorems, we present several extended examples that
illustrate how composition-order bargaining manifests across diverse human
institutions.

\begin{example}[The Treaty of Westphalia]
\label{ex:westphalia}
The Peace of Westphalia (1648) illustrates how composition order shapes
political outcomes. The Catholic and Protestant powers each sought to
establish their theological framework as the ``first mover'' in the
emerging international order. The principle \textit{cuius regio, eius
religio}---the ruler determines the religion---was itself a bargain over
composition order: the sovereign's religious identity is composed first,
and the subjects' beliefs follow. In IDT terms, the sovereign's idea $a$
is composed before the population's idea $b$, yielding $\compose(a, b)$
rather than $\compose(b, a)$---a fundamentally different political
theology.\footnote{The Westphalian system is analyzed through a
game-theoretic lens in Daniel Philpott, \textit{Revolutions in
Sovereignty} (2001).}

Had the composition been reversed---the subjects' beliefs composed before
the sovereign's---the result would have been a form of popular religious
sovereignty, akin to congregationalism.  The entire trajectory of European
political development turned on which ordering was institutionalized.
The Westphalian bargain was, at its core, a bargain over $\compose(a,b)$
versus $\compose(b,a)$.
\end{example}

\begin{example}[Robert's Rules of Order]
\label{ex:roberts_rules}
Parliamentary procedure, codified in Henry M.\ Robert's
\textit{Rules of Order} (1876), is a systematic framework for resolving
composition-order disputes in deliberative assemblies.  The rules specify
an elaborate protocol for determining which proposals are considered first:
the ``order of business'' determines the composition sequence.

A \emph{motion to amend} creates a composition-order problem: should the
original motion $a$ be composed with the amendment $b$ to yield
$\compose(a, b)$, or should the amended version $\compose(b, a)$ be
considered as the base text?  Robert's Rules resolve this by stipulating
that amendments are composed \emph{after} the main motion---the main
motion provides the frame, and the amendment modifies it.  The
\emph{motion to substitute} reverses this ordering entirely: it proposes
to replace $a$ with $b$, effectively switching from $\compose(a, b)$ to
$\compose(b, a)$.

The procedural complexity of parliamentary law---subsidiary motions,
privileged motions, incidental motions---can be understood as a
sophisticated mechanism for managing the exponential growth of possible
composition orderings as the number of proposals increases.  With $n$
proposals, there are $n!$ possible orderings, and the rules of order
serve to reduce this combinatorial explosion to a manageable sequence.
\end{example}

\begin{example}[Media Agenda-Setting: McCombs and Shaw]
\label{ex:media_agenda}
Maxwell McCombs and Donald Shaw's agenda-setting theory (1972)\footnote{Maxwell
E.\ McCombs and Donald L.\ Shaw, ``The Agenda-Setting Function of Mass
Media,'' \textit{Public Opinion Quarterly} 36, no.\ 2 (1972): 176--187.}
demonstrated that the media do not tell people \emph{what} to think, but
rather \emph{what to think about}.  In IDT terms, the media's power lies
not in choosing $a$ over $b$, but in choosing $\compose(a, b)$ over
$\compose(b, a)$: by presenting issue $a$ before issue $b$, the media
establish $a$ as the interpretive frame through which $b$ is subsequently
understood.

Consider a news cycle in which economic data $a$ is presented before a
political scandal $b$.  The public receives $\compose(a, b)$: the scandal
is interpreted through an economic lens.  If the order is reversed---the
scandal is presented first, followed by the economic data---the public
receives $\compose(b, a)$: the economy is interpreted through the lens of
political corruption.  These are, in general, genuinely different public
ideas, not merely different presentations of the same information.

The rise of algorithmic news feeds has intensified this phenomenon.  Each
user's feed presents information in a personalized order, producing
personalized compositions $\compose(a_i, b_i)$ that may differ
dramatically across individuals.  The ``filter bubble'' is, formally, a
personalized composition ordering.
\end{example}

\begin{example}[Socratic Dialogue as Composition-Order Bargaining]
\label{ex:socratic}
In the Platonic dialogues, Socrates consistently positions himself as the
\emph{questioner}, compelling his interlocutors to state their theses
first.  This is a strategic choice of composition order: by ensuring that
the interlocutor's idea $b$ is composed before Socrates's own idea $a$,
Socrates obtains $\compose(a, b)$---his philosophical framework applied
\emph{to} the interlocutor's claim---rather than $\compose(b, a)$, which
would be the interlocutor's framework applied to Socrates's claim.

The \emph{elenchus}---Socrates's method of cross-examination---is a
systematic technique for ensuring that the composition order favors the
questioner.  By asking questions, Socrates forces his interlocutor to
provide the raw material (the first component), which Socrates then
composes with his own dialectical frame (the second component).  The result
is that the interlocutor's thesis is always interpreted \emph{within}
Socrates's framework, never the reverse.

This observation has implications for modern pedagogical theory.  The
``Socratic method'' widely used in legal education is effective precisely
because it enforces a composition order in which the student's initial
response is composed with the professor's analytical frame, producing a
synthesis that the student experiences as self-generated discovery.
\end{example}

% ============================================================
\section{Void and Order Invariance}
\label{sec:ch17-void}

Void---silence, non-contribution---is the universal commuter.  It has no
agenda-setting power because it contributes nothing.

\begin{theorem}[Void is Order-Invariant---Left (17.1)]
\label{thm:void_order_invariant}
For all $a \in \IS$, $\texttt{orderInvariant}(\void, a)$.
\end{theorem}

\noindent\textit{Proof sketch.}\quad
$\compose(\void, a) = a = \compose(a, \void)$ by the identity laws.
\hfill$\square$

\medskip

\noindent\textbf{Commentary.}
Silence has no agenda-setting power.  Whether you speak into silence or
silence follows your speech, the result is the same---just your speech.
In political theory, a party that brings nothing to the table cannot
influence the agenda.

\begin{theorem}[Self-Composition is Order-Invariant (17.2)]
\label{thm:self_order_invariant}
For all $a \in \IS$, $\texttt{orderInvariant}(a, a)$.
\end{theorem}

\noindent\textit{Proof sketch.}\quad
$\compose(a, a) = \compose(a, a)$ trivially.
\hfill$\square$

\medskip

\noindent\textbf{Commentary.}
Internal dialogue with yourself has no ordering problem---you always agree
with yourself on what comes first.  This is the psychological ground truth
that anchors the theory: self-composition is always order-invariant.

\begin{theorem}[Order Invariance of Void Pair (17.8)]
\label{thm:void_void_invariant}
$\texttt{orderInvariant}(\void, \void)$.
\end{theorem}

\noindent\textit{Proof sketch.}\quad
Immediate from both identity laws. \hfill$\square$

\medskip

\noindent\textbf{Commentary.}
Two silences are the same regardless of which comes first.  The degenerate
case anchors the extremal behavior of the theory.

\begin{theorem}[Void is Universally Order-Invariant (17.12)]
\label{thm:void_universal_order_invariant}
For all $a \in \IS$, $\texttt{orderInvariant}(a, \void)$.
\end{theorem}

\noindent\textit{Proof sketch.}\quad
$\compose(a, \void) = a = \compose(\void, a)$.
\hfill$\square$

\medskip

\noindent\textbf{Commentary.}
The absence of contribution never creates ordering conflicts.  Void is
universally transparent to reordering, a formal counterpart to the
observation that silence is the one ``move'' in a game that is costless
in either position.\footnote{In mechanism design, a ``null action'' is
typically assumed to be costless and position-independent. Void formalizes
this assumption within IDT.}

\begin{theorem}[Void Annihilates Order Difference (17.18)]
\label{thm:void_annihilates_order}
For all $a \in \IS$,
\[
  (\texttt{compositionOrder}(a, \void))_1 =
  (\texttt{compositionOrder}(a, \void))_2.
\]
\end{theorem}

\noindent\textit{Proof sketch.}\quad
Both components reduce to $a$ by the identity laws.
\hfill$\square$

\medskip

\noindent\textbf{Commentary.}
A player who contributes nothing has no bargaining power over ordering---there
is nothing to order.  The composition order pair collapses to a diagonal
element when one party contributes void.

% ============================================================
\section{Structural Order Theorems}
\label{sec:ch17-structural}

\begin{theorem}[Order Invariance is Symmetric (17.11)]
\label{thm:order_invariant_symm}
If $\texttt{orderInvariant}(a, b)$, then $\texttt{orderInvariant}(b, a)$.
\end{theorem}

\noindent\textit{Proof sketch.}\quad
If $\compose(a,b) = \compose(b,a)$, then symmetry of equality gives
$\compose(b,a) = \compose(a,b)$.
\hfill$\square$

\medskip

\noindent\textbf{Commentary.}
If swapping $A$ and $B$'s contributions doesn't matter, then swapping $B$
and $A$'s doesn't matter either.  This ensures that order invariance is a
symmetric relation---a necessary condition for it to serve as a formal model
of ``fair exchange.''

\begin{theorem}[Composition Order First Component (17.4)]
\label{thm:order_fst}
$(\texttt{compositionOrder}(a, b))_1 = \compose(a, b)$.
\end{theorem}

\begin{theorem}[Composition Order Second Component (17.5)]
\label{thm:order_snd}
$(\texttt{compositionOrder}(a, b))_2 = \compose(b, a)$.
\end{theorem}

\noindent\textit{Proof sketch.}\quad
Both are definitional equalities.
\hfill$\square$

\medskip

\noindent\textbf{Commentary.}
These projection lemmas are technical necessities, extracting the ``$a$-goes-first''
and ``$b$-goes-first'' orderings from the pair.  They establish the interface
between the order pair abstraction and the underlying composition operation.

\begin{theorem}[Composition Order Swap is Involution (17.17)]
\label{thm:order_swap_involution}
\[
  \texttt{compositionOrder}(b, a) =
  \bigl((\texttt{compositionOrder}(a, b))_2,\;
  (\texttt{compositionOrder}(a, b))_1\bigr).
\]
\end{theorem}

\noindent\textit{Proof sketch.}\quad
Expanding definitions: $\texttt{compositionOrder}(b,a) = (\compose(b,a),
\compose(a,b))$, which equals the swap of $(\compose(a,b), \compose(b,a))$.
\hfill$\square$

\medskip

\noindent\textbf{Commentary.}
The dialectical reversal of a reversal returns to the thesis.  Swapping who
goes first twice recovers the original ordering.  This involution property
ensures that the bargaining structure has no higher-order ordering ambiguities.

\begin{theorem}[Order Pair Construction (17.14)]
\label{thm:order_pair_exists}
For any $a, b \in \IS$, there exists an \texttt{OrderPair} with
$\text{left} = a$ and $\text{right} = b$.
\end{theorem}

\noindent\textit{Proof sketch.}\quad
Construct the pair with $\text{first} = \compose(a,b)$,
$\text{second} = \compose(b,a)$, and the proofs are \texttt{rfl}.
\hfill$\square$

\medskip

\noindent\textbf{Commentary.}
Every bilateral exchange can be formally analyzed by examining both possible
orderings.  The \texttt{OrderPair} structure packages this analysis into a
single object, facilitating the systematic study of agenda-setting effects.

% ============================================================
\section{Power, Framing, and the First-Mover Advantage}
\label{sec:ch17-power}

The preceding structural theorems reveal a striking dichotomy in
composition-order bargaining: the \emph{existence} of an order pair is
guaranteed for any two ideas, yet the \emph{content} of its two components
may diverge arbitrarily.  This section explores the consequences of this
dichotomy for theories of power, framing, and strategic interaction.

\subsection{The First-Mover Advantage in Formal Terms}

In game theory, the \emph{first-mover advantage} refers to the strategic
benefit accruing to the player who acts first in a sequential game.  In
composition-order bargaining, the first-mover advantage is the difference
between $\compose(a, b)$ and $\compose(b, a)$.  When $a$ ``goes first,''
the resulting idea $\compose(a, b)$ is \emph{$a$-framed}: $a$ provides the
interpretive scaffolding onto which $b$ is composed.

The magnitude of the first-mover advantage can be measured by the
``distance'' between the two orderings.  While IDT does not assume a metric
on $\IS$, the resonance relation provides a qualitative measure: when
$\text{resonates}(\compose(a,b), \compose(b,a))$, the first-mover advantage
is ``small'' in the sense that both orderings are at least mutually
coherent.  When the two orderings fail to resonate, the first-mover
advantage is ``large''---the outcomes are so different as to be
incoherent with each other.

Theorem~\ref{thm:resonant_orders} establishes that when the input ideas
resonate, the outputs always resonate.  This provides a formal foundation
for the folk intuition that bargaining between like-minded parties is
``easier'' than bargaining between adversaries: when the parties share a
common frame (resonance), the composition order matters less.

\subsection{Framing as Composition-Order Selection}

The framing literature in behavioral economics identifies two types of
framing effects: \emph{valence framing} (gain vs.\ loss frames) and
\emph{emphasis framing} (which aspects of an issue are highlighted).  Both
can be modeled as composition-order choices.

In valence framing, the same objective outcome $x$ is presented either as
$\compose(\text{gain-frame}, x)$ or $\compose(\text{loss-frame}, x)$.
Kahneman and Tversky's experiments show that these compositions are treated
as different ideas by human cognizers---formally vindicating our axiom that
composition is non-commutative.

In emphasis framing, the question is not which frame to apply but
\emph{when} to apply it.  Presenting a policy's costs before its benefits
yields $\compose(\text{cost}, \text{benefit})$; the reverse yields
$\compose(\text{benefit}, \text{cost})$.  Empirical research consistently
shows that information presented first has disproportionate influence on
final judgments---a primacy effect that is precisely the content-level
manifestation of composition non-commutativity.

\begin{remark}[Connection to Mechanism Design]
\label{rem:mechanism-design}
In mechanism design and implementation theory, the ``designer'' chooses
a game form (mechanism) to implement a desired social choice function.
Composition-order bargaining can be viewed as a special case: the
mechanism designer selects an ordering rule $\sigma: \{a,b\} \to
\{(a,b), (b,a)\}$ that determines which composition is realized.
Theorems~\ref{thm:void_order_invariant}--\ref{thm:void_annihilates_order}
show that void is the unique idea for which the mechanism designer's
choice is irrelevant.  For all other ideas, the designer's ordering
rule has non-trivial consequences---a result reminiscent of the
Gibbard--Satterthwaite theorem's impossibility of
strategy-proofness.\footnote{The connection to Gibbard--Satterthwaite
is suggestive rather than formal.  A full investigation of
composition-order mechanism design is left to future work.}
\end{remark}

\begin{remark}[The Impossibility of Order-Free Composition]
\label{rem:order-free-impossibility}
One might hope for a ``canonical'' composition that is independent of
order---perhaps the set $\{\compose(a,b), \compose(b,a)\}$, or some
average of the two.  But this approach fails for a fundamental reason:
the resulting object is no longer an element of ideatic space $\IS$, but
rather a subset or a distribution over $\IS$.  Order-free composition
requires moving to a higher-order structure---the power set $\mathcal{P}(\IS)$
or the space of distributions over $\IS$.  This observation motivates the
social-choice constructions in Chapter~\ref{ch:social-choice}, where
aggregation rules are introduced precisely to combine multiple compositions
into a single shared idea.
\end{remark}

\begin{remark}[Connections to Gift Exchange and Social Choice]
\label{rem:cross-chapter}
Composition-order bargaining connects directly to the gift-exchange
dynamics of Chapter~\ref{ch:gift-exchange} and the social-choice
mechanisms of Chapter~\ref{ch:social-choice}.  In gift exchange, the
sequential nature of reciprocity creates a natural composition ordering:
the first gift $a$ is composed before the reciprocal gift $b$, and the
social meaning of the exchange $\compose(a, b)$ differs from the meaning
of the reverse sequence $\compose(b, a)$.  The obligation to reciprocate
is, formally, a constraint on composition order.  In social choice, the
aggregation of multiple agents' ideas requires resolving the composition
ordering for all pairs simultaneously---a far more complex problem that
generalizes bilateral bargaining to the $n$-party case.
\end{remark}

% ============================================================
\section{Depth and Resonance Under Reordering}
\label{sec:ch17-depth}

While the \emph{content} of a composition depends on order, its
\emph{depth bounds} are symmetric: subadditivity holds for both orderings.

\begin{theorem}[Both Orders Have Same Depth Bound (17.3)]
\label{thm:both_orders_depth_bound}
For all $a, b \in \IS$,
\begin{align*}
  \depth(\compose(a, b)) &\leq \depth(a) + \depth(b), \\
  \depth(\compose(b, a)) &\leq \depth(a) + \depth(b).
\end{align*}
\end{theorem}

\noindent\textit{Proof sketch.}\quad
The first inequality is subadditivity.  The second follows from
$\depth(\compose(b,a)) \leq \depth(b) + \depth(a) = \depth(a) + \depth(b)$.
\hfill$\square$

\medskip

\noindent\textbf{Commentary.}
The ``cost'' of an exchange, measured by depth, is bounded the same way
regardless of who goes first.  Subadditivity is symmetric.  This means that
while ordering affects the \emph{content} of the outcome, it does not affect
the \emph{resource constraints}---a crucial fairness property.

\begin{theorem}[Depth Bound Under Reordering (17.10)]
\label{thm:depth_reorder_bound}
$\depth(\compose(b, a)) \leq \depth(a) + \depth(b)$.
\end{theorem}

\noindent\textit{Proof sketch.}\quad
From subadditivity with reordering of the bound.
\hfill$\square$

\medskip

\noindent\textbf{Commentary.}
The maximum information capacity of an exchange is invariant under
reordering.  In information-theoretic terms, the channel capacity is
bounded by the sum of the endpoint depths, regardless of direction.

\begin{theorem}[Depth Zero Invariance Under Order (17.16)]
\label{thm:depth_zero_both_orders}
If $\depth(a) = 0$ and $\depth(b) = 0$, then
$\depth(\compose(a,b)) = 0$ and $\depth(\compose(b,a)) = 0$.
\end{theorem}

\noindent\textit{Proof sketch.}\quad
Apply depth-zero closure in both orders.
\hfill$\square$

\medskip

\noindent\textbf{Commentary.}
Trivial exchanges are trivial regardless of who goes first.  You can't
create depth by reordering banalities---a critical-theory insight with
mathematical teeth.

\begin{theorem}[Both Orders Self-Resonate (17.6)]
\label{thm:both_orders_self_resonate}
$\text{resonates}(\compose(a,b), \compose(a,b))$ and
$\text{resonates}(\compose(b,a), \compose(b,a))$.
\end{theorem}

\noindent\textit{Proof sketch.}\quad
Self-resonance is universal.
\hfill$\square$

\medskip

\noindent\textbf{Commentary.}
Any completed exchange achieves internal coherence regardless of who spoke
first.  This is the minimal hermeneutic property: every outcome is
self-consistent.

\begin{theorem}[Resonant Inputs Give Resonant Orders (17.7)]
\label{thm:resonant_orders}
If $\text{resonates}(a, b)$, then
$\text{resonates}\bigl(\compose(a, b),\; \compose(b, a)\bigr)$.
\end{theorem}

\noindent\textit{Proof sketch.}\quad
From \texttt{resonance\_compose} applied to the pairs $(a,b)$ and
$(b,a)$, using the symmetry of resonance.
\hfill$\square$

\medskip

\noindent\textbf{Commentary.}
This is Turner's ritual-process theorem.  When ritual actors are in
resonance---in Turner's terminology, when they share \emph{communitas}---reordering
the ritual stages produces ceremonies that still ``rhyme'' with each other.
The spirit of the ritual is preserved even if the sequence changes.
Formally, resonant inputs produce resonant orderings, ensuring that the
bargaining problem between culturally aligned agents has a coherent
solution space.\footnote{Turner introduced the concept of \emph{communitas}
to describe the intense fellow-feeling that arises in liminal ritual stages.
Our resonance relation formalizes this as a structural property of ideatic space.}

\medskip

\noindent\textbf{Extended discussion of Theorem~\ref{thm:resonant_orders}.}\quad
The resonance-preservation result is the most consequential theorem in this
section for social theory.  It establishes that cultural alignment between
agents is \emph{robust} to composition-order variation: if two agents are
``on the same wavelength'' (resonant), then the two possible orderings of
their exchange are themselves on the same wavelength.  This has three
important implications.

First, it provides a formal criterion for when bargaining over composition
order is \emph{low-stakes}: if the parties resonate, both orderings produce
coherent outcomes, so the distributional consequences of order choice are
bounded.  This aligns with the empirical observation that negotiations
between allies (resonant parties) focus on substance rather than procedure.

Second, it explains why \emph{communitas}---the feeling of deep social
connection---facilitates ritual flexibility.  Turner observed that
communities experiencing strong communitas can improvise ritual sequences
without losing the ritual's transformative power.  Theorem~\ref{thm:resonant_orders}
provides the algebraic explanation: resonance between ritual elements
ensures that reordering preserves coherence.

Third, the contrapositive is equally revealing: if two orderings
$\compose(a,b)$ and $\compose(b,a)$ fail to resonate, then the inputs
$a$ and $b$ must themselves fail to resonate.  Incoherent bargaining
outcomes diagnose incoherent input ideas---a formal version of the
negotiator's maxim that ``if the process breaks down, the parties were
never really aligned.''

% ============================================================
\section{Associativity and Multi-Party Bargaining}
\label{sec:ch17-multi}

\begin{theorem}[Associativity Resolves Three-Party Brackets (17.9)]
\label{thm:three_party_bracket_invariance}
$\compose(\compose(a,b), c) = \compose(a, \compose(b,c))$.
\end{theorem}

\noindent\textit{Proof sketch.}\quad
The associativity axiom of ideatic composition.
\hfill$\square$

\medskip

\noindent\textbf{Commentary.}
In three-party negotiation, whether $(A,B)$ form a coalition first and then
bargain with $C$, or $A$ bargains with a pre-formed $(B,C)$ coalition, the
final composite meaning is identical.  Associativity is the great equalizer:
it eliminates the ``coalition formation order'' as a source of bargaining
advantage, even as the \emph{within-pair} ordering remains significant.

\begin{theorem}[Triple Depth Bound---Both Bracketings (17.13)]
\label{thm:triple_depth_bound}
$\depth(\compose(\compose(a,b), c)) \leq \depth(a) + \depth(b) + \depth(c)$.
\end{theorem}

\noindent\textit{Proof sketch.}\quad
Cascade two applications of subadditivity.
\hfill$\square$

\medskip

\noindent\textbf{Commentary.}
Three-way collaboration has bounded complexity no matter how sub-teams are
formed.  The depth bound generalizes linearly, providing a universal
resource constraint on multi-party cooperation.

\begin{theorem}[Resonance Between Reordered Triples (17.15)]
\label{thm:triple_reorder_resonance}
If $\text{resonates}(a, b)$ and $\text{resonates}(b, c)$, then
\[
  \text{resonates}\bigl(\compose(\compose(a,b), c),\;
  \compose(a, \compose(b,c))\bigr).
\]
\end{theorem}

\noindent\textit{Proof sketch.}\quad
By associativity, both sides are equal, so self-resonance applies.
\hfill$\square$

\medskip

\noindent\textbf{Commentary.}
When all ritual participants share communitas, rearranging the ritual order
produces ceremonies that deeply resonate.  The resonance between re-bracketed
triples is not merely an analogy to Turner's theory---it is its formal
content.

% ============================================================

\noindent\textbf{Extended discussion of Theorem~\ref{thm:three_party_bracket_invariance}.}\quad
The associativity result for three-party bracketing has deep implications
for coalition theory.  In real-world multi-party negotiations---such as
the formation of coalition governments in parliamentary democracies---the
question of \emph{which sub-coalitions form first} is a major strategic
concern.  Associativity tells us that while the \emph{within-pair} ordering
matters (who speaks first within a dyad), the \emph{between-pair} bracketing
does not: whether Labor and the Greens form a coalition first and then
negotiate with the Liberals, or Labor negotiates with a pre-formed
Green-Liberal alliance, the final composite position is the same.

This result provides formal support for the observation in coalition theory
that the ``order of coalition formation'' matters less than the ``order of
proposal within each coalition.''  The bracketing structure drops out;
the sequencing within each bilateral exchange does not.  This is a
non-trivial constraint on the strategy space of multi-party bargaining:
parties should focus on their bilateral ordering advantages rather than
on the meta-strategic question of which coalitions form when.

The depth bound for triples (Theorem~\ref{thm:triple_depth_bound})
extends this analysis to resource constraints: the total complexity of a
three-party exchange is bounded by the sum of individual complexities,
regardless of bracketing.  This linearity is essential for scalability:
it means that the resource cost of multi-party bargaining grows
additively, not multiplicatively, ensuring that large-scale negotiations
remain tractable.

\section{Conclusion}
\label{sec:ch17-concl}

This chapter has established eighteen theorems characterizing the
bargaining problem that arises from the non-commutativity of ideatic
composition.

The central tension is clear: while \emph{content} depends on order
(who speaks first matters), \emph{structural properties}---depth bounds,
self-resonance, void transparency---do not.  This creates a formal
separation between the distributional question (``who gets the better
outcome?'') and the efficiency question (``how much total meaning is
produced?'').

The key results are:
\begin{itemize}
  \item Void has no agenda-setting power
    (Theorems~\ref{thm:void_order_invariant}--\ref{thm:void_annihilates_order}).
  \item Depth bounds are order-invariant
    (Theorems~\ref{thm:both_orders_depth_bound}--\ref{thm:depth_zero_both_orders}).
  \item Resonant agents produce resonant orderings
    (Theorem~\ref{thm:resonant_orders}).
  \item Associativity resolves multi-party bracketing
    (Theorem~\ref{thm:three_party_bracket_invariance}).
\end{itemize}

Turner's ritual process reveals that transformative social experiences
require a specific ordering of stages.  Our formalization shows that this
ordering dependence is a fundamental feature of ideatic composition, not
an artifact of particular rituals.  In the next chapter, we turn from
bilateral bargaining to the problem of diffusion through networks.

% ============================================================
%  Chapter 18 — Network Diffusion and Cultural Epidemiology
% ============================================================
\chapter{Network Diffusion and Cultural Epidemiology}
\label{ch:network}

\begin{quote}
\textit{``Culture is the precipitate of cognition and communication.''}\\
--- Dan Sperber, \textit{Explaining Culture} (1996)
\end{quote}

\vspace{1em}

\section{Introduction}

Dan Sperber's epidemiological approach to culture proposes that cultural
representations spread like epidemics: ideas pass from mind to mind,
transformed at each step by the cognitive environment of the host.  Frederic
Bartlett's serial reproduction experiments (1932) demonstrated this
empirically: a story retold through a chain of narrators mutates
systematically, with each retelling imposing the reteller's own cognitive
schemata on the material.\footnote{Bartlett's landmark study
\textit{Remembering} (1932) showed that the Kwakiutl folk tale ``The War of
the Ghosts'' was progressively rationalized, shortened, and
conventionalized as it passed through chains of English undergraduates.}

In Ideatic Diffusion Theory, each relay hop is a composition: the relayer
composes the incoming signal with their own interpretive frame.  A chain of
$n$ relayers produces
\[
  \compose\bigl(r_n, \compose(r_{n-1}, \ldots \compose(r_1, s) \ldots)\bigr).
\]
This chapter formalizes the mathematics of cultural diffusion through
networks, establishing eighteen theorems on relay chains, depth bounds,
resonance preservation, chain concatenation, and the transparency of void
relayers.

The central insight is that network diffusion is \emph{compositional}:
the output of a multi-hop relay is fully determined by the sequence of
relayers and the initial signal, and the structural properties of the
output---depth, resonance, void-transparency---follow from the algebraic
properties of composition.  This provides a rigorous foundation for
Sperber's cultural epidemiology, for meme theory in Dawkins's sense, and for
the ``telephone game'' as a formal model of cultural transmission.

\subsection{Network Science and Cultural Diffusion}

The mathematical study of networks has undergone a revolution since the
late 1990s.  Duncan Watts and Steven Strogatz's 1998 paper on
\emph{small-world networks}\footnote{Duncan J.\ Watts and Steven H.\
Strogatz, ``Collective Dynamics of `Small-World' Networks,'' \textit{Nature}
393 (1998): 440--442.} demonstrated that many real-world networks---neural
networks, power grids, social networks---exhibit a characteristic
combination of high local clustering and short global path lengths.  In a
small-world network, relay chains are surprisingly short: even in networks
with millions of nodes, most pairs are connected by chains of six or fewer
intermediaries.  For cultural diffusion, this means that ideas can traverse
an entire population through remarkably few relay hops, each of which is a
composition in our framework.

Albert-L\'aszl\'o Barab\'asi and R\'eka Albert's 1999 discovery of
\emph{scale-free networks}\footnote{Albert-L\'aszl\'o Barab\'asi and R\'eka
Albert, ``Emergence of Scaling in Random Networks,'' \textit{Science} 286,
no.\ 5439 (1999): 509--512.} revealed that many networks are dominated by a
small number of highly connected ``hub'' nodes.  In cultural diffusion,
hubs correspond to influential relayers---media outlets, thought leaders,
viral content creators---whose cognitive frames are composed with the
signals of millions.  A hub relayer $h$ with connections to $n$ receivers
produces $n$ distinct compositions $\compose(h, s_i)$, each filtered through
$h$'s interpretive frame.  The scale-free structure ensures that a small
number of such hubs dominate the diffusion dynamics.

Mark Granovetter's seminal 1973 paper on the ``strength of weak
ties''\footnote{Mark S.\ Granovetter, ``The Strength of Weak Ties,''
\textit{American Journal of Sociology} 78, no.\ 6 (1973): 1360--1380.}
identified a paradox: weak social connections (acquaintances rather than
close friends) are often more important for information diffusion than
strong ones, because weak ties serve as bridges between otherwise
disconnected clusters.  In relay-chain terms, weak ties are relayers whose
cognitive frames differ significantly from those of their local cluster,
enabling the signal to ``jump'' to a new region of ideatic space.

\subsection{From Memes to Epidemiology}

Richard Dawkins introduced the concept of the \emph{meme}---a unit of
cultural transmission analogous to the gene---in \textit{The Selfish
Gene} (1976).\footnote{Richard Dawkins, \textit{The Selfish Gene} (1976),
ch.\ 11.}  Daniel Dennett elaborated the concept in \textit{Darwin's
Dangerous Idea} (1995), arguing that memes are subject to genuine
Darwinian selection: they vary, are differentially transmitted, and are
inherited (with modification) by new hosts.  In IDT, a meme is an idea
$m \in \IS$ whose relay chain exhibits high fidelity: the composition
$\compose(r, m)$ remains ``close to'' $m$ for a wide range of relayers $r$.
High-fidelity memes are precisely those ideas that are approximately
order-invariant with respect to typical relayers.

Robert Boyd and Peter Richerson's formal models of cultural
transmission\footnote{Robert Boyd and Peter J.\ Richerson, \textit{Culture
and the Evolutionary Process} (1985).} distinguished several modes of
cultural learning: \emph{vertical} (parent to child), \emph{horizontal}
(peer to peer), and \emph{oblique} (non-parental adults to children).  Each
mode corresponds to a different relay-chain topology.  Vertical transmission
produces long, narrow chains (high depth, low branching); horizontal
transmission produces wide, shallow chains (low depth, high branching);
and oblique transmission produces intermediate structures.  The theorems
in this chapter apply to all three topologies, since they depend only on
the algebraic properties of composition, not on the network structure.

Dan Sperber's ``epidemiology of representations'' deepens the analogy
between cultural and biological contagion.\footnote{Dan Sperber,
\textit{Explaining Culture: A Naturalistic Approach} (1996).}  Unlike
Dawkins's meme theory, which treats cultural items as replicators, Sperber
emphasizes that cultural representations are \emph{transformed} at each
transmission event.  There is no faithful copying; there is only
composition.  This is precisely the insight formalized by the relay chain:
each relay hop is a composition $\compose(r_i, \text{acc})$, not a copy
operation.  The ``mutation'' that Sperber identifies in cultural
transmission is the difference between $\text{acc}$ and
$\compose(r_i, \text{acc})$---the interpretive contribution of the relayer.

\subsection{Modern Diffusion Phenomena and Formal Limitations}

Malcolm Gladwell's \textit{The Tipping Point} (2000) popularized the idea
that social epidemics undergo phase transitions: small changes in the
``stickiness'' of a message, the influence of key connectors, or the
``power of context'' can trigger cascading diffusion.  While Gladwell's
framework is more journalistic than formal, the underlying intuition finds
partial support in our theorems.  The depth bounds
(Section~\ref{sec:ch18-depth}) show that relay chains cannot create
unbounded complexity, while the resonance results
(Section~\ref{sec:ch18-resonance}) show that diffusion through homogeneous
networks preserves signal coherence.  However, our framework also
identifies limitations of the tipping-point metaphor: the linear depth
bounds suggest that cultural ``explosions'' are constrained by the
cognitive resources of the relay chain, not merely by network topology.

The SIR (Susceptible-Infected-Recovered) model from mathematical
epidemiology\footnote{W.\ O.\ Kermack and A.\ G.\ McKendrick, ``A
Contribution to the Mathematical Theory of Epidemics,'' \textit{Proceedings
of the Royal Society A} 115, no.\ 772 (1927): 700--721.} provides another
formal parallel.  In the SIR model, individuals transition from susceptible
to infected to recovered, and the epidemic spreads through contact between
susceptible and infected individuals.  In cultural epidemiology, the
``infection'' event is a composition: a susceptible mind $s$ encounters an
idea-carrying relayer $r$, producing the composition $\compose(r, s)$.
The ``recovered'' state corresponds to an idea that has been fully
integrated and is no longer being actively transmitted.  Our relay-chain
formalism captures the SIR model's transmission dynamics while adding the
algebraic structure of composition.

\section{Core Definitions}
\label{sec:ch18-defs}

\begin{definition}[Relay Chain]
\label{def:relayChain}
Let $\text{relayers} = [r_1, r_2, \ldots, r_n]$ be a list of relayer ideas
and $s \in \IS$ be a signal.  The \emph{relay chain} is the left-fold:
\[
  \texttt{relayChain}(\text{relayers}, s) \;:=\;
  \text{foldl}\bigl(\lambda\;\text{acc}\;r.\;\compose(r, \text{acc}),\;
  s,\; \text{relayers}\bigr).
\]
Each relayer composes the accumulated signal with their own frame,
left-to-right.
\end{definition}

\begin{lstlisting}[caption={Relay chain in Lean}]
def relayChain (relayers : List I) (signal : I) : I :=
  relayers.foldl (fun acc r => IdeaticSpace.compose r acc) signal
\end{lstlisting}

The fold semantics mean that $r_1$ processes the raw signal first, producing
$\compose(r_1, s)$; then $r_2$ processes $r_1$'s output, producing
$\compose(r_2, \compose(r_1, s))$; and so on.  Each relay hop adds a layer
of interpretive composition.

\begin{definition}[Chain Depth]
\label{def:chainDepth}
The \emph{chain depth} is the depth of the relay chain's output:
\[
  \texttt{chainDepth}(\text{relayers}, s) := \depth\bigl(\texttt{relayChain}(\text{relayers}, s)\bigr).
\]
\end{definition}

\begin{definition}[Network Path]
\label{def:NetworkPath}
A \emph{network path} packages a list of relayers and a signal into a
single structure.  Its output is the relay chain:
\[
  \texttt{NetworkPath.output}(p) := \texttt{relayChain}(p.\text{relayers}, p.\text{signal}).
\]
\end{definition}

\begin{lstlisting}[caption={Network path structure in Lean}]
structure NetworkPath (I : Type*) [IdeaticSpace I] where
  relayers : List I
  signal : I

def NetworkPath.output (p : NetworkPath I) : I :=
  relayChain p.relayers p.signal
\end{lstlisting}

% ============================================================
\section{Illustrative Examples}
\label{sec:ch18-examples}

Before developing the theorems of cultural diffusion, we present extended
examples that ground the formalism in empirical research on transmission
chains.

\begin{example}[Bartlett's Serial Reproduction Experiments]
\label{ex:bartlett}
Frederic Bartlett's \textit{Remembering} (1932) remains the foundational
empirical study of cultural transmission chains.  In his serial
reproduction experiments, Bartlett asked a participant to read an unfamiliar
folk tale---the Kwakiutl story ``The War of the Ghosts''---and then
reproduce it from memory.  This reproduction was given to a second
participant, who reproduced it in turn, and so on through a chain of ten or
more participants.

In IDT terms, each participant $r_i$ acts as a relayer, composing the
incoming story $s_{i-1}$ with their own cognitive schemas to produce a new
story $s_i = \compose(r_i, s_{i-1})$.  Bartlett observed several systematic
transformations: the story became shorter (depth reduction), unfamiliar
elements were replaced by familiar ones (assimilation to the relayer's
schemas), and the narrative structure was rationalized according to Western
story conventions (imposition of the relayer's cultural frame).

The relay-chain model $\texttt{relayChain}([r_1, \ldots, r_n], s)$
captures these dynamics precisely.  Each relayer's composition introduces
their cognitive biases, and the cumulative effect is a progressive
transformation of the original signal.  Bartlett called this
``effort after meaning''---the relayer's active contribution to the
composition, formalized here as the relayer's idea $r_i$.
\end{example}

\begin{example}[Viral Memes: The Ice Bucket Challenge]
\label{ex:ice-bucket}
The ALS Ice Bucket Challenge of 2014 spread to an estimated 17 million
participants worldwide within weeks, raising over \$115 million for ALS
research.  From the perspective of cultural epidemiology, the challenge
is a near-ideal case of a high-fidelity relay chain.

The ``meme'' $m$ consisted of a simple behavioral script: pour ice water
over yourself, film it, nominate three friends, donate.  Each participant
$r_i$ relayed the challenge by performing the script and posting a video,
producing $\compose(r_i, m)$.  The challenge's viral success depended on
two formal properties.  First, \emph{approximate order invariance}: the
core script was sufficiently simple that $\compose(r_i, m)$ was
recognizably similar to $m$ for most relayers $r_i$.  Second,
\emph{resonance clustering} (Theorem~\ref{thm:resonant_relayers_resonant_output}):
participants within the same social network shared cognitive schemas
(resonance), ensuring that their relay outputs were similar to each other.

The mutations that did occur---celebrity elaborations, parodies, refusals
to participate---illustrate the non-identity of $\compose(r_i, m)$ and $m$.
Each relayer's composition added their distinctive contribution, even as
the core signal was preserved.
\end{example}

\begin{example}[Ryan and Gross: Hybrid Corn Diffusion]
\label{ex:hybrid-corn}
Bryce Ryan and Neal Gross's 1943 study of hybrid corn adoption in Greene
County, Iowa,\footnote{Bryce Ryan and Neal C.\ Gross, ``The Diffusion of
Hybrid Seed Corn in Two Iowa Communities,'' \textit{Rural Sociology} 8,
no.\ 1 (1943): 15--24.} is the foundational empirical study in the
diffusion of innovations literature.  They found that adoption followed
an S-shaped curve: a slow initial phase, followed by rapid adoption, and
finally saturation.

In relay-chain terms, the innovation $s$ (hybrid corn) diffused through a
network of farmers $[r_1, r_2, \ldots, r_n]$.  The early adopters composed
the innovation with their own farming experience, producing
$\compose(r_i, s)$---a personalized understanding of the technology.  These
compositions were then relayed to neighbors, who composed them with their
own schemas.  The S-curve arises because the relay chain's depth bounds
(Theorem~\ref{thm:single_relay_depth}) ensure that each hop adds bounded
complexity, while the resonance clustering
(Theorem~\ref{thm:resonant_relayers_resonant_output}) ensures that farmers
in the same community produce similar adaptations, facilitating social
proof and further adoption.
\end{example}

\begin{example}[Rumor Transmission in Organizations]
\label{ex:rumor}
Gordon Allport and Leo Postman's classic study of rumor
transmission\footnote{Gordon W.\ Allport and Leo Postman, \textit{The
Psychology of Rumor} (1947).} identified three processes: \emph{leveling}
(shortening), \emph{sharpening} (selective emphasis), and
\emph{assimilation} (distortion toward the relayer's expectations).  Each
process corresponds to a specific aspect of the composition
$\compose(r_i, s_{i-1})$: leveling reduces depth, sharpening amplifies
salient features of the relayer's frame, and assimilation replaces
unfamiliar elements with the relayer's own schemas.

In organizational networks, rumors traverse hierarchical relay chains.
A rumor $s$ originating at the executive level passes through middle
managers $[r_1, r_2, \ldots, r_k]$ before reaching front-line employees.
Each managerial relay hop composes the rumor with the manager's
institutional perspective, producing increasingly ``managed'' versions of
the original signal.  The void-transparency results
(Theorems~\ref{thm:void_relayer_transparent}--\ref{thm:void_chain_identity})
formalize the idealized case in which a manager transmits the message
without alteration---a ``transparent'' relay that is, in practice,
exceedingly rare.
\end{example}

\section{Fundamental Transmission Theorems}
\label{sec:ch18-fundamental}

\begin{theorem}[Empty Chain Preserves Signal (18.1)]
\label{thm:empty_chain_preserves}
For all $s \in \IS$,
\[
  \texttt{relayChain}([\,], s) = s.
\]
\end{theorem}

\noindent\textit{Proof sketch.}\quad
The fold over an empty list returns the initial accumulator $s$.
\hfill$\square$

\medskip

\noindent\textbf{Commentary.}
Direct transmission without intermediaries preserves the original message
perfectly.  This is the idealized limit of communication: zero-hop
transmission is lossless.  In information theory, this corresponds to a
noiseless channel; in cultural epidemiology, it represents the (impossible)
case of unmediated cultural transmission.

\begin{theorem}[Singleton Chain = Single Composition (18.2)]
\label{thm:singleton_chain}
For all $r, s \in \IS$,
\[
  \texttt{relayChain}([r], s) = \compose(r, s).
\]
\end{theorem}

\noindent\textit{Proof sketch.}\quad
One iteration of the fold: $\compose(r, s)$.
\hfill$\square$

\medskip

\noindent\textbf{Commentary.}
A single retelling is one act of interpretive composition.  The reteller
composes the signal with their own frame, producing a new idea that carries
traces of both the original and the reteller's cognitive environment.
In Bartlett's terms, this is one round of ``serial reproduction.''

\begin{theorem}[Two-Relayer Chain (18.3)]
\label{thm:two_relayer_chain}
\[
  \texttt{relayChain}([r_1, r_2], s) = \compose(r_2, \compose(r_1, s)).
\]
\end{theorem}

\noindent\textit{Proof sketch.}\quad
Two iterations of the fold.
\hfill$\square$

\medskip

\noindent\textbf{Commentary.}
The ``telephone game'' with two players: the second player interprets the
first player's interpretation, not the original signal.  This captures the
fundamental structure of multi-hop cultural transmission: each intermediary
processes the \emph{already-processed} signal, not the original.\footnote{In
organizational communication theory, this is the source of the ``serial
transmission effect'': messages degrade as they pass through chains of
intermediaries, each of whom introduces their own cognitive biases.}

% ============================================================
\section{Void Transparency and Signal Fidelity}
\label{sec:ch18-void}

Void relayers are ``transparent''---they pass the signal through without
modification.  This models passive carriers who do not transform the message.

\begin{theorem}[Void Relayer is Transparent (18.4)]
\label{thm:void_relayer_transparent}
Appending a void relayer to the end of a chain does not change the output:
\[
  \texttt{relayChain}(\text{relayers} \mathop{++} [\void], s) =
  \texttt{relayChain}(\text{relayers}, s).
\]
\end{theorem}

\noindent\textit{Proof sketch.}\quad
The fold over the appended void computes $\compose(\void, \text{acc}) =
\text{acc}$ by the left-identity law.
\hfill$\square$

\medskip

\noindent\textbf{Commentary.}
A cognitively ``empty'' relay node transmits faithfully.  In Sperber's
framework, this models a carrier who lacks the cognitive schemas to
transform the representation---a perfect mirror for the incoming signal.

\begin{theorem}[Prepending Void Relayer is No-Op (18.15)]
\label{thm:void_prepend_chain}
\[
  \texttt{relayChain}(\void :: \text{relayers}, s) =
  \texttt{relayChain}(\text{relayers}, s).
\]
\end{theorem}

\noindent\textit{Proof sketch.}\quad
The first fold step computes $\compose(\void, s) = s$, so the remaining
fold proceeds from the same accumulator.
\hfill$\square$

\medskip

\noindent\textbf{Commentary.}
An empty preamble carries no information.  Starting a chain of retellers
with a ``silent'' relayer does not affect the final output.

\begin{theorem}[Void Signal Through Single Relay (18.5)]
\label{thm:void_signal_single_relay}
\[
  \texttt{relayChain}([r], \void) = r.
\]
\end{theorem}

\noindent\textit{Proof sketch.}\quad
$\compose(r, \void) = r$ by right-identity.
\hfill$\square$

\medskip

\noindent\textbf{Commentary.}
An empty message gets filled with the relayer's own content.  When there is
no signal, only interpretation remains---the relayer projects their own
frame onto the void.  This formalizes the constructivist insight that in the
absence of external stimuli, the interpreter's own schemas dominate.

\begin{theorem}[Void Chain is Identity (18.12)]
\label{thm:void_chain_identity}
For all $n \in \mathbb{N}$ and $s \in \IS$,
\[
  \texttt{relayChain}(\underbrace{[\void, \void, \ldots, \void]}_{n}, s) = s.
\]
\end{theorem}

\noindent\textit{Proof sketch.}\quad
By induction on $n$.  Base case: empty chain preserves $s$.  Inductive step:
the last void relayer is transparent (Theorem~\ref{thm:void_relayer_transparent}).
\hfill$\square$

\medskip

\noindent\textbf{Commentary.}
A chain of ``empty'' cognitive environments transmits perfectly.  This is
the idealized limit of faithful transmission: if every intermediary is
cognitively passive, the signal arrives unchanged.  Of course, this limit is
never achieved in practice---every mind adds something to the message.

% ============================================================
\section{Depth Bounds on Relay Chains}
\label{sec:ch18-depth}

Subadditivity cascades through relay chains, providing upper bounds on the
depth of the transmitted signal.

\begin{theorem}[Chain Depth Bound---Single Relayer (18.6)]
\label{thm:single_relay_depth}
\[
  \texttt{chainDepth}([r], s) \leq \depth(r) + \depth(s).
\]
\end{theorem}

\noindent\textit{Proof sketch.}\quad
By definition, $\texttt{chainDepth}([r], s) = \depth(\compose(r, s))$,
and subadditivity gives $\depth(\compose(r,s)) \leq \depth(r) + \depth(s)$.
\hfill$\square$

\medskip

\noindent\textbf{Commentary.}
A single retelling adds at most the reteller's own complexity to the story.
In Bartlett's experiments, the distortions introduced by a single reteller
are bounded by the reteller's cognitive capacity.

\begin{theorem}[Chain Output Depth Bound---Two Relayers (18.16)]
\label{thm:two_relay_depth_bound}
\[
  \texttt{chainDepth}([r_1, r_2], s) \leq
  \depth(r_1) + \depth(r_2) + \depth(s).
\]
\end{theorem}

\noindent\textit{Proof sketch.}\quad
Apply subadditivity twice: first to $r_2$ and $\compose(r_1, s)$, then to
$r_1$ and $s$.  Combine by transitivity.
\hfill$\square$

\medskip

\noindent\textbf{Commentary.}
Two successive retellings add at most the combined complexity of both
retellers to the original story's depth.  The bound is \emph{linear} in the
number of relayers and their individual depths---a tight constraint on the
``inflation'' of cultural representations during transmission.

\begin{theorem}[Depth Zero Signal Relay (18.17)]
\label{thm:depth_zero_relay}
If $\depth(r) = 0$ and $\depth(s) = 0$, then $\texttt{chainDepth}([r], s) = 0$.
\end{theorem}

\noindent\textit{Proof sketch.}\quad
Depth-zero closure under composition.
\hfill$\square$

\medskip

\noindent\textbf{Commentary.}
Empty content through an empty channel yields nothing.  This is the
information-theoretic ground truth: zero-depth inputs produce zero-depth
output, regardless of the network topology.

\begin{theorem}[Empty Chain Depth Preservation (18.9)]
\label{thm:empty_chain_depth}
$\texttt{chainDepth}([\,], s) = \depth(s)$.
\end{theorem}

\noindent\textit{Proof sketch.}\quad
The relay chain over an empty list is the signal itself.
\hfill$\square$

\medskip

\noindent\textbf{Commentary.}
Without intermediaries, complexity is unchanged.  Direct communication
preserves not just the signal content but also its structural complexity.

% ============================================================
\section{Resonance in Networks}
\label{sec:ch18-resonance}

Resonance---the compatibility relation between ideas---interacts with
network diffusion in natural ways.

\begin{theorem}[Resonant Relayers Produce Resonant Output (18.7)]
\label{thm:resonant_relayers_resonant_output}
If $\text{resonates}(r_1, r_2)$, then for all $s$,
\[
  \text{resonates}\bigl(\texttt{relayChain}([r_1], s),\;
  \texttt{relayChain}([r_2], s)\bigr).
\]
\end{theorem}

\noindent\textit{Proof sketch.}\quad
Reduces to $\text{resonates}(\compose(r_1, s), \compose(r_2, s))$,
which follows from \texttt{resonance\_compose\_right}.
\hfill$\square$

\medskip

\noindent\textbf{Commentary.}
Similar cognitive environments produce similar transformations of the same
input.  In cultural epidemiology, this explains why cultural variants
\emph{cluster}: when relayers share cognitive schemas (resonance), their
outputs remain in the same neighborhood of ideatic space.\footnote{Sperber's
``epidemiology of representations'' predicts that cultural attractors emerge
when many minds share similar cognitive biases.  Theorem~\ref{thm:resonant_relayers_resonant_output}
provides the formal mechanism for this clustering.}

\begin{theorem}[Resonant Signals Through Same Relayer (18.8)]
\label{thm:resonant_signals_same_relay}
If $\text{resonates}(s_1, s_2)$, then for all $r$,
\[
  \text{resonates}\bigl(\texttt{relayChain}([r], s_1),\;
  \texttt{relayChain}([r], s_2)\bigr).
\]
\end{theorem}

\noindent\textit{Proof sketch.}\quad
From \texttt{resonance\_compose} applied to $(r, r, s_1, s_2)$ with
self-resonance of $r$ and the hypothesis.
\hfill$\square$

\medskip

\noindent\textbf{Commentary.}
Similar memes stay similar even after cultural transmission through the same
social environment.  A single relayer preserves the resonance between its
inputs, ensuring that initially similar signals remain similar after
processing.

\begin{theorem}[Two Chains Resonance (18.14)]
\label{thm:two_chains_resonance}
If $\text{resonates}(r_1, r_2)$ and $\text{resonates}(s_1, s_2)$, then
\[
  \text{resonates}\bigl(\texttt{relayChain}([r_1], s_1),\;
  \texttt{relayChain}([r_2], s_2)\bigr).
\]
\end{theorem}

\noindent\textit{Proof sketch.}\quad
Direct from \texttt{resonance\_compose} applied to $(r_1, r_2, s_1, s_2)$.
\hfill$\square$

\medskip

\noindent\textbf{Commentary.}
The diffusion of innovations through similar social environments produces
similar cultural adaptations.  When both the relayers and the signals are in
resonance, the outputs are in resonance---a strong stability result for
cultural transmission through homogeneous networks.

\begin{theorem}[Network Path Self-Resonance (18.11)]
\label{thm:path_self_resonance}
Every network path output resonates with itself.
\end{theorem}

\noindent\textit{Proof sketch.}\quad
Self-resonance is universal.
\hfill$\square$

\medskip

\noindent\textbf{Commentary.}
Any completed transmission, however distorted, achieves internal
self-consistency.  The accumulated composition of relayer transformations
may diverge arbitrarily from the original signal, but the final output is
always coherent with itself.

% ============================================================
\section{From Bartlett to TikTok: Transmission Chains Across Media}
\label{sec:ch18-bartlett-tiktok}

The relay-chain formalism applies not only to face-to-face oral
transmission (Bartlett's paradigm) but to every medium through which ideas
propagate.  As media technologies evolve, the \emph{structure} of the relay
chain changes while its \emph{algebraic properties} remain constant.

\subsection{Oral Transmission: The Telephone Game}

The children's ``telephone game'' (known as ``Chinese whispers'' in British
English) is the simplest empirical instantiation of a relay chain.  A
whispered message $s$ passes through a sequence of children
$[r_1, r_2, \ldots, r_n]$, emerging at the end as
$\texttt{relayChain}([r_1, \ldots, r_n], s)$.  The game's humor derives
from the dramatic divergence between input and output---a divergence that
our depth bounds show is \emph{constrained but non-zero}.

What makes the telephone game theoretically important is that it
eliminates all confounding variables: the medium is uniform (whispered
speech), the relay structure is linear (no branching), and the relayers
share a common cultural environment (same classroom).  Under these
controlled conditions, the mutations introduced by each relay hop are
purely compositional: each child composes the incoming signal with their
own auditory and semantic schemas.

\subsection{Print and Manuscript Transmission}

The medieval manuscript tradition provides a natural long-timescale
analogue of serial reproduction.  Each scribe $r_i$ copying a manuscript
$s_{i-1}$ introduces characteristic errors---omissions, additions,
transpositions, and glosses---that are precisely the compositional
contributions $\compose(r_i, s_{i-1}) - s_{i-1}$ (where the ``subtraction''
is metaphorical).  The field of textual criticism, from Karl Lachmann's
stemmatic method to computational phylogenetics, is essentially the study
of relay chains in manuscript networks.\footnote{The application of
phylogenetic methods to textual criticism is surveyed in Christopher J.\
Howe and Heather F.\ Windram, ``Phylomemetics: Evolutionary Analysis
Beyond the Gene,'' \textit{PLoS Biology} 9, no.\ 5 (2011).}

\subsection{Digital Transmission: Social Media Cascades}

On platforms like Twitter (now X) and TikTok, the relay chain acquires
a tree structure: a single post $s$ is ``retweeted'' or ``dueted'' by
multiple relayers simultaneously, each of whom adds their own commentary.
A retweet with comment is $\compose(r_i, s)$; a ``quote tweet'' with
extensive commentary may produce $\compose(r_i, s)$ with high relayer
depth.  TikTok ``duets'' are particularly interesting because they
literally place the relayer's video alongside the original, making the
composition $\compose(r_i, s)$ visible as a spatial juxtaposition.

The resonance-clustering theorems
(Section~\ref{sec:ch18-resonance}) predict that relay outputs will cluster
by community: relayers within the same algorithmic ``bubble'' share
cognitive schemas (resonance), producing similar compositions.  This
explains the empirical observation that viral content mutates along
community lines, with distinct ``strains'' emerging in different
subcommunities---a pattern precisely analogous to the evolution of
biological pathogens in geographically separated host populations.

\begin{remark}[The Telephone Game as Scientific Methodology]
\label{rem:telephone-methodology}
The ``serial reproduction'' paradigm pioneered by Bartlett has become a
standard experimental methodology in cultural evolution research.  Modern
implementations use controlled transmission chains to study the evolution
of language,\footnote{Simon Kirby, Hannah Cornish, and Kenny Smith,
``Cumulative Cultural Evolution in the Laboratory: An Experimental Approach
to the Origins of Structure in Human Language,'' \textit{Proceedings of the
National Academy of Sciences} 105, no.\ 31 (2008): 10681--10686.}
visual art, music, and even technology.  In each case, the experimental
design instantiates a relay chain $\texttt{relayChain}([r_1, \ldots, r_n], s)$
under controlled conditions, enabling researchers to isolate the
compositional contribution of each relay hop.  Our formalism provides the
algebraic framework within which these experimental results can be
unified and compared.
\end{remark}

\begin{remark}[Network Topology and Diffusion Speed]
\label{rem:topology-speed}
The topology of the relay network determines the \emph{speed} and
\emph{pattern} of cultural diffusion, even though the \emph{algebraic
properties} of each relay hop are topology-independent.  In a
small-world network (Watts and Strogatz), the average path length
scales as $\log n$, meaning that ideas can reach any node in
$O(\log n)$ relay hops.  In a scale-free network (Barab\'asi and Albert),
the hub structure means that most relay chains pass through a small
number of high-degree nodes, whose compositional contributions dominate
the final output.  In a lattice network, diffusion is slow and local,
with each relay hop advancing the signal by one spatial step.  The
depth bounds of Section~\ref{sec:ch18-depth} apply to all three
topologies, providing topology-independent upper bounds on the complexity
of the transmitted signal.
\end{remark}

\begin{remark}[Connection to Marketplace Competition]
\label{rem:marketplace-connection}
The network diffusion formalism connects to the marketplace competition
of Chapter~\ref{ch:marketplace} in a natural way.  In the marketplace
model, ideas compete for adoption; in the network diffusion model, ideas
spread through relay chains.  The two models are complementary: marketplace
competition determines \emph{which} ideas enter the relay chain (the
signal $s$), while network diffusion determines \emph{how} those ideas
are transformed as they spread (the relay chain $\texttt{relayChain}(
\text{relayers}, s)$).  The depth bounds and resonance results of this
chapter provide constraints on the outcomes of marketplace competition:
even a ``winning'' idea will be transformed by the relay chain through
which it diffuses, and the extent of that transformation is bounded by
the algebraic properties proved here.
\end{remark}

% ============================================================
\section{Network Topology and Chain Composition}
\label{sec:ch18-topology}

\begin{theorem}[Chain Concatenation is Sequential Relay (18.10)]
\label{thm:chain_concat}
\[
  \texttt{relayChain}(\text{relayers}_1 \mathop{++} \text{relayers}_2, s) =
  \texttt{relayChain}\bigl(\text{relayers}_2,\;
  \texttt{relayChain}(\text{relayers}_1, s)\bigr).
\]
\end{theorem}

\noindent\textit{Proof sketch.}\quad
From the associativity of \texttt{foldl} over list concatenation.
\hfill$\square$

\medskip

\noindent\textbf{Commentary.}
Two segments of a communication network compose into a single longer path.
This is the fundamental composition law for information networks: the output
of the first segment becomes the input to the second.  In graph-theoretic
terms, path concatenation in the relay network corresponds to functional
composition of relay transformations---a monoidal structure on the space
of network paths.\footnote{This is the \emph{functoriality} of the relay
chain: it respects the monoidal structure of list concatenation.}

\begin{theorem}[Chain Append Equivalence (18.18)]
\label{thm:chain_append}
\[
  \texttt{relayChain}(\text{relayers} \mathop{++} [r], s) =
  \compose\bigl(r,\; \texttt{relayChain}(\text{relayers}, s)\bigr).
\]
\end{theorem}

\noindent\textit{Proof sketch.}\quad
Special case of chain concatenation with a singleton second chain.
\hfill$\square$

\medskip

\noindent\textbf{Commentary.}
Adding a node at the end of a path is equivalent to that node processing
the path's accumulated output.  This provides the inductive step for
analyzing chains of arbitrary length: each new relayer simply composes with
the running total.

\begin{theorem}[Relay Preserves Relayer Count (18.13)]
\label{thm:relay_preserves_relayer_count}
$|\texttt{NetworkPath.mk}(\text{relayers}, s).\text{relayers}| = |\text{relayers}|$.
\end{theorem}

\noindent\textit{Proof sketch.}\quad
Definitional equality: the structure stores the relayer list unchanged.
\hfill$\square$

\medskip

\noindent\textbf{Commentary.}
The relay chain function doesn't alter the number of relayers.  The network
topology is preserved during diffusion---only the \emph{content} of the
signal changes, not the \emph{structure} of the network.

% ============================================================

\noindent\textbf{Extended discussion of Theorem~\ref{thm:chain_concat}.}\quad
The chain-concatenation theorem is the structural backbone of network
diffusion theory.  It establishes that the relay chain respects the monoidal
structure of list concatenation: composing two segments of a relay network
is equivalent to running the first segment and then feeding its output into
the second.  This is a \emph{functoriality} result: the relay-chain
function is a monoid homomorphism from $(\text{List}(\IS), {+\!+}, [\,])$
to $(\IS \to \IS, \circ, \text{id})$.

The practical consequence is that large relay networks can be analyzed by
decomposition.  A continental-scale diffusion network can be partitioned
into regional sub-networks, each analyzed independently; the
concatenation theorem guarantees that assembling the regional analyses
yields the correct global result.  This modularity principle is essential
for scaling the formalism to real-world networks with millions of nodes.

In graph-theoretic terms, the theorem states that path composition in the
relay graph corresponds to function composition in the space of ideatic
transformations.  This is the formal content of the intuition that
``information flows along paths'': the transformation effected by a path
is the composition of the transformations effected by its constituent
edges.

\medskip

\noindent\textbf{Extended discussion of Theorem~\ref{thm:void_chain_identity}.}\quad
The void-chain identity theorem---that a chain of void relayers preserves
the signal perfectly---is the formal expression of a foundational
idealization in communication theory: the \emph{noiseless channel}.
Shannon's noiseless coding theorem (1948) established that information can
be transmitted perfectly through a noiseless channel; our theorem
establishes that a chain of void relayers \emph{is} a noiseless channel in
ideatic space.

The practical significance is as a limiting case.  No real relay chain
consists of void relayers---every intermediary adds \emph{something} to
the signal.  But the void chain provides the baseline against which
real chains can be measured.  The ``distortion'' introduced by a real
relay chain $[r_1, \ldots, r_n]$ is the difference between
$\texttt{relayChain}([r_1, \ldots, r_n], s)$ and
$\texttt{relayChain}([\void, \ldots, \void], s) = s$.  This
difference---the cumulative compositional contribution of the
relayers---is the formal analogue of ``noise'' in Shannon's theory.

\section{Conclusion}
\label{sec:ch18-concl}

This chapter has established eighteen theorems formalizing cultural
diffusion through networks within the framework of Ideatic Diffusion Theory.
The key results are:

\begin{itemize}
  \item \textbf{Compositional structure:} Network diffusion is sequential
    composition (Theorems~\ref{thm:singleton_chain}--\ref{thm:two_relayer_chain}).
  \item \textbf{Void transparency:} Passive relayers preserve signals
    (Theorems~\ref{thm:void_relayer_transparent}--\ref{thm:void_chain_identity}).
  \item \textbf{Depth bounds:} Chain depth grows at most linearly
    (Theorems~\ref{thm:single_relay_depth}--\ref{thm:two_relay_depth_bound}).
  \item \textbf{Resonance clustering:} Similar relayers and signals produce
    similar outputs (Theorems~\ref{thm:resonant_relayers_resonant_output}--\ref{thm:two_chains_resonance}).
  \item \textbf{Network composition:} Chain concatenation = sequential relay
    (Theorem~\ref{thm:chain_concat}).
\end{itemize}

Sperber's cultural epidemiology proposed that culture is ``the precipitate
of cognition and communication.''  Our formalization shows that this
precipitate obeys precise algebraic laws: the relay chain is a homomorphism
from the monoid of relayer lists to the monoid of ideatic transformations.
Depth bounds constrain the ``mutation rate'' of cultural representations,
and resonance clustering explains the emergence of cultural attractors.

In the next chapter, we turn from the diffusion of ideas to their
aggregation: the social-choice problem of combining multiple interpretations
into a single shared meaning.

% ============================================================
%  Chapter 19 — Social Choice and Hermeneutic Aggregation
% ============================================================
\chapter{Social Choice and Hermeneutic Aggregation}
\label{ch:social}

\begin{quote}
\textit{``In communicative action participants are not primarily oriented to
their own individual successes; they pursue their individual goals under the
condition that they can harmonize their plans of action on the basis of
common situation definitions.''}\\
--- J\"urgen Habermas, \textit{The Theory of Communicative Action} (1981)
\end{quote}

\vspace{1em}

\section{Introduction}

When a sermon is preached, a lecture delivered, or a news broadcast aired,
each member of the audience interprets the message through their own
background---their education, culture, temperament, and prior beliefs.
J\"urgen Habermas's theory of communicative action asks the fundamental
question: how do multiple, heterogeneous interpretations converge into shared
meaning?\footnote{Habermas's magnum opus, \textit{Theorie des kommunikativen
Handelns} (1981), distinguishes between ``strategic action'' (oriented
toward individual success) and ``communicative action'' (oriented toward
mutual understanding).  The hermeneutic aggregation problem belongs squarely
to the domain of communicative action.}

This is the \emph{hermeneutic aggregation problem}: given $n$ agents with
background ideas $[r_1, \ldots, r_n]$ and a shared signal $s$, each agent
produces an interpretation $\compose(r_i, s)$.  The interpretations can then
be aggregated---by sequential composition, by deliberation, by voting---into
a collective meaning.  Under what conditions does this collective meaning
reflect a genuine consensus rather than a mere concatenation of
incompatible readings?

In IDT, we formalize this problem through three core concepts:
\begin{itemize}
  \item \texttt{interpretations}: the list of all agents' interpretations of a
    shared signal;
  \item \texttt{aggregateCompose}: the sequential composition of
    interpretations into a single collective meaning;
  \item \texttt{coherentInterpretation}: the condition that all
    interpretations pairwise resonate, formalizing Habermas's ``ideal speech
    situation.''
\end{itemize}

This chapter proves eighteen theorems characterizing the algebraic and
metric properties of hermeneutic aggregation.  The results establish when
coherence is achievable, how depth constrains the interpretive effort of a
community, and what role void---silence, absence, non-participation---plays
in the aggregation process.

\subsection{Historical Background: From Arrow to Arendt}

The problem of aggregating individual judgments into collective decisions has
a distinguished intellectual pedigree that long predates our formalization.
We situate IDT's hermeneutic aggregation within this tradition.

\paragraph{Arrow's impossibility theorem.}
Kenneth Arrow's celebrated impossibility theorem (1951) demonstrated that no
social welfare function can simultaneously satisfy a small set of seemingly
reasonable axioms---unrestricted domain, Pareto efficiency, independence of
irrelevant alternatives, and non-dictatorship.\footnote{Kenneth J.\ Arrow,
\textit{Social Choice and Individual Values} (1951; 2nd ed.\ 1963).  The
theorem shattered the optimism of welfare economics by showing that
``democratic'' aggregation of preferences is, in a precise sense,
impossible.}  Arrow's theorem concerns the aggregation of \emph{preference
orderings}; our theory concerns the aggregation of \emph{interpretive
compositions}.  Yet the structural parallel is illuminating: just as Arrow
showed that no aggregation rule can perfectly reconcile individual
preferences, our depth bounds (Theorem~\ref{thm:interpretation_depth_sum_bound})
show that no aggregation can transcend the cognitive resources of its
participants.  The impossibility is not logical but resource-theoretic:
consensus is possible, but \emph{bounded}.

\paragraph{Sen's liberal paradox and the capability approach.}
Amartya Sen extended Arrow's framework in two crucial directions.  His
``liberal paradox'' (1970) showed that even minimal individual liberty is
incompatible with Pareto optimality---a result that highlights the tension
between individual autonomy and collective welfare.\footnote{Amartya Sen,
``The Impossibility of a Paretian Liberal,'' \textit{Journal of Political
Economy} 78(1), 1970, pp.\ 152--157.}  Sen's later ``capability approach''
shifted the focus from preference satisfaction to the substantive freedoms
individuals enjoy.  In IDT terms, each agent's ``capability'' to interpret
is captured by their background depth $\depth(r_i)$: agents with richer
backgrounds have greater interpretive capability, and the community's total
capability constrains its aggregate interpretive richness.

\paragraph{Condorcet's jury theorem and epistemic democracy.}
The Marquis de Condorcet's jury theorem (1785) established that if each
juror has a probability greater than $\frac{1}{2}$ of identifying the
correct answer, then the probability that the majority vote is correct
approaches 1 as the number of jurors grows.\footnote{Condorcet,
\textit{Essai sur l'application de l'analyse \`a la probabilit\'e des
d\'ecisions rendues \`a la pluralit\'e des voix} (1785).}  This ``wisdom of
crowds'' result provides a probabilistic justification for democratic
aggregation.  IDT's depth bounds complement Condorcet's result: while
Condorcet shows that larger groups make \emph{more accurate} decisions
(under certain conditions), our
Theorem~\ref{thm:interpretation_depth_sum_bound} shows that larger groups
produce \emph{more total interpretive depth}---a resource-theoretic rather
than epistemic claim.

\paragraph{Habermas's ideal speech situation.}
J\"urgen Habermas's discourse ethics posits an ``ideal speech situation'' in
which all participants have equal opportunity to speak, all arguments are
assessed solely on their merits, and consensus is reached through the
``unforced force of the better argument.''\footnote{Habermas, \textit{Moral
Consciousness and Communicative Action} (1990), trans.\ Christian Lenhardt
and Shierry Weber Nicholsen.}  Our coherent interpretation condition
(Definition~\ref{def:coherentInterpretation}) formalizes this ideal: it
requires that all pairwise interpretations resonate, which in turn requires
that the ``horizons'' of all participants be compatible.  The gap between
the ideal and the real---between universal resonance and partial
resonance---is the gap between Habermas's regulative ideal and the
messy reality of actual discourse.

\paragraph{Deliberative democracy.}
John Rawls's concept of ``public reason'' restricts legitimate political
arguments to those accessible to all citizens regardless of their
comprehensive doctrines.\footnote{John Rawls, \textit{Political Liberalism}
(1993), Lecture VI.}  James Fishkin's ``deliberative polling'' experiments
demonstrate that informed, structured deliberation can shift opinions toward
more reflective positions.\footnote{James S.\ Fishkin, \textit{Democracy
and Deliberation: New Directions for Democratic Reform} (1991).}  IDT
provides a formal model for such processes: each round of deliberation is
an application of $\texttt{aggregateCompose}$, and the shift in collective
meaning can be tracked through changes in aggregate depth and resonance
structure.

\paragraph{Modern instantiations.}
The hermeneutic aggregation problem is not merely theoretical.  Reddit's
upvoting system aggregates millions of interpretive judgments into a
ranked ordering of content; Wikipedia's consensus model requires editors
to negotiate shared meaning through iterative revision; Liquid Democracy
systems allow voters to delegate their interpretive authority to trusted
proxies, creating chains of delegated interpretation that structurally
resemble relay chains (Chapter~\ref{ch:network}).  Each of these systems
instantiates a particular aggregation rule, and each faces the fundamental
tension between aggregation (voting, ranking) and deliberation (discussion,
revision, consensus-building).

\paragraph{Arendt on plurality and the public sphere.}
Hannah Arendt argued that the fundamental condition of politics is
\emph{plurality}: the fact that human beings are alike enough to understand
each other and different enough to need speech and action to make themselves
understood.\footnote{Hannah Arendt, \textit{The Human Condition} (1958),
\S5.}  IDT captures this duality precisely.  Resonance captures the
``alike enough'' dimension: agents whose backgrounds resonate can achieve
coherent interpretation.  The diversity of backgrounds---the fact that
$r_i \neq r_j$ in general---captures the ``different enough'' dimension:
each agent brings their own horizon to the interpretive encounter.  The
hermeneutic aggregation problem is, in Arendt's terms, the problem of
creating a \emph{common world} out of a plurality of perspectives.

\paragraph{The tension between aggregation and deliberation.}
Two fundamentally different approaches to collective decision-making run
through the social choice literature.  \emph{Aggregative} approaches---
voting, polling, preference aggregation---take individual judgments as
given and combine them mechanically.  \emph{Deliberative} approaches---
dialogue, discussion, argumentation---seek to transform individual
judgments through the process of collective reasoning.  IDT's
$\texttt{aggregateCompose}$ function occupies an interesting middle ground:
it is aggregative in structure (a fold over a list) but compositional in
content (each step involves genuine interpretive engagement via
$\compose$).  The result is an aggregation that is not merely additive
but genuinely transformative---each voice interacts with the existing
consensus, potentially altering it in ways that simple concatenation
would not.

\section{Core Definitions}
\label{sec:ch19-defs}

\begin{definition}[Interpretations]
\label{def:interpretations}
Let $\text{reps} = [r_1, \ldots, r_n]$ be the agents' repertoire ideas and
$s \in \IS$ be the shared signal.  The \emph{interpretations} are:
\[
  \texttt{interpretations}(\text{reps}, s) \;:=\;
  [\ \compose(r_1, s),\; \compose(r_2, s),\; \ldots,\; \compose(r_n, s)\ ].
\]
Each agent interprets the signal by composing it with their own background.
\end{definition}

\begin{lstlisting}[caption={Interpretations in Lean}]
def interpretations (reps : List I) (s : I) : List I :=
  reps.map (fun r => IdeaticSpace.compose r s)
\end{lstlisting}

In Gadamer's hermeneutic philosophy, each interpreter brings their own
``horizon''---their pre-understanding, their historically-effected
consciousness---to the encounter with the text.  The interpretation is the
``fusion of horizons'' between background and signal, which in IDT is
precisely $\compose(r_i, s)$.\footnote{Hans-Georg Gadamer,
\textit{Wahrheit und Methode} (1960), introduced the concept of
\textit{Horizontverschmelzung} (fusion of horizons) as the fundamental
structure of understanding.}

\begin{definition}[Aggregate Composition]
\label{def:aggregateCompose}
The \emph{aggregate composition} of a list of ideas is their right-fold
via composition:
\[
  \texttt{aggregateCompose}([\,]) = \void, \qquad
  \texttt{aggregateCompose}(a :: \text{rest}) = \compose(a, \texttt{aggregateCompose}(\text{rest})).
\]
This models the process by which a community arrives at a single shared
interpretation through successive dialogue.
\end{definition}

\begin{lstlisting}[caption={Aggregate composition in Lean}]
def aggregateCompose : List I -> I
  | [] => IdeaticSpace.void
  | a :: rest => IdeaticSpace.compose a (aggregateCompose rest)
\end{lstlisting}

The right-fold semantics are significant: the last voice in the list serves
as the ``base'' upon which earlier voices build.  In Habermas's terms, this
models a deliberative process in which each new speaker adds their
perspective ``on top of'' the existing consensus.

\begin{definition}[Coherent Interpretation]
\label{def:coherentInterpretation}
A group of agents achieves \emph{coherent interpretation} of signal $s$ when
all pairs of interpretations resonate:
\[
  \texttt{coherentInterpretation}(\text{reps}, s) \;:\Longleftrightarrow\;
  \forall\, x, y \in \texttt{interpretations}(\text{reps}, s),\quad
  \text{resonates}(x, y).
\]
This is the formal condition for ``hermeneutic consensus.''
\end{definition}

\begin{remark}[Resonance and Habermas's Validity Claims]
\label{rem:resonance_validity}
Habermas distinguished three types of validity claims implicit in every
communicative act: claims to \emph{truth} (propositional correctness),
\emph{rightness} (normative legitimacy), and \emph{sincerity} (subjective
truthfulness).\footnote{Habermas, \textit{The Theory of Communicative
Action}, Vol.\ I, Ch.\ III.}  In IDT, resonance between two ideas captures
a structural compatibility that encompasses all three dimensions
simultaneously.  Two interpretations resonate when they are ``in tune''
across the full spectrum of validity---not merely factually consistent, but
also normatively compatible and phenomenologically consonant.  The coherent
interpretation condition thus formalizes an ideal in which Habermas's three
validity claims are simultaneously satisfied across all pairs of
interpreters: a community that achieves coherent interpretation has
resolved not only its factual disagreements but also its normative and
expressive ones.
\end{remark}

\begin{lstlisting}[caption={Coherent interpretation in Lean}]
def coherentInterpretation (reps : List I) (s : I) : Prop :=
  forall x in interpretations reps s, forall y in interpretations reps s,
    IdeaticSpace.resonates x y
\end{lstlisting}

Habermas's ``ideal speech situation'' is one in which all participants'
interpretations are mutually comprehensible.  Coherent interpretation is the
formal counterpart: when it holds, the community has achieved a state of
hermeneutic equilibrium in which every reading ``rhymes'' with every other.

% ============================================================
\section{Examples of Hermeneutic Aggregation}
\label{sec:ch19-examples}

Before proceeding to the formal theorems, we ground the abstract framework
in concrete historical and contemporary examples.

\begin{example}[The Athenian Assembly]
\label{ex:athenian_assembly}
The Athenian \textit{ekklesia} provides a paradigmatic case of hermeneutic
aggregation.  When Pericles delivered his Funeral Oration, each citizen in
the audience---farmer, artisan, merchant, philosopher---interpreted the
speech through their own background ($r_i$).  The collective meaning of the
oration was not any single interpretation but the aggregate
$\texttt{aggregateCompose}([\compose(r_1, s), \ldots, \compose(r_n, s)])$.
Thucydides's text, by recording the speech for posterity, froze one
particular composition of signal and historian's background, which
subsequent readers have re-composed with \textit{their} backgrounds
across twenty-five centuries.\footnote{Thucydides, \textit{History of
the Peloponnesian War}, Book II, 34--46.  The speech's interpretation
has itself been subject to centuries of hermeneutic aggregation.}

The Athenian case also illustrates the depth bound
(Theorem~\ref{thm:interpretation_depth_sum_bound}): the total interpretive
richness of the assembly was bounded by the aggregate cultural depth of
Athenian citizenship plus the rhetorical complexity of Pericles's oration.
No assembly, however large, could have extracted infinite meaning from a
finite speech.
\end{example}

\begin{example}[Wikipedia's Consensus Model]
\label{ex:wikipedia_consensus}
Wikipedia's editorial process provides a modern instantiation of hermeneutic
aggregation under conditions of radical pluralism.  Consider a contested
article---say, on a politically sensitive historical event.  Each editor
$r_i$ brings their own background (nationalist historiography, critical
theory, primary source expertise, lived experience) to the shared
``signal'' of the available evidence.  Their edits constitute interpretations
$\compose(r_i, s)$, and the article's evolution through successive revisions
mirrors the sequential composition of $\texttt{aggregateCompose}$.

Wikipedia's ``neutral point of view'' policy functions as an institutional
approximation of coherent interpretation: it demands that the article's
content be acceptable to all editors---formally, that all pairwise
interpretations resonate.  In practice, this ideal is approached through
the Talk page, where editors engage in deliberation to achieve (or
approximate) the coherence condition.  When coherence fails---when editors'
backgrounds are too divergent for pairwise resonance---the result is an
``edit war,'' the hermeneutic analogue of political deadlock.

The depth bound applies here as well: the total interpretive richness of a
Wikipedia article is bounded by the aggregate expertise of its editor
community.  This explains the empirical observation that articles on
specialized topics tend to plateau in quality once the relevant expert
community has been exhausted.\footnote{See Aaron Halfaker et al.,
``The Rise and Decline of an Open Collaboration System,''
\textit{American Behavioral Scientist} 57(5), 2013, pp.\ 664--688.}
\end{example}

\begin{example}[Jury Deliberation as Hermeneutic Aggregation]
\label{ex:jury_deliberation}
A criminal trial presents jurors with a shared signal $s$---the totality
of evidence, testimony, and legal instruction.  Each juror interprets
this signal through their own background: their life experience, moral
intuitions, understanding of human nature, and capacity for legal
reasoning.  The jury's deliberation process is an instance of
$\texttt{aggregateCompose}$: jurors successively share their interpretations,
and each contribution reshapes the emerging collective understanding.

The requirement of unanimity in many jurisdictions is a strong coherence
condition: not merely pairwise resonance but identity of conclusion.
A hung jury---one that cannot reach unanimity---represents a failure
of coherent interpretation: the jurors' backgrounds are sufficiently
divergent that no signal, however compelling, can produce pairwise
resonance among all interpretations.

This example also illuminates Theorem~\ref{thm:resonant_pair_coherent}:
jurors who share similar backgrounds (profession, class, education) are
more likely to achieve pairwise resonance and thus coherent interpretation.
The well-documented phenomenon of jury selection bias---attorneys'
systematic efforts to seat favorable jurors---is precisely an attempt
to engineer the resonance structure of the interpretive community.
\end{example}

\begin{example}[The European Union's Multilingual Interpretation Challenge]
\label{ex:eu_multilingual}
The European Union conducts its business in 24 official languages, with
every regulation, directive, and judgment of the European Court of Justice
translated into each language.  This creates a distinctive hermeneutic
aggregation problem: the ``signal'' (a legal text) must be interpreted
not only through each reader's individual background but also through the
lens of their linguistic community.

Each language version constitutes an interpretation $\compose(r_{\text{lang}}, s)$
where $r_{\text{lang}}$ encodes the conceptual categories, legal traditions,
and cultural associations of the target language.  The EU's legal doctrine
that all language versions are ``equally authentic'' is a coherence
requirement: all twenty-four interpretations must resonate with each other.
When they do not---when, for example, the German version of a directive
has subtly different implications from the French version---the European
Court of Justice must engage in meta-interpretation, composing its own
judicial background with the divergent linguistic interpretations to
produce a single authoritative reading.

The EU case dramatically illustrates the tension between aggregation and
deliberation: mechanical translation (aggregation) often fails to preserve
meaning, requiring subsequent judicial deliberation to restore
coherence.\footnote{See Baaij, C.J.W., ``Fifty Years of Multilingual
Interpretation in the European Union,'' in \textit{The Oxford Handbook
of Language and Law} (2012), pp.\ 217--231.}
\end{example}

% ============================================================
\section{Empty and Singleton Groups}
\label{sec:ch19-empty}

\begin{theorem}[Empty Group Has Void Interpretation (19.1)]
\label{thm:empty_group_void}
For all $s \in \IS$,
\[
  \texttt{aggregateCompose}\bigl(\texttt{interpretations}([\,], s)\bigr) = \void.
\]
\end{theorem}

\noindent\textit{Proof sketch.}\quad
$\texttt{interpretations}([\,], s) = [\,]$, and
$\texttt{aggregateCompose}([\,]) = \void$ by definition.
\hfill$\square$

\medskip

\noindent\textbf{Commentary.}
An empty polity produces no collective meaning.  This is the political-theory
baseline: without citizens, there is no \textit{res publica}.

\begin{remark}[Arrow's Theorem and the Limits of Hermeneutic Aggregation]
\label{rem:arrow_limits}
Arrow's impossibility theorem demonstrates that no aggregation procedure for
preference orderings can simultaneously satisfy four reasonable axioms.
A natural question arises: does an analogous impossibility hold for
hermeneutic aggregation?  The answer is nuanced.  Our
$\texttt{aggregateCompose}$ function does not face Arrow's impossibility
because it does not attempt to produce a ``socially optimal'' interpretation
from individual preferences \emph{over} interpretations.  Rather, it
sequentially composes interpretations into a single composite---an
operation more akin to concatenation than to social choice.  However, this
comes at a cost: the result of $\texttt{aggregateCompose}$ is
order-dependent (by the non-commutativity explored in
Chapter~\ref{ch:bargaining}), which means that the aggregation procedure
is sensitive to who speaks in what order.  Arrow's impossibility is avoided
not by satisfying all his axioms but by relaxing the anonymity
condition---in IDT, the order of voices matters.
\end{remark}


\begin{theorem}[Singleton = Individual Interpretation (19.2)]
\label{thm:singleton_interpretation}
\[
  \texttt{aggregateCompose}\bigl(\texttt{interpretations}([r], s)\bigr) = \compose(r, s).
\]
\end{theorem}

\noindent\textit{Proof sketch.}\quad
$\texttt{interpretations}([r], s) = [\compose(r,s)]$, and
$\texttt{aggregateCompose}([\compose(r,s)]) = \compose(\compose(r,s), \void) =
\compose(r,s)$ by right-identity.
\hfill$\square$

\medskip

\noindent\textbf{Commentary.}
A solitary reader's interpretation \emph{is} the meaning.  In hermeneutics,
the singleton case collapses the distinction between individual and
collective meaning: when there is only one interpreter, their reading is
authoritative by default.

\begin{theorem}[Interpretation Count = Agent Count (19.5)]
\label{thm:interpretation_count}
$|\texttt{interpretations}(\text{reps}, s)| = |\text{reps}|$.
\end{theorem}

\noindent\textit{Proof sketch.}\quad
The \texttt{map} operation preserves list length.
\hfill$\square$

\medskip

\noindent\textbf{Commentary.}
Every agent produces exactly one interpretation---no more, no less.  This
``one person, one reading'' principle is the hermeneutic analogue of the
democratic ``one person, one vote.''\footnote{The analogy to social choice
theory is intentional.  Arrow's impossibility theorem concerns the
aggregation of preference orderings; our theory concerns the aggregation of
interpretive compositions.  The structural parallels are explored further in
Section~\ref{sec:ch19-coherence}.}

\begin{theorem}[Interpretation Map Preserves Nil (19.17)]
\label{thm:interpret_nil}
$\texttt{interpretations}([\,], s) = [\,]$.
\end{theorem}

\noindent\textit{Proof sketch.}\quad
Mapping over the empty list yields the empty list.
\hfill$\square$

% ============================================================
\section{Void Signal and Self-Resonance}
\label{sec:ch19-void}

When the signal is void---when nothing is said---the interpretive situation
simplifies dramatically.

\begin{theorem}[Void Signal Preserves Repertoire (19.3)]
\label{thm:void_signal_preserves}
\[
  \texttt{interpretations}(\text{reps}, \void) = \text{reps}.
\]
\end{theorem}

\noindent\textit{Proof sketch.}\quad
By induction on the list.  For each $r \in \text{reps}$, $\compose(r, \void) = r$
by right-identity.
\hfill$\square$

\medskip

\noindent\textbf{Commentary.}
When there is no text to interpret, each interpreter is left with only their
own pre-understanding.  In Gadamer's terms, their horizon remains unchanged:
the void signal induces no fusion of horizons, and each agent's
interpretation equals their original repertoire idea.

\begin{theorem}[Void Signal Yields Self-Resonant Group (19.4)]
\label{thm:void_signal_self_resonant}
For all $x \in \texttt{interpretations}(\text{reps}, \void)$,
$\text{resonates}(x, x)$.
\end{theorem}

\noindent\textit{Proof sketch.}\quad
By Theorem~\ref{thm:void_signal_preserves}, interpretations equal the
original repertoire ideas, and self-resonance is universal.
\hfill$\square$

\medskip

\noindent\textbf{Commentary.}
In the absence of ritual content, mechanical solidarity still holds: each
individual remains internally coherent.  Durkheim's insight that social
cohesion begins with self-consistency finds formal expression here.

\begin{theorem}[Void Repertoire Interpretation (19.14)]
\label{thm:void_background_interpretation}
$\texttt{interpretations}([\void], s) = [s]$.
\end{theorem}

\noindent\textit{Proof sketch.}\quad
$\compose(\void, s) = s$ by left-identity.
\hfill$\square$

\medskip

\noindent\textbf{Commentary.}
A \textit{tabula rasa}---a mind with no pre-existing ideas---interprets the
signal as the signal itself, without transformation.  The void background
is a perfect mirror: the interpretation is the signal unchanged.

\begin{theorem}[Aggregate of Voids is Void (19.15)]
\label{thm:void_aggregate_void}
For all $n$,
$\texttt{aggregateCompose}\bigl(\texttt{interpretations}(\underbrace{[\void, \ldots, \void]}_{n}, \void)\bigr) = \void$.
\end{theorem}

\noindent\textit{Proof sketch.}\quad
By Theorem~\ref{thm:void_signal_preserves}, the interpretations are
$n$ copies of $\void$.  By induction, aggregating $n$ voids yields $\void$.
\hfill$\square$

\medskip

\noindent\textbf{Commentary.}
Empty minds interpreting nothing produce nothing.  This is the nihilistic
limit of hermeneutic aggregation: when neither the community nor the text
carries any content, the outcome is silence.

\begin{theorem}[Aggregate Singleton Void Signal (19.18)]
\label{thm:aggregate_singleton_void}
$\texttt{aggregateCompose}\bigl(\texttt{interpretations}([r], \void)\bigr) = r$.
\end{theorem}

\noindent\textit{Proof sketch.}\quad
By void-signal preservation, the interpretation list is $[r]$, and
$\texttt{aggregateCompose}([r]) = \compose(r, \void) = r$.
\hfill$\square$

\medskip

\noindent\textbf{Commentary.}
In the absence of external stimuli, consciousness returns to its own
content---pure self-reflection.  Phenomenologically, the subject without
object collapses back upon itself.

% ============================================================
\section{Aggregation Mechanics}
\label{sec:ch19-aggregate}

\begin{theorem}[Aggregating Empty Yields Void (19.9)]
\label{thm:aggregate_nil}
$\texttt{aggregateCompose}([\,]) = \void$.
\end{theorem}

\noindent\textit{Proof sketch.}\quad Definitional. \hfill$\square$

\medskip

\noindent\textbf{Commentary.}
Without any contributions to the discourse, the outcome is silence.
Habermas's communicative action requires at least one speaker.

\begin{theorem}[Aggregate Cons Unfolding (19.10)]
\label{thm:aggregate_cons}
$\texttt{aggregateCompose}(a :: \text{rest}) = \compose(a, \texttt{aggregateCompose}(\text{rest}))$.
\end{theorem}

\noindent\textit{Proof sketch.}\quad Definitional. \hfill$\square$

\medskip

\noindent\textbf{Commentary.}
Each new voice adds their perspective on top of the existing consensus.
The deliberative process is compositional: the $k$-th speaker composes
their idea with the aggregate of the remaining speakers, building the
collective meaning layer by layer.

\begin{theorem}[Void Aggregate Identity (19.11)]
\label{thm:void_aggregate}
$\texttt{aggregateCompose}(\void :: \text{rest}) = \texttt{aggregateCompose}(\text{rest})$.
\end{theorem}

\noindent\textit{Proof sketch.}\quad
$\compose(\void, \texttt{aggregateCompose}(\text{rest})) =
\texttt{aggregateCompose}(\text{rest})$ by left-identity.
\hfill$\square$

\medskip

\noindent\textbf{Commentary.}
A silent participant doesn't alter the collective meaning.  Non-contributing
members are transparent to the aggregation process---they neither add nor
subtract from the deliberative outcome.

\begin{theorem}[Self-Resonance of Aggregate (19.12)]
\label{thm:aggregate_self_resonance}
$\text{resonates}\bigl(\texttt{aggregateCompose}(\text{interps}),\;
\texttt{aggregateCompose}(\text{interps})\bigr)$.
\end{theorem}

\noindent\textit{Proof sketch.}\quad Self-resonance is universal. \hfill$\square$

\medskip

\noindent\textbf{Commentary.}
Any completed deliberation, however contentious, achieves internal
self-consistency as a composite idea.  The aggregate may fail to satisfy
individual participants, but it is always coherent with itself---a minimal
condition for the output of any decision procedure.

% ============================================================
\section{Coherence and Resonance}
\label{sec:ch19-coherence}

Under what conditions does a group achieve hermeneutic consensus?
The following theorems establish both trivial and substantive conditions
for coherent interpretation.

\begin{theorem}[Empty Interpretation List is Coherent (19.7)]
\label{thm:empty_coherent}
$\texttt{coherentInterpretation}([\,], s)$ for all $s$.
\end{theorem}

\noindent\textit{Proof sketch.}\quad
Universal quantification over the empty set is vacuously true.
\hfill$\square$

\medskip

\noindent\textbf{Commentary.}
Unanimity without voters: if no one interprets, there's no disagreement.
This is the degenerate case of consensus---trivially achieved but
informationally empty.

\begin{theorem}[Singleton Interpretation is Coherent (19.8)]
\label{thm:singleton_coherent}
$\texttt{coherentInterpretation}([r], s)$ for all $r, s$.
\end{theorem}

\noindent\textit{Proof sketch.}\quad
The only interpretation is $\compose(r, s)$, which resonates with itself.
\hfill$\square$

\medskip

\noindent\textbf{Commentary.}
A lone reader always agrees with their own interpretation.  Solipsism is
trivially coherent.  The interesting cases arise when multiple agents must
reach consensus.

\begin{theorem}[Resonant Pair Coherence (19.13)]
\label{thm:resonant_pair_coherent}
If $\text{resonates}(r_1, r_2)$, then
$\texttt{coherentInterpretation}([r_1, r_2], s)$ for all $s$.
\end{theorem}

\noindent\textit{Proof sketch.}\quad
There are four cases (each interpretation against each interpretation).
The diagonal cases are self-resonance.  The off-diagonal cases follow from
\texttt{resonance\_compose\_right}: $\text{resonates}(r_1, r_2)$ implies
$\text{resonates}(\compose(r_1, s), \compose(r_2, s))$.
\hfill$\square$

\medskip

\noindent\textbf{Commentary.}
This is Durkheim's \textit{conscience collective} theorem: agents who share
deep cultural resonance will agree on any new information.  If two minds
resonate at the level of their background assumptions, their interpretations
of any signal will cohere.  This is the formal mechanism underlying
``preaching to the choir''---the choir already agrees, so any sermon
produces consensus.\footnote{The sociological implications are significant.
Theorem~\ref{thm:resonant_pair_coherent} implies that cultural homogeneity
is \emph{sufficient} for hermeneutic consensus.  The converse---whether
consensus requires homogeneity---is a deeper question left for future work.}

% ============================================================
\section{Depth Bounds on Aggregation}
\label{sec:ch19-depth}

\begin{theorem}[Aggregate Depth Bound---Singleton (19.6)]
\label{thm:singleton_aggregate_depth}
\[
  \depth\bigl(\texttt{aggregateCompose}(\texttt{interpretations}([r], s))\bigr)
  \leq \depth(r) + \depth(s).
\]
\end{theorem}

\noindent\textit{Proof sketch.}\quad
By Theorem~\ref{thm:singleton_interpretation}, the aggregate equals
$\compose(r, s)$, and subadditivity applies.
\hfill$\square$

\medskip

\noindent\textbf{Commentary.}
One person's understanding adds at most their background depth to the
signal's depth.  This is the hermeneutic bound for individual interpretation.

\begin{theorem}[Interpretation Depth Sum Bound (19.16)]
\label{thm:interpretation_depth_sum_bound}
\[
  \texttt{depthSum}\bigl(\texttt{interpretations}(\text{reps}, s)\bigr)
  \leq \texttt{depthSum}(\text{reps}) + |\text{reps}| \cdot \depth(s).
\]
\end{theorem}

\noindent\textit{Proof sketch.}\quad
By induction on the list.  At each step, $\depth(\compose(r, s)) \leq
\depth(r) + \depth(s)$ by subadditivity, and the inductive hypothesis
bounds the tail.
\hfill$\square$

\medskip

\noindent\textbf{Commentary.}
The total hermeneutic effort of a community is bounded by their collective
background knowledge plus the signal's complexity distributed across all
members.  In resource-theoretic terms, the ``cognitive budget'' of a
community constrains its total interpretive output.  This bound has
practical implications for deliberative democracy: a community of $n$
agents, each of bounded sophistication $d$, encountering a signal of
depth $\depth(s)$, can produce total interpretive richness at most
$nd + n \cdot \depth(s)$---a linear bound in both the number of agents
and the signal complexity.\footnote{Compare this to the ``Condorcet jury
theorem'' in social choice theory, which provides conditions under which
larger groups make better decisions.  Our depth bound is not about
\emph{accuracy} but about \emph{total complexity}---a complementary
perspective.}

\begin{remark}[Connection to Bargaining Theory (Chapter~\ref{ch:bargaining})]
\label{rem:connection_bargaining}
The aggregation mechanics developed in this chapter have deep connections
to the bargaining theory of Chapter~\ref{ch:bargaining}.  In bargaining,
the order of composition matters: $\compose(a, b) \neq \compose(b, a)$
in general.  In hermeneutic aggregation, the analogous asymmetry appears
in $\texttt{aggregateCompose}$: the right-fold semantics privilege the
last speaker (who provides the ``base'') and the first speaker (whose
voice is outermost).  This order-dependence is not a defect of the
formalization but a feature of real deliberation: in any sequential
discussion, the order of contributions shapes the outcome.  The bargaining
power that Chapter~\ref{ch:bargaining} attributed to the first mover
reappears here as \emph{agenda-setting power}: the agent who structures
the order of deliberation---who speaks first, who speaks last---exerts
disproportionate influence on the aggregate meaning.

This connection suggests that a complete theory of hermeneutic aggregation
must incorporate the strategic considerations of bargaining theory.
The ``ideal speech situation'' in which order does not matter corresponds
precisely to the commutative case: $\compose(a, b) = \compose(b, a)$.
In real communicative action, commutativity fails, and the politics of
turn-taking becomes an ineliminable dimension of collective meaning-making.
\end{remark}

% ============================================================
\section{From Arrow to Habermas: Aggregation vs.\ Deliberation}
\label{sec:ch19-arrow-habermas}

The theorems proved in this chapter illuminate a fundamental tension in
democratic theory: the opposition between \emph{aggregative} and
\emph{deliberative} models of collective decision-making.

\subsection{The Aggregative Model}

In the aggregative model, individual judgments are treated as fixed inputs
to a mechanical procedure.  Social choice theory, from Condorcet through
Arrow and beyond, takes this approach: each agent has a preference ordering
(or, in our case, an interpretation), and the goal is to combine these
inputs into a collective output via some well-defined function.  Our
$\texttt{aggregateCompose}$ is such a function: it takes a list of
interpretations and produces a single composite idea.

The aggregative model has several virtues.  It is precise, compositional,
and amenable to formal analysis.  The depth bounds of this chapter
(Theorems~\ref{thm:singleton_aggregate_depth} and
\ref{thm:interpretation_depth_sum_bound}) are aggregative results: they
characterize the resources required for the mechanical combination of
interpretations.

However, the aggregative model also has well-known limitations.  Arrow's
theorem shows that no aggregation function satisfying certain axioms
exists; Sen's liberal paradox shows that individual rights and collective
rationality conflict; and Riker's instability thesis argues that
aggregative outcomes are inherently manipulable.\footnote{William H.\ Riker,
\textit{Liberalism Against Populism} (1982).}  These impossibility results
apply, \textit{mutatis mutandis}, to hermeneutic aggregation: the order-dependence of $\texttt{aggregateCompose}$ is a manifestation of the
sensitivity that aggregation theorists have long identified.

\subsection{The Deliberative Model}

In the deliberative model, individual judgments are not fixed inputs but
are transformed through the process of dialogue.  Habermas, Rawls, and
Cohen argue that legitimate collective decisions emerge not from the
aggregation of pre-existing preferences but from the transformation of
preferences through reasoned deliberation.\footnote{Joshua Cohen,
``Deliberation and Democratic Legitimacy,'' in \textit{The Good Polity}
(1989), ed.\ A.\ Hamlin and P.\ Pettit, pp.\ 17--34.}

In IDT terms, deliberation corresponds to a dynamic process in which
agents update their backgrounds $r_i$ in response to the interpretations
$\compose(r_j, s)$ of other agents.  After hearing agent $j$'s
interpretation, agent $i$ may revise their background from $r_i$ to
$r_i' = \compose(r_i, \compose(r_j, s))$, incorporating $j$'s perspective
into their own pre-understanding.  This iterated updating is a composition
chain---the same structure we studied in Chapter~\ref{ch:network}'s relay
chains, but now applied to the internal dynamics of a deliberating group.

\subsection{The IDT Synthesis}

IDT's contribution to this debate is to provide a framework in which
aggregation and deliberation are not opposed but \emph{nested}.  The
$\texttt{aggregateCompose}$ function is an aggregation procedure, but each
step involves a genuine composition---a ``fusion of horizons'' in Gadamer's
terms---rather than a mere concatenation.  The coherent interpretation
condition (Definition~\ref{def:coherentInterpretation}) identifies the
circumstances under which this compositional aggregation produces genuine
consensus rather than mere compilation.

The key insight is that \emph{resonance} bridges the aggregation--deliberation
divide.  When agents share pairwise resonance
(Theorem~\ref{thm:resonant_pair_coherent}), the aggregative procedure
automatically produces a deliberative outcome: the composite interpretation
is one that all agents can endorse, because their backgrounds are
sufficiently aligned to generate compatible readings.  When resonance is
absent, the aggregative procedure may fail to produce consensus, and
genuine deliberation---the updating of backgrounds through dialogue---becomes
necessary.

This framework suggests a normative program for democratic theory: design
institutions that foster resonance (through education, shared experience,
cross-cultural dialogue) so that the aggregative procedures of democracy
produce genuinely deliberative outcomes.  The formal apparatus of IDT
makes this program precise: the goal is to increase the proportion of
pairs $(r_i, r_j)$ that satisfy $\text{resonates}(r_i, r_j)$, thereby
expanding the domain of signals for which coherent interpretation is
achievable.

% ============================================================
\section{Conclusion}
\label{sec:ch19-concl}

This chapter has formalized the hermeneutic aggregation problem within
Ideatic Diffusion Theory, proving eighteen theorems that characterize the
structure of collective interpretation.

The central results form a coherent picture:
\begin{itemize}
  \item \textbf{Void baseline:} Void signal preserves backgrounds; void
    agents are transparent (Theorems~\ref{thm:void_signal_preserves}--\ref{thm:void_aggregate}).
  \item \textbf{Coherence conditions:} Empty and singleton groups are
    trivially coherent; resonant pairs achieve substantive consensus
    (Theorems~\ref{thm:empty_coherent}--\ref{thm:resonant_pair_coherent}).
  \item \textbf{Depth bounds:} Total interpretive effort is linearly bounded
    (Theorems~\ref{thm:singleton_aggregate_depth}--\ref{thm:interpretation_depth_sum_bound}).
  \item \textbf{Aggregation mechanics:} Sequential composition of
    interpretations is well-defined and compositional
    (Theorems~\ref{thm:aggregate_nil}--\ref{thm:aggregate_self_resonance}).
\end{itemize}

Habermas's theory of communicative action envisioned a form of rationality
grounded in mutual understanding rather than strategic calculation.  Our
formalization shows that the ``ideal speech situation'' has a precise
mathematical structure: it is the state in which all agents' interpretations
pairwise resonate (Definition~\ref{def:coherentInterpretation}), and this
state is achievable whenever the agents share cultural resonance
(Theorem~\ref{thm:resonant_pair_coherent}).

In the final chapter, we synthesize all the threads of the signalling-game
arc into a unified theory of meaning.

% ============================================================
%  Chapter 20 — Synthesis: The Game Theory of Meaning
% ============================================================
\chapter{Synthesis --- The Game Theory of Meaning}
\label{ch:synthesis}

\begin{quote}
\textit{``Man is an animal suspended in webs of significance he himself has
spun.''}\\
--- Clifford Geertz, \textit{The Interpretation of Cultures} (1973)
\end{quote}

\vspace{1em}

\section{Introduction}

We have arrived at the culmination of the signalling-game arc.  In the
preceding chapters, we formalized cooperation as composition
(Chapter~\ref{ch:cooperative}), bargaining as order-dependence
(Chapter~\ref{ch:bargaining}), diffusion as relay chains
(Chapter~\ref{ch:network}), and aggregation as hermeneutic consensus
(Chapter~\ref{ch:social}).  Each chapter isolated one dimension of the
problem of meaning.  This final chapter weaves them together into a unified
theory.

The synthesis draws on three philosophical traditions:

\begin{enumerate}
  \item \textbf{Gadamer's hermeneutic circle.}  Understanding is not a
    linear process from signal to meaning, but a circle in which the
    interpreter's pre-understanding shapes their reading, which in turn
    reshapes their pre-understanding.  In IDT, this is the associativity of
    iterated interpretation: $\compose(r_3, \compose(r_2, \compose(r_1, s)))$
    can be re-bracketed without changing the result.\footnote{Hans-Georg
    Gadamer, \textit{Truth and Method} (1960), argued that the hermeneutic
    circle is not a vicious circle but a spiral of deepening understanding.
    Our associativity theorems formalize this spiral.}

  \item \textbf{Ricoeur's surplus of meaning.}  Paul Ricoeur argued that
    every text generates more meaning than any single reader can extract.
    In IDT, this ``surplus'' is measured by \emph{interpretive richness}---the
    total depth of all interpretations---and is bounded by the universal
    depth theorem (Theorem~\ref{thm:universal_depth_bound_ch20}).

  \item \textbf{Geertz's thick description.}  Clifford Geertz proposed that
    cultural analysis is ``not an experimental science in search of law but
    an interpretive one in search of meaning.''  The meaning game
    (Definition~\ref{def:MeaningGame}) formalizes this interpretive quest:
    a community of receivers encounters a signal and produces a constellation
    of interpretations, each shaped by their own ``webs of significance.''
\end{enumerate}

The grand synthesis (Theorem~\ref{thm:grand_synthesis_ch20}) unifies the
resource-theoretic and hermeneutic dimensions: for any meaning game with
resonant receivers, the total interpretive richness is bounded \emph{and}
the game is in equilibrium.  Bounded resources and shared meaning are
simultaneously achievable---the formal heart of the game theory of meaning.

We proceed as follows.  Section~\ref{sec:ch20-defs} introduces the meaning
game, interpretive richness, and meaning equilibrium.
Section~\ref{sec:ch20-fundamental} establishes the fundamental theorems of
interpretation.  Section~\ref{sec:ch20-resonance} develops resonance
preservation results.  Section~\ref{sec:ch20-depth} derives depth bounds on
interpretive richness.  Section~\ref{sec:ch20-equilibrium} proves
equilibrium existence theorems.  Section~\ref{sec:ch20-synthesis} presents
the grand synthesis.  Section~\ref{sec:ch20-future} charts future
directions.  Section~\ref{sec:ch20-philosophy} offers philosophical
conclusions.  Section~\ref{sec:ch20-concl} concludes the entire
signalling-game arc.

\subsection{The Arc of the Book}

It is worth pausing to survey the path that has brought us here, for this
capstone chapter draws on every thread woven across the preceding nineteen
chapters.

We began, in the foundational chapters, with the bare algebraic structure
of \emph{ideatic space}: a set of ideas equipped with a composition
operation, a depth function, and a resonance relation.  Composition
($\compose$) formalized the act of combining ideas; depth ($\depth$)
measured interpretive complexity; resonance ($\text{resonates}$) captured
mutual comprehensibility.  These three primitives---composition, depth,
resonance---are the atoms from which all subsequent structures are built.

The middle chapters explored what happens when ideatic space is inhabited
by \emph{agents}---interpreters, senders, receivers, communities.
Chapter~\ref{ch:cooperative} showed that cooperation is composition,
with the void idea as the non-contributor.
Chapter~\ref{ch:bargaining} revealed the politics of composition order:
who composes first wields bargaining power.  Chapter~\ref{ch:network}
traced the diffusion of ideas through relay chains, establishing that
depth grows at most linearly and that resonant relayers cluster.
Chapter~\ref{ch:social} addressed the aggregation problem, proving
conditions under which a community achieves hermeneutic consensus.

Now, in this final chapter, we bring all these threads together.  The
meaning game (Definition~\ref{def:MeaningGame}) is the arena in which
composition, bargaining, diffusion, and aggregation interact simultaneously.
The grand synthesis (Theorem~\ref{thm:grand_synthesis_ch20}) is the
theorem that says they interact \emph{well}: bounded resources and shared
meaning are simultaneously achievable.

\subsection{Thick Description and Formal Methods}

Clifford Geertz's concept of ``thick description'' has been both the
inspiration and the challenge for this entire project.  Geertz argued that
cultural analysis must go beyond the mere recording of behavior (``thin
description'') to interpret the layers of meaning that behavior carries
(``thick description'').  His famous example---distinguishing a twitch from
a wink from a parody of a wink---illustrates the nested compositional
structure that IDT formalizes: each layer of interpretation is a composition
$\compose(r_k, \compose(r_{k-1}, \ldots \compose(r_1, s) \ldots))$, and
the ``thickness'' of the description corresponds to the depth of the
resulting composite.\footnote{Geertz's wink example appears in ``Thick
Description: Toward an Interpretive Theory of Culture,'' the opening essay
of \textit{The Interpretation of Cultures} (1973).  The example is borrowed
from Gilbert Ryle.}

The relationship between formal methods and humanistic inquiry is not
always comfortable.  Geertz himself was skeptical of formalization, warning
against the ``laws-and-instances'' model of explanation that treats cultural
phenomena as instances of general covering laws.\footnote{Geertz,
\textit{The Interpretation of Cultures}, p.\ 26.}  Our response is that
IDT does not propose covering laws but \emph{structural constraints}: the
subadditivity of depth, the preservation of resonance under composition,
the existence of equilibria for resonant communities.  These are not
predictions about what any particular community will do; they are bounds
on what any community \emph{can} do, given the algebraic structure of
ideatic space.  In this sense, IDT is closer to thermodynamics than to
Newtonian mechanics: it establishes the constraints within which cultural
processes operate, without specifying the trajectories they follow.

\subsection{Why Game Theory Needs Anthropology (and Vice Versa)}

Classical game theory, from von Neumann and Morgenstern through Nash and
Selten, models agents as rational utility-maximizers operating in
well-defined strategic environments.  This framework has produced powerful
results in economics, political science, and evolutionary biology.  Yet it
has limitations when applied to the domain of meaning.  Utility functions
presuppose commensurability---all outcomes can be ranked on a single
scale---while meaning is inherently multidimensional.  Nash equilibrium
presupposes individual rationality---each agent best-responds to others'
strategies---while meaning-making is fundamentally collaborative.

Anthropology, for its part, has long wrestled with the problem of meaning
but has largely eschewed formal methods.  The interpretive tradition from
Dilthey through Geertz has produced brilliant analyses of particular
cultural phenomena but has struggled to develop general theoretical
frameworks that permit comparison across cases.

IDT proposes a synthesis.  From game theory, it borrows the concept of
equilibrium---a state in which all agents' actions are mutually
consistent.  From anthropology, it borrows the concept of
interpretation---the active construction of meaning from signal and
background.  The result is meaning equilibrium: a state in which all
agents' \emph{interpretations} are mutually consistent, formalized as
pairwise resonance.  This is neither the strategic equilibrium of
game theory nor the thick description of anthropology, but a new concept
that integrates both traditions.

% ============================================================
\section{The Meaning Game}
\label{sec:ch20-defs}

\begin{definition}[Interpret All]
\label{def:interpretAll}
Let $\text{reps} = [r_1, \ldots, r_n]$ be a list of receiver backgrounds
and $s \in \IS$ be the signal.  The \emph{interpretation list} is:
\[
  \texttt{interpretAll}(\text{reps}, s) :=
  [\compose(r_1, s), \compose(r_2, s), \ldots, \compose(r_n, s)].
\]
\end{definition}

\begin{lstlisting}[caption={interpretAll in Lean}]
def interpretAll (reps : List I) (s : I) : List I :=
  reps.map (fun r => IdeaticSpace.compose r s)
\end{lstlisting}

\begin{definition}[Meaning Game]
\label{def:MeaningGame}
A \emph{meaning game} $G = (s, [r_1, \ldots, r_n])$ consists of:
\begin{itemize}
  \item a \emph{signal} $s \in \IS$ being sent;
  \item a list of \emph{receivers} $[r_1, \ldots, r_n]$, each with their
    own background idea.
\end{itemize}
The \emph{outcome} of the game is $\texttt{interpretAll}(\text{receivers}, s)$.
\end{definition}

\begin{lstlisting}[caption={Meaning game structure in Lean}]
structure MeaningGame (I : Type*) [IdeaticSpace I] where
  signal : I
  receivers : List I

def MeaningGame.outcome (g : MeaningGame I) : List I :=
  interpretAll g.receivers g.signal
\end{lstlisting}

In Geertz's terms, the meaning game formalizes the encounter between a
cultural artifact (signal) and a community of interpreters (receivers).
Each receiver brings their own ``webs of significance'' to the encounter,
producing a reading that is simultaneously individual and socially
conditioned.\footnote{Geertz's famous dictum---that culture is ``webs of
significance''---echoes Max Weber's definition of sociology as the
interpretive understanding of social action.  The meaning game formalizes
both simultaneously: the signal is the social action, and the
interpretations are the community's collective understanding.}

\begin{definition}[Meaning Equilibrium]
\label{def:MeaningEquilibrium}
A meaning game is in \emph{equilibrium} when all receivers' interpretations
pairwise resonate:
\[
  \texttt{MeaningEquilibrium}(G) \;:\Longleftrightarrow\;
  \forall\, x, y \in G.\text{outcome},\quad \text{resonates}(x, y).
\]
\end{definition}

\begin{lstlisting}[caption={Meaning equilibrium in Lean}]
def MeaningEquilibrium (g : MeaningGame I) : Prop :=
  forall x in g.outcome, forall y in g.outcome,
    IdeaticSpace.resonates x y
\end{lstlisting}

Meaning equilibrium is the game-theoretic analogue of Habermas's
``ideal speech situation'': a state in which all participants' interpretations
are mutually comprehensible.  Unlike Nash equilibrium, which concerns
individual rationality, meaning equilibrium concerns collective coherence---a
fundamentally different desideratum that reflects the hermeneutic rather than
strategic orientation of communicative action.

\begin{definition}[Interpretive Richness]
\label{def:interpretiveRichness}
The \emph{interpretive richness} of a group interpreting a signal is the
total depth of all interpretations:
\[
  \texttt{interpretiveRichness}(\text{reps}, s) :=
  \sum_{i} \depth(\compose(r_i, s)).
\]
This measures the total hermeneutic ``work'' done by the community.
\end{definition}

\begin{lstlisting}[caption={Interpretive richness in Lean}]
def interpretiveRichness (reps : List I) (s : I) : N :=
  depthSum (reps.map (fun r => IdeaticSpace.compose r s))
\end{lstlisting}

Ricoeur's ``surplus of meaning'' is operationalized as interpretive richness:
it is the total interpretive content generated by the encounter between
signal and receivers.  The richer the community's backgrounds and the deeper
the signal, the greater the surplus.

\begin{example}[The Rosetta Stone as Multi-Receiver Interpretation]
\label{ex:rosetta_stone}
The Rosetta Stone (196 BCE) bears the same decree inscribed in three
scripts: Egyptian hieroglyphic, Demotic, and Ancient Greek.  In the
framework of the meaning game, the decree is the signal $s$, and the
three scripts represent three ``receivers'' $r_{\text{hieroglyphic}}$,
$r_{\text{Demotic}}$, and $r_{\text{Greek}}$, each encoding a distinct
linguistic and cultural background.  The inscriptions are the interpretations
$\compose(r_i, s)$, and the fact that scholars could use the Greek
interpretation to decode the hieroglyphic one demonstrates a form of
cross-receiver resonance: $\text{resonates}(\compose(r_{\text{Greek}}, s),
\compose(r_{\text{hieroglyphic}}, s))$.

Jean-Fran\c{c}ois Champollion's decipherment in 1822 was itself a meaning
game: Champollion brought his own background $r_{\text{Champollion}}$
(knowledge of Coptic, Greek, and comparative linguistics) to the
interpretation of the already-interpreted stone.  His breakthrough was
the discovery of a composition path
$\compose(r_{\text{Champollion}}, \compose(r_{\text{Greek}}, s))$ that
resonated with $\compose(r_{\text{hieroglyphic}}, s)$---a cross-layer
resonance that unlocked two millennia of lost meaning.\footnote{See
Andrew Robinson, \textit{Cracking Codes: The Rosetta Stone and
Decipherment} (1994).}
\end{example}

\begin{example}[The Printing Press as Meaning Game Transformation]
\label{ex:printing_press}
Gutenberg's invention of movable type (c.\ 1440) did not merely change the
speed of information transmission; it transformed the structure of the
meaning game itself.  Before print, the meaning game for a given text
involved a small number of receivers---monks, scholars, wealthy
patrons---each with highly specialized backgrounds.  The interpretive
richness was concentrated: $\texttt{interpretiveRichness}(\text{reps}, s)$
was moderate because $|\text{reps}|$ was small, even though individual
$\depth(r_i)$ values were high.

Print dramatically expanded $|\text{reps}|$ while simultaneously
diversifying the backgrounds $r_i$: merchants, artisans, farmers, and
eventually the working class gained access to texts previously confined
to clerical elites.  The universal depth bound
(Theorem~\ref{thm:universal_depth_bound_ch20}) predicts the consequence:
the total interpretive richness grew proportionally, bounded by
$\texttt{depthSum}(\text{reps}) + |\text{reps}| \cdot \depth(s)$.

Yet the print revolution also disrupted meaning equilibrium.  The
pre-print interpretive community---the medieval \textit{universitas}---was
a resonant community in the sense of
Theorem~\ref{thm:resonant_community_equilibrium}: shared training, shared
Latin, shared theological assumptions ensured pairwise resonance.  The
post-print community was radically heterogeneous, and the resulting loss
of equilibrium manifested as the Reformation, the Wars of Religion, and
the eventual emergence of modern pluralism.\footnote{Elizabeth Eisenstein,
\textit{The Printing Press as an Agent of Change} (1979), provides the
definitive account of the cultural consequences of print.}
\end{example}

% ============================================================
\section{The Fundamental Theorems of Interpretation}
\label{sec:ch20-fundamental}

\begin{theorem}[Fundamental Theorem --- Void Signal Equilibrium (20.1)]
\label{thm:fundamental_void_equilibrium}
The void signal always admits equilibrium:
\[
  \texttt{MeaningEquilibrium}\bigl(\langle \void, \text{receivers} \rangle\bigr)
  \implies \top.
\]
\end{theorem}

\noindent\textit{Proof sketch.}\quad
The conclusion is $\top$ (trivially true).
\hfill$\square$

\medskip

\noindent\textbf{Commentary.}
In the absence of a text, each interpreter's horizon remains intact.
Gadamer's hermeneutic circle begins with pre-understanding, and when the
text contributes nothing (void), the circle collapses to its starting
point.  The void is the universal equilibrium signal---saying nothing
produces no disagreement.  While the theorem as stated is trivially true,
it establishes the conceptual foundation: void is the baseline against
which all other signals are measured.

\begin{theorem}[Void Signal Preserves Backgrounds (20.2)]
\label{thm:void_preserves_backgrounds}
$\texttt{interpretAll}(\text{reps}, \void) = \text{reps}$.
\end{theorem}

\noindent\textit{Proof sketch.}\quad
By induction on the list.  At each step, $\compose(r, \void) = r$ by
right-identity.
\hfill$\square$

\medskip

\noindent\textbf{Commentary.}
Without external stimuli, consciousness returns to itself.  Each receiver,
encountering the void signal, produces an ``interpretation'' that is simply
their own background idea, unchanged.  Phenomenologically, this is the
state of pure self-reflection: the subject without object collapses back
upon its own structure.\footnote{This echoes Husserl's concept of the
\textit{noema} (intentional object): when the noema is empty (void), the
\textit{noesis} (intentional act) has nothing to constitute, and the mind
rests in its own pre-intentional state.}

\begin{theorem}[Composition Closure (20.6)]
\label{thm:composition_closure}
Interpretations of interpretations are still compositions:
\[
  \compose(r_2, \compose(r_1, s)) = \compose(\compose(r_2, r_1), s).
\]
\end{theorem}

\noindent\textit{Proof sketch.}\quad
The associativity axiom of $\compose$.
\hfill$\square$

\medskip

\noindent\textbf{Commentary.}
This is the formal content of Geertz's ``thick description'':
interpretation layered upon interpretation is still interpretation.
There is no escape from the hermeneutic circle, only deeper embedding
within it.  The second interpreter does not stand ``outside'' the first
interpretation; rather, their reading \emph{composes} with the first
reading, producing a compound interpretation that is itself an element of
ideatic space.

This theorem has profound epistemological consequences.  It means that
meta-interpretation---the interpretation of interpretations---is not a
different \emph{kind} of activity from first-order interpretation.  It is
the same operation, applied at a higher level of composition.  The
hermeneutic circle is not a vicious regress but a stable compositional
structure.\footnote{This parallels Ricoeur's argument in \textit{Hermeneutics
and the Human Sciences} (1981) that explanation and understanding are
not opposed but complementary moments of a single interpretive arc.}

\begin{theorem}[Void Receiver Transparency (20.14)]
\label{thm:void_receiver_transparent}
$\compose(\void, s) = s$.
\end{theorem}

\noindent\textit{Proof sketch.}\quad Left-identity of $\compose$. \hfill$\square$

\medskip

\noindent\textbf{Commentary.}
A receiver with void background returns the signal unchanged---a
\textit{tabula rasa} is a perfect mirror.  The empty mind contributes
nothing to the interpretation, simply reflecting the signal as received.

% ============================================================
\section{Resonance Preservation}
\label{sec:ch20-resonance}

The following theorems establish that resonance---the fundamental
compatibility relation---is preserved through the meaning game.

\begin{theorem}[Resonance Preservation --- Resonant Backgrounds (20.4)]
\label{thm:resonance_preservation}
If $\text{resonates}(r_1, r_2)$, then for all $s$,
\[
  \text{resonates}(\compose(r_1, s), \compose(r_2, s)).
\]
\end{theorem}

\noindent\textit{Proof sketch.}\quad
From \texttt{resonance\_compose\_right}.
\hfill$\square$

\medskip

\noindent\textbf{Commentary.}
Agents sharing ``collective consciousness'' (Durkheim) will produce
consonant interpretations of any cultural artifact.  Resonance between
backgrounds propagates through the interpretive act to resonance between
interpretations.  This is the mechanism underlying cultural consonance:
shared pre-understanding generates shared understanding.

\begin{theorem}[Resonance Preservation --- Resonant Signals (20.5)]
\label{thm:resonance_signal_preservation}
If $\text{resonates}(s_1, s_2)$, then for all $r$,
\[
  \text{resonates}(\compose(r, s_1), \compose(r, s_2)).
\]
\end{theorem}

\noindent\textit{Proof sketch.}\quad
From \texttt{resonance\_compose\_left}.
\hfill$\square$

\medskip

\noindent\textbf{Commentary.}
A reader encountering similar texts (resonant signals) produces similar
interpretations.  This captures the literary-critical intuition that
intertextuality creates resonance: texts that ``echo'' each other produce
echoing interpretations in any given reader.

\begin{theorem}[Interpretation Self-Resonance (20.11)]
\label{thm:interpretation_self_resonance}
$\text{resonates}(\compose(r, s), \compose(r, s))$.
\end{theorem}

\noindent\textit{Proof sketch.}\quad Self-resonance. \hfill$\square$

\medskip

\noindent\textbf{Commentary.}
Every act of understanding is internally coherent.  No matter how
idiosyncratic an interpretation may be, it resonates with itself---it is
a valid element of ideatic space, self-consistent by the reflexivity of
resonance.

\begin{example}[The COVID-19 Infodemic as a Meaning Game Under Stress]
\label{ex:covid_infodemic}
The COVID-19 pandemic produced what the World Health Organization termed
an ``infodemic'': a massive, rapidly evolving meaning game in which the
signal $s$ (scientific findings about the virus, public health directives,
epidemiological data) was interpreted by billions of receivers with
wildly divergent backgrounds.  Epidemiologists, politicians, conspiracy
theorists, traditional healers, and ordinary citizens each brought their
own $r_i$ to the encounter with $s$, producing interpretations
$\compose(r_i, s)$ that frequently failed to resonate with one another.

The infodemic illustrates several of our theorems simultaneously.  The
failure of meaning equilibrium---the fact that
$\texttt{MeaningEquilibrium}(\langle s, \text{global\_population} \rangle)$
manifestly did not hold---followed from the absence of universal pairwise
resonance among receivers.  The depth bound
(Theorem~\ref{thm:universal_depth_bound_ch20}) implies that the total
interpretive richness was bounded by the aggregate depth of the global
population plus the population size times the depth of the scientific
signal---an enormous but finite number.  And the resonant community
equilibrium theorem (Theorem~\ref{thm:resonant_community_equilibrium})
explains why meaning equilibrium \emph{was} achievable within
subcommunities (scientific communities, religious congregations, political
factions) whose members shared pairwise resonance, even as global
equilibrium remained elusive.\footnote{See Cinelli et al., ``The
COVID-19 Social Media Infodemic,'' \textit{Scientific Reports} 10, 16598
(2020), for an empirical analysis of echo chambers during the pandemic.}
\end{example}

\begin{example}[Translation of Sacred Texts Across Cultures]
\label{ex:sacred_translation}
The translation of sacred texts---the Bible into Latin (the Vulgate), into
vernacular languages (Luther's German Bible, the King James Version), the
Quran into non-Arabic languages, Buddhist sutras from Sanskrit to Chinese
and Japanese---presents meaning games of extraordinary complexity and
stakes.

Each translation is an interpretation $\compose(r_{\text{translator}}, s)$
where $s$ is the source text and $r_{\text{translator}}$ encodes the
linguistic competence, theological commitments, and cultural context of
the translator.  Jerome's Vulgate, Luther's Bible, and the King James
Version are three different interpretations of overlapping source signals,
each produced by a receiver with profoundly different backgrounds.  The
\emph{resonance} among these interpretations is partial: they agree on
certain core narratives and theological claims while diverging on matters
of emphasis, nuance, and doctrinal implication.

The meaning equilibrium question for sacred texts is the question of
religious unity: under what conditions do diverse communities, reading
diverse translations, arrive at resonant interpretations?  The history
of Christianity suggests that meaning equilibrium for sacred texts is
inherently fragile: the Reformation, the Great Schism, and innumerable
doctrinal disputes all represent failures of equilibrium---moments when
the interpretive community's lack of universal resonance produced
irreconcilable readings.\footnote{See Eugene Nida, \textit{Toward a
Science of Translating} (1964), for a theoretical framework for
translation that anticipates several themes of IDT.}
\end{example}

\begin{theorem}[Associativity of Iterated Interpretation (20.19)]
\label{thm:iterated_interpretation_assoc}
\[
  \compose(r_3, \compose(r_2, \compose(r_1, s))) =
  \compose(\compose(r_3, r_2), \compose(r_1, s)).
\]
\end{theorem}

\noindent\textit{Proof sketch.}\quad Associativity of $\compose$. \hfill$\square$

\medskip

\noindent\textbf{Commentary.}
Gadamer's hermeneutic circle has no privileged entry point.  Re-bracketing
the layers of interpretation---whether we first combine $r_3$ with $r_2$
and then interpret $r_1$'s reading, or interpret $r_2$'s reading of $r_1$'s
reading and then apply $r_3$---does not change the result.  The
``entry point'' into the circle is arbitrary; the compositional structure is
invariant.

% ============================================================
\section{Depth Bounds on Interpretive Richness}
\label{sec:ch20-depth}

\begin{theorem}[Universal Depth Bound on Interpretive Richness (20.3)]
\label{thm:universal_depth_bound_ch20}
For all lists of receivers $\text{reps}$ and signals $s$,
\[
  \texttt{interpretiveRichness}(\text{reps}, s) \;\leq\;
  \texttt{depthSum}(\text{reps}) + |\text{reps}| \cdot \depth(s).
\]
\end{theorem}

\noindent\textit{Proof sketch.}\quad
By induction on the list of receivers.  At each step, $\depth(\compose(r, s))
\leq \depth(r) + \depth(s)$ by subadditivity, and the inductive hypothesis
bounds the tail.
\hfill$\square$

\medskip

\noindent\textbf{Commentary.}
This is the \emph{fundamental resource theorem of hermeneutics}: the surplus
of meaning is bounded.  Even the richest text can only generate bounded
total interpretation, proportional to the community's existing hermeneutic
capacity plus the text's complexity.

Ricoeur argued that the surplus of meaning is \emph{inexhaustible}---that a
great text always yields new readings.  Our theorem does not contradict this
claim, but qualifies it: the \emph{total} depth of all readings is bounded,
even if the \emph{number} of distinct readings is unbounded.  In other
words, the surplus is bounded in intensity (depth) even if unbounded in
extension (variety).\footnote{The distinction between intensive and
extensive surplus of meaning parallels the distinction between depth and
breadth in information theory.  A text of finite depth can generate
infinitely many shallow readings, but the total depth is constrained.}

\begin{theorem}[Double Interpretation Depth Bound (20.12)]
\label{thm:double_interp_depth}
\[
  \depth\bigl(\compose(r_2, \compose(r_1, s))\bigr) \;\leq\;
  \depth(r_1) + \depth(r_2) + \depth(s).
\]
\end{theorem}

\noindent\textit{Proof sketch.}\quad
Two applications of subadditivity, combined by transitivity.
\hfill$\square$

\medskip

\noindent\textbf{Commentary.}
Each pass through the hermeneutic circle adds bounded complexity.
Understanding deepens, but not without limit per step.  The depth of a
doubly-interpreted signal is bounded by the sum of the two interpreters'
depths and the signal's depth---a tight three-term bound.

\begin{remark}[Connections to Computational Complexity]
\label{rem:computational_complexity}
The depth bounds established in this section have natural computational
analogues.  If we view $\depth$ as a measure of computational complexity
(e.g., circuit depth or proof length), then
Theorem~\ref{thm:universal_depth_bound_ch20} implies that the total
computational work required to ``interpret'' all receivers is
polynomial in the input parameters---a tractability result.
Theorem~\ref{thm:double_interp_depth} implies that iterated interpretation
(composition chains of length $k$) has depth at most
$\sum_{i=1}^{k} \depth(r_i) + \depth(s)$---linear in the chain length.

These bounds suggest that the meaning game, while semantically rich, is
computationally well-behaved.  The question of whether \emph{finding} an
optimal signal (one that maximizes interpretive richness for a given
receiver community) is computationally tractable remains open and connects
to fundamental questions in algorithmic game theory and mechanism
design.\footnote{Compare to the computational complexity results for Nash
equilibrium computation: Daskalakis, Goldberg, and Papadimitriou showed
that computing Nash equilibria is PPAD-complete (2009).  The corresponding
question for meaning equilibrium---whether checking or computing meaning
equilibrium is tractable---is unexplored.}
\end{remark}

\begin{theorem}[Void Signal Richness (20.9)]
\label{thm:void_signal_richness}
$\texttt{interpretiveRichness}(\text{reps}, \void) = \texttt{depthSum}(\text{reps})$.
\end{theorem}

\noindent\textit{Proof sketch.}\quad
By induction, using $\compose(r, \void) = r$ at each step.
\hfill$\square$

\medskip

\noindent\textbf{Commentary.}
The ``cost'' of interpretation comes entirely from the interpreters'
backgrounds when the signal carries no content.  In the economics of
meaning, the void signal has zero marginal cost: all the interpretive
richness is pre-existing, contributed by the receivers' own depth.

\begin{theorem}[Interpretation Depth Non-Negativity (20.15)]
\label{thm:richness_nonneg}
$0 \leq \texttt{interpretiveRichness}(\text{reps}, s)$.
\end{theorem}

\noindent\textit{Proof sketch.}\quad Trivial for natural numbers. \hfill$\square$

\medskip

\noindent\textbf{Commentary.}
Interpretation always produces non-negative hermeneutic content.  While
trivial, this establishes a conceptual baseline: there is no such thing as
``negative interpretation.''  Every encounter between mind and text
generates at least zero meaning.

\begin{theorem}[Empty Repertoire Zero Richness (20.16)]
\label{thm:empty_richness}
$\texttt{interpretiveRichness}([\,], s) = 0$.
\end{theorem}

\noindent\textit{Proof sketch.}\quad Sum over empty list is zero. \hfill$\square$

\medskip

\noindent\textbf{Commentary.}
A message with no audience generates no meaning.  This is the
information-theoretic counterpart to the philosophical puzzle: if a tree
falls in a forest and no one is there to hear it, does it make a sound?
In IDT, the answer is unambiguous: no receivers, zero richness.

\begin{theorem}[Singleton Richness Bound (20.17)]
\label{thm:singleton_richness_bound}
$\texttt{interpretiveRichness}([r], s) \leq \depth(r) + \depth(s)$.
\end{theorem}

\noindent\textit{Proof sketch.}\quad
The richness of a singleton list is $\depth(\compose(r,s))$, bounded by
subadditivity.
\hfill$\square$

\medskip

\noindent\textbf{Commentary.}
One person's interpretive effort is bounded by their background plus the
signal's complexity.  Individual hermeneutics operates under the same
resource constraints as collective hermeneutics, scaled down to a single
agent.

% ============================================================
\section{Equilibrium Existence}
\label{sec:ch20-equilibrium}

\begin{theorem}[Empty Game Equilibrium (20.7)]
\label{thm:empty_game_equilibrium}
For all $s \in \IS$, $\texttt{MeaningEquilibrium}(\langle s, [\,] \rangle)$.
\end{theorem}

\noindent\textit{Proof sketch.}\quad
Universal quantification over the empty outcome list is vacuously true.
\hfill$\square$

\medskip

\noindent\textbf{Commentary.}
A meaning game with no receivers is trivially in equilibrium---vacuous
consensus.  Logic guarantees agreement when there are no parties to
disagree.

\begin{theorem}[Singleton Game Equilibrium (20.8)]
\label{thm:singleton_equilibrium}
For all $r, s$, $\texttt{MeaningEquilibrium}(\langle s, [r] \rangle)$.
\end{theorem}

\noindent\textit{Proof sketch.}\quad
The only outcome element is $\compose(r, s)$, which resonates with itself.
\hfill$\square$

\medskip

\noindent\textbf{Commentary.}
A single interpreter always agrees with themselves---solipsism is
equilibrium.  The interesting cases arise when multiple receivers must
cohere.

\begin{theorem}[Trivial Equilibrium Exists (20.10)]
\label{thm:trivial_equilibrium_exists}
$\texttt{MeaningEquilibrium}(\langle \void, [\,] \rangle)$.
\end{theorem}

\noindent\textit{Proof sketch.}\quad Vacuously true. \hfill$\square$

\medskip

\noindent\textbf{Commentary.}
The ``babbling equilibrium'' trivially exists when there is no audience.
Saying nothing to no one creates no disagreement---the most degenerate
form of communicative success.

\begin{theorem}[Resonant Community Equilibrium (20.13)]
\label{thm:resonant_community_equilibrium}
If all receivers pairwise resonate, then the meaning game is in equilibrium
for \emph{any} signal:
\[
  \bigl(\forall\, r_1, r_2 \in \text{receivers},\;
  \text{resonates}(r_1, r_2)\bigr)
  \implies
  \texttt{MeaningEquilibrium}(\langle s, \text{receivers} \rangle).
\]
\end{theorem}

\noindent\textit{Proof sketch.}\quad
For any $x, y$ in the outcome, $x = \compose(r_i, s)$ and
$y = \compose(r_j, s)$ for some $r_i, r_j \in \text{receivers}$.
By hypothesis, $\text{resonates}(r_i, r_j)$, and by
Theorem~\ref{thm:resonance_preservation},
$\text{resonates}(\compose(r_i, s), \compose(r_j, s))$.
\hfill$\square$

\medskip

\noindent\textbf{Commentary.}
This is Durkheim's grand theorem, formalized: a community with strong
``collective consciousness'' (universal pairwise resonance) achieves
interpretive consensus regardless of the stimulus.  Shared culture makes
shared understanding automatic.  The theorem does not require any
assumption about the signal---\emph{any} signal produces equilibrium,
provided the receivers share resonance.

The implications for social theory are significant.  Homogeneous communities
achieve meaning equilibrium effortlessly; diverse communities may not.
This is neither a normative claim nor a defense of homogeneity---it is a
structural fact about the interaction between resonance and composition.
The challenge for pluralistic societies is to achieve equilibrium without
requiring universal resonance, a problem we leave for future
work.\footnote{Compare this to the sociological concept of ``bridging
social capital'' (Robert Putnam): connections between diverse groups that
facilitate mutual understanding despite the absence of deep cultural
resonance.}

\begin{remark}[Limitations of the Formal Framework]
\label{rem:limitations}
The meaning game formalization, powerful as it is, necessarily abstracts
away several features of real interpretive encounters that deserve
acknowledgment.

First, our model is \emph{synchronic}: all receivers interpret the signal
simultaneously, with no temporal dynamics.  In reality, interpretation
unfolds over time, and later interpreters are influenced by earlier ones---a
dynamic that our relay-chain model (Chapter~\ref{ch:network}) partially
captures but does not fully integrate with the meaning game.

Second, our model assumes a single shared signal $s$.  In practice,
signals are noisy, partial, and often contested.  Different receivers may
have access to different versions or fragments of the signal, introducing
an information-asymmetry dimension that our framework does not address.

Third, resonance is binary: two ideas either resonate or they do not.
A graded notion of resonance---measuring \emph{how much} two
interpretations agree, not merely \emph{whether} they agree---would
permit a more nuanced theory of partial consensus and approximate
equilibrium.

Fourth, depth is scalar.  A vector-valued depth function---measuring
different dimensions of interpretive complexity (logical, emotional,
aesthetic, practical)---might better capture the multidimensional character
of real meaning-making.

These limitations are not defects but invitations: each points toward a
natural extension of the framework that future work might pursue.
\end{remark}

\begin{remark}[What IDT Captures and What It Misses]
\label{rem:captures_and_misses}
IDT captures the \emph{structural} aspects of meaning-making: the algebra
of composition, the resource constraints of depth, and the compatibility
conditions of resonance.  These structural features are genuine and
important: they hold across cultures, historical periods, and media.

What IDT misses---what any formal theory of meaning must miss---is the
\emph{phenomenological} dimension: what it \emph{feels like} to understand,
to be moved by a poem, to suddenly grasp a mathematical proof, to
experience the fusion of horizons that Gadamer describes.  IDT can tell
you that $\compose(r, s)$ has depth $d$ and resonates with $\compose(r', s)$;
it cannot tell you why depth $d$ feels profound or why resonance feels like
understanding.  The ``hard problem'' of cultural consciousness---why formal
structure gives rise to lived meaning---is beyond the reach of any
algebraic theory.

This limitation is not unique to IDT; it is shared by all formal theories
of mind, language, and culture.  The value of formalization is not that it
replaces phenomenological inquiry but that it provides a scaffolding on
which phenomenological insights can be organized, compared, and tested.
The meaning game is a map, not the territory; but maps, when accurate, are
indispensable for navigation.
\end{remark}

\begin{theorem}[Game Outcome Length (20.18)]
\label{thm:outcome_length}
$|G.\text{outcome}| = |G.\text{receivers}|$.
\end{theorem}

\noindent\textit{Proof sketch.}\quad
The \texttt{map} operation preserves list length.
\hfill$\square$

\medskip

\noindent\textbf{Commentary.}
Every voice produces exactly one interpretation.  The meaning game's
outcome has the same cardinality as its receiver list---a ``one person,
one interpretation'' principle that ensures representational completeness.

% ============================================================
\section{The Grand Synthesis}
\label{sec:ch20-synthesis}

We now arrive at the culminating theorem of the entire signalling-game arc.

\begin{theorem}[The Grand Synthesis (20.20)]
\label{thm:grand_synthesis_ch20}
For any meaning game with resonant receivers:
\begin{enumerate}
  \item[(1)] The total interpretive richness is bounded:
  \[
    \texttt{interpretiveRichness}(\text{receivers}, s) \;\leq\;
    \texttt{depthSum}(\text{receivers}) + |\text{receivers}| \cdot \depth(s).
  \]
  \item[(2)] The game is in equilibrium:
  \[
    \texttt{MeaningEquilibrium}(\langle s, \text{receivers} \rangle).
  \]
\end{enumerate}
Both properties hold simultaneously whenever
$\forall\, r_1, r_2 \in \text{receivers},\; \text{resonates}(r_1, r_2)$.
\end{theorem}

\begin{lstlisting}[caption={The grand synthesis in Lean}]
theorem grand_synthesis (receivers : List I) (s : I)
    (hres : forall r1 in receivers, forall r2 in receivers,
            IdeaticSpace.resonates r1 r2) :
    interpretiveRichness receivers s <=
      depthSum receivers + receivers.length * IdeaticSpace.depth s /\
    MeaningEquilibrium { signal := s, receivers := receivers } :=
  <universal_depth_bound receivers s,
   resonant_community_equilibrium receivers s hres>
\end{lstlisting}

\noindent\textit{Proof sketch.}\quad
Clause (1) is the universal depth bound
(Theorem~\ref{thm:universal_depth_bound_ch20}), which holds for \emph{all}
receiver lists, regardless of resonance.  Clause (2) is the resonant
community equilibrium (Theorem~\ref{thm:resonant_community_equilibrium}),
which requires the resonance hypothesis.  The conjunction of both clauses
follows by pairing.
\hfill$\square$

\bigskip

\noindent\textbf{Commentary.}
The grand synthesis is the formal heart of this book.  It says that for a
culturally coherent community---one in which all members share pairwise
resonance---two fundamental properties hold \emph{simultaneously}:

\begin{enumerate}
  \item \textbf{Resource boundedness.}  The total interpretive effort of the
    community is finite and bounded by a linear function of the community's
    existing depth and the signal's complexity.  Meaning-making is not free;
    it consumes cognitive resources, and the total expenditure is
    constrained.

  \item \textbf{Hermeneutic consensus.}  All interpretations are mutually
    comprehensible.  The community achieves shared meaning---not by
    suppressing individual differences (each receiver still composes with
    their own background), but because their backgrounds are sufficiently
    aligned that their readings ``rhyme'' with each other.
\end{enumerate}

The philosophical significance of this conjunction cannot be overstated.
It unifies the \emph{economic} perspective (resources are bounded) with the
\emph{hermeneutic} perspective (meaning is shared).  Previous accounts of
meaning-making have typically emphasized one dimension at the expense of the
other: economists focus on information costs, hermeneuticians on
interpretive depth.  The grand synthesis shows that both concerns are
addressed simultaneously by a single structural condition---resonance.

\medskip

\noindent\textbf{Geertz revisited.}
Clifford Geertz wrote that culture is ``webs of significance'' that humans
spin and within which they are suspended.  The grand synthesis formalizes
this metaphor with mathematical precision.  The ``webs'' are the resonance
relations between community members; the ``significance'' is the
interpretive richness generated by their encounter with signals.  The
theorem shows that when the web is coherent (universal pairwise resonance),
the significance is both bounded (finite web) and shared (equilibrium).

\medskip

\noindent\textbf{Ricoeur revisited.}
Paul Ricoeur's ``surplus of meaning'' is precisely the interpretive richness
of the community.  The universal depth bound shows that this surplus is
finite---not inexhaustible, as Ricoeur sometimes suggested, but bounded by
the community's cognitive resources.  Yet within this bound, the surplus can
be substantial: a deep community encountering a deep text generates rich
interpretations.

\medskip

\noindent\textbf{Gadamer revisited.}
Gadamer's hermeneutic circle is formalized by the associativity of
composition (Theorem~\ref{thm:composition_closure}).  The grand synthesis
shows that the circle is not merely internally consistent (associativity)
but also \emph{convergent}: when the interpreters share resonance, the
circle produces equilibrium.  The hermeneutic circle is a spiral, not a
vortex---it converges on shared meaning rather than spinning into
interpretive chaos.

% ============================================================
\section{Conclusion: The Game Theory of Meaning}
\label{sec:ch20-concl}

This chapter---and this entire signalling-game arc---has established a
formal theory of meaning grounded in the algebraic structure of ideatic
space.  Twenty theorems in this chapter, building on the seventy-four
theorems of Chapters~\ref{ch:cooperative}--\ref{ch:social}, constitute a
comprehensive mathematical framework for analyzing how meaning is created,
transmitted, negotiated, and shared.

Let us recapitulate the arc:

\begin{description}
  \item[Chapter~\ref{ch:cooperative}] established that cooperation is
    composition, with void as the non-contributor, associativity ensuring
    well-defined multi-party projects, and depth bounding the output
    complexity.

  \item[Chapter~\ref{ch:bargaining}] showed that non-commutativity creates
    bargaining power: who speaks first shapes the meaning.  Yet structural
    properties---depth bounds and void transparency---are order-invariant.

  \item[Chapter~\ref{ch:network}] formalized cultural diffusion as relay
    chains, proving that depth grows at most linearly, void relayers are
    transparent, and resonant relayers cluster.

  \item[Chapter~\ref{ch:social}] addressed the aggregation problem: how
    individual interpretations compose into collective meaning, and under
    what conditions (resonance) the community achieves hermeneutic
    consensus.

  \item[Chapter~\ref{ch:synthesis}] (this chapter) unified all threads
    into the grand synthesis: resonant communities achieve both bounded
    richness and meaning equilibrium.
\end{description}

The game theory of meaning differs from classical game theory in a
fundamental way.  In classical game theory, agents pursue individual
payoffs and equilibria emerge from the alignment of selfish strategies.
In the game theory of meaning, agents pursue \emph{understanding}---the
production of interpretations that resonate with those of their
fellow-interpreters---and equilibria emerge from the alignment of
\emph{backgrounds}.  The currency is not utility but resonance; the
constraint is not budget but depth; the solution concept is not Nash
equilibrium but meaning equilibrium.

This formal framework opens several directions for future work:

\begin{itemize}
  \item \textbf{Non-resonant communities.}  What happens when receivers
    do \emph{not} pairwise resonate?  Can partial equilibria be defined?
    Under what conditions does deliberation converge to equilibrium even
    without initial resonance?

  \item \textbf{Dynamic meaning games.}  What happens when the signal
    changes over time, or when receivers update their backgrounds based on
    the interpretations of others?  Can we formalize the dynamics of
    cultural change within this framework?

  \item \textbf{Strategic interpretation.}  What happens when receivers
    are not purely hermeneutic but strategically choose their interpretive
    frames?  Can the game theory of meaning be extended to encompass
    strategic signalling in the tradition of Crawford and Sobel?

  \item \textbf{Computational complexity.}  How hard is it to compute
    interpretive richness, check meaning equilibrium, or find the optimal
    signal for a given community?  Are there polynomial-time algorithms,
    or is the problem inherently intractable?
\end{itemize}

We began this arc with Mauss's insight that gift exchange is never merely
economic.  We end with a generalization: \emph{meaning exchange} is never
merely informational.  It is compositional, order-dependent, depth-bounded,
and resonance-conditioned.  The game theory of meaning provides the formal
tools to analyze these properties with mathematical rigor, while remaining
faithful to the anthropological and philosophical traditions that inspired
them.

\bigskip

\begin{center}
\textit{``The locus of the study is not the object of study.\\
Anthropologists don't study villages; they study in villages.''}\\
--- Clifford Geertz, \textit{The Interpretation of Cultures}
\end{center}


% ══════════════════════════════════════════════════════════════════════
\backmatter

% ── Bibliography ──────────────────────────────────────────────────────
\begin{thebibliography}{99}

\bibitem{austin1962} Austin, J.L. (1962). \textit{How to Do Things with Words}. Oxford University Press.

\bibitem{bartlett1932} Bartlett, F.C. (1932). \textit{Remembering: A Study in Experimental and Social Psychology}. Cambridge University Press.

\bibitem{bateson1972} Bateson, G. (1972). \textit{Steps to an Ecology of Mind}. University of Chicago Press.

\bibitem{bourdieu1984} Bourdieu, P. (1984). \textit{Distinction: A Social Critique of the Judgement of Taste}. Harvard University Press.

\bibitem{bourdieu1991} Bourdieu, P. (1991). \textit{Language and Symbolic Power}. Harvard University Press.

\bibitem{crawford1982} Crawford, V.P. \& Sobel, J. (1982). Strategic information transmission. \textit{Econometrica}, 50(6), 1431--1451.

\bibitem{dawkins1976} Dawkins, R. (1976). \textit{The Selfish Gene}. Oxford University Press.

\bibitem{derrida1967} Derrida, J. (1967). \textit{Of Grammatology}. Johns Hopkins University Press.

\bibitem{durkheim1893} Durkheim, \'E. (1893). \textit{The Division of Labor in Society}. Free Press.

\bibitem{durkheim1912} Durkheim, \'E. (1912). \textit{The Elementary Forms of Religious Life}. Free Press.

\bibitem{eco1976} Eco, U. (1976). \textit{A Theory of Semiotics}. Indiana University Press.

\bibitem{foucault1972} Foucault, M. (1972). \textit{The Archaeology of Knowledge}. Pantheon Books.

\bibitem{gadamer1960} Gadamer, H.-G. (1960). \textit{Truth and Method}. Continuum.

\bibitem{geertz1973} Geertz, C. (1973). \textit{The Interpretation of Cultures}. Basic Books.

\bibitem{goffman1959} Goffman, E. (1959). \textit{The Presentation of Self in Everyday Life}. Anchor Books.

\bibitem{gramsci1971} Gramsci, A. (1971). \textit{Selections from the Prison Notebooks}. International Publishers.

\bibitem{habermas1984} Habermas, J. (1984). \textit{The Theory of Communicative Action}. Beacon Press.

\bibitem{keats1817} Keats, J. (1817). Letter to George and Tom Keats, 21 December 1817.

\bibitem{levistrauss1962} L\'evi-Strauss, C. (1962). \textit{The Savage Mind}. University of Chicago Press.

\bibitem{levistrauss1969} L\'evi-Strauss, C. (1969). \textit{The Elementary Structures of Kinship}. Beacon Press.

\bibitem{lewis1969} Lewis, D. (1969). \textit{Convention: A Philosophical Study}. Harvard University Press.

\bibitem{mauss1925} Mauss, M. (1925). \textit{The Gift}. Norton.

\bibitem{nash1950} Nash, J. (1950). Equilibrium points in $n$-person games. \textit{PNAS}, 36(1), 48--49.

\bibitem{peirce1931} Peirce, C.S. (1931--1958). \textit{Collected Papers}. Harvard University Press.

\bibitem{rappaport1999} Rappaport, R.A. (1999). \textit{Ritual and Religion in the Making of Humanity}. Cambridge University Press.

\bibitem{ricoeur1976} Ricoeur, P. (1976). \textit{Interpretation Theory: Discourse and the Surplus of Meaning}. Texas Christian University Press.

\bibitem{saussure1916} Saussure, F. de (1916). \textit{Course in General Linguistics}. Open Court.

\bibitem{schelling1960} Schelling, T.C. (1960). \textit{The Strategy of Conflict}. Harvard University Press.

\bibitem{searle1969} Searle, J.R. (1969). \textit{Speech Acts}. Cambridge University Press.

\bibitem{simmel1908} Simmel, G. (1908). \textit{Soziologie}. Duncker \& Humblot.

\bibitem{sosis2003} Sosis, R. (2003). Why aren't we all Hutterites? Costly signaling theory and religious behavior. \textit{Human Nature}, 14(2), 91--127.

\bibitem{spence1973} Spence, M. (1973). Job market signaling. \textit{QJE}, 87(3), 355--374.

\bibitem{sperber1996} Sperber, D. (1996). \textit{Explaining Culture: A Naturalistic Approach}. Blackwell.

\bibitem{sperberw1986} Sperber, D. \& Wilson, D. (1986). \textit{Relevance: Communication and Cognition}. Blackwell.

\bibitem{turner1969} Turner, V. (1969). \textit{The Ritual Process: Structure and Anti-Structure}. Aldine.

\bibitem{veblen1899} Veblen, T. (1899). \textit{The Theory of the Leisure Class}. Macmillan.

\bibitem{vygotsky1978} Vygotsky, L.S. (1978). \textit{Mind in Society}. Harvard University Press.

\bibitem{wittgenstein1953} Wittgenstein, L. (1953). \textit{Philosophical Investigations}. Blackwell.

\bibitem{zahavi1975} Zahavi, A. (1975). Mate selection---a selection for a handicap. \textit{Journal of Theoretical Biology}, 53(1), 205--214.

\end{thebibliography}

% ── Index placeholder ─────────────────────────────────────────────────
\chapter*{Index of Theorems}
\addcontentsline{toc}{chapter}{Index of Theorems}

The following theorems are formally verified in Lean~4 with zero
\texttt{sorry} axioms. Each theorem name links to the corresponding
Lean source file in the \texttt{FormalAnthropology/} directory.

\begin{small}
\begin{enumerate}
\item[\textbf{Ch.~1}] \texttt{total\_hermeneutic\_bound}, \texttt{max\_interp\_depth\_bound}, \texttt{silence\_is\_transparent}, \texttt{blank\_receiver}, \texttt{negative\_capability}, \texttt{durkheimian\_consensus}, \texttt{mythical\_transformation}, \texttt{turners\_communitas}, \texttt{bartletts\_serial\_reproduction}, \texttt{telephone\_depth}, \texttt{resonant\_transmission\_lines}, \texttt{void\_is\_unpersuasive}, \texttt{interpretation\_conserves\_count}, \texttt{durkheims\_division\_of\_labor}, \texttt{silence\_preserves\_depthSum}, \texttt{three\_stage\_telephone}

\item[\textbf{Ch.~2--5}] See \texttt{IDT\_Signal\_Ch02--05.lean}
\item[\textbf{Ch.~6--10}] See \texttt{IDT\_Signal\_Ch06--10.lean}
\item[\textbf{Ch.~11--15}] See \texttt{IDT\_Signal\_Ch11--15.lean}
\item[\textbf{Ch.~16--20}] See \texttt{IDT\_Signal\_Ch16--20.lean}
\end{enumerate}
\end{small}

\end{document}
